% ==============================================================================
% Part VI: Full-Text Anthology across Five Academic Domains
% ==============================================================================
\part{Full-Text Anthology \\ 全篇范文与段落蓝图链拆解文库}\label{part:anthology}

\chapter{Computer Science \& Artificial Intelligence / 计算机科学与人工智能}\label{ch:anthology_18}
\Needspace{6\baselineskip}
本章精选该学科领域及对应文体的完整学术范文。每一篇范文不仅展示母语级发表标准的正文全貌，更清晰拆解全篇的“段落蓝图链条（Structural Blueprint Chain）”，并对每个自然段逐一剖析其核心交际意图、句际衔接逻辑公式（Cohesion Formula）以及每句话的句法成分与推进机制。

\section{FT01: 非平稳多变量时间序列预测完整小论文（DyTS-Net: Decoupling Spectral Phase and Continuous Scale）}\label{comp:FT01}
\begin{ConceptBox}{Composition Architecture \& Structural Blueprint Chain}
\textbf{Genre:} \texttt{AI\_Mini\_Paper} \quad \textbf{Paragraph Count:} 12\par
\textbf{Blueprint Chain:} \texttt{BP08 $\to$ BP01 $\to$ BP02 $\to$ BP03 $\to$ BP04 $\to$ BP07 $\to$ BP11 $\to$ BP12 $\to$ BP13 $\to$ BP18 $\to$ BP19 $\to$ BP24}
\end{ConceptBox}
\subsection*{Paragraph 1: 高密度四句式 Abstract 全段构建（Complete 4-Move Abstract） (\texttt{BP08})}
\noindent\textbf{Expressive Intent:}\enspace 我这一段想写出整篇小论文的完整 Abstract（摘要）：在 150--180 词内一气呵成交代背景、核心痛点、方法机制与量化实验结果。\par
\noindent\textbf{Cohesion Formula:}\enspace {\small\color{AcadSlate}Context (S1) $\to$ Gap (S2) $\to$ Method \& Mechanism (S3-S4) $\to$ Quantitative Proof (S5). 绑定逻辑规则：L001, L008, L010, L029.}\par\vspace{4pt}
\begin{ExampleBox}{Publication-Grade Authentic Text}
Recovering missing sensor observations is essential for reliable downstream monitoring in industrial and clinical time series. While diffusion-based imputation models generate plausible trajectories under random missingness, they suffer from high sampling latency and lose global periodic coherence when entire sensor blocks drop out consecutively. We present FlowImp, a conditional flow-matching framework that reconstructs contiguous missing blocks in three ODE solver steps. Specifically, FlowImp anchors straight transport paths on observed spectral priors and enforces cross-channel phase consistency via a lightweight Fourier regularizer. Across six multivariate benchmarks under 50\% block missingness, FlowImp reduces imputation MAE by 19.3\% compared with CSDI while accelerating inference by 24x.
\end{ExampleBox}
\begin{TipBox}{Rhetorical Move Breakdown \& Constituent Analysis}
\begin{itemize}[leftmargin=1.5em]
  \item \textbf{S1 [背景与任务重要性 Context]:} \textit{Recovering missing sensor observations is essential for reliable downstream monitoring in industrial and clinical time series.}\par\noindent{\small 对应步骤：S1 [背景与任务重要性 Context]：1句话锁定任务目标与核心应用场景（≤20词）。 | 句首锚定旧信息/逻辑视角，主干动词紧跟主语，复杂修饰置于宾语之后（符合尾重原则 L011 与旧新信息流 L001）。}
  \item \textbf{S2 [现有局限与痛点 Gap/Bottleneck]:} \textit{While diffusion-based imputation models generate plausible trajectories under random missingness, they suffer from high sampling latency and lose global periodic coherence when entire sensor blocks drop out consecutively.}\par\noindent{\small 对应步骤：S2 [现有局限与痛点 Gap/Bottleneck]：While existing [Paradigms] achieve [Merit], they struggle with [Core Bottleneck] because [Underlying Cause]. | 句首锚定旧信息/逻辑视角，主干动词紧跟主语，复杂修饰置于宾语之后（符合尾重原则 L011 与旧新信息流 L001）。}
  \item \textbf{S3-S4 [本文方法与核心机制 Approach \& Mechanism]:} \textit{We present FlowImp, a conditional flow-matching framework that reconstructs contiguous missing blocks in three ODE solver steps.}\par\noindent{\small 对应步骤：S3-S4 [本文方法与核心机制 Approach \& Mechanism]：In this work, we present [Model], which [Core Idea]. Specifically, [Module 1]..., while [Module 2]... | 句首锚定旧信息/逻辑视角，主干动词紧跟主语，复杂修饰置于宾语之后（符合尾重原则 L011 与旧新信息流 L001）。}
  \item \textbf{S5 [量化实证结果 Quantitative Results]:} \textit{Specifically, FlowImp anchors straight transport paths on observed spectral priors and enforces cross-channel phase consistency via a lightweight Fourier regularizer.}\par\noindent{\small 对应步骤：S5 [量化实证结果 Quantitative Results]：Experiments on [Benchmarks/Hardware] show that [Model] outperforms [Baselines] by [X\%] in [Metric] and achieves [Efficiency/Real-world Result]. | 句首锚定旧信息/逻辑视角，主干动词紧跟主语，复杂修饰置于宾语之后（符合尾重原则 L011 与旧新信息流 L001）。}
  \item \textbf{S5 [量化实证结果 Quantitative Results]:} \textit{Across six multivariate benchmarks under 50\% block missingness, FlowImp reduces imputation MAE by 19.3\% compared with CSDI while accelerating inference by 24x.}\par\noindent{\small 对应步骤：S5 [量化实证结果 Quantitative Results]：Experiments on [Benchmarks/Hardware] show that [Model] outperforms [Baselines] by [X\%] in [Metric] and achieves [Efficiency/Real-world Result]. | 句首锚定旧信息/逻辑视角，主干动词紧跟主语，复杂修饰置于宾语之后（符合尾重原则 L011 与旧新信息流 L001）。}
\end{itemize}
\end{TipBox}
\medskip
\subsection*{Paragraph 2: 确立研究背景与真实工程价值（拒绝空喊口号） (\texttt{BP01})}
\noindent\textbf{Expressive Intent:}\enspace 我这一段想作为 Introduction 第一段：说明多变量时间序列预测或具身智能控制在实际工业系统中的关键作用，以及高可靠性的刚需。\par
\noindent\textbf{Cohesion Formula:}\enspace {\small\color{AcadSlate}Given-to-New 链：Core Task (S1) $\to$ Practical Deployment (S2) $\to$ Real-world Disturbances (S3) $\to$ Consequence of Failure (S4). 绑定逻辑规则：L001, L011, L029.}\par\vspace{4pt}
\begin{ExampleBox}{Publication-Grade Authentic Text}
Multivariate time series forecasting underpins safety-critical scheduling across smart power grids, cloud clusters, and automated manufacturing. In these operational settings, control systems rely on multi-horizon predictions to balance supply fluctuations and prevent cascading overload. However, field sensors routinely operate under abrupt regime shifts, irregular sampling intervals, and high-frequency measurement noise. Failing to accommodate such non-stationary disturbances degrades long-range forecast reliability and triggers costly control oscillations.
\end{ExampleBox}
\begin{TipBox}{Rhetorical Move Breakdown \& Constituent Analysis}
\begin{itemize}[leftmargin=1.5em]
  \item \textbf{S1 [任务锚定 Topic Anchor]:} \textit{Multivariate time series forecasting underpins safety-critical scheduling across smart power grids, cloud clusters, and automated manufacturing.}\par\noindent{\small 对应步骤：S1 [任务锚定 Topic Anchor]：直接用紧凑主谓宾点明研究对象及其在真实系统中的核心功能（≤16词，禁用 In the era of AI）。 | 句首锚定旧信息/逻辑视角，主干动词紧跟主语，复杂修饰置于宾语之后（符合尾重原则 L011 与旧新信息流 L001）。}
  \item \textbf{S2 [应用场景展开 Concrete Scenarios]:} \textit{In these operational settings, control systems rely on multi-horizon predictions to balance supply fluctuations and prevent cascading overload.}\par\noindent{\small 对应步骤：S2 [应用场景展开 Concrete Scenarios]：列举 2--3 个具体工程场景（如电网负荷调度、灵巧机械手装配）及其决策依赖。 | 句首锚定旧信息/逻辑视角，主干动词紧跟主语，复杂修饰置于宾语之后（符合尾重原则 L011 与旧新信息流 L001）。}
  \item \textbf{S3 [苛刻工况约束 Operational Bottleneck]:} \textit{However, field sensors routinely operate under abrupt regime shifts, irregular sampling intervals, and high-frequency measurement noise.}\par\noindent{\small 对应步骤：S3 [苛刻工况约束 Operational Bottleneck]：指出真实部署环境下不可避免的物理/统计挑战（如非平稳漂移、传感器缺失、毫秒级延迟约束）。 | 句首锚定旧信息/逻辑视角，主干动词紧跟主语，复杂修饰置于宾语之后（符合尾重原则 L011 与旧新信息流 L001）。}
  \item \textbf{S4 [研究必要性收束 Imperative Hook]:} \textit{Failing to accommodate such non-stationary disturbances degrades long-range forecast reliability and triggers costly control oscillations.}\par\noindent{\small 对应步骤：S4 [研究必要性收束 Imperative Hook]：点明若无法克服上述约束会造成什么具体系统后果，自然引出下一段对现有方法的审视。 | 句首锚定旧信息/逻辑视角，主干动词紧跟主语，复杂修饰置于宾语之后（符合尾重原则 L011 与旧新信息流 L001）。}
\end{itemize}
\end{TipBox}
\medskip
\subsection*{Paragraph 3: 总结现有主流范式及其隐含假设（欲抑先扬铺垫段） (\texttt{BP02})}
\noindent\textbf{Expressive Intent:}\enspace 我这一段想客观概括前人主流模型（如 Transformer / Diffusion / MLP）取得了什么进展，但点出它们共同依赖的一个脆弱前提假设。\par
\noindent\textbf{Cohesion Formula:}\enspace {\small\color{AcadSlate}先让步肯定（While/Recent paradigms...）$\to$ 后釜底抽薪揭示隐含前提（Yet these architectures implicitly assume...）。绑定逻辑规则：L008, L009.}\par\vspace{4pt}
\begin{ExampleBox}{Publication-Grade Authentic Text}
Recent deep forecasters have substantially reduced long-horizon error by replacing point-wise attention with sub-series patch tokenization and frequency-domain decomposition. These architectures effectively capture recurring seasonal patterns when training and test sequences share stationary statistical moments. Yet they implicitly assume that inter-channel correlations and spectral energy distributions remain time-invariant across the prediction window. When external shocks alter underlying system dynamics, this stationarity assumption collapses, causing frozen spectral filters to project outdated trends.
\end{ExampleBox}
\begin{TipBox}{Rhetorical Move Breakdown \& Constituent Analysis}
\begin{itemize}[leftmargin=1.5em]
  \item \textbf{S1 [流派归类 Taxonomic Opening]:} \textit{Recent deep forecasters have substantially reduced long-horizon error by replacing point-wise attention with sub-series patch tokenization and frequency-domain decomposition.}\par\noindent{\small 对应步骤：S1 [流派归类 Taxonomic Opening]：按核心机制归纳现有主流范式及其已解决的子问题。 | 句首锚定旧信息/逻辑视角，主干动词紧跟主语，复杂修饰置于宾语之后（符合尾重原则 L011 与旧新信息流 L001）。}
  \item \textbf{S2 [代表性机制展开 Mechanism Recap]:} \textit{These architectures effectively capture recurring seasonal patterns when training and test sequences share stationary statistical moments.}\par\noindent{\small 对应步骤：S2 [代表性机制展开 Mechanism Recap]：简述这些方法如何通过特定架构（如 Patch 分块、频域分解或扩散去噪）提升表征能力。 | 句首锚定旧信息/逻辑视角，主干动词紧跟主语，复杂修饰置于宾语之后（符合尾重原则 L011 与旧新信息流 L001）。}
  \item \textbf{S3 [转折揭示隐含假设 Pivot to Hidden Assumption]:} \textit{Yet they implicitly assume that inter-channel correlations and spectral energy distributions remain time-invariant across the prediction window.}\par\noindent{\small 对应步骤：S3 [转折揭示隐含假设 Pivot to Hidden Assumption]：用 Yet / However 锋利转折，指出它们暗中依赖的理想化假设（如平稳分布、静态通道相关性、仿真动力学无偏差）。 | 句首锚定旧信息/逻辑视角，主干动词紧跟主语，复杂修饰置于宾语之后（符合尾重原则 L011 与旧新信息流 L001）。}
  \item \textbf{S4 [假设破裂后果 Breakdown Under Stress]:} \textit{When external shocks alter underlying system dynamics, this stationarity assumption collapses, causing frozen spectral filters to project outdated trends.}\par\noindent{\small 对应步骤：S4 [假设破裂后果 Breakdown Under Stress]：说明当真实数据违背该假设时，现有架构为何会遭遇性能断崖。 | 句首锚定旧信息/逻辑视角，主干动词紧跟主语，复杂修饰置于宾语之后（符合尾重原则 L011 与旧新信息流 L001）。}
\end{itemize}
\end{TipBox}
\medskip
\subsection*{Paragraph 4: 尖锐指出核心技术瓶颈与失效机理（Gap Identification） (\texttt{BP03})}
\noindent\textbf{Expressive Intent:}\enspace 我这一段想深入剖析现有方法在特定难题（如非平稳突变、长程自回归误差累积、Sim-to-Real 接触失配）上为什么会失效，把机理讲透。\par
\noindent\textbf{Cohesion Formula:}\enspace {\small\color{AcadSlate}因果链式推进（A induces B $\to$ This B compounds into C $\to$ Naive fix D sacrifices E）。绑定逻辑规则：L003, L004, L024/L025.}\par\vspace{4pt}
\begin{ExampleBox}{Publication-Grade Authentic Text}
Bridging non-stationary robustness and amplitude fidelity remains a core bottleneck in long-horizon forecasting. While instance normalization strips time-varying mean and variance shifts to stabilize attention weights, it simultaneously erases physical scale cues required to predict extreme load peaks. Conversely, retaining raw unnormalized inputs exposes self-attention layers to severe covariate shift, causing validation error to diverge over extended horizons. Resolving this tension requires a mechanism that disentangles intrinsic temporal dynamics from transient scale drift without discarding magnitude information.
\end{ExampleBox}
\begin{TipBox}{Rhetorical Move Breakdown \& Constituent Analysis}
\begin{itemize}[leftmargin=1.5em]
  \item \textbf{S1 [核心痛点立论 Bottleneck Claim]:} \textit{Bridging non-stationary robustness and amplitude fidelity remains a core bottleneck in long-horizon forecasting.}\par\noindent{\small 对应步骤：S1 [核心痛点立论 Bottleneck Claim]：一句话直接命名该核心瓶颈（如 Compound error accumulation / Sim-to-real contact mismatch）。 | 句首锚定旧信息/逻辑视角，主干动词紧跟主语，复杂修饰置于宾语之后（符合尾重原则 L011 与旧新信息流 L001）。}
  \item \textbf{S2 [第一层失效机理 Primary Failure Mechanism]:} \textit{While instance normalization strips time-varying mean and variance shifts to stabilize attention weights, it simultaneously erases physical scale cues required to predict extreme load peaks.}\par\noindent{\small 对应步骤：S2 [第一层失效机理 Primary Failure Mechanism]：从数学或物理传导链条解释误差是如何产生并放大的。 | 句首锚定旧信息/逻辑视角，主干动词紧跟主语，复杂修饰置于宾语之后（符合尾重原则 L011 与旧新信息流 L001）。}
  \item \textbf{S3 [第二层耦合困境 Secondary Dilemma]:} \textit{Conversely, retaining raw unnormalized inputs exposes self-attention layers to severe covariate shift, causing validation error to diverge over extended horizons.}\par\noindent{\small 对应步骤：S3 [第二层耦合困境 Secondary Dilemma]：指出为什么简单的修补手段（如单纯加大模型、简单归一化或域随机化）不仅无效还会带来副作用。 | 句首锚定旧信息/逻辑视角，主干动词紧跟主语，复杂修饰置于宾语之后（符合尾重原则 L011 与旧新信息流 L001）。}
  \item \textbf{S4 [凝练核心研究问题 Research Question Formulation]:} \textit{Resolving this tension requires a mechanism that disentangles intrinsic temporal dynamics from transient scale drift without discarding magnitude information.}\par\noindent{\small 对应步骤：S4 [凝练核心研究问题 Research Question Formulation]：用 Therefore, a fundamental question arises: how can we... 引出本文目标。 | 句首锚定旧信息/逻辑视角，主干动词紧跟主语，复杂修饰置于宾语之后（符合尾重原则 L011 与旧新信息流 L001）。}
\end{itemize}
\end{TipBox}
\medskip
\subsection*{Paragraph 5: 亮出本文核心洞察与破局方法总览（Key Insight \& Method Overview） (\texttt{BP04})}
\noindent\textbf{Expressive Intent:}\enspace 我这一段想在指出前人痛点后，正式亮出我们的方法（To address this dilemma, we propose...），并用最通俗清晰的直觉告诉审稿人为什么我们的设计能破局。\par
\noindent\textbf{Cohesion Formula:}\enspace {\small\color{AcadSlate}Problem Anchor $\to$ Key Insight $\to$ Dual-Module Mechanism $\to$ Pareto Resolution. 绑定逻辑规则：L002, L023.}\par\vspace{4pt}
\begin{ExampleBox}{Publication-Grade Authentic Text}
To resolve this dilemma, we propose DyTS-Net, a dual-stream forecasting architecture that decouples stationary structural modeling from dynamic scale trajectory estimation. Our key insight is that while high-frequency phase relationships benefit from normalized shift-invariant attention, macro-scale amplitude evolution follows a low-dimensional continuous dynamical system. Specifically, DyTS-Net first extracts scale-invariant temporal tokens via spectral patch attention, and concurrently models non-stationary mean-variance evolution through a lightweight neural ODE branch. Recombining both streams at the decoding stage preserves sharp peak magnitudes without sacrificing training stability.
\end{ExampleBox}
\begin{TipBox}{Rhetorical Move Breakdown \& Constituent Analysis}
\begin{itemize}[leftmargin=1.5em]
  \item \textbf{S1 [方法登场与定位 Method Introduction]:} \textit{To resolve this dilemma, we propose DyTS-Net, a dual-stream forecasting architecture that decouples stationary structural modeling from dynamic scale trajectory estimation.}\par\noindent{\small 对应步骤：S1 [方法登场与定位 Method Introduction]：To address this dilemma, we propose [Model Name], a [Core Paradigm] framework that [Primary Goal]. | 句首锚定旧信息/逻辑视角，主干动词紧跟主语，复杂修饰置于宾语之后（符合尾重原则 L011 与旧新信息流 L001）。}
  \item \textbf{S2 [核心直觉洞察 Key Insight]:} \textit{Our key insight is that while high-frequency phase relationships benefit from normalized shift-invariant attention, macro-scale amplitude evolution follows a low-dimensional continuous dynamical system.}\par\noindent{\small 对应步骤：S2 [核心直觉洞察 Key Insight]：Our key insight is that [Physical/Statistical Principle]，点破破局的底层支点。 | 句首锚定旧信息/逻辑视角，主干动词紧跟主语，复杂修饰置于宾语之后（符合尾重原则 L011 与旧新信息流 L001）。}
  \item \textbf{S3 [模块协同主干 Module Synergy]:} \textit{Specifically, DyTS-Net first extracts scale-invariant temporal tokens via spectral patch attention, and concurrently models non-stationary mean-variance evolution through a lightweight neural ODE branch.}\par\noindent{\small 对应步骤：S3 [模块协同主干 Module Synergy]：用 First / Specifically 与 Concurrently / Furthermore 简述两大核心模块如何分工咬合。 | 句首锚定旧信息/逻辑视角，主干动词紧跟主语，复杂修饰置于宾语之后（符合尾重原则 L011 与旧新信息流 L001）。}
  \item \textbf{S4 [闭环收益总结 Closed-Loop Benefit]:} \textit{Recombining both streams at the decoding stage preserves sharp peak magnitudes without sacrificing training stability.}\par\noindent{\small 对应步骤：S4 [闭环收益总结 Closed-Loop Benefit]：说明这种协同设计如何同时兼顾原先冲突的两个指标（如兼顾平稳表征与幅值保真，或兼顾多模态多样性与单步实时控制）。 | 句首锚定旧信息/逻辑视角，主干动词紧跟主语，复杂修饰置于宾语之后（符合尾重原则 L011 与旧新信息流 L001）。}
\end{itemize}
\end{TipBox}
\medskip
\subsection*{Paragraph 6: 凝练论文三点核心贡献（Contributions Summary） (\texttt{BP07})}
\noindent\textbf{Expressive Intent:}\enspace 我这一段想在 Introduction 最末尾清晰列出本文的三点核心贡献（理论/视角贡献 + 核心算法/架构贡献 + 实验与开源实证贡献）。\par
\noindent\textbf{Cohesion Formula:}\enspace {\small\color{AcadSlate}Parallel Grammatical Structure（严格平行结构：全部用 We + 强及物动词 identify / propose / demonstrate 开头）。绑定逻辑规则：L002, L010.}\par\vspace{4pt}
\begin{ExampleBox}{Publication-Grade Authentic Text}
In summary, our main contributions are threefold: (1) We formalize the amplitude-stationarity trade-off in long-horizon forecasting and prove that standard instance normalization bounds peak-load recoverability. (2) We propose DyTS-Net, a hybrid architecture that couples shift-invariant spectral attention with a continuous neural ODE scale tracker. (3) We demonstrate across eight public benchmarks and industrial wind-power telemetry that DyTS-Net lowers long-range MSE by up to 17.4\% while cutting inference latency nearly in half.
\end{ExampleBox}
\begin{TipBox}{Rhetorical Move Breakdown \& Constituent Analysis}
\begin{itemize}[leftmargin=1.5em]
  \item \textbf{S1 [总起引导句 Lead-in]:} \textit{In summary, our main contributions are threefold: (1) We formalize the amplitude-stationarity trade-off in long-horizon forecasting and prove that standard instance normalization bounds peak-load recoverability.}\par\noindent{\small 对应步骤：S1 [总起引导句 Lead-in]：In summary, our main contributions are threefold: | 句首锚定旧信息/逻辑视角，主干动词紧跟主语，复杂修饰置于宾语之后（符合尾重原则 L011 与旧新信息流 L001）。}
  \item \textbf{S2 [贡献 1:} \textit{(2) We propose DyTS-Net, a hybrid architecture that couples shift-invariant spectral attention with a continuous neural ODE scale tracker.}\par\noindent{\small 对应步骤：S2 [贡献 1：新视角/问题形式化 Conceptual/Formulation Contribution]：用强动词名物化或动词原形开头（We identify and formalize... / We establish...），点明打破了什么传统局限。 | 句首锚定旧信息/逻辑视角，主干动词紧跟主语，复杂修饰置于宾语之后（符合尾重原则 L011 与旧新信息流 L001）。}
  \item \textbf{S3 [贡献 2:} \textit{(3) We demonstrate across eight public benchmarks and industrial wind-power telemetry that DyTS-Net lowers long-range MSE by up to 17.4\% while cutting inference latency nearly in half.}\par\noindent{\small 对应步骤：S3 [贡献 2：核心方法与模块 Algorithmic/Architectural Contribution]：We propose [Model/Mechanism], which integrates [A] and [B] to achieve [Theoretical/Functional Guarantee]. | 句首锚定旧信息/逻辑视角，主干动词紧跟主语，复杂修饰置于宾语之后（符合尾重原则 L011 与旧新信息流 L001）。}
\end{itemize}
\end{TipBox}
\medskip
\subsection*{Paragraph 7: 问题形式化定义与符号约定段（Problem Formulation \& Notation） (\texttt{BP11})}
\noindent\textbf{Expressive Intent:}\enspace 我这一段想在 Methodology 开篇（Section 3.1 Problem Formulation），严谨、简洁地定义输入张量、输出目标、约束条件与优化目标。\par
\noindent\textbf{Cohesion Formula:}\enspace {\small\color{AcadSlate}Input Space $\to$ Auxiliary Conditions $\to$ Target Mapping $\to$ Crucial Distributional Constraint. 绑定逻辑规则：L014, L011.}\par\vspace{4pt}
\begin{ExampleBox}{Publication-Grade Authentic Text}
Consider a multivariate time series $\mathbf{X}_{1:L} = [x_1, \dots, x_L]^\top \in \mathbb{R}^{L \times C}$ recording $C$ correlated sensors over lookback horizon $L$. Given an binary observation mask $\mathbf{M} \in \{0, 1\}^{L \times C}$ indicating irregular sensor dropouts, our objective is to learn a generative mapping $p_\theta(\mathbf{Y}_{L+1:L+H} \mid \mathbf{X}_{1:L} \odot \mathbf{M})$ over future horizon $H$. Crucially, unlike standard settings that assume time-invariant marginal distributions, we explicitly account for non-stationary regime shifts where $\mathbb{E}[x_t]$ and $\text{Cov}(x_t)$ evolve continuously between the lookback and forecasting windows.
\end{ExampleBox}
\begin{TipBox}{Rhetorical Move Breakdown \& Constituent Analysis}
\begin{itemize}[leftmargin=1.5em]
  \item \textbf{S1 [任务总体形式化 Task Definition]:} \textit{Consider a multivariate time series $\mathbf{X}_{1:L} = [x_1, \dots, x_L]^\top \in \mathbb{R}^{L \times C}$ recording $C$ correlated sensors over lookback horizon $L$.}\par\noindent{\small 对应步骤：S1 [任务总体形式化 Task Definition]：Consider a [System/Dataset] observed over a lookback window of length $L$, denoted as $\mathbf{X} = [x_1, \dots, x_L] \in \mathbb{R}^{L \times C}$, where $C$ represents [Meaning]. | 句首锚定旧信息/逻辑视角，主干动词紧跟主语，复杂修饰置于宾语之后（符合尾重原则 L011 与旧新信息流 L001）。}
  \item \textbf{S2 [补充条件或模态输入 Auxiliary/Condition Inputs]:} \textit{Given an binary observation mask $\mathbf{M} \in \{0, 1\}^{L \times C}$ indicating irregular sensor dropouts, our objective is to learn a generative mapping $p_\theta(\mathbf{Y}_{L+1:L+H} \mid \mathbf{X}_{1:L} \odot \mathbf{M})$ over future horizon $H$.}\par\noindent{\small 对应步骤：S2 [补充条件或模态输入 Auxiliary/Condition Inputs]：At each step $t$, the system also receives [Condition/Observation] $z_t \in \mathcal{Z}$ under [Noise/Missingness/Dynamic Regime]. | 句首锚定旧信息/逻辑视角，主干动词紧跟主语，复杂修饰置于宾语之后（符合尾重原则 L011 与旧新信息流 L001）。}
  \item \textbf{S3 [预测/控制目标 Objective]:} \textit{Crucially, unlike standard settings that assume time-invariant marginal distributions, we explicitly account for non-stationary regime shifts where $\mathbb{E}[x_t]$ and $\text{Cov}(x_t)$ evolve continuously between the lookback and forecasting windows.}\par\noindent{\small 对应步骤：S3 [预测/控制目标 Objective]：Our goal is to learn a parameterized mapping $f_\theta(\cdot)$ that predicts future states $\hat{\mathbf{Y}} \in \mathbb{R}^{H \times C}$ over horizon $H$ (or action trajectory $a_{t:t+K}$). | 句首锚定旧信息/逻辑视角，主干动词紧跟主语，复杂修饰置于宾语之后（符合尾重原则 L011 与旧新信息流 L001）。}
\end{itemize}
\end{TipBox}
\medskip
\subsection*{Paragraph 8: 整体模型架构数据流俯瞰段（Framework Overview \& Pipeline Walkthrough） (\texttt{BP12})}
\noindent\textbf{Expressive Intent:}\enspace 我这一段想配合论文里的 Figure 2（模型总架构图），用一段话把输入张量从左到右流经各个模块直到输出预测结果的完整全貌讲清楚。\par
\noindent\textbf{Cohesion Formula:}\enspace {\small\color{AcadSlate}严格的链式数据流（$\mathbf{X} \to \mathbf{Z}^{(1)} \to \mathbf{Z}^{(2)} \to \hat{\mathbf{Y}}$），每一句句首承接上一句的输出张量。绑定逻辑规则：L001, L003, L004.}\par\vspace{4pt}
\begin{ExampleBox}{Publication-Grade Authentic Text}
Figure 2 illustrates the overall architecture of DyTS-Net, which maps raw sensor streams to multi-horizon forecasts through three sequential stages. Given an input sequence $\mathbf{X}$, the Adaptive De-stationary Tokenizer first separates local trend-scale statistics from normalized sub-series patches $\mathbf{Z}_0$. These normalized patches then pass through stacked Spectral-Spatial Mixing blocks, which jointly filter high-frequency channel noise in the Fourier domain and propagate cross-sensor phase dependencies via a sparse directed graph. Finally, the Scale-Guided Decoder fuses the refined structural representations with the projected neural-ODE scale trajectory to reconstruct the multi-horizon output $\hat{\mathbf{Y}}$.
\end{ExampleBox}
\begin{TipBox}{Rhetorical Move Breakdown \& Constituent Analysis}
\begin{itemize}[leftmargin=1.5em]
  \item \textbf{S1 [总图锚定与设计哲学 Visual Anchor \& Philosophy]:} \textit{Figure 2 illustrates the overall architecture of DyTS-Net, which maps raw sensor streams to multi-horizon forecasts through three sequential stages.}\par\noindent{\small 对应步骤：S1 [总图锚定与设计哲学 Visual Anchor \& Philosophy]：Figure 2 illustrates the overall architecture of [Model], which processes [Input] through three sequential stages: [Stage 1], [Stage 2], and [Stage 3]. | 句首锚定旧信息/逻辑视角，主干动词紧跟主语，复杂修饰置于宾语之后（符合尾重原则 L011 与旧新信息流 L001）。}
  \item \textbf{S2 [阶段 1 数据流 Stage 1 Flow]:} \textit{Given an input sequence $\mathbf{X}$, the Adaptive De-stationary Tokenizer first separates local trend-scale statistics from normalized sub-series patches $\mathbf{Z}_0$.}\par\noindent{\small 对应步骤：S2 [阶段 1 数据流 Stage 1 Flow]：Given raw input $\mathbf{X}$, the [Module 1] first transforms [Raw Input] into [Intermediate Representation $\mathbf{Z}^{(1)}$] to [Functional Purpose 1]. | 句首锚定旧信息/逻辑视角，主干动词紧跟主语，复杂修饰置于宾语之后（符合尾重原则 L011 与旧新信息流 L001）。}
  \item \textbf{S3 [阶段 2 数据流 Stage 2 Flow]:} \textit{These normalized patches then pass through stacked Spectral-Spatial Mixing blocks, which jointly filter high-frequency channel noise in the Fourier domain and propagate cross-sensor phase dependencies via a sparse directed graph.}\par\noindent{\small 对应步骤：S3 [阶段 2 数据流 Stage 2 Flow]：Next, these [Summary Noun: representations $\mathbf{Z}^{(1)}$] are fed into [Module 2], which [Mechanism 2] to produce [Refined State $\mathbf{Z}^{(2)}$]. | 句首锚定旧信息/逻辑视角，主干动词紧跟主语，复杂修饰置于宾语之后（符合尾重原则 L011 与旧新信息流 L001）。}
  \item \textbf{S4 [阶段 3 解码与闭环输出 Stage 3 Decoding]:} \textit{Finally, the Scale-Guided Decoder fuses the refined structural representations with the projected neural-ODE scale trajectory to reconstruct the multi-horizon output $\hat{\mathbf{Y}}$.}\par\noindent{\small 对应步骤：S4 [阶段 3 解码与闭环输出 Stage 3 Decoding]：Finally, the [Module 3] integrates $\mathbf{Z}^{(2)}$ with [Prior/Scale Cue] to output the final [Prediction/Action Trajectory]. | 句首锚定旧信息/逻辑视角，主干动词紧跟主语，复杂修饰置于宾语之后（符合尾重原则 L011 与旧新信息流 L001）。}
\end{itemize}
\end{TipBox}
\medskip
\subsection*{Paragraph 9: 核心子模块：直觉动机与数学设计无缝衔接段（Intuition-before-Math） (\texttt{BP13})}
\noindent\textbf{Expressive Intent:}\enspace 我这一段想详细介绍我们提出的某个核心创新模块：先用两句话讲清楚为什么直接用传统算子会出问题，再顺理成章推导出我们的新公式。\par
\noindent\textbf{Cohesion Formula:}\enspace {\small\color{AcadSlate}Pathology of Standard Operator $\to$ Physical Intuition $\to$ Equation Integration $\to$ Theoretical Guarantee. 绑定逻辑规则：L014, L003.}\par\vspace{4pt}
\begin{ExampleBox}{Publication-Grade Authentic Text}
Standard dot-product attention measures token similarity purely by instantaneous amplitude alignment, often matching noisy spikes across unrelated time steps. To enforce temporal coherence, we observe that genuine physical interactions between industrial sensors exhibit consistent phase lag across dominant frequencies. Formally, given query and key projections $Q, K \in \mathbb{R}^{N \times d}$, we modulate the attention logits via a phase-alignment mask: $\mathcal{A}(Q, K, V) = \text{softmax}\big(\frac{QK^\top}{\sqrt{d}} - \lambda \|\Delta \Phi_{Q,K}\|_1\big)V$, where $\Delta \Phi_{Q,K} \in [-\pi, \pi]^{N \times N}$ denotes the instantaneous Hilbert phase difference matrix and $\lambda > 0$ controls phase selectivity. Penalizing large phase discrepancies prevents spurious attention weights between out-of-phase noise bursts.
\end{ExampleBox}
\begin{TipBox}{Rhetorical Move Breakdown \& Constituent Analysis}
\begin{itemize}[leftmargin=1.5em]
  \item \textbf{S1 [传统算子缺陷 Motivation]:} \textit{Standard dot-product attention measures token similarity purely by instantaneous amplitude alignment, often matching noisy spikes across unrelated time steps.}\par\noindent{\small 对应步骤：S1 [传统算子缺陷 Motivation]：Standard [Conventional Operator, e.g., softmax self-attention] computes [Operation], which suffers from [Pathology] when applied to [Domain Condition]. | 句首锚定旧信息/逻辑视角，主干动词紧跟主语，复杂修饰置于宾语之后（符合尾重原则 L011 与旧新信息流 L001）。}
  \item \textbf{S2 [物理/统计直觉 Physical Intuition]:} \textit{To enforce temporal coherence, we observe that genuine physical interactions between industrial sensors exhibit consistent phase lag across dominant frequencies.}\par\noindent{\small 对应步骤：S2 [物理/统计直觉 Physical Intuition]：To overcome this pathology, we observe that [Key Domain Property] can be effectively regularized by [Intuitive Operation]. | 句首锚定旧信息/逻辑视角，主干动词紧跟主语，复杂修饰置于宾语之后（符合尾重原则 L011 与旧新信息流 L001）。}
  \item \textbf{S3 [数学公式定义 Mathematical Formulation]:} \textit{Formally, given query and key projections $Q, K \in \mathbb{R}^{N \times d}$, we modulate the attention logits via a phase-alignment mask: $\mathcal{A}(Q, K, V) = \text{softmax}\big(\frac{QK^\top}{\sqrt{d}} - \lambda \|\Delta \Phi_{Q,K}\|_1\big)V$, where $\Delta \Phi_{Q,K} \in [-\pi, \pi]^{N \times N}$ denotes the instantaneous Hilbert phase difference matrix and $\lambda > 0$ controls phase selectivity.}\par\noindent{\small 对应步骤：S3 [数学公式定义 Mathematical Formulation]：Formally, given input representation $H \in \mathbb{R}^{N \times d}$, we formulate the [Module Name] as: [Equation], where [Symbol Explanations]. | 句首锚定旧信息/逻辑视角，主干动词紧跟主语，复杂修饰置于宾语之后（符合尾重原则 L011 与旧新信息流 L001）。}
  \item \textbf{S4 [理论/计算性质 Guarantee]:} \textit{Penalizing large phase discrepancies prevents spurious attention weights between out-of-phase noise bursts.}\par\noindent{\small 对应步骤：S4 [理论/计算性质 Guarantee]：By construction, this formulation guarantees [Desirable Property, e.g., $O(N)$ complexity / translation invariance]. | 句首锚定旧信息/逻辑视角，主干动词紧跟主语，复杂修饰置于宾语之后（符合尾重原则 L011 与旧新信息流 L001）。}
\end{itemize}
\end{TipBox}
\medskip
\subsection*{Paragraph 10: 主实验结果整体对比与核心胜出原因剖析段（Main Results Analysis） (\texttt{BP18})}
\noindent\textbf{Expressive Intent:}\enspace 我这一段想解读主实验大表（Table 1）：既报告全局胜率，又挑出最有代表性的对比案例，解释为什么竞争对手在那里跌倒而我们稳赢。\par
\noindent\textbf{Cohesion Formula:}\enspace {\small\color{AcadSlate}Global Win Rate $\to$ Advantage over Family A $\to$ Resilience at Long Horizon vs. Family B $\to$ Mechanistic Takeaway. 绑定逻辑规则：L015, L013, L010.}\par\vspace{4pt}
\begin{ExampleBox}{Publication-Grade Authentic Text}
Table 1 reports the multivariate forecasting results, where DyTS-Net secures the lowest error in 25 out of 28 dataset-horizon configurations. Compared with channel-independent PatchTST, DyTS-Net reduces average MSE by 14.2\% on the 862-channel Traffic benchmark, where modeling spatial-temporal congestion propagation across adjacent road sensors is critical. Furthermore, while variate-mixing iTransformer performs competitively at $H=96$, its MSE escalates by 68\% at $H=720$ on volatile Weather series as unregularized cross-channel attention overfits transient wind gusts. By decoupling stationary phase coupling from neural-ODE scale projection, DyTS-Net suppresses long-horizon drift and maintains an 11.6\% margin over the strongest baseline at $H=720$.
\end{ExampleBox}
\begin{TipBox}{Rhetorical Move Breakdown \& Constituent Analysis}
\begin{itemize}[leftmargin=1.5em]
  \item \textbf{S1 [全局统治力概述 Overall Dominance]:} \textit{Table 1 reports the multivariate forecasting results, where DyTS-Net secures the lowest error in 25 out of 28 dataset-horizon configurations.}\par\noindent{\small 对应步骤：S1 [全局统治力概述 Overall Dominance]：Table 1 summarizes the multivariate forecasting results, where [Our Model] achieves the lowest error in [X] out of [Y] dataset-horizon configurations. | 句首锚定旧信息/逻辑视角，主干动词紧跟主语，复杂修饰置于宾语之后（符合尾重原则 L011 与旧新信息流 L001）。}
  \item \textbf{S2 [对比流派 1 的差距与原因 Contrast vs. Family 1]:} \textit{Compared with channel-independent PatchTST, DyTS-Net reduces average MSE by 14.2\% on the 862-channel Traffic benchmark, where modeling spatial-temporal congestion propagation across adjacent road sensors is critical.}\par\noindent{\small 对应步骤：S2 [对比流派 1 的差距与原因 Contrast vs. Family 1]：Compared with channel-independent baselines such as [Baseline A], our method reduces average MSE by [X\%] on [Dataset], where strong cross-sensor coupling is critical. | 句首锚定旧信息/逻辑视角，主干动词紧跟主语，复杂修饰置于宾语之后（符合尾重原则 L011 与旧新信息流 L001）。}
  \item \textbf{S3 [对比流派 2 在困难场景下的崩溃 Contrast vs. Family 2 under Hard Regime]:} \textit{Furthermore, while variate-mixing iTransformer performs competitively at $H=96$, its MSE escalates by 68\% at $H=720$ on volatile Weather series as unregularized cross-channel attention overfits transient wind gusts.}\par\noindent{\small 对应步骤：S3 [对比流派 2 在困难场景下的崩溃 Contrast vs. Family 2 under Hard Regime]：Furthermore, while [Baseline B] performs competitively at short horizons ($H=96$), its error escalates by [Y\%] at $H=720$ due to [Failure Mechanism]. | 句首锚定旧信息/逻辑视角，主干动词紧跟主语，复杂修饰置于宾语之后（符合尾重原则 L011 与旧新信息流 L001）。}
  \item \textbf{S4 [回扣本文核心机制 Mechanistic Confirmation]:} \textit{By decoupling stationary phase coupling from neural-ODE scale projection, DyTS-Net suppresses long-horizon drift and maintains an 11.6\% margin over the strongest baseline at $H=720$.}\par\noindent{\small 对应步骤：S4 [回扣本文核心机制 Mechanistic Confirmation]：By contrast, our [Core Module] bounds long-range drift, maintaining a [Z\%] margin over the second-best baseline at $H=720$. | 句首锚定旧信息/逻辑视角，主干动词紧跟主语，复杂修饰置于宾语之后（符合尾重原则 L011 与旧新信息流 L001）。}
\end{itemize}
\end{TipBox}
\medskip
\subsection*{Paragraph 11: 消融实验逐组件归因论证段（Counterfactual Ablation Study） (\texttt{BP19})}
\noindent\textbf{Expressive Intent:}\enspace 我这一段想汇报消融实验（Table 2）：证明我们的模块 A、模块 B 和正则项 C 每一个都不可或缺，且各自负责解决不同的特定子问题。\par
\noindent\textbf{Cohesion Formula:}\enspace {\small\color{AcadSlate}Ablation Overview $\to$ w/o Module A (Specific Failure) $\to$ w/o Module B (Specific Failure) $\to$ Cross-Module Synergy. 绑定逻辑规则：L016, L007.}\par\vspace{4pt}
\begin{ExampleBox}{Publication-Grade Authentic Text}
To isolate the contribution of each design choice, Table 2 compares full DyTS-Net against three ablated variants across ETTm2 and ECL. Removing the neural-ODE scale tracker (w/o Scale-ODE) triggers the sharpest degradation during seasonal transition months on ETTm2, raising 720-step MSE by 19.4\% as the decoder fails to track shifting baseline temperatures. Conversely, replacing phase-aligned graph attention with standard dot-product attention (w/o Phase-Attn) degrades performance most severely on the 321-channel ECL grid (+13.8\% MSE), where unaligned high-frequency load spikes corrupt cross-substation aggregation. Jointly deploying both modules yields complementary gains, as accurate scale tracking prevents amplitude distortion from biasing phase estimation.
\end{ExampleBox}
\begin{TipBox}{Rhetorical Move Breakdown \& Constituent Analysis}
\begin{itemize}[leftmargin=1.5em]
  \item \textbf{S1 [消融实验总括 Ablation Setup]:} \textit{To isolate the contribution of each design choice, Table 2 compares full DyTS-Net against three ablated variants across ETTm2 and ECL.}\par\noindent{\small 对应步骤：S1 [消融实验总括 Ablation Setup]：To isolate the contribution of each architectural component, we evaluate four ablated variants in Table 2 by progressively removing [Mod A], [Mod B], and [Loss C]. | 句首锚定旧信息/逻辑视角，主干动词紧跟主语，复杂修饰置于宾语之后（符合尾重原则 L011 与旧新信息流 L001）。}
  \item \textbf{S2 [移除模块 A 的冲击与归因 Impact of Removing Mod A]:} \textit{Removing the neural-ODE scale tracker (w/o Scale-ODE) triggers the sharpest degradation during seasonal transition months on ETTm2, raising 720-step MSE by 19.4\% as the decoder fails to track shifting baseline temperatures.}\par\noindent{\small 对应步骤：S2 [移除模块 A 的冲击与归因 Impact of Removing Mod A]：Removing [Mod A] (w/o A) causes the largest degradation on [Non-stationary/Volatile Subset], increasing MSE by [X\%] because the model loses [Specific Capability A]. | 句首锚定旧信息/逻辑视角，主干动词紧跟主语，复杂修饰置于宾语之后（符合尾重原则 L011 与旧新信息流 L001）。}
  \item \textbf{S3 [移除模块 B 的冲击与归因 Impact of Removing Mod B]:} \textit{Conversely, replacing phase-aligned graph attention with standard dot-product attention (w/o Phase-Attn) degrades performance most severely on the 321-channel ECL grid (+13.8\% MSE), where unaligned high-frequency load spikes corrupt cross-substation aggregation.}\par\noindent{\small 对应步骤：S3 [移除模块 B 的冲击与归因 Impact of Removing Mod B]：Similarly, replacing [Mod B] with standard [Naive Operator] degrades performance by [Y\%] on high-dimensional benchmarks, confirming its role in [Specific Capability B]. | 句首锚定旧信息/逻辑视角，主干动词紧跟主语，复杂修饰置于宾语之后（符合尾重原则 L011 与旧新信息流 L001）。}
  \item \textbf{S4 [协同效应总结 Synergistic Conclusion]:} \textit{Jointly deploying both modules yields complementary gains, as accurate scale tracking prevents amplitude distortion from biasing phase estimation.}\par\noindent{\small 对应步骤：S4 [协同效应总结 Synergistic Conclusion]：Combining all components yields super-additive gains, as [Mod A] stabilizes input statistics for [Mod B]'s relational reasoning. | 句首锚定旧信息/逻辑视角，主干动词紧跟主语，复杂修饰置于宾语之后（符合尾重原则 L011 与旧新信息流 L001）。}
\end{itemize}
\end{TipBox}
\medskip
\subsection*{Paragraph 12: 全文总结与学术影响收尾段（Conclusion \& Broader Outlook） (\texttt{BP24})}
\noindent\textbf{Expressive Intent:}\enspace 我这一段想写整篇小论文的 Conclusion：不简单复制粘贴 Abstract，而是从更高维度总结本文解决的科学矛盾、方法精髓、实证高度与领域启示。\par
\noindent\textbf{Cohesion Formula:}\enspace {\small\color{AcadSlate}Tension Resolved $\to$ Methodological Insight $\to$ Empirical Confirmation $\to$ Generalizable Blueprint. 绑定逻辑规则：L013, L023, L029.}\par\vspace{4pt}
\begin{ExampleBox}{Publication-Grade Authentic Text}
In this work, we resolved the long-standing tension between shift-invariant normalization and physical amplitude preservation in non-stationary time series forecasting. By decoupling frequency-domain phase attention from continuous neural-ODE scale tracking, DyTS-Net captures long-range seasonal dependencies while accurately reconstructing extreme peak magnitudes. Empirical evaluations across eight benchmarks and industrial wind-power telemetry demonstrated MSE reductions of up to 17.4\% alongside a 45\% reduction in computational overhead. Beyond forecasting, this structural-scale disentanglement principle provides a generalizable foundation for anomaly detection and closed-loop predictive control in non-stationary cyber-physical systems.
\end{ExampleBox}
\begin{TipBox}{Rhetorical Move Breakdown \& Constituent Analysis}
\begin{itemize}[leftmargin=1.5em]
  \item \textbf{S1 [核心科学问题闭合 Scientific Closure]:} \textit{In this work, we resolved the long-standing tension between shift-invariant normalization and physical amplitude preservation in non-stationary time series forecasting.}\par\noindent{\small 对应步骤：S1 [核心科学问题闭合 Scientific Closure]：In this paper, we addressed the fundamental tension between [Goal A] and [Goal B] in [Domain]. | 句首锚定旧信息/逻辑视角，主干动词紧跟主语，复杂修饰置于宾语之后（符合尾重原则 L011 与旧新信息流 L001）。}
  \item \textbf{S2 [方法论精髓回顾 Methodological Synthesis]:} \textit{By decoupling frequency-domain phase attention from continuous neural-ODE scale tracking, DyTS-Net captures long-range seasonal dependencies while accurately reconstructing extreme peak magnitudes.}\par\noindent{\small 对应步骤：S2 [方法论精髓回顾 Methodological Synthesis]：By introducing [Model/Mechanism], we demonstrated that [Core Insight] enables [Capability] without incurring [Traditional Penalty]. | 句首锚定旧信息/逻辑视角，主干动词紧跟主语，复杂修饰置于宾语之后（符合尾重原则 L011 与旧新信息流 L001）。}
  \item \textbf{S3 [实证结果高度概括 Empirical Recap]:} \textit{Empirical evaluations across eight benchmarks and industrial wind-power telemetry demonstrated MSE reductions of up to 17.4\% alongside a 45\% reduction in computational overhead.}\par\noindent{\small 对应步骤：S3 [实证结果高度概括 Empirical Recap]：Comprehensive evaluations across [Benchmarks/Hardware] confirmed consistent error reductions of up to [X\%] alongside substantial inference speedups. | 句首锚定旧信息/逻辑视角，主干动词紧跟主语，复杂修饰置于宾语之后（符合尾重原则 L011 与旧新信息流 L001）。}
  \item \textbf{S4 [跨领域迁移展望 Broader Outlook]:} \textit{Beyond forecasting, this structural-scale disentanglement principle provides a generalizable foundation for anomaly detection and closed-loop predictive control in non-stationary cyber-physical systems.}\par\noindent{\small 对应步骤：S4 [跨领域迁移展望 Broader Outlook]：Beyond [Current Task], our [Core Principle] offers a generalizable blueprint for broader sequential decision-making tasks under non-stationary environments. | 句首锚定旧信息/逻辑视角，主干动词紧跟主语，复杂修饰置于宾语之后（符合尾重原则 L011 与旧新信息流 L001）。}
\end{itemize}
\end{TipBox}
\medskip
\clearpage
\section{FT02: 基于条件流匹配的块缺失时间序列生成与插补完整小论文（FlowImp: Manifold-Consistent Time Series Imputation）}\label{comp:FT02}
\begin{ConceptBox}{Composition Architecture \& Structural Blueprint Chain}
\textbf{Genre:} \texttt{AI\_Mini\_Paper} \quad \textbf{Paragraph Count:} 11\par
\textbf{Blueprint Chain:} \texttt{BP08 $\to$ BP01 $\to$ BP09 $\to$ BP10 $\to$ BP11 $\to$ BP15 $\to$ BP17 $\to$ BP18 $\to$ BP20 $\to$ BP23 $\to$ BP24}
\end{ConceptBox}
\subsection*{Paragraph 1: 高密度四句式 Abstract 全段构建（Complete 4-Move Abstract） (\texttt{BP08})}
\noindent\textbf{Expressive Intent:}\enspace 我这一段想写出整篇小论文的完整 Abstract（摘要）：在 150--180 词内一气呵成交代背景、核心痛点、方法机制与量化实验结果。\par
\noindent\textbf{Cohesion Formula:}\enspace {\small\color{AcadSlate}Context (S1) $\to$ Gap (S2) $\to$ Method \& Mechanism (S3-S4) $\to$ Quantitative Proof (S5). 绑定逻辑规则：L001, L008, L010, L029.}\par\vspace{4pt}
\begin{ExampleBox}{Publication-Grade Authentic Text}
Recovering missing sensor observations is essential for reliable downstream monitoring in industrial and clinical time series. While diffusion-based imputation models generate plausible trajectories under random missingness, they suffer from high sampling latency and lose global periodic coherence when entire sensor blocks drop out consecutively. We present FlowImp, a conditional flow-matching framework that reconstructs contiguous missing blocks in three ODE solver steps. Specifically, FlowImp anchors straight transport paths on observed spectral priors and enforces cross-channel phase consistency via a lightweight Fourier regularizer. Across six multivariate benchmarks under 50\% block missingness, FlowImp reduces imputation MAE by 19.3\% compared with CSDI while accelerating inference by 24x.
\end{ExampleBox}
\begin{TipBox}{Rhetorical Move Breakdown \& Constituent Analysis}
\begin{itemize}[leftmargin=1.5em]
  \item \textbf{S1 [背景与任务重要性 Context]:} \textit{Recovering missing sensor observations is essential for reliable downstream monitoring in industrial and clinical time series.}\par\noindent{\small 对应步骤：S1 [背景与任务重要性 Context]：1句话锁定任务目标与核心应用场景（≤20词）。 | 句首锚定旧信息/逻辑视角，主干动词紧跟主语，复杂修饰置于宾语之后（符合尾重原则 L011 与旧新信息流 L001）。}
  \item \textbf{S2 [现有局限与痛点 Gap/Bottleneck]:} \textit{While diffusion-based imputation models generate plausible trajectories under random missingness, they suffer from high sampling latency and lose global periodic coherence when entire sensor blocks drop out consecutively.}\par\noindent{\small 对应步骤：S2 [现有局限与痛点 Gap/Bottleneck]：While existing [Paradigms] achieve [Merit], they struggle with [Core Bottleneck] because [Underlying Cause]. | 句首锚定旧信息/逻辑视角，主干动词紧跟主语，复杂修饰置于宾语之后（符合尾重原则 L011 与旧新信息流 L001）。}
  \item \textbf{S3-S4 [本文方法与核心机制 Approach \& Mechanism]:} \textit{We present FlowImp, a conditional flow-matching framework that reconstructs contiguous missing blocks in three ODE solver steps.}\par\noindent{\small 对应步骤：S3-S4 [本文方法与核心机制 Approach \& Mechanism]：In this work, we present [Model], which [Core Idea]. Specifically, [Module 1]..., while [Module 2]... | 句首锚定旧信息/逻辑视角，主干动词紧跟主语，复杂修饰置于宾语之后（符合尾重原则 L011 与旧新信息流 L001）。}
  \item \textbf{S5 [量化实证结果 Quantitative Results]:} \textit{Specifically, FlowImp anchors straight transport paths on observed spectral priors and enforces cross-channel phase consistency via a lightweight Fourier regularizer.}\par\noindent{\small 对应步骤：S5 [量化实证结果 Quantitative Results]：Experiments on [Benchmarks/Hardware] show that [Model] outperforms [Baselines] by [X\%] in [Metric] and achieves [Efficiency/Real-world Result]. | 句首锚定旧信息/逻辑视角，主干动词紧跟主语，复杂修饰置于宾语之后（符合尾重原则 L011 与旧新信息流 L001）。}
  \item \textbf{S5 [量化实证结果 Quantitative Results]:} \textit{Across six multivariate benchmarks under 50\% block missingness, FlowImp reduces imputation MAE by 19.3\% compared with CSDI while accelerating inference by 24x.}\par\noindent{\small 对应步骤：S5 [量化实证结果 Quantitative Results]：Experiments on [Benchmarks/Hardware] show that [Model] outperforms [Baselines] by [X\%] in [Metric] and achieves [Efficiency/Real-world Result]. | 句首锚定旧信息/逻辑视角，主干动词紧跟主语，复杂修饰置于宾语之后（符合尾重原则 L011 与旧新信息流 L001）。}
\end{itemize}
\end{TipBox}
\medskip
\subsection*{Paragraph 2: 确立研究背景与真实工程价值（拒绝空喊口号） (\texttt{BP01})}
\noindent\textbf{Expressive Intent:}\enspace 我这一段想作为 Introduction 第一段：说明多变量时间序列预测或具身智能控制在实际工业系统中的关键作用，以及高可靠性的刚需。\par
\noindent\textbf{Cohesion Formula:}\enspace {\small\color{AcadSlate}Given-to-New 链：Core Task (S1) $\to$ Practical Deployment (S2) $\to$ Real-world Disturbances (S3) $\to$ Consequence of Failure (S4). 绑定逻辑规则：L001, L011, L029.}\par\vspace{4pt}
\begin{ExampleBox}{Publication-Grade Authentic Text}
Multivariate time series forecasting underpins safety-critical scheduling across smart power grids, cloud clusters, and automated manufacturing. In these operational settings, control systems rely on multi-horizon predictions to balance supply fluctuations and prevent cascading overload. However, field sensors routinely operate under abrupt regime shifts, irregular sampling intervals, and high-frequency measurement noise. Failing to accommodate such non-stationary disturbances degrades long-range forecast reliability and triggers costly control oscillations.
\end{ExampleBox}
\begin{TipBox}{Rhetorical Move Breakdown \& Constituent Analysis}
\begin{itemize}[leftmargin=1.5em]
  \item \textbf{S1 [任务锚定 Topic Anchor]:} \textit{Multivariate time series forecasting underpins safety-critical scheduling across smart power grids, cloud clusters, and automated manufacturing.}\par\noindent{\small 对应步骤：S1 [任务锚定 Topic Anchor]：直接用紧凑主谓宾点明研究对象及其在真实系统中的核心功能（≤16词，禁用 In the era of AI）。 | 句首锚定旧信息/逻辑视角，主干动词紧跟主语，复杂修饰置于宾语之后（符合尾重原则 L011 与旧新信息流 L001）。}
  \item \textbf{S2 [应用场景展开 Concrete Scenarios]:} \textit{In these operational settings, control systems rely on multi-horizon predictions to balance supply fluctuations and prevent cascading overload.}\par\noindent{\small 对应步骤：S2 [应用场景展开 Concrete Scenarios]：列举 2--3 个具体工程场景（如电网负荷调度、灵巧机械手装配）及其决策依赖。 | 句首锚定旧信息/逻辑视角，主干动词紧跟主语，复杂修饰置于宾语之后（符合尾重原则 L011 与旧新信息流 L001）。}
  \item \textbf{S3 [苛刻工况约束 Operational Bottleneck]:} \textit{However, field sensors routinely operate under abrupt regime shifts, irregular sampling intervals, and high-frequency measurement noise.}\par\noindent{\small 对应步骤：S3 [苛刻工况约束 Operational Bottleneck]：指出真实部署环境下不可避免的物理/统计挑战（如非平稳漂移、传感器缺失、毫秒级延迟约束）。 | 句首锚定旧信息/逻辑视角，主干动词紧跟主语，复杂修饰置于宾语之后（符合尾重原则 L011 与旧新信息流 L001）。}
  \item \textbf{S4 [研究必要性收束 Imperative Hook]:} \textit{Failing to accommodate such non-stationary disturbances degrades long-range forecast reliability and triggers costly control oscillations.}\par\noindent{\small 对应步骤：S4 [研究必要性收束 Imperative Hook]：点明若无法克服上述约束会造成什么具体系统后果，自然引出下一段对现有方法的审视。 | 句首锚定旧信息/逻辑视角，主干动词紧跟主语，复杂修饰置于宾语之后（符合尾重原则 L011 与旧新信息流 L001）。}
\end{itemize}
\end{TipBox}
\medskip
\subsection*{Paragraph 3: 相关工作分类综述段（拒绝 A did X, B did Y 流水账） (\texttt{BP09})}
\noindent\textbf{Expressive Intent:}\enspace 我这一段想写 Related Work 的一个子方向：把过去三年的 6--8 篇代表性论文按技术路线分成两类进行归纳，并指出它们的共同盲区。\par
\noindent\textbf{Cohesion Formula:}\enspace {\small\color{AcadSlate}Two-Branch Taxonomy $\to$ Family A $\to$ Contrastive Family B $\to$ Shared Limitation. 绑定逻辑规则：L009, L008.}\par\vspace{4pt}
\begin{ExampleBox}{Publication-Grade Authentic Text}
Pre-trained time series foundation models generally follow either continuous patch reconstruction or discrete vocabulary tokenization. Continuous reconstruction frameworks (e.g., PatchTST, MOMENT) pre-train encoders via masked patch regression, preserving fine-grained numerical resolution across continuous domains. Conversely, discrete tokenization approaches (e.g., Chronos) quantize scaled observations into fixed bins to leverage cross-entropy language modeling objectives. Despite their zero-shot transferability, both paradigms treat univariate channels independently during pre-training, ignoring physical cross-variable conservation laws that govern multivariate engineering systems.
\end{ExampleBox}
\begin{TipBox}{Rhetorical Move Breakdown \& Constituent Analysis}
\begin{itemize}[leftmargin=1.5em]
  \item \textbf{S1 [分类总纲 Taxonomic Thesis]:} \textit{Pre-trained time series foundation models generally follow either continuous patch reconstruction or discrete vocabulary tokenization.}\par\noindent{\small 对应步骤：S1 [分类总纲 Taxonomic Thesis]：Existing methods for [Task] can be categorized into two main families: [Family A] and [Family B]. | 句首锚定旧信息/逻辑视角，主干动词紧跟主语，复杂修饰置于宾语之后（符合尾重原则 L011 与旧新信息流 L001）。}
  \item \textbf{S2 [流派 A 归纳与代表作 Family A Synthesis]:} \textit{Continuous reconstruction frameworks (e.g., PatchTST, MOMENT) pre-train encoders via masked patch regression, preserving fine-grained numerical resolution across continuous domains.}\par\noindent{\small 对应步骤：S2 [流派 A 归纳与代表作 Family A Synthesis]：[Family A] approaches (e.g., [Models]) focus on [Mechanism A], effectively addressing [Sub-problem 1]. | 句首锚定旧信息/逻辑视角，主干动词紧跟主语，复杂修饰置于宾语之后（符合尾重原则 L011 与旧新信息流 L001）。}
  \item \textbf{S3 [流派 B 归纳与对比 Family B Synthesis]:} \textit{Conversely, discrete tokenization approaches (e.g., Chronos) quantize scaled observations into fixed bins to leverage cross-entropy language modeling objectives.}\par\noindent{\small 对应步骤：S3 [流派 B 归纳与对比 Family B Synthesis]：By contrast, [Family B] architectures (e.g., [Models]) reformulate the problem via [Mechanism B] to reduce [Cost/Complexity]. | 句首锚定旧信息/逻辑视角，主干动词紧跟主语，复杂修饰置于宾语之后（符合尾重原则 L011 与旧新信息流 L001）。}
  \item \textbf{S4 [共同盲区提炼 Shared Blind Spot]:} \textit{Despite their zero-shot transferability, both paradigms treat univariate channels independently during pre-training, ignoring physical cross-variable conservation laws that govern multivariate engineering systems.}\par\noindent{\small 对应步骤：S4 [共同盲区提炼 Shared Blind Spot]：Despite their respective strengths, both families overlook [Critical Factor], which degrades performance under [Target Scenario]. | 句首锚定旧信息/逻辑视角，主干动词紧跟主语，复杂修饰置于宾语之后（符合尾重原则 L011 与旧新信息流 L001）。}
\end{itemize}
\end{TipBox}
\medskip
\subsection*{Paragraph 4: 对比现有流派优劣并精准定位本文差异（Contrastive Positioning） (\texttt{BP10})}
\noindent\textbf{Expressive Intent:}\enspace 我这一段想在 Related Work 小节末尾，把我们的方法与最接近的 2 个竞争方案做精准对比，防止审稿人说我们只是对某篇旧文的微调（Incremental）。\par
\noindent\textbf{Cohesion Formula:}\enspace {\small\color{AcadSlate}Acknowledge Overlap $\to$ Pinpoint Structural Divergence $\to$ Explain Functional Gain. 绑定逻辑规则：L008, L019.}\par\vspace{4pt}
\begin{ExampleBox}{Publication-Grade Authentic Text}
Among spectral forecasters, FEDformer and TimesNet are most closely related to our approach. While both methods transform temporal representations via fast Fourier transforms, they apply static top-k frequency selection globally across the entire lookback window, which blurs localized transient spikes. In contrast, our wavelet-coupled attention dynamically adapts filter bandwidths over sliding sub-intervals, isolating non-stationary burst events without corrupting global seasonal bases. Consequently, our model captures localized fault transients without requiring pre-defined frequency thresholds.
\end{ExampleBox}
\begin{TipBox}{Rhetorical Move Breakdown \& Constituent Analysis}
\begin{itemize}[leftmargin=1.5em]
  \item \textbf{S1 [锁定最相近工作 Closest Prior Work]:} \textit{Among spectral forecasters, FEDformer and TimesNet are most closely related to our approach.}\par\noindent{\small 对应步骤：S1 [锁定最相近工作 Closest Prior Work]：Among prior studies, [Closest Method A] and [Closest Method B] are most closely related to our approach. | 句首锚定旧信息/逻辑视角，主干动词紧跟主语，复杂修饰置于宾语之后（符合尾重原则 L011 与旧新信息流 L001）。}
  \item \textbf{S2 [剖析其局限性 Structural Difference 1]:} \textit{While both methods transform temporal representations via fast Fourier transforms, they apply static top-k frequency selection globally across the entire lookback window, which blurs localized transient spikes.}\par\noindent{\small 对应步骤：S2 [剖析其局限性 Structural Difference 1]：While [Method A] also employs [Shared Concept], it relies on [Limitation 1], restricting its applicability to [Narrow Setting]. | 句首锚定旧信息/逻辑视角，主干动词紧跟主语，复杂修饰置于宾语之后（符合尾重原则 L011 与旧新信息流 L001）。}
  \item \textbf{S3 [剖析本文本质突破 Core Departure]:} \textit{In contrast, our wavelet-coupled attention dynamically adapts filter bandwidths over sliding sub-intervals, isolating non-stationary burst events without corrupting global seasonal bases.}\par\noindent{\small 对应步骤：S3 [剖析本文本质突破 Core Departure]：In contrast, our framework introduces [Core Innovation], which fundamentally eliminates [Bottleneck] by [Mechanism]. | 句首锚定旧信息/逻辑视角，主干动词紧跟主语，复杂修饰置于宾语之后（符合尾重原则 L011 与旧新信息流 L001）。}
  \item \textbf{S4 [实证或理论优势锁定 Empirical/Theoretical Edge]:} \textit{Consequently, our model captures localized fault transients without requiring pre-defined frequency thresholds.}\par\noindent{\small 对应步骤：S4 [实证或理论优势锁定 Empirical/Theoretical Edge]：Consequently, our approach achieves [Capability] without requiring [Expensive Requirement of Prior Work]. | 句首锚定旧信息/逻辑视角，主干动词紧跟主语，复杂修饰置于宾语之后（符合尾重原则 L011 与旧新信息流 L001）。}
\end{itemize}
\end{TipBox}
\medskip
\subsection*{Paragraph 5: 问题形式化定义与符号约定段（Problem Formulation \& Notation） (\texttt{BP11})}
\noindent\textbf{Expressive Intent:}\enspace 我这一段想在 Methodology 开篇（Section 3.1 Problem Formulation），严谨、简洁地定义输入张量、输出目标、约束条件与优化目标。\par
\noindent\textbf{Cohesion Formula:}\enspace {\small\color{AcadSlate}Input Space $\to$ Auxiliary Conditions $\to$ Target Mapping $\to$ Crucial Distributional Constraint. 绑定逻辑规则：L014, L011.}\par\vspace{4pt}
\begin{ExampleBox}{Publication-Grade Authentic Text}
Consider a multivariate time series $\mathbf{X}_{1:L} = [x_1, \dots, x_L]^\top \in \mathbb{R}^{L \times C}$ recording $C$ correlated sensors over lookback horizon $L$. Given an binary observation mask $\mathbf{M} \in \{0, 1\}^{L \times C}$ indicating irregular sensor dropouts, our objective is to learn a generative mapping $p_\theta(\mathbf{Y}_{L+1:L+H} \mid \mathbf{X}_{1:L} \odot \mathbf{M})$ over future horizon $H$. Crucially, unlike standard settings that assume time-invariant marginal distributions, we explicitly account for non-stationary regime shifts where $\mathbb{E}[x_t]$ and $\text{Cov}(x_t)$ evolve continuously between the lookback and forecasting windows.
\end{ExampleBox}
\begin{TipBox}{Rhetorical Move Breakdown \& Constituent Analysis}
\begin{itemize}[leftmargin=1.5em]
  \item \textbf{S1 [任务总体形式化 Task Definition]:} \textit{Consider a multivariate time series $\mathbf{X}_{1:L} = [x_1, \dots, x_L]^\top \in \mathbb{R}^{L \times C}$ recording $C$ correlated sensors over lookback horizon $L$.}\par\noindent{\small 对应步骤：S1 [任务总体形式化 Task Definition]：Consider a [System/Dataset] observed over a lookback window of length $L$, denoted as $\mathbf{X} = [x_1, \dots, x_L] \in \mathbb{R}^{L \times C}$, where $C$ represents [Meaning]. | 句首锚定旧信息/逻辑视角，主干动词紧跟主语，复杂修饰置于宾语之后（符合尾重原则 L011 与旧新信息流 L001）。}
  \item \textbf{S2 [补充条件或模态输入 Auxiliary/Condition Inputs]:} \textit{Given an binary observation mask $\mathbf{M} \in \{0, 1\}^{L \times C}$ indicating irregular sensor dropouts, our objective is to learn a generative mapping $p_\theta(\mathbf{Y}_{L+1:L+H} \mid \mathbf{X}_{1:L} \odot \mathbf{M})$ over future horizon $H$.}\par\noindent{\small 对应步骤：S2 [补充条件或模态输入 Auxiliary/Condition Inputs]：At each step $t$, the system also receives [Condition/Observation] $z_t \in \mathcal{Z}$ under [Noise/Missingness/Dynamic Regime]. | 句首锚定旧信息/逻辑视角，主干动词紧跟主语，复杂修饰置于宾语之后（符合尾重原则 L011 与旧新信息流 L001）。}
  \item \textbf{S3 [预测/控制目标 Objective]:} \textit{Crucially, unlike standard settings that assume time-invariant marginal distributions, we explicitly account for non-stationary regime shifts where $\mathbb{E}[x_t]$ and $\text{Cov}(x_t)$ evolve continuously between the lookback and forecasting windows.}\par\noindent{\small 对应步骤：S3 [预测/控制目标 Objective]：Our goal is to learn a parameterized mapping $f_\theta(\cdot)$ that predicts future states $\hat{\mathbf{Y}} \in \mathbb{R}^{H \times C}$ over horizon $H$ (or action trajectory $a_{t:t+K}$). | 句首锚定旧信息/逻辑视角，主干动词紧跟主语，复杂修饰置于宾语之后（符合尾重原则 L011 与旧新信息流 L001）。}
\end{itemize}
\end{TipBox}
\medskip
\subsection*{Paragraph 6: 生成模型专属：条件控制与采样去噪轨迹阐述段（Generative Modeling Mechanism） (\texttt{BP15})}
\noindent\textbf{Expressive Intent:}\enspace 我这一段想阐述生成模型（Diffusion / Flow Matching / Autoregressive）如何在满足外部物理约束或条件观测的前提下，生成高质量且多样化的序列/轨迹。\par
\noindent\textbf{Cohesion Formula:}\enspace {\small\color{AcadSlate}Guidance Artifact Problem $\to$ Intrinsic Velocity Conditioning $\to$ Tangent Projection Step $\to$ Fast Exact Convergence. 绑定逻辑规则：L026, L014.}\par\vspace{4pt}
\begin{ExampleBox}{Publication-Grade Authentic Text}
Applying post-hoc gradient guidance during diffusion sampling often pushes intermediate noisy states off the data manifold, introducing unphysical high-frequency oscillations. To enforce physical boundary constraints without manifold deviation, we incorporate observation consistency directly into the conditional flow-matching velocity field. At flow time $\tau \in [0, 1]$, our network predicts a residual velocity $v_\theta(x_\tau, \tau \mid c)$ that projects the transport trajectory onto the null space of the linear constraint operator. This projection guarantees exact boundary satisfaction at terminal time $\tau=1$ while maintaining straight transport paths amenable to 4-step Euler integration.
\end{ExampleBox}
\begin{TipBox}{Rhetorical Move Breakdown \& Constituent Analysis}
\begin{itemize}[leftmargin=1.5em]
  \item \textbf{S1 [无条件/朴素引导缺陷 Naive Guidance Bottleneck]:} \textit{Applying post-hoc gradient guidance during diffusion sampling often pushes intermediate noisy states off the data manifold, introducing unphysical high-frequency oscillations.}\par\noindent{\small 对应步骤：S1 [无条件/朴素引导缺陷 Naive Guidance Bottleneck]：Applying post-hoc classifier guidance to pre-trained diffusion samplers often pushes intermediate noisy states off the data manifold, inducing high-frequency artifacts. | 句首锚定旧信息/逻辑视角，主干动词紧跟主语，复杂修饰置于宾语之后（符合尾重原则 L011 与旧新信息流 L001）。}
  \item \textbf{S2 [本文流形内条件注入机制 Manifold-Consistent Conditioning]:} \textit{To enforce physical boundary constraints without manifold deviation, we incorporate observation consistency directly into the conditional flow-matching velocity field.}\par\noindent{\small 对应步骤：S2 [本文流形内条件注入机制 Manifold-Consistent Conditioning]：To enforce hard boundary constraints without manifold deviation, we embed [Condition/Constraint] directly into the velocity field / denoising kernel. | 句首锚定旧信息/逻辑视角，主干动词紧跟主语，复杂修饰置于宾语之后（符合尾重原则 L011 与旧新信息流 L001）。}
  \item \textbf{S3 [具体轨迹更新方程 Trajectory Update Formulation]:} \textit{At flow time $\tau \in [0, 1]$, our network predicts a residual velocity $v_\theta(x_\tau, \tau \mid c)$ that projects the transport trajectory onto the null space of the linear constraint operator.}\par\noindent{\small 对应步骤：S3 [具体轨迹更新方程 Trajectory Update Formulation]：At solver step $\tau \in [0, 1]$, the conditional velocity field $v_\theta(x_\tau, \tau \mid c)$ projects the unconstrained flow onto the tangent space of feasible trajectories. | 句首锚定旧信息/逻辑视角，主干动词紧跟主语，复杂修饰置于宾语之后（符合尾重原则 L011 与旧新信息流 L001）。}
  \item \textbf{S4 [采样效率与分布保真收益 Sampling \& Fidelity Guarantee]:} \textit{This projection guarantees exact boundary satisfaction at terminal time $\tau=1$ while maintaining straight transport paths amenable to 4-step Euler integration.}\par\noindent{\small 对应步骤：S4 [采样效率与分布保真收益 Sampling \& Fidelity Guarantee]：This tangent projection guarantees exact boundary satisfaction at $\tau=1$ while permitting straight 4-step Euler integration. | 句首锚定旧信息/逻辑视角，主干动词紧跟主语，复杂修饰置于宾语之后（符合尾重原则 L011 与旧新信息流 L001）。}
\end{itemize}
\end{TipBox}
\medskip
\subsection*{Paragraph 7: 实验设置、数据集特性与强基线选型说明段（Experimental Setup \& Baselines） (\texttt{BP17})}
\noindent\textbf{Expressive Intent:}\enspace 我这一段想在 Experiments 开篇说明：我们为什么选这几类数据集和这几类基线模型，证明我们的实验设计是全面、公平且直击痛点的。\par
\noindent\textbf{Cohesion Formula:}\enspace {\small\color{AcadSlate}Dataset Diversity $\to$ Challenging Regime Highlights $\to$ Categorized Strong Baselines $\to$ Strict Fairness Protocol. 绑定逻辑规则：L009, L010.}\par\vspace{4pt}
\begin{ExampleBox}{Publication-Grade Authentic Text}
To evaluate robustness across diverse non-stationary regimes, we conduct experiments on seven public benchmarks spanning energy grid load (ETT, ECL), meteorological telemetry (Weather), and urban traffic flow (Traffic). Notably, Weather contains abrupt atmospheric pressure transitions that test non-stationary adaptability, whereas Traffic spans 862 synchronized sensors that challenge high-dimensional scalability. We benchmark our approach against 11 state-of-the-art methods across three paradigms: channel-independent patch Transformers (PatchTST), variate-mixing architectures (iTransformer, Crossformer), and frequency-domain forecasters (TimesNet, FEDformer). All baselines are reproduced under identical lookback windows ($L=96$) and prediction horizons ($H \in \{96, 192, 336, 720\}$) using standard MSE and MAE metrics.
\end{ExampleBox}
\begin{TipBox}{Rhetorical Move Breakdown \& Constituent Analysis}
\begin{itemize}[leftmargin=1.5em]
  \item \textbf{S1 [数据集选型逻辑 Dataset Selection Rationale]:} \textit{To evaluate robustness across diverse non-stationary regimes, we conduct experiments on seven public benchmarks spanning energy grid load (ETT, ECL), meteorological telemetry (Weather), and urban traffic flow (Traffic).}\par\noindent{\small 对应步骤：S1 [数据集选型逻辑 Dataset Selection Rationale]：To evaluate performance across diverse non-stationary regimes, we conduct experiments on [N] benchmarks spanning [Domain 1], [Domain 2], and [Domain 3]. | 句首锚定旧信息/逻辑视角，主干动词紧跟主语，复杂修饰置于宾语之后（符合尾重原则 L011 与旧新信息流 L001）。}
  \item \textbf{S2 [极端/困难特性突出 Challenging Dataset Properties]:} \textit{Notably, Weather contains abrupt atmospheric pressure transitions that test non-stationary adaptability, whereas Traffic spans 862 synchronized sensors that challenge high-dimensional scalability.}\par\noindent{\small 对应步骤：S2 [极端/困难特性突出 Challenging Dataset Properties]：Notably, [Dataset A] exhibits [Hard Property 1, e.g., strong covariate shift], whereas [Dataset B] tests scalability across [Hard Property 2, e.g., 862 correlated channels]. | 句首锚定旧信息/逻辑视角，主干动词紧跟主语，复杂修饰置于宾语之后（符合尾重原则 L011 与旧新信息流 L001）。}
  \item \textbf{S3 [基线流派全覆盖 Baseline Family Coverage]:} \textit{We benchmark our approach against 11 state-of-the-art methods across three paradigms: channel-independent patch Transformers (PatchTST), variate-mixing architectures (iTransformer, Crossformer), and frequency-domain forecasters (TimesNet, FEDformer).}\par\noindent{\small 对应步骤：S3 [基线流派全覆盖 Baseline Family Coverage]：We compare our method against [K] competitive baselines grouped into three families: [Family 1], [Family 2], and [Family 3]. | 句首锚定旧信息/逻辑视角，主干动词紧跟主语，复杂修饰置于宾语之后（符合尾重原则 L011 与旧新信息流 L001）。}
  \item \textbf{S4 [公平评测协议 Fair Evaluation Protocol]:} \textit{All baselines are reproduced under identical lookback windows ($L=96$) and prediction horizons ($H \in \{96, 192, 336, 720\}$) using standard MSE and MAE metrics.}\par\noindent{\small 对应步骤：S4 [公平评测协议 Fair Evaluation Protocol]：All models are evaluated under identical lookback windows ($L=96$), prediction horizons ($H \in \{96, 192, 336, 720\}$), and hyperparameter search budgets. | 句首锚定旧信息/逻辑视角，主干动词紧跟主语，复杂修饰置于宾语之后（符合尾重原则 L011 与旧新信息流 L001）。}
\end{itemize}
\end{TipBox}
\medskip
\subsection*{Paragraph 8: 主实验结果整体对比与核心胜出原因剖析段（Main Results Analysis） (\texttt{BP18})}
\noindent\textbf{Expressive Intent:}\enspace 我这一段想解读主实验大表（Table 1）：既报告全局胜率，又挑出最有代表性的对比案例，解释为什么竞争对手在那里跌倒而我们稳赢。\par
\noindent\textbf{Cohesion Formula:}\enspace {\small\color{AcadSlate}Global Win Rate $\to$ Advantage over Family A $\to$ Resilience at Long Horizon vs. Family B $\to$ Mechanistic Takeaway. 绑定逻辑规则：L015, L013, L010.}\par\vspace{4pt}
\begin{ExampleBox}{Publication-Grade Authentic Text}
Table 1 reports the multivariate forecasting results, where DyTS-Net secures the lowest error in 25 out of 28 dataset-horizon configurations. Compared with channel-independent PatchTST, DyTS-Net reduces average MSE by 14.2\% on the 862-channel Traffic benchmark, where modeling spatial-temporal congestion propagation across adjacent road sensors is critical. Furthermore, while variate-mixing iTransformer performs competitively at $H=96$, its MSE escalates by 68\% at $H=720$ on volatile Weather series as unregularized cross-channel attention overfits transient wind gusts. By decoupling stationary phase coupling from neural-ODE scale projection, DyTS-Net suppresses long-horizon drift and maintains an 11.6\% margin over the strongest baseline at $H=720$.
\end{ExampleBox}
\begin{TipBox}{Rhetorical Move Breakdown \& Constituent Analysis}
\begin{itemize}[leftmargin=1.5em]
  \item \textbf{S1 [全局统治力概述 Overall Dominance]:} \textit{Table 1 reports the multivariate forecasting results, where DyTS-Net secures the lowest error in 25 out of 28 dataset-horizon configurations.}\par\noindent{\small 对应步骤：S1 [全局统治力概述 Overall Dominance]：Table 1 summarizes the multivariate forecasting results, where [Our Model] achieves the lowest error in [X] out of [Y] dataset-horizon configurations. | 句首锚定旧信息/逻辑视角，主干动词紧跟主语，复杂修饰置于宾语之后（符合尾重原则 L011 与旧新信息流 L001）。}
  \item \textbf{S2 [对比流派 1 的差距与原因 Contrast vs. Family 1]:} \textit{Compared with channel-independent PatchTST, DyTS-Net reduces average MSE by 14.2\% on the 862-channel Traffic benchmark, where modeling spatial-temporal congestion propagation across adjacent road sensors is critical.}\par\noindent{\small 对应步骤：S2 [对比流派 1 的差距与原因 Contrast vs. Family 1]：Compared with channel-independent baselines such as [Baseline A], our method reduces average MSE by [X\%] on [Dataset], where strong cross-sensor coupling is critical. | 句首锚定旧信息/逻辑视角，主干动词紧跟主语，复杂修饰置于宾语之后（符合尾重原则 L011 与旧新信息流 L001）。}
  \item \textbf{S3 [对比流派 2 在困难场景下的崩溃 Contrast vs. Family 2 under Hard Regime]:} \textit{Furthermore, while variate-mixing iTransformer performs competitively at $H=96$, its MSE escalates by 68\% at $H=720$ on volatile Weather series as unregularized cross-channel attention overfits transient wind gusts.}\par\noindent{\small 对应步骤：S3 [对比流派 2 在困难场景下的崩溃 Contrast vs. Family 2 under Hard Regime]：Furthermore, while [Baseline B] performs competitively at short horizons ($H=96$), its error escalates by [Y\%] at $H=720$ due to [Failure Mechanism]. | 句首锚定旧信息/逻辑视角，主干动词紧跟主语，复杂修饰置于宾语之后（符合尾重原则 L011 与旧新信息流 L001）。}
  \item \textbf{S4 [回扣本文核心机制 Mechanistic Confirmation]:} \textit{By decoupling stationary phase coupling from neural-ODE scale projection, DyTS-Net suppresses long-horizon drift and maintains an 11.6\% margin over the strongest baseline at $H=720$.}\par\noindent{\small 对应步骤：S4 [回扣本文核心机制 Mechanistic Confirmation]：By contrast, our [Core Module] bounds long-range drift, maintaining a [Z\%] margin over the second-best baseline at $H=720$. | 句首锚定旧信息/逻辑视角，主干动词紧跟主语，复杂修饰置于宾语之后（符合尾重原则 L011 与旧新信息流 L001）。}
\end{itemize}
\end{TipBox}
\medskip
\subsection*{Paragraph 9: 超参数敏感性与极端噪声鲁棒性分析段（Sensitivity \& Robustness Analysis） (\texttt{BP20})}
\noindent\textbf{Expressive Intent:}\enspace 我这一段想回应审稿人最常问的：你的模型是不是对超参数很敏感？在缺失值或传感器强噪声下会不会突然崩盘？\par
\noindent\textbf{Cohesion Formula:}\enspace {\small\color{AcadSlate}Goal $\to$ Wide Hyperparameter Plateau $\to$ Stress-Test Comparison $\to$ Mechanistic Explanation. 绑定逻辑规则：L010, L019.}\par\vspace{4pt}
\begin{ExampleBox}{Publication-Grade Authentic Text}
Figure 4 evaluates hyperparameter sensitivity and resilience against severe observation corruption. Varying the phase regularization weight $\lambda$ across $[0.05, 0.50]$ alters validation MSE by less than 1.4\%, indicating a wide, tuning-friendly operating plateau around $\lambda=0.15$. Under stress tests where 30\% of sensor readings are randomly masked or corrupted by impulse noise ($\text{SNR} = 5\text{ dB}$), point-wise Transformer baselines suffer a 34.2\% error surge, whereas DyTS-Net degrades by only 4.8\%. This robustness stems from our spectral energy gate, which automatically attenuates high-entropy corrupted channels prior to cross-variate message passing.
\end{ExampleBox}
\begin{TipBox}{Rhetorical Move Breakdown \& Constituent Analysis}
\begin{itemize}[leftmargin=1.5em]
  \item \textbf{S1 [分析目的 Objective]:} \textit{Figure 4 evaluates hyperparameter sensitivity and resilience against severe observation corruption.}\par\noindent{\small 对应步骤：S1 [分析目的 Objective]：We further examine sensitivity to key hyperparameters and resilience against severe sensor corruption in Figure 4. | 句首锚定旧信息/逻辑视角，主干动词紧跟主语，复杂修饰置于宾语之后（符合尾重原则 L011 与旧新信息流 L001）。}
  \item \textbf{S2 [超参数平坦区间 Plateau Stability]:} \textit{Varying the phase regularization weight $\lambda$ across $[0.05, 0.50]$ alters validation MSE by less than 1.4\%, indicating a wide, tuning-friendly operating plateau around $\lambda=0.15$.}\par\noindent{\small 对应步骤：S2 [超参数平坦区间 Plateau Stability]：Varying [Hyperparameter $\lambda$] across two orders of magnitude ($[0.01, 1.0]$) alters validation MSE by less than [1.8\%], demonstrating a wide stable operating plateau. | 句首锚定旧信息/逻辑视角，主干动词紧跟主语，复杂修饰置于宾语之后（符合尾重原则 L011 与旧新信息流 L001）。}
  \item \textbf{S3 [极端边缘工况压力测试 Stress Test under Corruption]:} \textit{Under stress tests where 30\% of sensor readings are randomly masked or corrupted by impulse noise ($\text{SNR} = 5\text{ dB}$), point-wise Transformer baselines suffer a 34.2\% error surge, whereas DyTS-Net degrades by only 4.8\%.}\par\noindent{\small 对应步骤：S3 [极端边缘工况压力测试 Stress Test under Corruption]：Under synthetic sensor dropouts up to [40\%] and Gaussian noise ($\sigma=0.3$), baseline accuracy plummets by over [35\%], whereas our model degrades by less than [6.5\%]. | 句首锚定旧信息/逻辑视角，主干动词紧跟主语，复杂修饰置于宾语之后（符合尾重原则 L011 与旧新信息流 L001）。}
  \item \textbf{S4 [鲁棒性物理归因 Robustness Mechanism]:} \textit{This robustness stems from our spectral energy gate, which automatically attenuates high-entropy corrupted channels prior to cross-variate message passing.}\par\noindent{\small 对应步骤：S4 [鲁棒性物理归因 Robustness Mechanism]：This resilience stems from [Specific Mechanism], which filters out corrupted channels before global aggregation. | 句首锚定旧信息/逻辑视角，主干动词紧跟主语，复杂修饰置于宾语之后（符合尾重原则 L011 与旧新信息流 L001）。}
\end{itemize}
\end{TipBox}
\medskip
\subsection*{Paragraph 10: 局限性（Limitations）客观陈述与未来扩展段（Constructive Limitations） (\texttt{BP23})}
\noindent\textbf{Expressive Intent:}\enspace 我这一段想写论文末尾或 NeurIPS/ICLR Checklist 要求的 Limitations 段落：诚实指出当前方法的两个客观边界，但绝不贬损论文的核心贡献。\par
\noindent\textbf{Cohesion Formula:}\enspace {\small\color{AcadSlate}Two Objective Boundaries $\to$ [Boundary 1 + Technical Remedy] $\to$ [Boundary 2 + Stochastic Extension]. 绑定逻辑规则：L017.}\par\vspace{4pt}
\begin{ExampleBox}{Publication-Grade Authentic Text}
Despite its strong empirical performance, our framework exhibits two theoretical boundaries worth noting. First, while our spectral branch scales quasi-linearly with sequence length, computing pairwise Hilbert phase differences scales quadratically with channel count ($O(C^2)$), which increases memory footprint when $C > 2,000$; incorporating anchor-based low-rank graph sparsification provides a direct remedy. Second, our continuous neural-ODE scale tracker assumes piece-wise smooth drift and may smooth over instantaneous step discontinuities caused by sensor recalibration. Extending the continuous drift prior to jump-diffusion stochastic differential equations offers a principled direction for future investigation.
\end{ExampleBox}
\begin{TipBox}{Rhetorical Move Breakdown \& Constituent Analysis}
\begin{itemize}[leftmargin=1.5em]
  \item \textbf{S1 [边界 1:} \textit{Despite its strong empirical performance, our framework exhibits two theoretical boundaries worth noting.}\par\noindent{\small 对应步骤：S1 [边界 1：极端高维或极度稀疏场景 Boundary 1]：Despite its effectiveness across standard benchmarks, our framework currently exhibits two practical boundaries. | 句首锚定旧信息/逻辑视角，主干动词紧跟主语，复杂修饰置于宾语之后（符合尾重原则 L011 与旧新信息流 L001）。}
  \item \textbf{S2 [边界 1 机理分析与解法 Mechanism \& Solution 1]:} \textit{First, while our spectral branch scales quasi-linearly with sequence length, computing pairwise Hilbert phase differences scales quadratically with channel count ($O(C^2)$), which increases memory footprint when $C > 2,000$; incorporating anchor-based low-rank graph sparsification provides a direct remedy.}\par\noindent{\small 对应步骤：S2 [边界 1 机理分析与解法 Mechanism \& Solution 1]：First, when the variate dimension scales beyond $C > 5,000$ unclustered sensors, constructing the cross-channel phase graph incurs noticeable memory overhead; integrating hierarchical sensor clustering offers a natural remedy. | 句首锚定旧信息/逻辑视角，主干动词紧跟主语，复杂修饰置于宾语之后（符合尾重原则 L011 与旧新信息流 L001）。}
  \item \textbf{S3 [边界 2:} \textit{Second, our continuous neural-ODE scale tracker assumes piece-wise smooth drift and may smooth over instantaneous step discontinuities caused by sensor recalibration.}\par\noindent{\small 对应步骤：S3 [边界 2：特定假设限制 Boundary 2]：Second, our continuous scale tracker assumes piece-wise differentiable drift and may under-estimate instantaneous discontinuous jumps caused by hardware resets. | 句首锚定旧信息/逻辑视角，主干动词紧跟主语，复杂修饰置于宾语之后（符合尾重原则 L011 与旧新信息流 L001）。}
  \item \textbf{S4 [未来演进方向 Future Extension]:} \textit{Extending the continuous drift prior to jump-diffusion stochastic differential equations offers a principled direction for future investigation.}\par\noindent{\small 对应步骤：S4 [未来演进方向 Future Extension]：Coupling our neural ODE tracker with jump-diffusion stochastic processes represents a promising avenue for future work. | 句首锚定旧信息/逻辑视角，主干动词紧跟主语，复杂修饰置于宾语之后（符合尾重原则 L011 与旧新信息流 L001）。}
\end{itemize}
\end{TipBox}
\medskip
\subsection*{Paragraph 11: 全文总结与学术影响收尾段（Conclusion \& Broader Outlook） (\texttt{BP24})}
\noindent\textbf{Expressive Intent:}\enspace 我这一段想写整篇小论文的 Conclusion：不简单复制粘贴 Abstract，而是从更高维度总结本文解决的科学矛盾、方法精髓、实证高度与领域启示。\par
\noindent\textbf{Cohesion Formula:}\enspace {\small\color{AcadSlate}Tension Resolved $\to$ Methodological Insight $\to$ Empirical Confirmation $\to$ Generalizable Blueprint. 绑定逻辑规则：L013, L023, L029.}\par\vspace{4pt}
\begin{ExampleBox}{Publication-Grade Authentic Text}
In this work, we resolved the long-standing tension between shift-invariant normalization and physical amplitude preservation in non-stationary time series forecasting. By decoupling frequency-domain phase attention from continuous neural-ODE scale tracking, DyTS-Net captures long-range seasonal dependencies while accurately reconstructing extreme peak magnitudes. Empirical evaluations across eight benchmarks and industrial wind-power telemetry demonstrated MSE reductions of up to 17.4\% alongside a 45\% reduction in computational overhead. Beyond forecasting, this structural-scale disentanglement principle provides a generalizable foundation for anomaly detection and closed-loop predictive control in non-stationary cyber-physical systems.
\end{ExampleBox}
\begin{TipBox}{Rhetorical Move Breakdown \& Constituent Analysis}
\begin{itemize}[leftmargin=1.5em]
  \item \textbf{S1 [核心科学问题闭合 Scientific Closure]:} \textit{In this work, we resolved the long-standing tension between shift-invariant normalization and physical amplitude preservation in non-stationary time series forecasting.}\par\noindent{\small 对应步骤：S1 [核心科学问题闭合 Scientific Closure]：In this paper, we addressed the fundamental tension between [Goal A] and [Goal B] in [Domain]. | 句首锚定旧信息/逻辑视角，主干动词紧跟主语，复杂修饰置于宾语之后（符合尾重原则 L011 与旧新信息流 L001）。}
  \item \textbf{S2 [方法论精髓回顾 Methodological Synthesis]:} \textit{By decoupling frequency-domain phase attention from continuous neural-ODE scale tracking, DyTS-Net captures long-range seasonal dependencies while accurately reconstructing extreme peak magnitudes.}\par\noindent{\small 对应步骤：S2 [方法论精髓回顾 Methodological Synthesis]：By introducing [Model/Mechanism], we demonstrated that [Core Insight] enables [Capability] without incurring [Traditional Penalty]. | 句首锚定旧信息/逻辑视角，主干动词紧跟主语，复杂修饰置于宾语之后（符合尾重原则 L011 与旧新信息流 L001）。}
  \item \textbf{S3 [实证结果高度概括 Empirical Recap]:} \textit{Empirical evaluations across eight benchmarks and industrial wind-power telemetry demonstrated MSE reductions of up to 17.4\% alongside a 45\% reduction in computational overhead.}\par\noindent{\small 对应步骤：S3 [实证结果高度概括 Empirical Recap]：Comprehensive evaluations across [Benchmarks/Hardware] confirmed consistent error reductions of up to [X\%] alongside substantial inference speedups. | 句首锚定旧信息/逻辑视角，主干动词紧跟主语，复杂修饰置于宾语之后（符合尾重原则 L011 与旧新信息流 L001）。}
  \item \textbf{S4 [跨领域迁移展望 Broader Outlook]:} \textit{Beyond forecasting, this structural-scale disentanglement principle provides a generalizable foundation for anomaly detection and closed-loop predictive control in non-stationary cyber-physical systems.}\par\noindent{\small 对应步骤：S4 [跨领域迁移展望 Broader Outlook]：Beyond [Current Task], our [Core Principle] offers a generalizable blueprint for broader sequential decision-making tasks under non-stationary environments. | 句首锚定旧信息/逻辑视角，主干动词紧跟主语，复杂修饰置于宾语之后（符合尾重原则 L011 与旧新信息流 L001）。}
\end{itemize}
\end{TipBox}
\medskip
\clearpage
\section{FT03: 时频双域解耦的长序列异常检测与预测完整小论文（WavePatch: Multi-Scale Wavelet-Temporal Mixing）}\label{comp:FT03}
\begin{ConceptBox}{Composition Architecture \& Structural Blueprint Chain}
\textbf{Genre:} \texttt{AI\_Mini\_Paper} \quad \textbf{Paragraph Count:} 11\par
\textbf{Blueprint Chain:} \texttt{BP08 $\to$ BP02 $\to$ BP03 $\to$ BP05 $\to$ BP06 $\to$ BP12 $\to$ BP14 $\to$ BP18 $\to$ BP19 $\to$ BP21 $\to$ BP24}
\end{ConceptBox}
\subsection*{Paragraph 1: 高密度四句式 Abstract 全段构建（Complete 4-Move Abstract） (\texttt{BP08})}
\noindent\textbf{Expressive Intent:}\enspace 我这一段想写出整篇小论文的完整 Abstract（摘要）：在 150--180 词内一气呵成交代背景、核心痛点、方法机制与量化实验结果。\par
\noindent\textbf{Cohesion Formula:}\enspace {\small\color{AcadSlate}Context (S1) $\to$ Gap (S2) $\to$ Method \& Mechanism (S3-S4) $\to$ Quantitative Proof (S5). 绑定逻辑规则：L001, L008, L010, L029.}\par\vspace{4pt}
\begin{ExampleBox}{Publication-Grade Authentic Text}
Recovering missing sensor observations is essential for reliable downstream monitoring in industrial and clinical time series. While diffusion-based imputation models generate plausible trajectories under random missingness, they suffer from high sampling latency and lose global periodic coherence when entire sensor blocks drop out consecutively. We present FlowImp, a conditional flow-matching framework that reconstructs contiguous missing blocks in three ODE solver steps. Specifically, FlowImp anchors straight transport paths on observed spectral priors and enforces cross-channel phase consistency via a lightweight Fourier regularizer. Across six multivariate benchmarks under 50\% block missingness, FlowImp reduces imputation MAE by 19.3\% compared with CSDI while accelerating inference by 24x.
\end{ExampleBox}
\begin{TipBox}{Rhetorical Move Breakdown \& Constituent Analysis}
\begin{itemize}[leftmargin=1.5em]
  \item \textbf{S1 [背景与任务重要性 Context]:} \textit{Recovering missing sensor observations is essential for reliable downstream monitoring in industrial and clinical time series.}\par\noindent{\small 对应步骤：S1 [背景与任务重要性 Context]：1句话锁定任务目标与核心应用场景（≤20词）。 | 句首锚定旧信息/逻辑视角，主干动词紧跟主语，复杂修饰置于宾语之后（符合尾重原则 L011 与旧新信息流 L001）。}
  \item \textbf{S2 [现有局限与痛点 Gap/Bottleneck]:} \textit{While diffusion-based imputation models generate plausible trajectories under random missingness, they suffer from high sampling latency and lose global periodic coherence when entire sensor blocks drop out consecutively.}\par\noindent{\small 对应步骤：S2 [现有局限与痛点 Gap/Bottleneck]：While existing [Paradigms] achieve [Merit], they struggle with [Core Bottleneck] because [Underlying Cause]. | 句首锚定旧信息/逻辑视角，主干动词紧跟主语，复杂修饰置于宾语之后（符合尾重原则 L011 与旧新信息流 L001）。}
  \item \textbf{S3-S4 [本文方法与核心机制 Approach \& Mechanism]:} \textit{We present FlowImp, a conditional flow-matching framework that reconstructs contiguous missing blocks in three ODE solver steps.}\par\noindent{\small 对应步骤：S3-S4 [本文方法与核心机制 Approach \& Mechanism]：In this work, we present [Model], which [Core Idea]. Specifically, [Module 1]..., while [Module 2]... | 句首锚定旧信息/逻辑视角，主干动词紧跟主语，复杂修饰置于宾语之后（符合尾重原则 L011 与旧新信息流 L001）。}
  \item \textbf{S5 [量化实证结果 Quantitative Results]:} \textit{Specifically, FlowImp anchors straight transport paths on observed spectral priors and enforces cross-channel phase consistency via a lightweight Fourier regularizer.}\par\noindent{\small 对应步骤：S5 [量化实证结果 Quantitative Results]：Experiments on [Benchmarks/Hardware] show that [Model] outperforms [Baselines] by [X\%] in [Metric] and achieves [Efficiency/Real-world Result]. | 句首锚定旧信息/逻辑视角，主干动词紧跟主语，复杂修饰置于宾语之后（符合尾重原则 L011 与旧新信息流 L001）。}
  \item \textbf{S5 [量化实证结果 Quantitative Results]:} \textit{Across six multivariate benchmarks under 50\% block missingness, FlowImp reduces imputation MAE by 19.3\% compared with CSDI while accelerating inference by 24x.}\par\noindent{\small 对应步骤：S5 [量化实证结果 Quantitative Results]：Experiments on [Benchmarks/Hardware] show that [Model] outperforms [Baselines] by [X\%] in [Metric] and achieves [Efficiency/Real-world Result]. | 句首锚定旧信息/逻辑视角，主干动词紧跟主语，复杂修饰置于宾语之后（符合尾重原则 L011 与旧新信息流 L001）。}
\end{itemize}
\end{TipBox}
\medskip
\subsection*{Paragraph 2: 总结现有主流范式及其隐含假设（欲抑先扬铺垫段） (\texttt{BP02})}
\noindent\textbf{Expressive Intent:}\enspace 我这一段想客观概括前人主流模型（如 Transformer / Diffusion / MLP）取得了什么进展，但点出它们共同依赖的一个脆弱前提假设。\par
\noindent\textbf{Cohesion Formula:}\enspace {\small\color{AcadSlate}先让步肯定（While/Recent paradigms...）$\to$ 后釜底抽薪揭示隐含前提（Yet these architectures implicitly assume...）。绑定逻辑规则：L008, L009.}\par\vspace{4pt}
\begin{ExampleBox}{Publication-Grade Authentic Text}
Recent deep forecasters have substantially reduced long-horizon error by replacing point-wise attention with sub-series patch tokenization and frequency-domain decomposition. These architectures effectively capture recurring seasonal patterns when training and test sequences share stationary statistical moments. Yet they implicitly assume that inter-channel correlations and spectral energy distributions remain time-invariant across the prediction window. When external shocks alter underlying system dynamics, this stationarity assumption collapses, causing frozen spectral filters to project outdated trends.
\end{ExampleBox}
\begin{TipBox}{Rhetorical Move Breakdown \& Constituent Analysis}
\begin{itemize}[leftmargin=1.5em]
  \item \textbf{S1 [流派归类 Taxonomic Opening]:} \textit{Recent deep forecasters have substantially reduced long-horizon error by replacing point-wise attention with sub-series patch tokenization and frequency-domain decomposition.}\par\noindent{\small 对应步骤：S1 [流派归类 Taxonomic Opening]：按核心机制归纳现有主流范式及其已解决的子问题。 | 句首锚定旧信息/逻辑视角，主干动词紧跟主语，复杂修饰置于宾语之后（符合尾重原则 L011 与旧新信息流 L001）。}
  \item \textbf{S2 [代表性机制展开 Mechanism Recap]:} \textit{These architectures effectively capture recurring seasonal patterns when training and test sequences share stationary statistical moments.}\par\noindent{\small 对应步骤：S2 [代表性机制展开 Mechanism Recap]：简述这些方法如何通过特定架构（如 Patch 分块、频域分解或扩散去噪）提升表征能力。 | 句首锚定旧信息/逻辑视角，主干动词紧跟主语，复杂修饰置于宾语之后（符合尾重原则 L011 与旧新信息流 L001）。}
  \item \textbf{S3 [转折揭示隐含假设 Pivot to Hidden Assumption]:} \textit{Yet they implicitly assume that inter-channel correlations and spectral energy distributions remain time-invariant across the prediction window.}\par\noindent{\small 对应步骤：S3 [转折揭示隐含假设 Pivot to Hidden Assumption]：用 Yet / However 锋利转折，指出它们暗中依赖的理想化假设（如平稳分布、静态通道相关性、仿真动力学无偏差）。 | 句首锚定旧信息/逻辑视角，主干动词紧跟主语，复杂修饰置于宾语之后（符合尾重原则 L011 与旧新信息流 L001）。}
  \item \textbf{S4 [假设破裂后果 Breakdown Under Stress]:} \textit{When external shocks alter underlying system dynamics, this stationarity assumption collapses, causing frozen spectral filters to project outdated trends.}\par\noindent{\small 对应步骤：S4 [假设破裂后果 Breakdown Under Stress]：说明当真实数据违背该假设时，现有架构为何会遭遇性能断崖。 | 句首锚定旧信息/逻辑视角，主干动词紧跟主语，复杂修饰置于宾语之后（符合尾重原则 L011 与旧新信息流 L001）。}
\end{itemize}
\end{TipBox}
\medskip
\subsection*{Paragraph 3: 尖锐指出核心技术瓶颈与失效机理（Gap Identification） (\texttt{BP03})}
\noindent\textbf{Expressive Intent:}\enspace 我这一段想深入剖析现有方法在特定难题（如非平稳突变、长程自回归误差累积、Sim-to-Real 接触失配）上为什么会失效，把机理讲透。\par
\noindent\textbf{Cohesion Formula:}\enspace {\small\color{AcadSlate}因果链式推进（A induces B $\to$ This B compounds into C $\to$ Naive fix D sacrifices E）。绑定逻辑规则：L003, L004, L024/L025.}\par\vspace{4pt}
\begin{ExampleBox}{Publication-Grade Authentic Text}
Bridging non-stationary robustness and amplitude fidelity remains a core bottleneck in long-horizon forecasting. While instance normalization strips time-varying mean and variance shifts to stabilize attention weights, it simultaneously erases physical scale cues required to predict extreme load peaks. Conversely, retaining raw unnormalized inputs exposes self-attention layers to severe covariate shift, causing validation error to diverge over extended horizons. Resolving this tension requires a mechanism that disentangles intrinsic temporal dynamics from transient scale drift without discarding magnitude information.
\end{ExampleBox}
\begin{TipBox}{Rhetorical Move Breakdown \& Constituent Analysis}
\begin{itemize}[leftmargin=1.5em]
  \item \textbf{S1 [核心痛点立论 Bottleneck Claim]:} \textit{Bridging non-stationary robustness and amplitude fidelity remains a core bottleneck in long-horizon forecasting.}\par\noindent{\small 对应步骤：S1 [核心痛点立论 Bottleneck Claim]：一句话直接命名该核心瓶颈（如 Compound error accumulation / Sim-to-real contact mismatch）。 | 句首锚定旧信息/逻辑视角，主干动词紧跟主语，复杂修饰置于宾语之后（符合尾重原则 L011 与旧新信息流 L001）。}
  \item \textbf{S2 [第一层失效机理 Primary Failure Mechanism]:} \textit{While instance normalization strips time-varying mean and variance shifts to stabilize attention weights, it simultaneously erases physical scale cues required to predict extreme load peaks.}\par\noindent{\small 对应步骤：S2 [第一层失效机理 Primary Failure Mechanism]：从数学或物理传导链条解释误差是如何产生并放大的。 | 句首锚定旧信息/逻辑视角，主干动词紧跟主语，复杂修饰置于宾语之后（符合尾重原则 L011 与旧新信息流 L001）。}
  \item \textbf{S3 [第二层耦合困境 Secondary Dilemma]:} \textit{Conversely, retaining raw unnormalized inputs exposes self-attention layers to severe covariate shift, causing validation error to diverge over extended horizons.}\par\noindent{\small 对应步骤：S3 [第二层耦合困境 Secondary Dilemma]：指出为什么简单的修补手段（如单纯加大模型、简单归一化或域随机化）不仅无效还会带来副作用。 | 句首锚定旧信息/逻辑视角，主干动词紧跟主语，复杂修饰置于宾语之后（符合尾重原则 L011 与旧新信息流 L001）。}
  \item \textbf{S4 [凝练核心研究问题 Research Question Formulation]:} \textit{Resolving this tension requires a mechanism that disentangles intrinsic temporal dynamics from transient scale drift without discarding magnitude information.}\par\noindent{\small 对应步骤：S4 [凝练核心研究问题 Research Question Formulation]：用 Therefore, a fundamental question arises: how can we... 引出本文目标。 | 句首锚定旧信息/逻辑视角，主干动词紧跟主语，复杂修饰置于宾语之后（符合尾重原则 L011 与旧新信息流 L001）。}
\end{itemize}
\end{TipBox}
\medskip
\subsection*{Paragraph 4: 挑战与机制一对一映射段（Challenge-to-Mechanism Mapping） (\texttt{BP05})}
\noindent\textbf{Expressive Intent:}\enspace 我这一段想说明：把我们的核心思想落地时遇到了两个具体工程/数学难点，我们分别设计了什么精巧机制去攻克它们。\par
\noindent\textbf{Cohesion Formula:}\enspace {\small\color{AcadSlate}Challenge 1 $\to$ Mechanism 1; Challenge 2 $\to$ Mechanism 2 $\to$ Joint System Property. 绑定逻辑规则：L003, L011.}\par\vspace{4pt}
\begin{ExampleBox}{Publication-Grade Authentic Text}
Instantiating closed-loop visuomotor diffusion on physical robots entails two technical challenges. First, because raw depth camera streams contain occlusion artifacts around reflective object edges, we introduce a tactile-guided cross-attention gate that reweights visual tokens upon contact onset. Second, to eliminate the 200 ms sampling bottleneck of iterative denoising, we formulate a trajectory consistency distillation objective that maps Gaussian noise directly to collision-free action chunks in two forward steps. Together, these designs sustain 50 Hz reactive control under partial visual occlusion.
\end{ExampleBox}
\begin{TipBox}{Rhetorical Move Breakdown \& Constituent Analysis}
\begin{itemize}[leftmargin=1.5em]
  \item \textbf{S1 [落地难点总起 Implementation Hurdle]:} \textit{Instantiating closed-loop visuomotor diffusion on physical robots entails two technical challenges.}\par\noindent{\small 对应步骤：S1 [落地难点总起 Implementation Hurdle]：Translating this principle into a scalable architecture entails two technical challenges. | 句首锚定旧信息/逻辑视角，主干动词紧跟主语，复杂修饰置于宾语之后（符合尾重原则 L011 与旧新信息流 L001）。}
  \item \textbf{S2 [难点 1 + 靶向设计 1]:} \textit{First, because raw depth camera streams contain occlusion artifacts around reflective object edges, we introduce a tactile-guided cross-attention gate that reweights visual tokens upon contact onset.}\par\noindent{\small 对应步骤：S2 [难点 1 + 靶向设计 1]：First, because [Challenge 1], we design [Mechanism 1] to [Specific Function]. | 句首锚定旧信息/逻辑视角，主干动词紧跟主语，复杂修饰置于宾语之后（符合尾重原则 L011 与旧新信息流 L001）。}
  \item \textbf{S3 [难点 2 + 靶向设计 2]:} \textit{Second, to eliminate the 200 ms sampling bottleneck of iterative denoising, we formulate a trajectory consistency distillation objective that maps Gaussian noise directly to collision-free action chunks in two forward steps.}\par\noindent{\small 对应步骤：S3 [难点 2 + 靶向设计 2]：Second, to prevent [Challenge 2] under [Constraint], we introduce [Mechanism 2] that [Specific Function]. | 句首锚定旧信息/逻辑视角，主干动词紧跟主语，复杂修饰置于宾语之后（符合尾重原则 L011 与旧新信息流 L001）。}
  \item \textbf{S4 [整体计算/泛化特性 System Property]:} \textit{Together, these designs sustain 50 Hz reactive control under partial visual occlusion.}\par\noindent{\small 对应步骤：S4 [整体计算/泛化特性 System Property]：Together, these components enable [Model] to operate at [Efficiency/Robustness Metric]. | 句首锚定旧信息/逻辑视角，主干动词紧跟主语，复杂修饰置于宾语之后（符合尾重原则 L011 与旧新信息流 L001）。}
\end{itemize}
\end{TipBox}
\medskip
\subsection*{Paragraph 5: 实验证言与性能突破预览段（Quantitative Empirical Preview） (\texttt{BP06})}
\noindent\textbf{Expressive Intent:}\enspace 我这一段想在 Introduction 结尾前，用三句话硬核总结我们在哪些基准数据集/真实硬件上做了实验，取得了多大幅度的领先。\par
\noindent\textbf{Cohesion Formula:}\enspace {\small\color{AcadSlate}Evaluation Scope $\to$ Accuracy Gain over Strongest Baselines $\to$ Efficiency/Generalization Gain. 绑定逻辑规则：L010, L015.}\par\vspace{4pt}
\begin{ExampleBox}{Publication-Grade Authentic Text}
We evaluate DyTS-Net across eight long-horizon forecasting benchmarks and a real-world wind-turbine telemetry system spanning 1,400 sensors. Across all prediction horizons up to 720 steps, DyTS-Net reduces MSE by 12.8\% over PatchTST and 17.4\% over iTransformer, with the largest gains emerging during sudden weather-front transitions. Furthermore, our architecture requires 45\% fewer FLOPs than full channel-mixing Transformers and transfers zero-shot to newly deployed turbine units without fine-tuning.
\end{ExampleBox}
\begin{TipBox}{Rhetorical Move Breakdown \& Constituent Analysis}
\begin{itemize}[leftmargin=1.5em]
  \item \textbf{S1 [评测广度 Scope of Evaluation]:} \textit{We evaluate DyTS-Net across eight long-horizon forecasting benchmarks and a real-world wind-turbine telemetry system spanning 1,400 sensors.}\par\noindent{\small 对应步骤：S1 [评测广度 Scope of Evaluation]：We evaluate [Model] across [N] public benchmarks and [Real-world Hardware/Scenario], spanning [Tasks]. | 句首锚定旧信息/逻辑视角，主干动词紧跟主语，复杂修饰置于宾语之后（符合尾重原则 L011 与旧新信息流 L001）。}
  \item \textbf{S2 [核心精度突破 Primary Accuracy Gain]:} \textit{Across all prediction horizons up to 720 steps, DyTS-Net reduces MSE by 12.8\% over PatchTST and 17.4\% over iTransformer, with the largest gains emerging during sudden weather-front transitions.}\par\noindent{\small 对应步骤：S2 [核心精度突破 Primary Accuracy Gain]：实验主指标对比（必须包含具体百分比、最强基线名称及困难场景表现）。 | 句首锚定旧信息/逻辑视角，主干动词紧跟主语，复杂修饰置于宾语之后（符合尾重原则 L011 与旧新信息流 L001）。}
  \item \textbf{S3 [效率或泛化突破 Efficiency / Zero-Shot Gain]:} \textit{Furthermore, our architecture requires 45\% fewer FLOPs than full channel-mixing Transformers and transfers zero-shot to newly deployed turbine units without fine-tuning.}\par\noindent{\small 对应步骤：S3 [效率或泛化突破 Efficiency / Zero-Shot Gain]：报告计算速度、显存占用或零样本跨场景泛化收益。 | 句首锚定旧信息/逻辑视角，主干动词紧跟主语，复杂修饰置于宾语之后（符合尾重原则 L011 与旧新信息流 L001）。}
\end{itemize}
\end{TipBox}
\medskip
\subsection*{Paragraph 6: 整体模型架构数据流俯瞰段（Framework Overview \& Pipeline Walkthrough） (\texttt{BP12})}
\noindent\textbf{Expressive Intent:}\enspace 我这一段想配合论文里的 Figure 2（模型总架构图），用一段话把输入张量从左到右流经各个模块直到输出预测结果的完整全貌讲清楚。\par
\noindent\textbf{Cohesion Formula:}\enspace {\small\color{AcadSlate}严格的链式数据流（$\mathbf{X} \to \mathbf{Z}^{(1)} \to \mathbf{Z}^{(2)} \to \hat{\mathbf{Y}}$），每一句句首承接上一句的输出张量。绑定逻辑规则：L001, L003, L004.}\par\vspace{4pt}
\begin{ExampleBox}{Publication-Grade Authentic Text}
Figure 2 illustrates the overall architecture of DyTS-Net, which maps raw sensor streams to multi-horizon forecasts through three sequential stages. Given an input sequence $\mathbf{X}$, the Adaptive De-stationary Tokenizer first separates local trend-scale statistics from normalized sub-series patches $\mathbf{Z}_0$. These normalized patches then pass through stacked Spectral-Spatial Mixing blocks, which jointly filter high-frequency channel noise in the Fourier domain and propagate cross-sensor phase dependencies via a sparse directed graph. Finally, the Scale-Guided Decoder fuses the refined structural representations with the projected neural-ODE scale trajectory to reconstruct the multi-horizon output $\hat{\mathbf{Y}}$.
\end{ExampleBox}
\begin{TipBox}{Rhetorical Move Breakdown \& Constituent Analysis}
\begin{itemize}[leftmargin=1.5em]
  \item \textbf{S1 [总图锚定与设计哲学 Visual Anchor \& Philosophy]:} \textit{Figure 2 illustrates the overall architecture of DyTS-Net, which maps raw sensor streams to multi-horizon forecasts through three sequential stages.}\par\noindent{\small 对应步骤：S1 [总图锚定与设计哲学 Visual Anchor \& Philosophy]：Figure 2 illustrates the overall architecture of [Model], which processes [Input] through three sequential stages: [Stage 1], [Stage 2], and [Stage 3]. | 句首锚定旧信息/逻辑视角，主干动词紧跟主语，复杂修饰置于宾语之后（符合尾重原则 L011 与旧新信息流 L001）。}
  \item \textbf{S2 [阶段 1 数据流 Stage 1 Flow]:} \textit{Given an input sequence $\mathbf{X}$, the Adaptive De-stationary Tokenizer first separates local trend-scale statistics from normalized sub-series patches $\mathbf{Z}_0$.}\par\noindent{\small 对应步骤：S2 [阶段 1 数据流 Stage 1 Flow]：Given raw input $\mathbf{X}$, the [Module 1] first transforms [Raw Input] into [Intermediate Representation $\mathbf{Z}^{(1)}$] to [Functional Purpose 1]. | 句首锚定旧信息/逻辑视角，主干动词紧跟主语，复杂修饰置于宾语之后（符合尾重原则 L011 与旧新信息流 L001）。}
  \item \textbf{S3 [阶段 2 数据流 Stage 2 Flow]:} \textit{These normalized patches then pass through stacked Spectral-Spatial Mixing blocks, which jointly filter high-frequency channel noise in the Fourier domain and propagate cross-sensor phase dependencies via a sparse directed graph.}\par\noindent{\small 对应步骤：S3 [阶段 2 数据流 Stage 2 Flow]：Next, these [Summary Noun: representations $\mathbf{Z}^{(1)}$] are fed into [Module 2], which [Mechanism 2] to produce [Refined State $\mathbf{Z}^{(2)}$]. | 句首锚定旧信息/逻辑视角，主干动词紧跟主语，复杂修饰置于宾语之后（符合尾重原则 L011 与旧新信息流 L001）。}
  \item \textbf{S4 [阶段 3 解码与闭环输出 Stage 3 Decoding]:} \textit{Finally, the Scale-Guided Decoder fuses the refined structural representations with the projected neural-ODE scale trajectory to reconstruct the multi-horizon output $\hat{\mathbf{Y}}$.}\par\noindent{\small 对应步骤：S4 [阶段 3 解码与闭环输出 Stage 3 Decoding]：Finally, the [Module 3] integrates $\mathbf{Z}^{(2)}$ with [Prior/Scale Cue] to output the final [Prediction/Action Trajectory]. | 句首锚定旧信息/逻辑视角，主干动词紧跟主语，复杂修饰置于宾语之后（符合尾重原则 L011 与旧新信息流 L001）。}
\end{itemize}
\end{TipBox}
\medskip
\subsection*{Paragraph 7: 时间序列专属：多尺度频域-时域耦合机制阐述段（Time Series Mechanism） (\texttt{BP14})}
\noindent\textbf{Expressive Intent:}\enspace 我这一段想专门阐述时间序列模型中“如何在频域捕获全局周期性，同时在时域保留局部突发异常（Transient Anomalies）”的双分支耦合机制。\par
\noindent\textbf{Cohesion Formula:}\enspace {\small\color{AcadSlate}Time-Frequency Dilemma $\to$ Spectral Global Branch $\to$ Temporal Local Branch $\to$ Uncertainty-Gated Fusion. 绑定逻辑规则：L023, L024.}\par\vspace{4pt}
\begin{ExampleBox}{Publication-Grade Authentic Text}
Modeling telemetry streams solely in the time domain struggles to span multi-day periodicities, whereas pure Fourier projection spreads localized transient spikes across the entire spectrum. To reconcile both scales, our dual-domain block processes input patches along two complementary pathways. The spectral pathway applies real-valued RFFT to isolate top-$k$ harmonic bases governing macro-periodicity, while the parallel temporal pathway employs dilated depthwise convolutions to retain sharp step changes and localized anomalies. An energy-ratio gate then fuses both representations at each time step, dynamically shifting weight toward the temporal pathway whenever local high-frequency residual energy spikes.
\end{ExampleBox}
\begin{TipBox}{Rhetorical Move Breakdown \& Constituent Analysis}
\begin{itemize}[leftmargin=1.5em]
  \item \textbf{S1 [单域局限性 Single-Domain Limitation]:} \textit{Modeling telemetry streams solely in the time domain struggles to span multi-day periodicities, whereas pure Fourier projection spreads localized transient spikes across the entire spectrum.}\par\noindent{\small 对应步骤：S1 [单域局限性 Single-Domain Limitation]：Modeling non-stationary time series solely in the time domain struggles to capture long-range periodicities, whereas pure Fourier projection blurs localized transient anomalies. | 句首锚定旧信息/逻辑视角，主干动词紧跟主语，复杂修饰置于宾语之后（符合尾重原则 L011 与旧新信息流 L001）。}
  \item \textbf{S2 [频域分支设计 Frequency Branch]:} \textit{To reconcile both scales, our dual-domain block processes input patches along two complementary pathways.}\par\noindent{\small 对应步骤：S2 [频域分支设计 Frequency Branch]：To capture global seasonal structure, our spectral branch applies real-valued FFT $\mathcal{F}(\cdot)$ along the temporal dimension and reweights dominant harmonic components. | 句首锚定旧信息/逻辑视角，主干动词紧跟主语，复杂修饰置于宾语之后（符合尾重原则 L011 与旧新信息流 L001）。}
  \item \textbf{S3 [时域分支设计 Temporal Branch]:} \textit{The spectral pathway applies real-valued RFFT to isolate top-$k$ harmonic bases governing macro-periodicity, while the parallel temporal pathway employs dilated depthwise convolutions to retain sharp step changes and localized anomalies.}\par\noindent{\small 对应步骤：S3 [时域分支设计 Temporal Branch]：In parallel, a multi-scale depthwise convolution branch operates directly on local patches to preserve sharp step changes and localized spikes. | 句首锚定旧信息/逻辑视角，主干动词紧跟主语，复杂修饰置于宾语之后（符合尾重原则 L011 与旧新信息流 L001）。}
  \item \textbf{S4 [自适应门控融合 Adaptive Fusion]:} \textit{An energy-ratio gate then fuses both representations at each time step, dynamically shifting weight toward the temporal pathway whenever local high-frequency residual energy spikes.}\par\noindent{\small 对应步骤：S4 [自适应门控融合 Adaptive Fusion]：An uncertainty-aware gate then fuses both branches step-by-step, assigning higher weight to the temporal branch whenever local residual variance exceeds a threshold. | 句首锚定旧信息/逻辑视角，主干动词紧跟主语，复杂修饰置于宾语之后（符合尾重原则 L011 与旧新信息流 L001）。}
\end{itemize}
\end{TipBox}
\medskip
\subsection*{Paragraph 8: 主实验结果整体对比与核心胜出原因剖析段（Main Results Analysis） (\texttt{BP18})}
\noindent\textbf{Expressive Intent:}\enspace 我这一段想解读主实验大表（Table 1）：既报告全局胜率，又挑出最有代表性的对比案例，解释为什么竞争对手在那里跌倒而我们稳赢。\par
\noindent\textbf{Cohesion Formula:}\enspace {\small\color{AcadSlate}Global Win Rate $\to$ Advantage over Family A $\to$ Resilience at Long Horizon vs. Family B $\to$ Mechanistic Takeaway. 绑定逻辑规则：L015, L013, L010.}\par\vspace{4pt}
\begin{ExampleBox}{Publication-Grade Authentic Text}
Table 1 reports the multivariate forecasting results, where DyTS-Net secures the lowest error in 25 out of 28 dataset-horizon configurations. Compared with channel-independent PatchTST, DyTS-Net reduces average MSE by 14.2\% on the 862-channel Traffic benchmark, where modeling spatial-temporal congestion propagation across adjacent road sensors is critical. Furthermore, while variate-mixing iTransformer performs competitively at $H=96$, its MSE escalates by 68\% at $H=720$ on volatile Weather series as unregularized cross-channel attention overfits transient wind gusts. By decoupling stationary phase coupling from neural-ODE scale projection, DyTS-Net suppresses long-horizon drift and maintains an 11.6\% margin over the strongest baseline at $H=720$.
\end{ExampleBox}
\begin{TipBox}{Rhetorical Move Breakdown \& Constituent Analysis}
\begin{itemize}[leftmargin=1.5em]
  \item \textbf{S1 [全局统治力概述 Overall Dominance]:} \textit{Table 1 reports the multivariate forecasting results, where DyTS-Net secures the lowest error in 25 out of 28 dataset-horizon configurations.}\par\noindent{\small 对应步骤：S1 [全局统治力概述 Overall Dominance]：Table 1 summarizes the multivariate forecasting results, where [Our Model] achieves the lowest error in [X] out of [Y] dataset-horizon configurations. | 句首锚定旧信息/逻辑视角，主干动词紧跟主语，复杂修饰置于宾语之后（符合尾重原则 L011 与旧新信息流 L001）。}
  \item \textbf{S2 [对比流派 1 的差距与原因 Contrast vs. Family 1]:} \textit{Compared with channel-independent PatchTST, DyTS-Net reduces average MSE by 14.2\% on the 862-channel Traffic benchmark, where modeling spatial-temporal congestion propagation across adjacent road sensors is critical.}\par\noindent{\small 对应步骤：S2 [对比流派 1 的差距与原因 Contrast vs. Family 1]：Compared with channel-independent baselines such as [Baseline A], our method reduces average MSE by [X\%] on [Dataset], where strong cross-sensor coupling is critical. | 句首锚定旧信息/逻辑视角，主干动词紧跟主语，复杂修饰置于宾语之后（符合尾重原则 L011 与旧新信息流 L001）。}
  \item \textbf{S3 [对比流派 2 在困难场景下的崩溃 Contrast vs. Family 2 under Hard Regime]:} \textit{Furthermore, while variate-mixing iTransformer performs competitively at $H=96$, its MSE escalates by 68\% at $H=720$ on volatile Weather series as unregularized cross-channel attention overfits transient wind gusts.}\par\noindent{\small 对应步骤：S3 [对比流派 2 在困难场景下的崩溃 Contrast vs. Family 2 under Hard Regime]：Furthermore, while [Baseline B] performs competitively at short horizons ($H=96$), its error escalates by [Y\%] at $H=720$ due to [Failure Mechanism]. | 句首锚定旧信息/逻辑视角，主干动词紧跟主语，复杂修饰置于宾语之后（符合尾重原则 L011 与旧新信息流 L001）。}
  \item \textbf{S4 [回扣本文核心机制 Mechanistic Confirmation]:} \textit{By decoupling stationary phase coupling from neural-ODE scale projection, DyTS-Net suppresses long-horizon drift and maintains an 11.6\% margin over the strongest baseline at $H=720$.}\par\noindent{\small 对应步骤：S4 [回扣本文核心机制 Mechanistic Confirmation]：By contrast, our [Core Module] bounds long-range drift, maintaining a [Z\%] margin over the second-best baseline at $H=720$. | 句首锚定旧信息/逻辑视角，主干动词紧跟主语，复杂修饰置于宾语之后（符合尾重原则 L011 与旧新信息流 L001）。}
\end{itemize}
\end{TipBox}
\medskip
\subsection*{Paragraph 9: 消融实验逐组件归因论证段（Counterfactual Ablation Study） (\texttt{BP19})}
\noindent\textbf{Expressive Intent:}\enspace 我这一段想汇报消融实验（Table 2）：证明我们的模块 A、模块 B 和正则项 C 每一个都不可或缺，且各自负责解决不同的特定子问题。\par
\noindent\textbf{Cohesion Formula:}\enspace {\small\color{AcadSlate}Ablation Overview $\to$ w/o Module A (Specific Failure) $\to$ w/o Module B (Specific Failure) $\to$ Cross-Module Synergy. 绑定逻辑规则：L016, L007.}\par\vspace{4pt}
\begin{ExampleBox}{Publication-Grade Authentic Text}
To isolate the contribution of each design choice, Table 2 compares full DyTS-Net against three ablated variants across ETTm2 and ECL. Removing the neural-ODE scale tracker (w/o Scale-ODE) triggers the sharpest degradation during seasonal transition months on ETTm2, raising 720-step MSE by 19.4\% as the decoder fails to track shifting baseline temperatures. Conversely, replacing phase-aligned graph attention with standard dot-product attention (w/o Phase-Attn) degrades performance most severely on the 321-channel ECL grid (+13.8\% MSE), where unaligned high-frequency load spikes corrupt cross-substation aggregation. Jointly deploying both modules yields complementary gains, as accurate scale tracking prevents amplitude distortion from biasing phase estimation.
\end{ExampleBox}
\begin{TipBox}{Rhetorical Move Breakdown \& Constituent Analysis}
\begin{itemize}[leftmargin=1.5em]
  \item \textbf{S1 [消融实验总括 Ablation Setup]:} \textit{To isolate the contribution of each design choice, Table 2 compares full DyTS-Net against three ablated variants across ETTm2 and ECL.}\par\noindent{\small 对应步骤：S1 [消融实验总括 Ablation Setup]：To isolate the contribution of each architectural component, we evaluate four ablated variants in Table 2 by progressively removing [Mod A], [Mod B], and [Loss C]. | 句首锚定旧信息/逻辑视角，主干动词紧跟主语，复杂修饰置于宾语之后（符合尾重原则 L011 与旧新信息流 L001）。}
  \item \textbf{S2 [移除模块 A 的冲击与归因 Impact of Removing Mod A]:} \textit{Removing the neural-ODE scale tracker (w/o Scale-ODE) triggers the sharpest degradation during seasonal transition months on ETTm2, raising 720-step MSE by 19.4\% as the decoder fails to track shifting baseline temperatures.}\par\noindent{\small 对应步骤：S2 [移除模块 A 的冲击与归因 Impact of Removing Mod A]：Removing [Mod A] (w/o A) causes the largest degradation on [Non-stationary/Volatile Subset], increasing MSE by [X\%] because the model loses [Specific Capability A]. | 句首锚定旧信息/逻辑视角，主干动词紧跟主语，复杂修饰置于宾语之后（符合尾重原则 L011 与旧新信息流 L001）。}
  \item \textbf{S3 [移除模块 B 的冲击与归因 Impact of Removing Mod B]:} \textit{Conversely, replacing phase-aligned graph attention with standard dot-product attention (w/o Phase-Attn) degrades performance most severely on the 321-channel ECL grid (+13.8\% MSE), where unaligned high-frequency load spikes corrupt cross-substation aggregation.}\par\noindent{\small 对应步骤：S3 [移除模块 B 的冲击与归因 Impact of Removing Mod B]：Similarly, replacing [Mod B] with standard [Naive Operator] degrades performance by [Y\%] on high-dimensional benchmarks, confirming its role in [Specific Capability B]. | 句首锚定旧信息/逻辑视角，主干动词紧跟主语，复杂修饰置于宾语之后（符合尾重原则 L011 与旧新信息流 L001）。}
  \item \textbf{S4 [协同效应总结 Synergistic Conclusion]:} \textit{Jointly deploying both modules yields complementary gains, as accurate scale tracking prevents amplitude distortion from biasing phase estimation.}\par\noindent{\small 对应步骤：S4 [协同效应总结 Synergistic Conclusion]：Combining all components yields super-additive gains, as [Mod A] stabilizes input statistics for [Mod B]'s relational reasoning. | 句首锚定旧信息/逻辑视角，主干动词紧跟主语，复杂修饰置于宾语之后（符合尾重原则 L011 与旧新信息流 L001）。}
\end{itemize}
\end{TipBox}
\medskip
\subsection*{Paragraph 10: 极端/反常情况与失败案例坦诚剖析段（Failure Case \& Outlier Analysis） (\texttt{BP21})}
\noindent\textbf{Expressive Intent:}\enspace 我这一段想通过可视化案例（Case Study / Qualitative Analysis），既展示模型在复杂跳变下的精准跟踪，又坦诚剖析一个模型预测失准的极端案例。\par
\noindent\textbf{Cohesion Formula:}\enspace {\small\color{AcadSlate}Visual Setup $\to$ Success under Shift $\to$ Honest Failure Case under Unobserved Exogenous Shock $\to$ Actionable Insight. 绑定逻辑规则：L017, L018.}\par\vspace{4pt}
\begin{ExampleBox}{Publication-Grade Authentic Text}
Figure 5 visualizes 336-step forecasts during an abrupt voltage drop on the ETTm1 transformer load benchmark. Whereas PatchTST and iTransformer exhibit a 14-step phase lag and severe post-recovery overshoot, DyTS-Net tracks both the sharp downward ramp and subsequent stabilization within 5\% amplitude error. Nevertheless, the failure case in Panel (d) reveals that when an unscheduled manual grid curtailment occurs without preceding spectral precursors, our endogenous model experiences a 6-step delayed response. Anticipating such zero-precursor manual interventions requires integrating exogenous dispatch logs alongside historical telemetry.
\end{ExampleBox}
\begin{TipBox}{Rhetorical Move Breakdown \& Constituent Analysis}
\begin{itemize}[leftmargin=1.5em]
  \item \textbf{S1 [定性对比切入 Qualitative Setup]:} \textit{Figure 5 visualizes 336-step forecasts during an abrupt voltage drop on the ETTm1 transformer load benchmark.}\par\noindent{\small 对应步骤：S1 [定性对比切入 Qualitative Setup]：Figure 5 visualizes predicted trajectories during an abrupt regime shift on [Dataset/Task]. | 句首锚定旧信息/逻辑视角，主干动词紧跟主语，复杂修饰置于宾语之后（符合尾重原则 L011 与旧新信息流 L001）。}
  \item \textbf{S2 [成功案例亮点 Success Highlight]:} \textit{Whereas PatchTST and iTransformer exhibit a 14-step phase lag and severe post-recovery overshoot, DyTS-Net tracks both the sharp downward ramp and subsequent stabilization within 5\% amplitude error.}\par\noindent{\small 对应步骤：S2 [成功案例亮点 Success Highlight]：While [Baseline] lags behind the sudden step change by [K steps] and overshoots, our model rapidly locks onto the new amplitude regime within [M steps]. | 句首锚定旧信息/逻辑视角，主干动词紧跟主语，复杂修饰置于宾语之后（符合尾重原则 L011 与旧新信息流 L001）。}
  \item \textbf{S3 [失败案例坦诚剖析 Honest Failure Mode]:} \textit{Nevertheless, the failure case in Panel (d) reveals that when an unscheduled manual grid curtailment occurs without preceding spectral precursors, our endogenous model experiences a 6-step delayed response.}\par\noindent{\small 对应步骤：S3 [失败案例坦诚剖析 Honest Failure Mode]：Nevertheless, inspecting the bottom-right failure case reveals that when an external manual override occurs without historical precursors, all purely endogenous forecasters---including ours---exhibit a transient lag. | 句首锚定旧信息/逻辑视角，主干动词紧跟主语，复杂修饰置于宾语之后（符合尾重原则 L011 与旧新信息流 L001）。}
  \item \textbf{S4 [改进启示 Constructive Takeaway]:} \textit{Anticipating such zero-precursor manual interventions requires integrating exogenous dispatch logs alongside historical telemetry.}\par\noindent{\small 对应步骤：S4 [改进启示 Constructive Takeaway]：Incorporating exogenous event logs or textual maintenance records is necessary to anticipate such unscheduled interventions. | 句首锚定旧信息/逻辑视角，主干动词紧跟主语，复杂修饰置于宾语之后（符合尾重原则 L011 与旧新信息流 L001）。}
\end{itemize}
\end{TipBox}
\medskip
\subsection*{Paragraph 11: 全文总结与学术影响收尾段（Conclusion \& Broader Outlook） (\texttt{BP24})}
\noindent\textbf{Expressive Intent:}\enspace 我这一段想写整篇小论文的 Conclusion：不简单复制粘贴 Abstract，而是从更高维度总结本文解决的科学矛盾、方法精髓、实证高度与领域启示。\par
\noindent\textbf{Cohesion Formula:}\enspace {\small\color{AcadSlate}Tension Resolved $\to$ Methodological Insight $\to$ Empirical Confirmation $\to$ Generalizable Blueprint. 绑定逻辑规则：L013, L023, L029.}\par\vspace{4pt}
\begin{ExampleBox}{Publication-Grade Authentic Text}
In this work, we resolved the long-standing tension between shift-invariant normalization and physical amplitude preservation in non-stationary time series forecasting. By decoupling frequency-domain phase attention from continuous neural-ODE scale tracking, DyTS-Net captures long-range seasonal dependencies while accurately reconstructing extreme peak magnitudes. Empirical evaluations across eight benchmarks and industrial wind-power telemetry demonstrated MSE reductions of up to 17.4\% alongside a 45\% reduction in computational overhead. Beyond forecasting, this structural-scale disentanglement principle provides a generalizable foundation for anomaly detection and closed-loop predictive control in non-stationary cyber-physical systems.
\end{ExampleBox}
\begin{TipBox}{Rhetorical Move Breakdown \& Constituent Analysis}
\begin{itemize}[leftmargin=1.5em]
  \item \textbf{S1 [核心科学问题闭合 Scientific Closure]:} \textit{In this work, we resolved the long-standing tension between shift-invariant normalization and physical amplitude preservation in non-stationary time series forecasting.}\par\noindent{\small 对应步骤：S1 [核心科学问题闭合 Scientific Closure]：In this paper, we addressed the fundamental tension between [Goal A] and [Goal B] in [Domain]. | 句首锚定旧信息/逻辑视角，主干动词紧跟主语，复杂修饰置于宾语之后（符合尾重原则 L011 与旧新信息流 L001）。}
  \item \textbf{S2 [方法论精髓回顾 Methodological Synthesis]:} \textit{By decoupling frequency-domain phase attention from continuous neural-ODE scale tracking, DyTS-Net captures long-range seasonal dependencies while accurately reconstructing extreme peak magnitudes.}\par\noindent{\small 对应步骤：S2 [方法论精髓回顾 Methodological Synthesis]：By introducing [Model/Mechanism], we demonstrated that [Core Insight] enables [Capability] without incurring [Traditional Penalty]. | 句首锚定旧信息/逻辑视角，主干动词紧跟主语，复杂修饰置于宾语之后（符合尾重原则 L011 与旧新信息流 L001）。}
  \item \textbf{S3 [实证结果高度概括 Empirical Recap]:} \textit{Empirical evaluations across eight benchmarks and industrial wind-power telemetry demonstrated MSE reductions of up to 17.4\% alongside a 45\% reduction in computational overhead.}\par\noindent{\small 对应步骤：S3 [实证结果高度概括 Empirical Recap]：Comprehensive evaluations across [Benchmarks/Hardware] confirmed consistent error reductions of up to [X\%] alongside substantial inference speedups. | 句首锚定旧信息/逻辑视角，主干动词紧跟主语，复杂修饰置于宾语之后（符合尾重原则 L011 与旧新信息流 L001）。}
  \item \textbf{S4 [跨领域迁移展望 Broader Outlook]:} \textit{Beyond forecasting, this structural-scale disentanglement principle provides a generalizable foundation for anomaly detection and closed-loop predictive control in non-stationary cyber-physical systems.}\par\noindent{\small 对应步骤：S4 [跨领域迁移展望 Broader Outlook]：Beyond [Current Task], our [Core Principle] offers a generalizable blueprint for broader sequential decision-making tasks under non-stationary environments. | 句首锚定旧信息/逻辑视角，主干动词紧跟主语，复杂修饰置于宾语之后（符合尾重原则 L011 与旧新信息流 L001）。}
\end{itemize}
\end{TipBox}
\medskip
\clearpage
\section{FT04: 具身智能视觉-力触觉闭环操作与 Sim-to-Real 自适应完整小论文（TacFlow: Contact-Aware Visuomotor Policy）}\label{comp:FT04}
\begin{ConceptBox}{Composition Architecture \& Structural Blueprint Chain}
\textbf{Genre:} \texttt{AI\_Mini\_Paper} \quad \textbf{Paragraph Count:} 11\par
\textbf{Blueprint Chain:} \texttt{BP08 $\to$ BP01 $\to$ BP03 $\to$ BP05 $\to$ BP07 $\to$ BP11 $\to$ BP16 $\to$ BP18 $\to$ BP19 $\to$ BP22 $\to$ BP24}
\end{ConceptBox}
\subsection*{Paragraph 1: 高密度四句式 Abstract 全段构建（Complete 4-Move Abstract） (\texttt{BP08})}
\noindent\textbf{Expressive Intent:}\enspace 我这一段想写出整篇小论文的完整 Abstract（摘要）：在 150--180 词内一气呵成交代背景、核心痛点、方法机制与量化实验结果。\par
\noindent\textbf{Cohesion Formula:}\enspace {\small\color{AcadSlate}Context (S1) $\to$ Gap (S2) $\to$ Method \& Mechanism (S3-S4) $\to$ Quantitative Proof (S5). 绑定逻辑规则：L001, L008, L010, L029.}\par\vspace{4pt}
\begin{ExampleBox}{Publication-Grade Authentic Text}
Recovering missing sensor observations is essential for reliable downstream monitoring in industrial and clinical time series. While diffusion-based imputation models generate plausible trajectories under random missingness, they suffer from high sampling latency and lose global periodic coherence when entire sensor blocks drop out consecutively. We present FlowImp, a conditional flow-matching framework that reconstructs contiguous missing blocks in three ODE solver steps. Specifically, FlowImp anchors straight transport paths on observed spectral priors and enforces cross-channel phase consistency via a lightweight Fourier regularizer. Across six multivariate benchmarks under 50\% block missingness, FlowImp reduces imputation MAE by 19.3\% compared with CSDI while accelerating inference by 24x.
\end{ExampleBox}
\begin{TipBox}{Rhetorical Move Breakdown \& Constituent Analysis}
\begin{itemize}[leftmargin=1.5em]
  \item \textbf{S1 [背景与任务重要性 Context]:} \textit{Recovering missing sensor observations is essential for reliable downstream monitoring in industrial and clinical time series.}\par\noindent{\small 对应步骤：S1 [背景与任务重要性 Context]：1句话锁定任务目标与核心应用场景（≤20词）。 | 句首锚定旧信息/逻辑视角，主干动词紧跟主语，复杂修饰置于宾语之后（符合尾重原则 L011 与旧新信息流 L001）。}
  \item \textbf{S2 [现有局限与痛点 Gap/Bottleneck]:} \textit{While diffusion-based imputation models generate plausible trajectories under random missingness, they suffer from high sampling latency and lose global periodic coherence when entire sensor blocks drop out consecutively.}\par\noindent{\small 对应步骤：S2 [现有局限与痛点 Gap/Bottleneck]：While existing [Paradigms] achieve [Merit], they struggle with [Core Bottleneck] because [Underlying Cause]. | 句首锚定旧信息/逻辑视角，主干动词紧跟主语，复杂修饰置于宾语之后（符合尾重原则 L011 与旧新信息流 L001）。}
  \item \textbf{S3-S4 [本文方法与核心机制 Approach \& Mechanism]:} \textit{We present FlowImp, a conditional flow-matching framework that reconstructs contiguous missing blocks in three ODE solver steps.}\par\noindent{\small 对应步骤：S3-S4 [本文方法与核心机制 Approach \& Mechanism]：In this work, we present [Model], which [Core Idea]. Specifically, [Module 1]..., while [Module 2]... | 句首锚定旧信息/逻辑视角，主干动词紧跟主语，复杂修饰置于宾语之后（符合尾重原则 L011 与旧新信息流 L001）。}
  \item \textbf{S5 [量化实证结果 Quantitative Results]:} \textit{Specifically, FlowImp anchors straight transport paths on observed spectral priors and enforces cross-channel phase consistency via a lightweight Fourier regularizer.}\par\noindent{\small 对应步骤：S5 [量化实证结果 Quantitative Results]：Experiments on [Benchmarks/Hardware] show that [Model] outperforms [Baselines] by [X\%] in [Metric] and achieves [Efficiency/Real-world Result]. | 句首锚定旧信息/逻辑视角，主干动词紧跟主语，复杂修饰置于宾语之后（符合尾重原则 L011 与旧新信息流 L001）。}
  \item \textbf{S5 [量化实证结果 Quantitative Results]:} \textit{Across six multivariate benchmarks under 50\% block missingness, FlowImp reduces imputation MAE by 19.3\% compared with CSDI while accelerating inference by 24x.}\par\noindent{\small 对应步骤：S5 [量化实证结果 Quantitative Results]：Experiments on [Benchmarks/Hardware] show that [Model] outperforms [Baselines] by [X\%] in [Metric] and achieves [Efficiency/Real-world Result]. | 句首锚定旧信息/逻辑视角，主干动词紧跟主语，复杂修饰置于宾语之后（符合尾重原则 L011 与旧新信息流 L001）。}
\end{itemize}
\end{TipBox}
\medskip
\subsection*{Paragraph 2: 确立研究背景与真实工程价值（拒绝空喊口号） (\texttt{BP01})}
\noindent\textbf{Expressive Intent:}\enspace 我这一段想作为 Introduction 第一段：说明多变量时间序列预测或具身智能控制在实际工业系统中的关键作用，以及高可靠性的刚需。\par
\noindent\textbf{Cohesion Formula:}\enspace {\small\color{AcadSlate}Given-to-New 链：Core Task (S1) $\to$ Practical Deployment (S2) $\to$ Real-world Disturbances (S3) $\to$ Consequence of Failure (S4). 绑定逻辑规则：L001, L011, L029.}\par\vspace{4pt}
\begin{ExampleBox}{Publication-Grade Authentic Text}
Multivariate time series forecasting underpins safety-critical scheduling across smart power grids, cloud clusters, and automated manufacturing. In these operational settings, control systems rely on multi-horizon predictions to balance supply fluctuations and prevent cascading overload. However, field sensors routinely operate under abrupt regime shifts, irregular sampling intervals, and high-frequency measurement noise. Failing to accommodate such non-stationary disturbances degrades long-range forecast reliability and triggers costly control oscillations.
\end{ExampleBox}
\begin{TipBox}{Rhetorical Move Breakdown \& Constituent Analysis}
\begin{itemize}[leftmargin=1.5em]
  \item \textbf{S1 [任务锚定 Topic Anchor]:} \textit{Multivariate time series forecasting underpins safety-critical scheduling across smart power grids, cloud clusters, and automated manufacturing.}\par\noindent{\small 对应步骤：S1 [任务锚定 Topic Anchor]：直接用紧凑主谓宾点明研究对象及其在真实系统中的核心功能（≤16词，禁用 In the era of AI）。 | 句首锚定旧信息/逻辑视角，主干动词紧跟主语，复杂修饰置于宾语之后（符合尾重原则 L011 与旧新信息流 L001）。}
  \item \textbf{S2 [应用场景展开 Concrete Scenarios]:} \textit{In these operational settings, control systems rely on multi-horizon predictions to balance supply fluctuations and prevent cascading overload.}\par\noindent{\small 对应步骤：S2 [应用场景展开 Concrete Scenarios]：列举 2--3 个具体工程场景（如电网负荷调度、灵巧机械手装配）及其决策依赖。 | 句首锚定旧信息/逻辑视角，主干动词紧跟主语，复杂修饰置于宾语之后（符合尾重原则 L011 与旧新信息流 L001）。}
  \item \textbf{S3 [苛刻工况约束 Operational Bottleneck]:} \textit{However, field sensors routinely operate under abrupt regime shifts, irregular sampling intervals, and high-frequency measurement noise.}\par\noindent{\small 对应步骤：S3 [苛刻工况约束 Operational Bottleneck]：指出真实部署环境下不可避免的物理/统计挑战（如非平稳漂移、传感器缺失、毫秒级延迟约束）。 | 句首锚定旧信息/逻辑视角，主干动词紧跟主语，复杂修饰置于宾语之后（符合尾重原则 L011 与旧新信息流 L001）。}
  \item \textbf{S4 [研究必要性收束 Imperative Hook]:} \textit{Failing to accommodate such non-stationary disturbances degrades long-range forecast reliability and triggers costly control oscillations.}\par\noindent{\small 对应步骤：S4 [研究必要性收束 Imperative Hook]：点明若无法克服上述约束会造成什么具体系统后果，自然引出下一段对现有方法的审视。 | 句首锚定旧信息/逻辑视角，主干动词紧跟主语，复杂修饰置于宾语之后（符合尾重原则 L011 与旧新信息流 L001）。}
\end{itemize}
\end{TipBox}
\medskip
\subsection*{Paragraph 3: 尖锐指出核心技术瓶颈与失效机理（Gap Identification） (\texttt{BP03})}
\noindent\textbf{Expressive Intent:}\enspace 我这一段想深入剖析现有方法在特定难题（如非平稳突变、长程自回归误差累积、Sim-to-Real 接触失配）上为什么会失效，把机理讲透。\par
\noindent\textbf{Cohesion Formula:}\enspace {\small\color{AcadSlate}因果链式推进（A induces B $\to$ This B compounds into C $\to$ Naive fix D sacrifices E）。绑定逻辑规则：L003, L004, L024/L025.}\par\vspace{4pt}
\begin{ExampleBox}{Publication-Grade Authentic Text}
Bridging non-stationary robustness and amplitude fidelity remains a core bottleneck in long-horizon forecasting. While instance normalization strips time-varying mean and variance shifts to stabilize attention weights, it simultaneously erases physical scale cues required to predict extreme load peaks. Conversely, retaining raw unnormalized inputs exposes self-attention layers to severe covariate shift, causing validation error to diverge over extended horizons. Resolving this tension requires a mechanism that disentangles intrinsic temporal dynamics from transient scale drift without discarding magnitude information.
\end{ExampleBox}
\begin{TipBox}{Rhetorical Move Breakdown \& Constituent Analysis}
\begin{itemize}[leftmargin=1.5em]
  \item \textbf{S1 [核心痛点立论 Bottleneck Claim]:} \textit{Bridging non-stationary robustness and amplitude fidelity remains a core bottleneck in long-horizon forecasting.}\par\noindent{\small 对应步骤：S1 [核心痛点立论 Bottleneck Claim]：一句话直接命名该核心瓶颈（如 Compound error accumulation / Sim-to-real contact mismatch）。 | 句首锚定旧信息/逻辑视角，主干动词紧跟主语，复杂修饰置于宾语之后（符合尾重原则 L011 与旧新信息流 L001）。}
  \item \textbf{S2 [第一层失效机理 Primary Failure Mechanism]:} \textit{While instance normalization strips time-varying mean and variance shifts to stabilize attention weights, it simultaneously erases physical scale cues required to predict extreme load peaks.}\par\noindent{\small 对应步骤：S2 [第一层失效机理 Primary Failure Mechanism]：从数学或物理传导链条解释误差是如何产生并放大的。 | 句首锚定旧信息/逻辑视角，主干动词紧跟主语，复杂修饰置于宾语之后（符合尾重原则 L011 与旧新信息流 L001）。}
  \item \textbf{S3 [第二层耦合困境 Secondary Dilemma]:} \textit{Conversely, retaining raw unnormalized inputs exposes self-attention layers to severe covariate shift, causing validation error to diverge over extended horizons.}\par\noindent{\small 对应步骤：S3 [第二层耦合困境 Secondary Dilemma]：指出为什么简单的修补手段（如单纯加大模型、简单归一化或域随机化）不仅无效还会带来副作用。 | 句首锚定旧信息/逻辑视角，主干动词紧跟主语，复杂修饰置于宾语之后（符合尾重原则 L011 与旧新信息流 L001）。}
  \item \textbf{S4 [凝练核心研究问题 Research Question Formulation]:} \textit{Resolving this tension requires a mechanism that disentangles intrinsic temporal dynamics from transient scale drift without discarding magnitude information.}\par\noindent{\small 对应步骤：S4 [凝练核心研究问题 Research Question Formulation]：用 Therefore, a fundamental question arises: how can we... 引出本文目标。 | 句首锚定旧信息/逻辑视角，主干动词紧跟主语，复杂修饰置于宾语之后（符合尾重原则 L011 与旧新信息流 L001）。}
\end{itemize}
\end{TipBox}
\medskip
\subsection*{Paragraph 4: 挑战与机制一对一映射段（Challenge-to-Mechanism Mapping） (\texttt{BP05})}
\noindent\textbf{Expressive Intent:}\enspace 我这一段想说明：把我们的核心思想落地时遇到了两个具体工程/数学难点，我们分别设计了什么精巧机制去攻克它们。\par
\noindent\textbf{Cohesion Formula:}\enspace {\small\color{AcadSlate}Challenge 1 $\to$ Mechanism 1; Challenge 2 $\to$ Mechanism 2 $\to$ Joint System Property. 绑定逻辑规则：L003, L011.}\par\vspace{4pt}
\begin{ExampleBox}{Publication-Grade Authentic Text}
Instantiating closed-loop visuomotor diffusion on physical robots entails two technical challenges. First, because raw depth camera streams contain occlusion artifacts around reflective object edges, we introduce a tactile-guided cross-attention gate that reweights visual tokens upon contact onset. Second, to eliminate the 200 ms sampling bottleneck of iterative denoising, we formulate a trajectory consistency distillation objective that maps Gaussian noise directly to collision-free action chunks in two forward steps. Together, these designs sustain 50 Hz reactive control under partial visual occlusion.
\end{ExampleBox}
\begin{TipBox}{Rhetorical Move Breakdown \& Constituent Analysis}
\begin{itemize}[leftmargin=1.5em]
  \item \textbf{S1 [落地难点总起 Implementation Hurdle]:} \textit{Instantiating closed-loop visuomotor diffusion on physical robots entails two technical challenges.}\par\noindent{\small 对应步骤：S1 [落地难点总起 Implementation Hurdle]：Translating this principle into a scalable architecture entails two technical challenges. | 句首锚定旧信息/逻辑视角，主干动词紧跟主语，复杂修饰置于宾语之后（符合尾重原则 L011 与旧新信息流 L001）。}
  \item \textbf{S2 [难点 1 + 靶向设计 1]:} \textit{First, because raw depth camera streams contain occlusion artifacts around reflective object edges, we introduce a tactile-guided cross-attention gate that reweights visual tokens upon contact onset.}\par\noindent{\small 对应步骤：S2 [难点 1 + 靶向设计 1]：First, because [Challenge 1], we design [Mechanism 1] to [Specific Function]. | 句首锚定旧信息/逻辑视角，主干动词紧跟主语，复杂修饰置于宾语之后（符合尾重原则 L011 与旧新信息流 L001）。}
  \item \textbf{S3 [难点 2 + 靶向设计 2]:} \textit{Second, to eliminate the 200 ms sampling bottleneck of iterative denoising, we formulate a trajectory consistency distillation objective that maps Gaussian noise directly to collision-free action chunks in two forward steps.}\par\noindent{\small 对应步骤：S3 [难点 2 + 靶向设计 2]：Second, to prevent [Challenge 2] under [Constraint], we introduce [Mechanism 2] that [Specific Function]. | 句首锚定旧信息/逻辑视角，主干动词紧跟主语，复杂修饰置于宾语之后（符合尾重原则 L011 与旧新信息流 L001）。}
  \item \textbf{S4 [整体计算/泛化特性 System Property]:} \textit{Together, these designs sustain 50 Hz reactive control under partial visual occlusion.}\par\noindent{\small 对应步骤：S4 [整体计算/泛化特性 System Property]：Together, these components enable [Model] to operate at [Efficiency/Robustness Metric]. | 句首锚定旧信息/逻辑视角，主干动词紧跟主语，复杂修饰置于宾语之后（符合尾重原则 L011 与旧新信息流 L001）。}
\end{itemize}
\end{TipBox}
\medskip
\subsection*{Paragraph 5: 凝练论文三点核心贡献（Contributions Summary） (\texttt{BP07})}
\noindent\textbf{Expressive Intent:}\enspace 我这一段想在 Introduction 最末尾清晰列出本文的三点核心贡献（理论/视角贡献 + 核心算法/架构贡献 + 实验与开源实证贡献）。\par
\noindent\textbf{Cohesion Formula:}\enspace {\small\color{AcadSlate}Parallel Grammatical Structure（严格平行结构：全部用 We + 强及物动词 identify / propose / demonstrate 开头）。绑定逻辑规则：L002, L010.}\par\vspace{4pt}
\begin{ExampleBox}{Publication-Grade Authentic Text}
In summary, our main contributions are threefold: (1) We formalize the amplitude-stationarity trade-off in long-horizon forecasting and prove that standard instance normalization bounds peak-load recoverability. (2) We propose DyTS-Net, a hybrid architecture that couples shift-invariant spectral attention with a continuous neural ODE scale tracker. (3) We demonstrate across eight public benchmarks and industrial wind-power telemetry that DyTS-Net lowers long-range MSE by up to 17.4\% while cutting inference latency nearly in half.
\end{ExampleBox}
\begin{TipBox}{Rhetorical Move Breakdown \& Constituent Analysis}
\begin{itemize}[leftmargin=1.5em]
  \item \textbf{S1 [总起引导句 Lead-in]:} \textit{In summary, our main contributions are threefold: (1) We formalize the amplitude-stationarity trade-off in long-horizon forecasting and prove that standard instance normalization bounds peak-load recoverability.}\par\noindent{\small 对应步骤：S1 [总起引导句 Lead-in]：In summary, our main contributions are threefold: | 句首锚定旧信息/逻辑视角，主干动词紧跟主语，复杂修饰置于宾语之后（符合尾重原则 L011 与旧新信息流 L001）。}
  \item \textbf{S2 [贡献 1:} \textit{(2) We propose DyTS-Net, a hybrid architecture that couples shift-invariant spectral attention with a continuous neural ODE scale tracker.}\par\noindent{\small 对应步骤：S2 [贡献 1：新视角/问题形式化 Conceptual/Formulation Contribution]：用强动词名物化或动词原形开头（We identify and formalize... / We establish...），点明打破了什么传统局限。 | 句首锚定旧信息/逻辑视角，主干动词紧跟主语，复杂修饰置于宾语之后（符合尾重原则 L011 与旧新信息流 L001）。}
  \item \textbf{S3 [贡献 2:} \textit{(3) We demonstrate across eight public benchmarks and industrial wind-power telemetry that DyTS-Net lowers long-range MSE by up to 17.4\% while cutting inference latency nearly in half.}\par\noindent{\small 对应步骤：S3 [贡献 2：核心方法与模块 Algorithmic/Architectural Contribution]：We propose [Model/Mechanism], which integrates [A] and [B] to achieve [Theoretical/Functional Guarantee]. | 句首锚定旧信息/逻辑视角，主干动词紧跟主语，复杂修饰置于宾语之后（符合尾重原则 L011 与旧新信息流 L001）。}
\end{itemize}
\end{TipBox}
\medskip
\subsection*{Paragraph 6: 问题形式化定义与符号约定段（Problem Formulation \& Notation） (\texttt{BP11})}
\noindent\textbf{Expressive Intent:}\enspace 我这一段想在 Methodology 开篇（Section 3.1 Problem Formulation），严谨、简洁地定义输入张量、输出目标、约束条件与优化目标。\par
\noindent\textbf{Cohesion Formula:}\enspace {\small\color{AcadSlate}Input Space $\to$ Auxiliary Conditions $\to$ Target Mapping $\to$ Crucial Distributional Constraint. 绑定逻辑规则：L014, L011.}\par\vspace{4pt}
\begin{ExampleBox}{Publication-Grade Authentic Text}
Consider a multivariate time series $\mathbf{X}_{1:L} = [x_1, \dots, x_L]^\top \in \mathbb{R}^{L \times C}$ recording $C$ correlated sensors over lookback horizon $L$. Given an binary observation mask $\mathbf{M} \in \{0, 1\}^{L \times C}$ indicating irregular sensor dropouts, our objective is to learn a generative mapping $p_\theta(\mathbf{Y}_{L+1:L+H} \mid \mathbf{X}_{1:L} \odot \mathbf{M})$ over future horizon $H$. Crucially, unlike standard settings that assume time-invariant marginal distributions, we explicitly account for non-stationary regime shifts where $\mathbb{E}[x_t]$ and $\text{Cov}(x_t)$ evolve continuously between the lookback and forecasting windows.
\end{ExampleBox}
\begin{TipBox}{Rhetorical Move Breakdown \& Constituent Analysis}
\begin{itemize}[leftmargin=1.5em]
  \item \textbf{S1 [任务总体形式化 Task Definition]:} \textit{Consider a multivariate time series $\mathbf{X}_{1:L} = [x_1, \dots, x_L]^\top \in \mathbb{R}^{L \times C}$ recording $C$ correlated sensors over lookback horizon $L$.}\par\noindent{\small 对应步骤：S1 [任务总体形式化 Task Definition]：Consider a [System/Dataset] observed over a lookback window of length $L$, denoted as $\mathbf{X} = [x_1, \dots, x_L] \in \mathbb{R}^{L \times C}$, where $C$ represents [Meaning]. | 句首锚定旧信息/逻辑视角，主干动词紧跟主语，复杂修饰置于宾语之后（符合尾重原则 L011 与旧新信息流 L001）。}
  \item \textbf{S2 [补充条件或模态输入 Auxiliary/Condition Inputs]:} \textit{Given an binary observation mask $\mathbf{M} \in \{0, 1\}^{L \times C}$ indicating irregular sensor dropouts, our objective is to learn a generative mapping $p_\theta(\mathbf{Y}_{L+1:L+H} \mid \mathbf{X}_{1:L} \odot \mathbf{M})$ over future horizon $H$.}\par\noindent{\small 对应步骤：S2 [补充条件或模态输入 Auxiliary/Condition Inputs]：At each step $t$, the system also receives [Condition/Observation] $z_t \in \mathcal{Z}$ under [Noise/Missingness/Dynamic Regime]. | 句首锚定旧信息/逻辑视角，主干动词紧跟主语，复杂修饰置于宾语之后（符合尾重原则 L011 与旧新信息流 L001）。}
  \item \textbf{S3 [预测/控制目标 Objective]:} \textit{Crucially, unlike standard settings that assume time-invariant marginal distributions, we explicitly account for non-stationary regime shifts where $\mathbb{E}[x_t]$ and $\text{Cov}(x_t)$ evolve continuously between the lookback and forecasting windows.}\par\noindent{\small 对应步骤：S3 [预测/控制目标 Objective]：Our goal is to learn a parameterized mapping $f_\theta(\cdot)$ that predicts future states $\hat{\mathbf{Y}} \in \mathbb{R}^{H \times C}$ over horizon $H$ (or action trajectory $a_{t:t+K}$). | 句首锚定旧信息/逻辑视角，主干动词紧跟主语，复杂修饰置于宾语之后（符合尾重原则 L011 与旧新信息流 L001）。}
\end{itemize}
\end{TipBox}
\medskip
\subsection*{Paragraph 7: 具身智能专属：多模态感知-动作闭环与 Sim-to-Real 阐述段（Embodied AI Policy） (\texttt{BP16})}
\noindent\textbf{Expressive Intent:}\enspace 我这一段想阐述具身智能机器人策略（Embodied Policy）如何融合视觉、本体感觉（Proprioception）或触觉，并克服仿真到真实硬件的动力学差距（Sim-to-Real Gap）。\par
\noindent\textbf{Cohesion Formula:}\enspace {\small\color{AcadSlate}Visual Occlusion Bottleneck $\to$ Asynchronous Multimodal Alignment $\to$ Action Chunking $\to$ Online Sim-to-Real System Identification. 绑定逻辑规则：L025, L003.}\par\vspace{4pt}
\begin{ExampleBox}{Publication-Grade Authentic Text}
Relying solely on third-person RGB-D observations leaves visuomotor policies vulnerable to severe self-occlusion once the gripper approaches the target object. To preserve observability during contact-rich manipulation, our encoder aligns 30 Hz visual spatial tokens with 500 Hz joint proprioception and wrist force-torque histories via cross-modal rotary embeddings. Rather than issuing single-step commands prone to high-frequency jitter, the policy predicts 16-step SE(3) end-effector action chunks smoothed through exponential temporal ensembling. At deployment, a recurrent history adapter estimates latent surface friction and payload mass online, bridging the sim-to-real dynamics gap in zero shot.
\end{ExampleBox}
\begin{TipBox}{Rhetorical Move Breakdown \& Constituent Analysis}
\begin{itemize}[leftmargin=1.5em]
  \item \textbf{S1 [单模态开环局限 Open-Loop Visual Limitation]:} \textit{Relying solely on third-person RGB-D observations leaves visuomotor policies vulnerable to severe self-occlusion once the gripper approaches the target object.}\par\noindent{\small 对应步骤：S1 [单模态开环局限 Open-Loop Visual Limitation]：Relying solely on RGB-D point clouds leaves manipulation policies vulnerable to self-occlusion by the robot gripper during fine-grained contact. | 句首锚定旧信息/逻辑视角，主干动词紧跟主语，复杂修饰置于宾语之后（符合尾重原则 L011 与旧新信息流 L001）。}
  \item \textbf{S2 [多模态时空对齐 Multi-Modal Alignment]:} \textit{To preserve observability during contact-rich manipulation, our encoder aligns 30 Hz visual spatial tokens with 500 Hz joint proprioception and wrist force-torque histories via cross-modal rotary embeddings.}\par\noindent{\small 对应步骤：S2 [多模态时空对齐 Multi-Modal Alignment]：To maintain state observability during contact, our encoder aligns asynchronous 30 Hz visual tokens with 500 Hz joint proprioception and wrist force-torque histories. | 句首锚定旧信息/逻辑视角，主干动词紧跟主语，复杂修饰置于宾语之后（符合尾重原则 L011 与旧新信息流 L001）。}
  \item \textbf{S3 [动作块预测与平滑执行 Action Chunking \& Temporal Ensembling]:} \textit{Rather than issuing single-step commands prone to high-frequency jitter, the policy predicts 16-step SE(3) end-effector action chunks smoothed through exponential temporal ensembling.}\par\noindent{\small 对应步骤：S3 [动作块预测与平滑执行 Action Chunking \& Temporal Ensembling]：Rather than predicting single-step torques prone to jitter, the policy outputs overlapping $K$-step SE(3) end-effector action chunks smoothed via exponential temporal ensembling. | 句首锚定旧信息/逻辑视角，主干动词紧跟主语，复杂修饰置于宾语之后（符合尾重原则 L011 与旧新信息流 L001）。}
  \item \textbf{S4 [Sim-to-Real 自适应机制 Sim-to-Real Adaptation]:} \textit{At deployment, a recurrent history adapter estimates latent surface friction and payload mass online, bridging the sim-to-real dynamics gap in zero shot.}\par\noindent{\small 对应步骤：S4 [Sim-to-Real 自适应机制 Sim-to-Real Adaptation]：During physical deployment, a lightweight history adapter infers latent friction and payload parameters online, bridging the sim-to-real dynamics gap without policy fine-tuning. | 句首锚定旧信息/逻辑视角，主干动词紧跟主语，复杂修饰置于宾语之后（符合尾重原则 L011 与旧新信息流 L001）。}
\end{itemize}
\end{TipBox}
\medskip
\subsection*{Paragraph 8: 主实验结果整体对比与核心胜出原因剖析段（Main Results Analysis） (\texttt{BP18})}
\noindent\textbf{Expressive Intent:}\enspace 我这一段想解读主实验大表（Table 1）：既报告全局胜率，又挑出最有代表性的对比案例，解释为什么竞争对手在那里跌倒而我们稳赢。\par
\noindent\textbf{Cohesion Formula:}\enspace {\small\color{AcadSlate}Global Win Rate $\to$ Advantage over Family A $\to$ Resilience at Long Horizon vs. Family B $\to$ Mechanistic Takeaway. 绑定逻辑规则：L015, L013, L010.}\par\vspace{4pt}
\begin{ExampleBox}{Publication-Grade Authentic Text}
Table 1 reports the multivariate forecasting results, where DyTS-Net secures the lowest error in 25 out of 28 dataset-horizon configurations. Compared with channel-independent PatchTST, DyTS-Net reduces average MSE by 14.2\% on the 862-channel Traffic benchmark, where modeling spatial-temporal congestion propagation across adjacent road sensors is critical. Furthermore, while variate-mixing iTransformer performs competitively at $H=96$, its MSE escalates by 68\% at $H=720$ on volatile Weather series as unregularized cross-channel attention overfits transient wind gusts. By decoupling stationary phase coupling from neural-ODE scale projection, DyTS-Net suppresses long-horizon drift and maintains an 11.6\% margin over the strongest baseline at $H=720$.
\end{ExampleBox}
\begin{TipBox}{Rhetorical Move Breakdown \& Constituent Analysis}
\begin{itemize}[leftmargin=1.5em]
  \item \textbf{S1 [全局统治力概述 Overall Dominance]:} \textit{Table 1 reports the multivariate forecasting results, where DyTS-Net secures the lowest error in 25 out of 28 dataset-horizon configurations.}\par\noindent{\small 对应步骤：S1 [全局统治力概述 Overall Dominance]：Table 1 summarizes the multivariate forecasting results, where [Our Model] achieves the lowest error in [X] out of [Y] dataset-horizon configurations. | 句首锚定旧信息/逻辑视角，主干动词紧跟主语，复杂修饰置于宾语之后（符合尾重原则 L011 与旧新信息流 L001）。}
  \item \textbf{S2 [对比流派 1 的差距与原因 Contrast vs. Family 1]:} \textit{Compared with channel-independent PatchTST, DyTS-Net reduces average MSE by 14.2\% on the 862-channel Traffic benchmark, where modeling spatial-temporal congestion propagation across adjacent road sensors is critical.}\par\noindent{\small 对应步骤：S2 [对比流派 1 的差距与原因 Contrast vs. Family 1]：Compared with channel-independent baselines such as [Baseline A], our method reduces average MSE by [X\%] on [Dataset], where strong cross-sensor coupling is critical. | 句首锚定旧信息/逻辑视角，主干动词紧跟主语，复杂修饰置于宾语之后（符合尾重原则 L011 与旧新信息流 L001）。}
  \item \textbf{S3 [对比流派 2 在困难场景下的崩溃 Contrast vs. Family 2 under Hard Regime]:} \textit{Furthermore, while variate-mixing iTransformer performs competitively at $H=96$, its MSE escalates by 68\% at $H=720$ on volatile Weather series as unregularized cross-channel attention overfits transient wind gusts.}\par\noindent{\small 对应步骤：S3 [对比流派 2 在困难场景下的崩溃 Contrast vs. Family 2 under Hard Regime]：Furthermore, while [Baseline B] performs competitively at short horizons ($H=96$), its error escalates by [Y\%] at $H=720$ due to [Failure Mechanism]. | 句首锚定旧信息/逻辑视角，主干动词紧跟主语，复杂修饰置于宾语之后（符合尾重原则 L011 与旧新信息流 L001）。}
  \item \textbf{S4 [回扣本文核心机制 Mechanistic Confirmation]:} \textit{By decoupling stationary phase coupling from neural-ODE scale projection, DyTS-Net suppresses long-horizon drift and maintains an 11.6\% margin over the strongest baseline at $H=720$.}\par\noindent{\small 对应步骤：S4 [回扣本文核心机制 Mechanistic Confirmation]：By contrast, our [Core Module] bounds long-range drift, maintaining a [Z\%] margin over the second-best baseline at $H=720$. | 句首锚定旧信息/逻辑视角，主干动词紧跟主语，复杂修饰置于宾语之后（符合尾重原则 L011 与旧新信息流 L001）。}
\end{itemize}
\end{TipBox}
\medskip
\subsection*{Paragraph 9: 消融实验逐组件归因论证段（Counterfactual Ablation Study） (\texttt{BP19})}
\noindent\textbf{Expressive Intent:}\enspace 我这一段想汇报消融实验（Table 2）：证明我们的模块 A、模块 B 和正则项 C 每一个都不可或缺，且各自负责解决不同的特定子问题。\par
\noindent\textbf{Cohesion Formula:}\enspace {\small\color{AcadSlate}Ablation Overview $\to$ w/o Module A (Specific Failure) $\to$ w/o Module B (Specific Failure) $\to$ Cross-Module Synergy. 绑定逻辑规则：L016, L007.}\par\vspace{4pt}
\begin{ExampleBox}{Publication-Grade Authentic Text}
To isolate the contribution of each design choice, Table 2 compares full DyTS-Net against three ablated variants across ETTm2 and ECL. Removing the neural-ODE scale tracker (w/o Scale-ODE) triggers the sharpest degradation during seasonal transition months on ETTm2, raising 720-step MSE by 19.4\% as the decoder fails to track shifting baseline temperatures. Conversely, replacing phase-aligned graph attention with standard dot-product attention (w/o Phase-Attn) degrades performance most severely on the 321-channel ECL grid (+13.8\% MSE), where unaligned high-frequency load spikes corrupt cross-substation aggregation. Jointly deploying both modules yields complementary gains, as accurate scale tracking prevents amplitude distortion from biasing phase estimation.
\end{ExampleBox}
\begin{TipBox}{Rhetorical Move Breakdown \& Constituent Analysis}
\begin{itemize}[leftmargin=1.5em]
  \item \textbf{S1 [消融实验总括 Ablation Setup]:} \textit{To isolate the contribution of each design choice, Table 2 compares full DyTS-Net against three ablated variants across ETTm2 and ECL.}\par\noindent{\small 对应步骤：S1 [消融实验总括 Ablation Setup]：To isolate the contribution of each architectural component, we evaluate four ablated variants in Table 2 by progressively removing [Mod A], [Mod B], and [Loss C]. | 句首锚定旧信息/逻辑视角，主干动词紧跟主语，复杂修饰置于宾语之后（符合尾重原则 L011 与旧新信息流 L001）。}
  \item \textbf{S2 [移除模块 A 的冲击与归因 Impact of Removing Mod A]:} \textit{Removing the neural-ODE scale tracker (w/o Scale-ODE) triggers the sharpest degradation during seasonal transition months on ETTm2, raising 720-step MSE by 19.4\% as the decoder fails to track shifting baseline temperatures.}\par\noindent{\small 对应步骤：S2 [移除模块 A 的冲击与归因 Impact of Removing Mod A]：Removing [Mod A] (w/o A) causes the largest degradation on [Non-stationary/Volatile Subset], increasing MSE by [X\%] because the model loses [Specific Capability A]. | 句首锚定旧信息/逻辑视角，主干动词紧跟主语，复杂修饰置于宾语之后（符合尾重原则 L011 与旧新信息流 L001）。}
  \item \textbf{S3 [移除模块 B 的冲击与归因 Impact of Removing Mod B]:} \textit{Conversely, replacing phase-aligned graph attention with standard dot-product attention (w/o Phase-Attn) degrades performance most severely on the 321-channel ECL grid (+13.8\% MSE), where unaligned high-frequency load spikes corrupt cross-substation aggregation.}\par\noindent{\small 对应步骤：S3 [移除模块 B 的冲击与归因 Impact of Removing Mod B]：Similarly, replacing [Mod B] with standard [Naive Operator] degrades performance by [Y\%] on high-dimensional benchmarks, confirming its role in [Specific Capability B]. | 句首锚定旧信息/逻辑视角，主干动词紧跟主语，复杂修饰置于宾语之后（符合尾重原则 L011 与旧新信息流 L001）。}
  \item \textbf{S4 [协同效应总结 Synergistic Conclusion]:} \textit{Jointly deploying both modules yields complementary gains, as accurate scale tracking prevents amplitude distortion from biasing phase estimation.}\par\noindent{\small 对应步骤：S4 [协同效应总结 Synergistic Conclusion]：Combining all components yields super-additive gains, as [Mod A] stabilizes input statistics for [Mod B]'s relational reasoning. | 句首锚定旧信息/逻辑视角，主干动词紧跟主语，复杂修饰置于宾语之后（符合尾重原则 L011 与旧新信息流 L001）。}
\end{itemize}
\end{TipBox}
\medskip
\subsection*{Paragraph 10: 计算开销、推理延迟与精度帕累托权衡段（Efficiency-Accuracy Pareto Analysis） (\texttt{BP22})}
\noindent\textbf{Expressive Intent:}\enspace 我这一段想说明我们的模型不仅准，而且算得快、显存省，在“精度-延迟/显存”帕累托前沿（Pareto Frontier）上显著优于现有大模型。\par
\noindent\textbf{Cohesion Formula:}\enspace {\small\color{AcadSlate}Hardware Benchmark $\to$ Complexity Reduction \& Memory Gain $\to$ Real-Time Latency Margin $\to$ Pareto Conclusion. 绑定逻辑规则：L010, L023.}\par\vspace{4pt}
\begin{ExampleBox}{Publication-Grade Authentic Text}
Figure 6 plots the MSE-latency Pareto frontier measured on an onboard NVIDIA Orin AGX module. By replacing $O(L^2)$ point-wise self-attention with $O(L \log L)$ spectral filtering and compact patch tokens, DyTS-Net reduces peak GPU memory by 62\% and training wall-clock time by 3.4x compared with iTransformer. At inference time, generating a 16-step visuomotor trajectory requires only 6.8 ms---well below the 20 ms cycle budget of a 50 Hz industrial robot controller---while outperforming a 200 ms DDPM baseline by 11.4\% in task success rate.
\end{ExampleBox}
\begin{TipBox}{Rhetorical Move Breakdown \& Constituent Analysis}
\begin{itemize}[leftmargin=1.5em]
  \item \textbf{S1 [硬件与测试基准 Hardware Setup]:} \textit{Figure 6 plots the MSE-latency Pareto frontier measured on an onboard NVIDIA Orin AGX module.}\par\noindent{\small 对应步骤：S1 [硬件与测试基准 Hardware Setup]：Figure 6 compares the accuracy-efficiency Pareto frontier across all baselines measured on a single NVIDIA RTX 4090 GPU. | 句首锚定旧信息/逻辑视角，主干动词紧跟主语，复杂修饰置于宾语之后（符合尾重原则 L011 与旧新信息流 L001）。}
  \item \textbf{S2 [复杂度理论与实测对比 Complexity \& Runtime]:} \textit{By replacing $O(L^2)$ point-wise self-attention with $O(L \log L)$ spectral filtering and compact patch tokens, DyTS-Net reduces peak GPU memory by 62\% and training wall-clock time by 3.4x compared with iTransformer.}\par\noindent{\small 对应步骤：S2 [复杂度理论与实测对比 Complexity \& Runtime]：By replacing $O(L^2)$ point-wise attention with $O(L \log L)$ spectral filtering and patch aggregation, our model cuts training time per epoch by [X\%] and peak VRAM by [Y\%]. | 句首锚定旧信息/逻辑视角，主干动词紧跟主语，复杂修饰置于宾语之后（符合尾重原则 L011 与旧新信息流 L001）。}
  \item \textbf{S3 [实时闭环部署可行性 Real-Time Viability]:} \textit{At inference time, generating a 16-step visuomotor trajectory requires only 6.8 ms---well below the 20 ms cycle budget of a 50 Hz industrial robot controller---while outperforming a 200 ms DDPM baseline by 11.4\% in task success rate.}\par\noindent{\small 对应步骤：S3 [实时闭环部署可行性 Real-Time Viability]：At inference time, generating a 720-step forecast (or 16-step robot action chunk) requires only [4.2 ms], well below the [20 ms] real-time control deadline. | 句首锚定旧信息/逻辑视角，主干动词紧跟主语，复杂修饰置于宾语之后（符合尾重原则 L011 与旧新信息流 L001）。}
\end{itemize}
\end{TipBox}
\medskip
\subsection*{Paragraph 11: 全文总结与学术影响收尾段（Conclusion \& Broader Outlook） (\texttt{BP24})}
\noindent\textbf{Expressive Intent:}\enspace 我这一段想写整篇小论文的 Conclusion：不简单复制粘贴 Abstract，而是从更高维度总结本文解决的科学矛盾、方法精髓、实证高度与领域启示。\par
\noindent\textbf{Cohesion Formula:}\enspace {\small\color{AcadSlate}Tension Resolved $\to$ Methodological Insight $\to$ Empirical Confirmation $\to$ Generalizable Blueprint. 绑定逻辑规则：L013, L023, L029.}\par\vspace{4pt}
\begin{ExampleBox}{Publication-Grade Authentic Text}
In this work, we resolved the long-standing tension between shift-invariant normalization and physical amplitude preservation in non-stationary time series forecasting. By decoupling frequency-domain phase attention from continuous neural-ODE scale tracking, DyTS-Net captures long-range seasonal dependencies while accurately reconstructing extreme peak magnitudes. Empirical evaluations across eight benchmarks and industrial wind-power telemetry demonstrated MSE reductions of up to 17.4\% alongside a 45\% reduction in computational overhead. Beyond forecasting, this structural-scale disentanglement principle provides a generalizable foundation for anomaly detection and closed-loop predictive control in non-stationary cyber-physical systems.
\end{ExampleBox}
\begin{TipBox}{Rhetorical Move Breakdown \& Constituent Analysis}
\begin{itemize}[leftmargin=1.5em]
  \item \textbf{S1 [核心科学问题闭合 Scientific Closure]:} \textit{In this work, we resolved the long-standing tension between shift-invariant normalization and physical amplitude preservation in non-stationary time series forecasting.}\par\noindent{\small 对应步骤：S1 [核心科学问题闭合 Scientific Closure]：In this paper, we addressed the fundamental tension between [Goal A] and [Goal B] in [Domain]. | 句首锚定旧信息/逻辑视角，主干动词紧跟主语，复杂修饰置于宾语之后（符合尾重原则 L011 与旧新信息流 L001）。}
  \item \textbf{S2 [方法论精髓回顾 Methodological Synthesis]:} \textit{By decoupling frequency-domain phase attention from continuous neural-ODE scale tracking, DyTS-Net captures long-range seasonal dependencies while accurately reconstructing extreme peak magnitudes.}\par\noindent{\small 对应步骤：S2 [方法论精髓回顾 Methodological Synthesis]：By introducing [Model/Mechanism], we demonstrated that [Core Insight] enables [Capability] without incurring [Traditional Penalty]. | 句首锚定旧信息/逻辑视角，主干动词紧跟主语，复杂修饰置于宾语之后（符合尾重原则 L011 与旧新信息流 L001）。}
  \item \textbf{S3 [实证结果高度概括 Empirical Recap]:} \textit{Empirical evaluations across eight benchmarks and industrial wind-power telemetry demonstrated MSE reductions of up to 17.4\% alongside a 45\% reduction in computational overhead.}\par\noindent{\small 对应步骤：S3 [实证结果高度概括 Empirical Recap]：Comprehensive evaluations across [Benchmarks/Hardware] confirmed consistent error reductions of up to [X\%] alongside substantial inference speedups. | 句首锚定旧信息/逻辑视角，主干动词紧跟主语，复杂修饰置于宾语之后（符合尾重原则 L011 与旧新信息流 L001）。}
  \item \textbf{S4 [跨领域迁移展望 Broader Outlook]:} \textit{Beyond forecasting, this structural-scale disentanglement principle provides a generalizable foundation for anomaly detection and closed-loop predictive control in non-stationary cyber-physical systems.}\par\noindent{\small 对应步骤：S4 [跨领域迁移展望 Broader Outlook]：Beyond [Current Task], our [Core Principle] offers a generalizable blueprint for broader sequential decision-making tasks under non-stationary environments. | 句首锚定旧信息/逻辑视角，主干动词紧跟主语，复杂修饰置于宾语之后（符合尾重原则 L011 与旧新信息流 L001）。}
\end{itemize}
\end{TipBox}
\medskip
\clearpage
\section{FT05: 基于整流蒸馏的单步实时机械臂轨迹生成完整小论文（RectiPolicy: Real-Time Action Chunk Generation）}\label{comp:FT05}
\begin{ConceptBox}{Composition Architecture \& Structural Blueprint Chain}
\textbf{Genre:} \texttt{AI\_Mini\_Paper} \quad \textbf{Paragraph Count:} 11\par
\textbf{Blueprint Chain:} \texttt{BP08 $\to$ BP02 $\to$ BP04 $\to$ BP05 $\to$ BP10 $\to$ BP15 $\to$ BP16 $\to$ BP18 $\to$ BP22 $\to$ BP23 $\to$ BP24}
\end{ConceptBox}
\subsection*{Paragraph 1: 高密度四句式 Abstract 全段构建（Complete 4-Move Abstract） (\texttt{BP08})}
\noindent\textbf{Expressive Intent:}\enspace 我这一段想写出整篇小论文的完整 Abstract（摘要）：在 150--180 词内一气呵成交代背景、核心痛点、方法机制与量化实验结果。\par
\noindent\textbf{Cohesion Formula:}\enspace {\small\color{AcadSlate}Context (S1) $\to$ Gap (S2) $\to$ Method \& Mechanism (S3-S4) $\to$ Quantitative Proof (S5). 绑定逻辑规则：L001, L008, L010, L029.}\par\vspace{4pt}
\begin{ExampleBox}{Publication-Grade Authentic Text}
Recovering missing sensor observations is essential for reliable downstream monitoring in industrial and clinical time series. While diffusion-based imputation models generate plausible trajectories under random missingness, they suffer from high sampling latency and lose global periodic coherence when entire sensor blocks drop out consecutively. We present FlowImp, a conditional flow-matching framework that reconstructs contiguous missing blocks in three ODE solver steps. Specifically, FlowImp anchors straight transport paths on observed spectral priors and enforces cross-channel phase consistency via a lightweight Fourier regularizer. Across six multivariate benchmarks under 50\% block missingness, FlowImp reduces imputation MAE by 19.3\% compared with CSDI while accelerating inference by 24x.
\end{ExampleBox}
\begin{TipBox}{Rhetorical Move Breakdown \& Constituent Analysis}
\begin{itemize}[leftmargin=1.5em]
  \item \textbf{S1 [背景与任务重要性 Context]:} \textit{Recovering missing sensor observations is essential for reliable downstream monitoring in industrial and clinical time series.}\par\noindent{\small 对应步骤：S1 [背景与任务重要性 Context]：1句话锁定任务目标与核心应用场景（≤20词）。 | 句首锚定旧信息/逻辑视角，主干动词紧跟主语，复杂修饰置于宾语之后（符合尾重原则 L011 与旧新信息流 L001）。}
  \item \textbf{S2 [现有局限与痛点 Gap/Bottleneck]:} \textit{While diffusion-based imputation models generate plausible trajectories under random missingness, they suffer from high sampling latency and lose global periodic coherence when entire sensor blocks drop out consecutively.}\par\noindent{\small 对应步骤：S2 [现有局限与痛点 Gap/Bottleneck]：While existing [Paradigms] achieve [Merit], they struggle with [Core Bottleneck] because [Underlying Cause]. | 句首锚定旧信息/逻辑视角，主干动词紧跟主语，复杂修饰置于宾语之后（符合尾重原则 L011 与旧新信息流 L001）。}
  \item \textbf{S3-S4 [本文方法与核心机制 Approach \& Mechanism]:} \textit{We present FlowImp, a conditional flow-matching framework that reconstructs contiguous missing blocks in three ODE solver steps.}\par\noindent{\small 对应步骤：S3-S4 [本文方法与核心机制 Approach \& Mechanism]：In this work, we present [Model], which [Core Idea]. Specifically, [Module 1]..., while [Module 2]... | 句首锚定旧信息/逻辑视角，主干动词紧跟主语，复杂修饰置于宾语之后（符合尾重原则 L011 与旧新信息流 L001）。}
  \item \textbf{S5 [量化实证结果 Quantitative Results]:} \textit{Specifically, FlowImp anchors straight transport paths on observed spectral priors and enforces cross-channel phase consistency via a lightweight Fourier regularizer.}\par\noindent{\small 对应步骤：S5 [量化实证结果 Quantitative Results]：Experiments on [Benchmarks/Hardware] show that [Model] outperforms [Baselines] by [X\%] in [Metric] and achieves [Efficiency/Real-world Result]. | 句首锚定旧信息/逻辑视角，主干动词紧跟主语，复杂修饰置于宾语之后（符合尾重原则 L011 与旧新信息流 L001）。}
  \item \textbf{S5 [量化实证结果 Quantitative Results]:} \textit{Across six multivariate benchmarks under 50\% block missingness, FlowImp reduces imputation MAE by 19.3\% compared with CSDI while accelerating inference by 24x.}\par\noindent{\small 对应步骤：S5 [量化实证结果 Quantitative Results]：Experiments on [Benchmarks/Hardware] show that [Model] outperforms [Baselines] by [X\%] in [Metric] and achieves [Efficiency/Real-world Result]. | 句首锚定旧信息/逻辑视角，主干动词紧跟主语，复杂修饰置于宾语之后（符合尾重原则 L011 与旧新信息流 L001）。}
\end{itemize}
\end{TipBox}
\medskip
\subsection*{Paragraph 2: 总结现有主流范式及其隐含假设（欲抑先扬铺垫段） (\texttt{BP02})}
\noindent\textbf{Expressive Intent:}\enspace 我这一段想客观概括前人主流模型（如 Transformer / Diffusion / MLP）取得了什么进展，但点出它们共同依赖的一个脆弱前提假设。\par
\noindent\textbf{Cohesion Formula:}\enspace {\small\color{AcadSlate}先让步肯定（While/Recent paradigms...）$\to$ 后釜底抽薪揭示隐含前提（Yet these architectures implicitly assume...）。绑定逻辑规则：L008, L009.}\par\vspace{4pt}
\begin{ExampleBox}{Publication-Grade Authentic Text}
Recent deep forecasters have substantially reduced long-horizon error by replacing point-wise attention with sub-series patch tokenization and frequency-domain decomposition. These architectures effectively capture recurring seasonal patterns when training and test sequences share stationary statistical moments. Yet they implicitly assume that inter-channel correlations and spectral energy distributions remain time-invariant across the prediction window. When external shocks alter underlying system dynamics, this stationarity assumption collapses, causing frozen spectral filters to project outdated trends.
\end{ExampleBox}
\begin{TipBox}{Rhetorical Move Breakdown \& Constituent Analysis}
\begin{itemize}[leftmargin=1.5em]
  \item \textbf{S1 [流派归类 Taxonomic Opening]:} \textit{Recent deep forecasters have substantially reduced long-horizon error by replacing point-wise attention with sub-series patch tokenization and frequency-domain decomposition.}\par\noindent{\small 对应步骤：S1 [流派归类 Taxonomic Opening]：按核心机制归纳现有主流范式及其已解决的子问题。 | 句首锚定旧信息/逻辑视角，主干动词紧跟主语，复杂修饰置于宾语之后（符合尾重原则 L011 与旧新信息流 L001）。}
  \item \textbf{S2 [代表性机制展开 Mechanism Recap]:} \textit{These architectures effectively capture recurring seasonal patterns when training and test sequences share stationary statistical moments.}\par\noindent{\small 对应步骤：S2 [代表性机制展开 Mechanism Recap]：简述这些方法如何通过特定架构（如 Patch 分块、频域分解或扩散去噪）提升表征能力。 | 句首锚定旧信息/逻辑视角，主干动词紧跟主语，复杂修饰置于宾语之后（符合尾重原则 L011 与旧新信息流 L001）。}
  \item \textbf{S3 [转折揭示隐含假设 Pivot to Hidden Assumption]:} \textit{Yet they implicitly assume that inter-channel correlations and spectral energy distributions remain time-invariant across the prediction window.}\par\noindent{\small 对应步骤：S3 [转折揭示隐含假设 Pivot to Hidden Assumption]：用 Yet / However 锋利转折，指出它们暗中依赖的理想化假设（如平稳分布、静态通道相关性、仿真动力学无偏差）。 | 句首锚定旧信息/逻辑视角，主干动词紧跟主语，复杂修饰置于宾语之后（符合尾重原则 L011 与旧新信息流 L001）。}
  \item \textbf{S4 [假设破裂后果 Breakdown Under Stress]:} \textit{When external shocks alter underlying system dynamics, this stationarity assumption collapses, causing frozen spectral filters to project outdated trends.}\par\noindent{\small 对应步骤：S4 [假设破裂后果 Breakdown Under Stress]：说明当真实数据违背该假设时，现有架构为何会遭遇性能断崖。 | 句首锚定旧信息/逻辑视角，主干动词紧跟主语，复杂修饰置于宾语之后（符合尾重原则 L011 与旧新信息流 L001）。}
\end{itemize}
\end{TipBox}
\medskip
\subsection*{Paragraph 3: 亮出本文核心洞察与破局方法总览（Key Insight \& Method Overview） (\texttt{BP04})}
\noindent\textbf{Expressive Intent:}\enspace 我这一段想在指出前人痛点后，正式亮出我们的方法（To address this dilemma, we propose...），并用最通俗清晰的直觉告诉审稿人为什么我们的设计能破局。\par
\noindent\textbf{Cohesion Formula:}\enspace {\small\color{AcadSlate}Problem Anchor $\to$ Key Insight $\to$ Dual-Module Mechanism $\to$ Pareto Resolution. 绑定逻辑规则：L002, L023.}\par\vspace{4pt}
\begin{ExampleBox}{Publication-Grade Authentic Text}
To resolve this dilemma, we propose DyTS-Net, a dual-stream forecasting architecture that decouples stationary structural modeling from dynamic scale trajectory estimation. Our key insight is that while high-frequency phase relationships benefit from normalized shift-invariant attention, macro-scale amplitude evolution follows a low-dimensional continuous dynamical system. Specifically, DyTS-Net first extracts scale-invariant temporal tokens via spectral patch attention, and concurrently models non-stationary mean-variance evolution through a lightweight neural ODE branch. Recombining both streams at the decoding stage preserves sharp peak magnitudes without sacrificing training stability.
\end{ExampleBox}
\begin{TipBox}{Rhetorical Move Breakdown \& Constituent Analysis}
\begin{itemize}[leftmargin=1.5em]
  \item \textbf{S1 [方法登场与定位 Method Introduction]:} \textit{To resolve this dilemma, we propose DyTS-Net, a dual-stream forecasting architecture that decouples stationary structural modeling from dynamic scale trajectory estimation.}\par\noindent{\small 对应步骤：S1 [方法登场与定位 Method Introduction]：To address this dilemma, we propose [Model Name], a [Core Paradigm] framework that [Primary Goal]. | 句首锚定旧信息/逻辑视角，主干动词紧跟主语，复杂修饰置于宾语之后（符合尾重原则 L011 与旧新信息流 L001）。}
  \item \textbf{S2 [核心直觉洞察 Key Insight]:} \textit{Our key insight is that while high-frequency phase relationships benefit from normalized shift-invariant attention, macro-scale amplitude evolution follows a low-dimensional continuous dynamical system.}\par\noindent{\small 对应步骤：S2 [核心直觉洞察 Key Insight]：Our key insight is that [Physical/Statistical Principle]，点破破局的底层支点。 | 句首锚定旧信息/逻辑视角，主干动词紧跟主语，复杂修饰置于宾语之后（符合尾重原则 L011 与旧新信息流 L001）。}
  \item \textbf{S3 [模块协同主干 Module Synergy]:} \textit{Specifically, DyTS-Net first extracts scale-invariant temporal tokens via spectral patch attention, and concurrently models non-stationary mean-variance evolution through a lightweight neural ODE branch.}\par\noindent{\small 对应步骤：S3 [模块协同主干 Module Synergy]：用 First / Specifically 与 Concurrently / Furthermore 简述两大核心模块如何分工咬合。 | 句首锚定旧信息/逻辑视角，主干动词紧跟主语，复杂修饰置于宾语之后（符合尾重原则 L011 与旧新信息流 L001）。}
  \item \textbf{S4 [闭环收益总结 Closed-Loop Benefit]:} \textit{Recombining both streams at the decoding stage preserves sharp peak magnitudes without sacrificing training stability.}\par\noindent{\small 对应步骤：S4 [闭环收益总结 Closed-Loop Benefit]：说明这种协同设计如何同时兼顾原先冲突的两个指标（如兼顾平稳表征与幅值保真，或兼顾多模态多样性与单步实时控制）。 | 句首锚定旧信息/逻辑视角，主干动词紧跟主语，复杂修饰置于宾语之后（符合尾重原则 L011 与旧新信息流 L001）。}
\end{itemize}
\end{TipBox}
\medskip
\subsection*{Paragraph 4: 挑战与机制一对一映射段（Challenge-to-Mechanism Mapping） (\texttt{BP05})}
\noindent\textbf{Expressive Intent:}\enspace 我这一段想说明：把我们的核心思想落地时遇到了两个具体工程/数学难点，我们分别设计了什么精巧机制去攻克它们。\par
\noindent\textbf{Cohesion Formula:}\enspace {\small\color{AcadSlate}Challenge 1 $\to$ Mechanism 1; Challenge 2 $\to$ Mechanism 2 $\to$ Joint System Property. 绑定逻辑规则：L003, L011.}\par\vspace{4pt}
\begin{ExampleBox}{Publication-Grade Authentic Text}
Instantiating closed-loop visuomotor diffusion on physical robots entails two technical challenges. First, because raw depth camera streams contain occlusion artifacts around reflective object edges, we introduce a tactile-guided cross-attention gate that reweights visual tokens upon contact onset. Second, to eliminate the 200 ms sampling bottleneck of iterative denoising, we formulate a trajectory consistency distillation objective that maps Gaussian noise directly to collision-free action chunks in two forward steps. Together, these designs sustain 50 Hz reactive control under partial visual occlusion.
\end{ExampleBox}
\begin{TipBox}{Rhetorical Move Breakdown \& Constituent Analysis}
\begin{itemize}[leftmargin=1.5em]
  \item \textbf{S1 [落地难点总起 Implementation Hurdle]:} \textit{Instantiating closed-loop visuomotor diffusion on physical robots entails two technical challenges.}\par\noindent{\small 对应步骤：S1 [落地难点总起 Implementation Hurdle]：Translating this principle into a scalable architecture entails two technical challenges. | 句首锚定旧信息/逻辑视角，主干动词紧跟主语，复杂修饰置于宾语之后（符合尾重原则 L011 与旧新信息流 L001）。}
  \item \textbf{S2 [难点 1 + 靶向设计 1]:} \textit{First, because raw depth camera streams contain occlusion artifacts around reflective object edges, we introduce a tactile-guided cross-attention gate that reweights visual tokens upon contact onset.}\par\noindent{\small 对应步骤：S2 [难点 1 + 靶向设计 1]：First, because [Challenge 1], we design [Mechanism 1] to [Specific Function]. | 句首锚定旧信息/逻辑视角，主干动词紧跟主语，复杂修饰置于宾语之后（符合尾重原则 L011 与旧新信息流 L001）。}
  \item \textbf{S3 [难点 2 + 靶向设计 2]:} \textit{Second, to eliminate the 200 ms sampling bottleneck of iterative denoising, we formulate a trajectory consistency distillation objective that maps Gaussian noise directly to collision-free action chunks in two forward steps.}\par\noindent{\small 对应步骤：S3 [难点 2 + 靶向设计 2]：Second, to prevent [Challenge 2] under [Constraint], we introduce [Mechanism 2] that [Specific Function]. | 句首锚定旧信息/逻辑视角，主干动词紧跟主语，复杂修饰置于宾语之后（符合尾重原则 L011 与旧新信息流 L001）。}
  \item \textbf{S4 [整体计算/泛化特性 System Property]:} \textit{Together, these designs sustain 50 Hz reactive control under partial visual occlusion.}\par\noindent{\small 对应步骤：S4 [整体计算/泛化特性 System Property]：Together, these components enable [Model] to operate at [Efficiency/Robustness Metric]. | 句首锚定旧信息/逻辑视角，主干动词紧跟主语，复杂修饰置于宾语之后（符合尾重原则 L011 与旧新信息流 L001）。}
\end{itemize}
\end{TipBox}
\medskip
\subsection*{Paragraph 5: 对比现有流派优劣并精准定位本文差异（Contrastive Positioning） (\texttt{BP10})}
\noindent\textbf{Expressive Intent:}\enspace 我这一段想在 Related Work 小节末尾，把我们的方法与最接近的 2 个竞争方案做精准对比，防止审稿人说我们只是对某篇旧文的微调（Incremental）。\par
\noindent\textbf{Cohesion Formula:}\enspace {\small\color{AcadSlate}Acknowledge Overlap $\to$ Pinpoint Structural Divergence $\to$ Explain Functional Gain. 绑定逻辑规则：L008, L019.}\par\vspace{4pt}
\begin{ExampleBox}{Publication-Grade Authentic Text}
Among spectral forecasters, FEDformer and TimesNet are most closely related to our approach. While both methods transform temporal representations via fast Fourier transforms, they apply static top-k frequency selection globally across the entire lookback window, which blurs localized transient spikes. In contrast, our wavelet-coupled attention dynamically adapts filter bandwidths over sliding sub-intervals, isolating non-stationary burst events without corrupting global seasonal bases. Consequently, our model captures localized fault transients without requiring pre-defined frequency thresholds.
\end{ExampleBox}
\begin{TipBox}{Rhetorical Move Breakdown \& Constituent Analysis}
\begin{itemize}[leftmargin=1.5em]
  \item \textbf{S1 [锁定最相近工作 Closest Prior Work]:} \textit{Among spectral forecasters, FEDformer and TimesNet are most closely related to our approach.}\par\noindent{\small 对应步骤：S1 [锁定最相近工作 Closest Prior Work]：Among prior studies, [Closest Method A] and [Closest Method B] are most closely related to our approach. | 句首锚定旧信息/逻辑视角，主干动词紧跟主语，复杂修饰置于宾语之后（符合尾重原则 L011 与旧新信息流 L001）。}
  \item \textbf{S2 [剖析其局限性 Structural Difference 1]:} \textit{While both methods transform temporal representations via fast Fourier transforms, they apply static top-k frequency selection globally across the entire lookback window, which blurs localized transient spikes.}\par\noindent{\small 对应步骤：S2 [剖析其局限性 Structural Difference 1]：While [Method A] also employs [Shared Concept], it relies on [Limitation 1], restricting its applicability to [Narrow Setting]. | 句首锚定旧信息/逻辑视角，主干动词紧跟主语，复杂修饰置于宾语之后（符合尾重原则 L011 与旧新信息流 L001）。}
  \item \textbf{S3 [剖析本文本质突破 Core Departure]:} \textit{In contrast, our wavelet-coupled attention dynamically adapts filter bandwidths over sliding sub-intervals, isolating non-stationary burst events without corrupting global seasonal bases.}\par\noindent{\small 对应步骤：S3 [剖析本文本质突破 Core Departure]：In contrast, our framework introduces [Core Innovation], which fundamentally eliminates [Bottleneck] by [Mechanism]. | 句首锚定旧信息/逻辑视角，主干动词紧跟主语，复杂修饰置于宾语之后（符合尾重原则 L011 与旧新信息流 L001）。}
  \item \textbf{S4 [实证或理论优势锁定 Empirical/Theoretical Edge]:} \textit{Consequently, our model captures localized fault transients without requiring pre-defined frequency thresholds.}\par\noindent{\small 对应步骤：S4 [实证或理论优势锁定 Empirical/Theoretical Edge]：Consequently, our approach achieves [Capability] without requiring [Expensive Requirement of Prior Work]. | 句首锚定旧信息/逻辑视角，主干动词紧跟主语，复杂修饰置于宾语之后（符合尾重原则 L011 与旧新信息流 L001）。}
\end{itemize}
\end{TipBox}
\medskip
\subsection*{Paragraph 6: 生成模型专属：条件控制与采样去噪轨迹阐述段（Generative Modeling Mechanism） (\texttt{BP15})}
\noindent\textbf{Expressive Intent:}\enspace 我这一段想阐述生成模型（Diffusion / Flow Matching / Autoregressive）如何在满足外部物理约束或条件观测的前提下，生成高质量且多样化的序列/轨迹。\par
\noindent\textbf{Cohesion Formula:}\enspace {\small\color{AcadSlate}Guidance Artifact Problem $\to$ Intrinsic Velocity Conditioning $\to$ Tangent Projection Step $\to$ Fast Exact Convergence. 绑定逻辑规则：L026, L014.}\par\vspace{4pt}
\begin{ExampleBox}{Publication-Grade Authentic Text}
Applying post-hoc gradient guidance during diffusion sampling often pushes intermediate noisy states off the data manifold, introducing unphysical high-frequency oscillations. To enforce physical boundary constraints without manifold deviation, we incorporate observation consistency directly into the conditional flow-matching velocity field. At flow time $\tau \in [0, 1]$, our network predicts a residual velocity $v_\theta(x_\tau, \tau \mid c)$ that projects the transport trajectory onto the null space of the linear constraint operator. This projection guarantees exact boundary satisfaction at terminal time $\tau=1$ while maintaining straight transport paths amenable to 4-step Euler integration.
\end{ExampleBox}
\begin{TipBox}{Rhetorical Move Breakdown \& Constituent Analysis}
\begin{itemize}[leftmargin=1.5em]
  \item \textbf{S1 [无条件/朴素引导缺陷 Naive Guidance Bottleneck]:} \textit{Applying post-hoc gradient guidance during diffusion sampling often pushes intermediate noisy states off the data manifold, introducing unphysical high-frequency oscillations.}\par\noindent{\small 对应步骤：S1 [无条件/朴素引导缺陷 Naive Guidance Bottleneck]：Applying post-hoc classifier guidance to pre-trained diffusion samplers often pushes intermediate noisy states off the data manifold, inducing high-frequency artifacts. | 句首锚定旧信息/逻辑视角，主干动词紧跟主语，复杂修饰置于宾语之后（符合尾重原则 L011 与旧新信息流 L001）。}
  \item \textbf{S2 [本文流形内条件注入机制 Manifold-Consistent Conditioning]:} \textit{To enforce physical boundary constraints without manifold deviation, we incorporate observation consistency directly into the conditional flow-matching velocity field.}\par\noindent{\small 对应步骤：S2 [本文流形内条件注入机制 Manifold-Consistent Conditioning]：To enforce hard boundary constraints without manifold deviation, we embed [Condition/Constraint] directly into the velocity field / denoising kernel. | 句首锚定旧信息/逻辑视角，主干动词紧跟主语，复杂修饰置于宾语之后（符合尾重原则 L011 与旧新信息流 L001）。}
  \item \textbf{S3 [具体轨迹更新方程 Trajectory Update Formulation]:} \textit{At flow time $\tau \in [0, 1]$, our network predicts a residual velocity $v_\theta(x_\tau, \tau \mid c)$ that projects the transport trajectory onto the null space of the linear constraint operator.}\par\noindent{\small 对应步骤：S3 [具体轨迹更新方程 Trajectory Update Formulation]：At solver step $\tau \in [0, 1]$, the conditional velocity field $v_\theta(x_\tau, \tau \mid c)$ projects the unconstrained flow onto the tangent space of feasible trajectories. | 句首锚定旧信息/逻辑视角，主干动词紧跟主语，复杂修饰置于宾语之后（符合尾重原则 L011 与旧新信息流 L001）。}
  \item \textbf{S4 [采样效率与分布保真收益 Sampling \& Fidelity Guarantee]:} \textit{This projection guarantees exact boundary satisfaction at terminal time $\tau=1$ while maintaining straight transport paths amenable to 4-step Euler integration.}\par\noindent{\small 对应步骤：S4 [采样效率与分布保真收益 Sampling \& Fidelity Guarantee]：This tangent projection guarantees exact boundary satisfaction at $\tau=1$ while permitting straight 4-step Euler integration. | 句首锚定旧信息/逻辑视角，主干动词紧跟主语，复杂修饰置于宾语之后（符合尾重原则 L011 与旧新信息流 L001）。}
\end{itemize}
\end{TipBox}
\medskip
\subsection*{Paragraph 7: 具身智能专属：多模态感知-动作闭环与 Sim-to-Real 阐述段（Embodied AI Policy） (\texttt{BP16})}
\noindent\textbf{Expressive Intent:}\enspace 我这一段想阐述具身智能机器人策略（Embodied Policy）如何融合视觉、本体感觉（Proprioception）或触觉，并克服仿真到真实硬件的动力学差距（Sim-to-Real Gap）。\par
\noindent\textbf{Cohesion Formula:}\enspace {\small\color{AcadSlate}Visual Occlusion Bottleneck $\to$ Asynchronous Multimodal Alignment $\to$ Action Chunking $\to$ Online Sim-to-Real System Identification. 绑定逻辑规则：L025, L003.}\par\vspace{4pt}
\begin{ExampleBox}{Publication-Grade Authentic Text}
Relying solely on third-person RGB-D observations leaves visuomotor policies vulnerable to severe self-occlusion once the gripper approaches the target object. To preserve observability during contact-rich manipulation, our encoder aligns 30 Hz visual spatial tokens with 500 Hz joint proprioception and wrist force-torque histories via cross-modal rotary embeddings. Rather than issuing single-step commands prone to high-frequency jitter, the policy predicts 16-step SE(3) end-effector action chunks smoothed through exponential temporal ensembling. At deployment, a recurrent history adapter estimates latent surface friction and payload mass online, bridging the sim-to-real dynamics gap in zero shot.
\end{ExampleBox}
\begin{TipBox}{Rhetorical Move Breakdown \& Constituent Analysis}
\begin{itemize}[leftmargin=1.5em]
  \item \textbf{S1 [单模态开环局限 Open-Loop Visual Limitation]:} \textit{Relying solely on third-person RGB-D observations leaves visuomotor policies vulnerable to severe self-occlusion once the gripper approaches the target object.}\par\noindent{\small 对应步骤：S1 [单模态开环局限 Open-Loop Visual Limitation]：Relying solely on RGB-D point clouds leaves manipulation policies vulnerable to self-occlusion by the robot gripper during fine-grained contact. | 句首锚定旧信息/逻辑视角，主干动词紧跟主语，复杂修饰置于宾语之后（符合尾重原则 L011 与旧新信息流 L001）。}
  \item \textbf{S2 [多模态时空对齐 Multi-Modal Alignment]:} \textit{To preserve observability during contact-rich manipulation, our encoder aligns 30 Hz visual spatial tokens with 500 Hz joint proprioception and wrist force-torque histories via cross-modal rotary embeddings.}\par\noindent{\small 对应步骤：S2 [多模态时空对齐 Multi-Modal Alignment]：To maintain state observability during contact, our encoder aligns asynchronous 30 Hz visual tokens with 500 Hz joint proprioception and wrist force-torque histories. | 句首锚定旧信息/逻辑视角，主干动词紧跟主语，复杂修饰置于宾语之后（符合尾重原则 L011 与旧新信息流 L001）。}
  \item \textbf{S3 [动作块预测与平滑执行 Action Chunking \& Temporal Ensembling]:} \textit{Rather than issuing single-step commands prone to high-frequency jitter, the policy predicts 16-step SE(3) end-effector action chunks smoothed through exponential temporal ensembling.}\par\noindent{\small 对应步骤：S3 [动作块预测与平滑执行 Action Chunking \& Temporal Ensembling]：Rather than predicting single-step torques prone to jitter, the policy outputs overlapping $K$-step SE(3) end-effector action chunks smoothed via exponential temporal ensembling. | 句首锚定旧信息/逻辑视角，主干动词紧跟主语，复杂修饰置于宾语之后（符合尾重原则 L011 与旧新信息流 L001）。}
  \item \textbf{S4 [Sim-to-Real 自适应机制 Sim-to-Real Adaptation]:} \textit{At deployment, a recurrent history adapter estimates latent surface friction and payload mass online, bridging the sim-to-real dynamics gap in zero shot.}\par\noindent{\small 对应步骤：S4 [Sim-to-Real 自适应机制 Sim-to-Real Adaptation]：During physical deployment, a lightweight history adapter infers latent friction and payload parameters online, bridging the sim-to-real dynamics gap without policy fine-tuning. | 句首锚定旧信息/逻辑视角，主干动词紧跟主语，复杂修饰置于宾语之后（符合尾重原则 L011 与旧新信息流 L001）。}
\end{itemize}
\end{TipBox}
\medskip
\subsection*{Paragraph 8: 主实验结果整体对比与核心胜出原因剖析段（Main Results Analysis） (\texttt{BP18})}
\noindent\textbf{Expressive Intent:}\enspace 我这一段想解读主实验大表（Table 1）：既报告全局胜率，又挑出最有代表性的对比案例，解释为什么竞争对手在那里跌倒而我们稳赢。\par
\noindent\textbf{Cohesion Formula:}\enspace {\small\color{AcadSlate}Global Win Rate $\to$ Advantage over Family A $\to$ Resilience at Long Horizon vs. Family B $\to$ Mechanistic Takeaway. 绑定逻辑规则：L015, L013, L010.}\par\vspace{4pt}
\begin{ExampleBox}{Publication-Grade Authentic Text}
Table 1 reports the multivariate forecasting results, where DyTS-Net secures the lowest error in 25 out of 28 dataset-horizon configurations. Compared with channel-independent PatchTST, DyTS-Net reduces average MSE by 14.2\% on the 862-channel Traffic benchmark, where modeling spatial-temporal congestion propagation across adjacent road sensors is critical. Furthermore, while variate-mixing iTransformer performs competitively at $H=96$, its MSE escalates by 68\% at $H=720$ on volatile Weather series as unregularized cross-channel attention overfits transient wind gusts. By decoupling stationary phase coupling from neural-ODE scale projection, DyTS-Net suppresses long-horizon drift and maintains an 11.6\% margin over the strongest baseline at $H=720$.
\end{ExampleBox}
\begin{TipBox}{Rhetorical Move Breakdown \& Constituent Analysis}
\begin{itemize}[leftmargin=1.5em]
  \item \textbf{S1 [全局统治力概述 Overall Dominance]:} \textit{Table 1 reports the multivariate forecasting results, where DyTS-Net secures the lowest error in 25 out of 28 dataset-horizon configurations.}\par\noindent{\small 对应步骤：S1 [全局统治力概述 Overall Dominance]：Table 1 summarizes the multivariate forecasting results, where [Our Model] achieves the lowest error in [X] out of [Y] dataset-horizon configurations. | 句首锚定旧信息/逻辑视角，主干动词紧跟主语，复杂修饰置于宾语之后（符合尾重原则 L011 与旧新信息流 L001）。}
  \item \textbf{S2 [对比流派 1 的差距与原因 Contrast vs. Family 1]:} \textit{Compared with channel-independent PatchTST, DyTS-Net reduces average MSE by 14.2\% on the 862-channel Traffic benchmark, where modeling spatial-temporal congestion propagation across adjacent road sensors is critical.}\par\noindent{\small 对应步骤：S2 [对比流派 1 的差距与原因 Contrast vs. Family 1]：Compared with channel-independent baselines such as [Baseline A], our method reduces average MSE by [X\%] on [Dataset], where strong cross-sensor coupling is critical. | 句首锚定旧信息/逻辑视角，主干动词紧跟主语，复杂修饰置于宾语之后（符合尾重原则 L011 与旧新信息流 L001）。}
  \item \textbf{S3 [对比流派 2 在困难场景下的崩溃 Contrast vs. Family 2 under Hard Regime]:} \textit{Furthermore, while variate-mixing iTransformer performs competitively at $H=96$, its MSE escalates by 68\% at $H=720$ on volatile Weather series as unregularized cross-channel attention overfits transient wind gusts.}\par\noindent{\small 对应步骤：S3 [对比流派 2 在困难场景下的崩溃 Contrast vs. Family 2 under Hard Regime]：Furthermore, while [Baseline B] performs competitively at short horizons ($H=96$), its error escalates by [Y\%] at $H=720$ due to [Failure Mechanism]. | 句首锚定旧信息/逻辑视角，主干动词紧跟主语，复杂修饰置于宾语之后（符合尾重原则 L011 与旧新信息流 L001）。}
  \item \textbf{S4 [回扣本文核心机制 Mechanistic Confirmation]:} \textit{By decoupling stationary phase coupling from neural-ODE scale projection, DyTS-Net suppresses long-horizon drift and maintains an 11.6\% margin over the strongest baseline at $H=720$.}\par\noindent{\small 对应步骤：S4 [回扣本文核心机制 Mechanistic Confirmation]：By contrast, our [Core Module] bounds long-range drift, maintaining a [Z\%] margin over the second-best baseline at $H=720$. | 句首锚定旧信息/逻辑视角，主干动词紧跟主语，复杂修饰置于宾语之后（符合尾重原则 L011 与旧新信息流 L001）。}
\end{itemize}
\end{TipBox}
\medskip
\subsection*{Paragraph 9: 计算开销、推理延迟与精度帕累托权衡段（Efficiency-Accuracy Pareto Analysis） (\texttt{BP22})}
\noindent\textbf{Expressive Intent:}\enspace 我这一段想说明我们的模型不仅准，而且算得快、显存省，在“精度-延迟/显存”帕累托前沿（Pareto Frontier）上显著优于现有大模型。\par
\noindent\textbf{Cohesion Formula:}\enspace {\small\color{AcadSlate}Hardware Benchmark $\to$ Complexity Reduction \& Memory Gain $\to$ Real-Time Latency Margin $\to$ Pareto Conclusion. 绑定逻辑规则：L010, L023.}\par\vspace{4pt}
\begin{ExampleBox}{Publication-Grade Authentic Text}
Figure 6 plots the MSE-latency Pareto frontier measured on an onboard NVIDIA Orin AGX module. By replacing $O(L^2)$ point-wise self-attention with $O(L \log L)$ spectral filtering and compact patch tokens, DyTS-Net reduces peak GPU memory by 62\% and training wall-clock time by 3.4x compared with iTransformer. At inference time, generating a 16-step visuomotor trajectory requires only 6.8 ms---well below the 20 ms cycle budget of a 50 Hz industrial robot controller---while outperforming a 200 ms DDPM baseline by 11.4\% in task success rate.
\end{ExampleBox}
\begin{TipBox}{Rhetorical Move Breakdown \& Constituent Analysis}
\begin{itemize}[leftmargin=1.5em]
  \item \textbf{S1 [硬件与测试基准 Hardware Setup]:} \textit{Figure 6 plots the MSE-latency Pareto frontier measured on an onboard NVIDIA Orin AGX module.}\par\noindent{\small 对应步骤：S1 [硬件与测试基准 Hardware Setup]：Figure 6 compares the accuracy-efficiency Pareto frontier across all baselines measured on a single NVIDIA RTX 4090 GPU. | 句首锚定旧信息/逻辑视角，主干动词紧跟主语，复杂修饰置于宾语之后（符合尾重原则 L011 与旧新信息流 L001）。}
  \item \textbf{S2 [复杂度理论与实测对比 Complexity \& Runtime]:} \textit{By replacing $O(L^2)$ point-wise self-attention with $O(L \log L)$ spectral filtering and compact patch tokens, DyTS-Net reduces peak GPU memory by 62\% and training wall-clock time by 3.4x compared with iTransformer.}\par\noindent{\small 对应步骤：S2 [复杂度理论与实测对比 Complexity \& Runtime]：By replacing $O(L^2)$ point-wise attention with $O(L \log L)$ spectral filtering and patch aggregation, our model cuts training time per epoch by [X\%] and peak VRAM by [Y\%]. | 句首锚定旧信息/逻辑视角，主干动词紧跟主语，复杂修饰置于宾语之后（符合尾重原则 L011 与旧新信息流 L001）。}
  \item \textbf{S3 [实时闭环部署可行性 Real-Time Viability]:} \textit{At inference time, generating a 16-step visuomotor trajectory requires only 6.8 ms---well below the 20 ms cycle budget of a 50 Hz industrial robot controller---while outperforming a 200 ms DDPM baseline by 11.4\% in task success rate.}\par\noindent{\small 对应步骤：S3 [实时闭环部署可行性 Real-Time Viability]：At inference time, generating a 720-step forecast (or 16-step robot action chunk) requires only [4.2 ms], well below the [20 ms] real-time control deadline. | 句首锚定旧信息/逻辑视角，主干动词紧跟主语，复杂修饰置于宾语之后（符合尾重原则 L011 与旧新信息流 L001）。}
\end{itemize}
\end{TipBox}
\medskip
\subsection*{Paragraph 10: 局限性（Limitations）客观陈述与未来扩展段（Constructive Limitations） (\texttt{BP23})}
\noindent\textbf{Expressive Intent:}\enspace 我这一段想写论文末尾或 NeurIPS/ICLR Checklist 要求的 Limitations 段落：诚实指出当前方法的两个客观边界，但绝不贬损论文的核心贡献。\par
\noindent\textbf{Cohesion Formula:}\enspace {\small\color{AcadSlate}Two Objective Boundaries $\to$ [Boundary 1 + Technical Remedy] $\to$ [Boundary 2 + Stochastic Extension]. 绑定逻辑规则：L017.}\par\vspace{4pt}
\begin{ExampleBox}{Publication-Grade Authentic Text}
Despite its strong empirical performance, our framework exhibits two theoretical boundaries worth noting. First, while our spectral branch scales quasi-linearly with sequence length, computing pairwise Hilbert phase differences scales quadratically with channel count ($O(C^2)$), which increases memory footprint when $C > 2,000$; incorporating anchor-based low-rank graph sparsification provides a direct remedy. Second, our continuous neural-ODE scale tracker assumes piece-wise smooth drift and may smooth over instantaneous step discontinuities caused by sensor recalibration. Extending the continuous drift prior to jump-diffusion stochastic differential equations offers a principled direction for future investigation.
\end{ExampleBox}
\begin{TipBox}{Rhetorical Move Breakdown \& Constituent Analysis}
\begin{itemize}[leftmargin=1.5em]
  \item \textbf{S1 [边界 1:} \textit{Despite its strong empirical performance, our framework exhibits two theoretical boundaries worth noting.}\par\noindent{\small 对应步骤：S1 [边界 1：极端高维或极度稀疏场景 Boundary 1]：Despite its effectiveness across standard benchmarks, our framework currently exhibits two practical boundaries. | 句首锚定旧信息/逻辑视角，主干动词紧跟主语，复杂修饰置于宾语之后（符合尾重原则 L011 与旧新信息流 L001）。}
  \item \textbf{S2 [边界 1 机理分析与解法 Mechanism \& Solution 1]:} \textit{First, while our spectral branch scales quasi-linearly with sequence length, computing pairwise Hilbert phase differences scales quadratically with channel count ($O(C^2)$), which increases memory footprint when $C > 2,000$; incorporating anchor-based low-rank graph sparsification provides a direct remedy.}\par\noindent{\small 对应步骤：S2 [边界 1 机理分析与解法 Mechanism \& Solution 1]：First, when the variate dimension scales beyond $C > 5,000$ unclustered sensors, constructing the cross-channel phase graph incurs noticeable memory overhead; integrating hierarchical sensor clustering offers a natural remedy. | 句首锚定旧信息/逻辑视角，主干动词紧跟主语，复杂修饰置于宾语之后（符合尾重原则 L011 与旧新信息流 L001）。}
  \item \textbf{S3 [边界 2:} \textit{Second, our continuous neural-ODE scale tracker assumes piece-wise smooth drift and may smooth over instantaneous step discontinuities caused by sensor recalibration.}\par\noindent{\small 对应步骤：S3 [边界 2：特定假设限制 Boundary 2]：Second, our continuous scale tracker assumes piece-wise differentiable drift and may under-estimate instantaneous discontinuous jumps caused by hardware resets. | 句首锚定旧信息/逻辑视角，主干动词紧跟主语，复杂修饰置于宾语之后（符合尾重原则 L011 与旧新信息流 L001）。}
  \item \textbf{S4 [未来演进方向 Future Extension]:} \textit{Extending the continuous drift prior to jump-diffusion stochastic differential equations offers a principled direction for future investigation.}\par\noindent{\small 对应步骤：S4 [未来演进方向 Future Extension]：Coupling our neural ODE tracker with jump-diffusion stochastic processes represents a promising avenue for future work. | 句首锚定旧信息/逻辑视角，主干动词紧跟主语，复杂修饰置于宾语之后（符合尾重原则 L011 与旧新信息流 L001）。}
\end{itemize}
\end{TipBox}
\medskip
\subsection*{Paragraph 11: 全文总结与学术影响收尾段（Conclusion \& Broader Outlook） (\texttt{BP24})}
\noindent\textbf{Expressive Intent:}\enspace 我这一段想写整篇小论文的 Conclusion：不简单复制粘贴 Abstract，而是从更高维度总结本文解决的科学矛盾、方法精髓、实证高度与领域启示。\par
\noindent\textbf{Cohesion Formula:}\enspace {\small\color{AcadSlate}Tension Resolved $\to$ Methodological Insight $\to$ Empirical Confirmation $\to$ Generalizable Blueprint. 绑定逻辑规则：L013, L023, L029.}\par\vspace{4pt}
\begin{ExampleBox}{Publication-Grade Authentic Text}
In this work, we resolved the long-standing tension between shift-invariant normalization and physical amplitude preservation in non-stationary time series forecasting. By decoupling frequency-domain phase attention from continuous neural-ODE scale tracking, DyTS-Net captures long-range seasonal dependencies while accurately reconstructing extreme peak magnitudes. Empirical evaluations across eight benchmarks and industrial wind-power telemetry demonstrated MSE reductions of up to 17.4\% alongside a 45\% reduction in computational overhead. Beyond forecasting, this structural-scale disentanglement principle provides a generalizable foundation for anomaly detection and closed-loop predictive control in non-stationary cyber-physical systems.
\end{ExampleBox}
\begin{TipBox}{Rhetorical Move Breakdown \& Constituent Analysis}
\begin{itemize}[leftmargin=1.5em]
  \item \textbf{S1 [核心科学问题闭合 Scientific Closure]:} \textit{In this work, we resolved the long-standing tension between shift-invariant normalization and physical amplitude preservation in non-stationary time series forecasting.}\par\noindent{\small 对应步骤：S1 [核心科学问题闭合 Scientific Closure]：In this paper, we addressed the fundamental tension between [Goal A] and [Goal B] in [Domain]. | 句首锚定旧信息/逻辑视角，主干动词紧跟主语，复杂修饰置于宾语之后（符合尾重原则 L011 与旧新信息流 L001）。}
  \item \textbf{S2 [方法论精髓回顾 Methodological Synthesis]:} \textit{By decoupling frequency-domain phase attention from continuous neural-ODE scale tracking, DyTS-Net captures long-range seasonal dependencies while accurately reconstructing extreme peak magnitudes.}\par\noindent{\small 对应步骤：S2 [方法论精髓回顾 Methodological Synthesis]：By introducing [Model/Mechanism], we demonstrated that [Core Insight] enables [Capability] without incurring [Traditional Penalty]. | 句首锚定旧信息/逻辑视角，主干动词紧跟主语，复杂修饰置于宾语之后（符合尾重原则 L011 与旧新信息流 L001）。}
  \item \textbf{S3 [实证结果高度概括 Empirical Recap]:} \textit{Empirical evaluations across eight benchmarks and industrial wind-power telemetry demonstrated MSE reductions of up to 17.4\% alongside a 45\% reduction in computational overhead.}\par\noindent{\small 对应步骤：S3 [实证结果高度概括 Empirical Recap]：Comprehensive evaluations across [Benchmarks/Hardware] confirmed consistent error reductions of up to [X\%] alongside substantial inference speedups. | 句首锚定旧信息/逻辑视角，主干动词紧跟主语，复杂修饰置于宾语之后（符合尾重原则 L011 与旧新信息流 L001）。}
  \item \textbf{S4 [跨领域迁移展望 Broader Outlook]:} \textit{Beyond forecasting, this structural-scale disentanglement principle provides a generalizable foundation for anomaly detection and closed-loop predictive control in non-stationary cyber-physical systems.}\par\noindent{\small 对应步骤：S4 [跨领域迁移展望 Broader Outlook]：Beyond [Current Task], our [Core Principle] offers a generalizable blueprint for broader sequential decision-making tasks under non-stationary environments. | 句首锚定旧信息/逻辑视角，主干动词紧跟主语，复杂修饰置于宾语之后（符合尾重原则 L011 与旧新信息流 L001）。}
\end{itemize}
\end{TipBox}
\medskip
\clearpage
\chapter{Biomedical \& Life Sciences / 生物医学与生命科学}\label{ch:anthology_19}
\Needspace{6\baselineskip}
本章精选该学科领域及对应文体的完整学术范文。每一篇范文不仅展示母语级发表标准的正文全貌，更清晰拆解全篇的“段落蓝图链条（Structural Blueprint Chain）”，并对每个自然段逐一剖析其核心交际意图、句际衔接逻辑公式（Cohesion Formula）以及每句话的句法成分与推进机制。

\section{FT06: 面向高维传感器网络的稀疏相位图注意力完整小论文（PhaseGraph: Scalable Spatio-Temporal Forecasting）}\label{comp:FT06}
\begin{ConceptBox}{Composition Architecture \& Structural Blueprint Chain}
\textbf{Genre:} \texttt{AI\_Mini\_Paper} \quad \textbf{Paragraph Count:} 11\par
\textbf{Blueprint Chain:} \texttt{BP08 $\to$ BP01 $\to$ BP09 $\to$ BP11 $\to$ BP12 $\to$ BP13 $\to$ BP17 $\to$ BP18 $\to$ BP19 $\to$ BP20 $\to$ BP24}
\end{ConceptBox}
\subsection*{Paragraph 1: 高密度四句式 Abstract 全段构建（Complete 4-Move Abstract） (\texttt{BP08})}
\noindent\textbf{Expressive Intent:}\enspace 我这一段想写出整篇小论文的完整 Abstract（摘要）：在 150--180 词内一气呵成交代背景、核心痛点、方法机制与量化实验结果。\par
\noindent\textbf{Cohesion Formula:}\enspace {\small\color{AcadSlate}Context (S1) $\to$ Gap (S2) $\to$ Method \& Mechanism (S3-S4) $\to$ Quantitative Proof (S5). 绑定逻辑规则：L001, L008, L010, L029.}\par\vspace{4pt}
\begin{ExampleBox}{Publication-Grade Authentic Text}
Recovering missing sensor observations is essential for reliable downstream monitoring in industrial and clinical time series. While diffusion-based imputation models generate plausible trajectories under random missingness, they suffer from high sampling latency and lose global periodic coherence when entire sensor blocks drop out consecutively. We present FlowImp, a conditional flow-matching framework that reconstructs contiguous missing blocks in three ODE solver steps. Specifically, FlowImp anchors straight transport paths on observed spectral priors and enforces cross-channel phase consistency via a lightweight Fourier regularizer. Across six multivariate benchmarks under 50\% block missingness, FlowImp reduces imputation MAE by 19.3\% compared with CSDI while accelerating inference by 24x.
\end{ExampleBox}
\begin{TipBox}{Rhetorical Move Breakdown \& Constituent Analysis}
\begin{itemize}[leftmargin=1.5em]
  \item \textbf{S1 [背景与任务重要性 Context]:} \textit{Recovering missing sensor observations is essential for reliable downstream monitoring in industrial and clinical time series.}\par\noindent{\small 对应步骤：S1 [背景与任务重要性 Context]：1句话锁定任务目标与核心应用场景（≤20词）。 | 句首锚定旧信息/逻辑视角，主干动词紧跟主语，复杂修饰置于宾语之后（符合尾重原则 L011 与旧新信息流 L001）。}
  \item \textbf{S2 [现有局限与痛点 Gap/Bottleneck]:} \textit{While diffusion-based imputation models generate plausible trajectories under random missingness, they suffer from high sampling latency and lose global periodic coherence when entire sensor blocks drop out consecutively.}\par\noindent{\small 对应步骤：S2 [现有局限与痛点 Gap/Bottleneck]：While existing [Paradigms] achieve [Merit], they struggle with [Core Bottleneck] because [Underlying Cause]. | 句首锚定旧信息/逻辑视角，主干动词紧跟主语，复杂修饰置于宾语之后（符合尾重原则 L011 与旧新信息流 L001）。}
  \item \textbf{S3-S4 [本文方法与核心机制 Approach \& Mechanism]:} \textit{We present FlowImp, a conditional flow-matching framework that reconstructs contiguous missing blocks in three ODE solver steps.}\par\noindent{\small 对应步骤：S3-S4 [本文方法与核心机制 Approach \& Mechanism]：In this work, we present [Model], which [Core Idea]. Specifically, [Module 1]..., while [Module 2]... | 句首锚定旧信息/逻辑视角，主干动词紧跟主语，复杂修饰置于宾语之后（符合尾重原则 L011 与旧新信息流 L001）。}
  \item \textbf{S5 [量化实证结果 Quantitative Results]:} \textit{Specifically, FlowImp anchors straight transport paths on observed spectral priors and enforces cross-channel phase consistency via a lightweight Fourier regularizer.}\par\noindent{\small 对应步骤：S5 [量化实证结果 Quantitative Results]：Experiments on [Benchmarks/Hardware] show that [Model] outperforms [Baselines] by [X\%] in [Metric] and achieves [Efficiency/Real-world Result]. | 句首锚定旧信息/逻辑视角，主干动词紧跟主语，复杂修饰置于宾语之后（符合尾重原则 L011 与旧新信息流 L001）。}
  \item \textbf{S5 [量化实证结果 Quantitative Results]:} \textit{Across six multivariate benchmarks under 50\% block missingness, FlowImp reduces imputation MAE by 19.3\% compared with CSDI while accelerating inference by 24x.}\par\noindent{\small 对应步骤：S5 [量化实证结果 Quantitative Results]：Experiments on [Benchmarks/Hardware] show that [Model] outperforms [Baselines] by [X\%] in [Metric] and achieves [Efficiency/Real-world Result]. | 句首锚定旧信息/逻辑视角，主干动词紧跟主语，复杂修饰置于宾语之后（符合尾重原则 L011 与旧新信息流 L001）。}
\end{itemize}
\end{TipBox}
\medskip
\subsection*{Paragraph 2: 确立研究背景与真实工程价值（拒绝空喊口号） (\texttt{BP01})}
\noindent\textbf{Expressive Intent:}\enspace 我这一段想作为 Introduction 第一段：说明多变量时间序列预测或具身智能控制在实际工业系统中的关键作用，以及高可靠性的刚需。\par
\noindent\textbf{Cohesion Formula:}\enspace {\small\color{AcadSlate}Given-to-New 链：Core Task (S1) $\to$ Practical Deployment (S2) $\to$ Real-world Disturbances (S3) $\to$ Consequence of Failure (S4). 绑定逻辑规则：L001, L011, L029.}\par\vspace{4pt}
\begin{ExampleBox}{Publication-Grade Authentic Text}
Multivariate time series forecasting underpins safety-critical scheduling across smart power grids, cloud clusters, and automated manufacturing. In these operational settings, control systems rely on multi-horizon predictions to balance supply fluctuations and prevent cascading overload. However, field sensors routinely operate under abrupt regime shifts, irregular sampling intervals, and high-frequency measurement noise. Failing to accommodate such non-stationary disturbances degrades long-range forecast reliability and triggers costly control oscillations.
\end{ExampleBox}
\begin{TipBox}{Rhetorical Move Breakdown \& Constituent Analysis}
\begin{itemize}[leftmargin=1.5em]
  \item \textbf{S1 [任务锚定 Topic Anchor]:} \textit{Multivariate time series forecasting underpins safety-critical scheduling across smart power grids, cloud clusters, and automated manufacturing.}\par\noindent{\small 对应步骤：S1 [任务锚定 Topic Anchor]：直接用紧凑主谓宾点明研究对象及其在真实系统中的核心功能（≤16词，禁用 In the era of AI）。 | 句首锚定旧信息/逻辑视角，主干动词紧跟主语，复杂修饰置于宾语之后（符合尾重原则 L011 与旧新信息流 L001）。}
  \item \textbf{S2 [应用场景展开 Concrete Scenarios]:} \textit{In these operational settings, control systems rely on multi-horizon predictions to balance supply fluctuations and prevent cascading overload.}\par\noindent{\small 对应步骤：S2 [应用场景展开 Concrete Scenarios]：列举 2--3 个具体工程场景（如电网负荷调度、灵巧机械手装配）及其决策依赖。 | 句首锚定旧信息/逻辑视角，主干动词紧跟主语，复杂修饰置于宾语之后（符合尾重原则 L011 与旧新信息流 L001）。}
  \item \textbf{S3 [苛刻工况约束 Operational Bottleneck]:} \textit{However, field sensors routinely operate under abrupt regime shifts, irregular sampling intervals, and high-frequency measurement noise.}\par\noindent{\small 对应步骤：S3 [苛刻工况约束 Operational Bottleneck]：指出真实部署环境下不可避免的物理/统计挑战（如非平稳漂移、传感器缺失、毫秒级延迟约束）。 | 句首锚定旧信息/逻辑视角，主干动词紧跟主语，复杂修饰置于宾语之后（符合尾重原则 L011 与旧新信息流 L001）。}
  \item \textbf{S4 [研究必要性收束 Imperative Hook]:} \textit{Failing to accommodate such non-stationary disturbances degrades long-range forecast reliability and triggers costly control oscillations.}\par\noindent{\small 对应步骤：S4 [研究必要性收束 Imperative Hook]：点明若无法克服上述约束会造成什么具体系统后果，自然引出下一段对现有方法的审视。 | 句首锚定旧信息/逻辑视角，主干动词紧跟主语，复杂修饰置于宾语之后（符合尾重原则 L011 与旧新信息流 L001）。}
\end{itemize}
\end{TipBox}
\medskip
\subsection*{Paragraph 3: 相关工作分类综述段（拒绝 A did X, B did Y 流水账） (\texttt{BP09})}
\noindent\textbf{Expressive Intent:}\enspace 我这一段想写 Related Work 的一个子方向：把过去三年的 6--8 篇代表性论文按技术路线分成两类进行归纳，并指出它们的共同盲区。\par
\noindent\textbf{Cohesion Formula:}\enspace {\small\color{AcadSlate}Two-Branch Taxonomy $\to$ Family A $\to$ Contrastive Family B $\to$ Shared Limitation. 绑定逻辑规则：L009, L008.}\par\vspace{4pt}
\begin{ExampleBox}{Publication-Grade Authentic Text}
Pre-trained time series foundation models generally follow either continuous patch reconstruction or discrete vocabulary tokenization. Continuous reconstruction frameworks (e.g., PatchTST, MOMENT) pre-train encoders via masked patch regression, preserving fine-grained numerical resolution across continuous domains. Conversely, discrete tokenization approaches (e.g., Chronos) quantize scaled observations into fixed bins to leverage cross-entropy language modeling objectives. Despite their zero-shot transferability, both paradigms treat univariate channels independently during pre-training, ignoring physical cross-variable conservation laws that govern multivariate engineering systems.
\end{ExampleBox}
\begin{TipBox}{Rhetorical Move Breakdown \& Constituent Analysis}
\begin{itemize}[leftmargin=1.5em]
  \item \textbf{S1 [分类总纲 Taxonomic Thesis]:} \textit{Pre-trained time series foundation models generally follow either continuous patch reconstruction or discrete vocabulary tokenization.}\par\noindent{\small 对应步骤：S1 [分类总纲 Taxonomic Thesis]：Existing methods for [Task] can be categorized into two main families: [Family A] and [Family B]. | 句首锚定旧信息/逻辑视角，主干动词紧跟主语，复杂修饰置于宾语之后（符合尾重原则 L011 与旧新信息流 L001）。}
  \item \textbf{S2 [流派 A 归纳与代表作 Family A Synthesis]:} \textit{Continuous reconstruction frameworks (e.g., PatchTST, MOMENT) pre-train encoders via masked patch regression, preserving fine-grained numerical resolution across continuous domains.}\par\noindent{\small 对应步骤：S2 [流派 A 归纳与代表作 Family A Synthesis]：[Family A] approaches (e.g., [Models]) focus on [Mechanism A], effectively addressing [Sub-problem 1]. | 句首锚定旧信息/逻辑视角，主干动词紧跟主语，复杂修饰置于宾语之后（符合尾重原则 L011 与旧新信息流 L001）。}
  \item \textbf{S3 [流派 B 归纳与对比 Family B Synthesis]:} \textit{Conversely, discrete tokenization approaches (e.g., Chronos) quantize scaled observations into fixed bins to leverage cross-entropy language modeling objectives.}\par\noindent{\small 对应步骤：S3 [流派 B 归纳与对比 Family B Synthesis]：By contrast, [Family B] architectures (e.g., [Models]) reformulate the problem via [Mechanism B] to reduce [Cost/Complexity]. | 句首锚定旧信息/逻辑视角，主干动词紧跟主语，复杂修饰置于宾语之后（符合尾重原则 L011 与旧新信息流 L001）。}
  \item \textbf{S4 [共同盲区提炼 Shared Blind Spot]:} \textit{Despite their zero-shot transferability, both paradigms treat univariate channels independently during pre-training, ignoring physical cross-variable conservation laws that govern multivariate engineering systems.}\par\noindent{\small 对应步骤：S4 [共同盲区提炼 Shared Blind Spot]：Despite their respective strengths, both families overlook [Critical Factor], which degrades performance under [Target Scenario]. | 句首锚定旧信息/逻辑视角，主干动词紧跟主语，复杂修饰置于宾语之后（符合尾重原则 L011 与旧新信息流 L001）。}
\end{itemize}
\end{TipBox}
\medskip
\subsection*{Paragraph 4: 问题形式化定义与符号约定段（Problem Formulation \& Notation） (\texttt{BP11})}
\noindent\textbf{Expressive Intent:}\enspace 我这一段想在 Methodology 开篇（Section 3.1 Problem Formulation），严谨、简洁地定义输入张量、输出目标、约束条件与优化目标。\par
\noindent\textbf{Cohesion Formula:}\enspace {\small\color{AcadSlate}Input Space $\to$ Auxiliary Conditions $\to$ Target Mapping $\to$ Crucial Distributional Constraint. 绑定逻辑规则：L014, L011.}\par\vspace{4pt}
\begin{ExampleBox}{Publication-Grade Authentic Text}
Consider a multivariate time series $\mathbf{X}_{1:L} = [x_1, \dots, x_L]^\top \in \mathbb{R}^{L \times C}$ recording $C$ correlated sensors over lookback horizon $L$. Given an binary observation mask $\mathbf{M} \in \{0, 1\}^{L \times C}$ indicating irregular sensor dropouts, our objective is to learn a generative mapping $p_\theta(\mathbf{Y}_{L+1:L+H} \mid \mathbf{X}_{1:L} \odot \mathbf{M})$ over future horizon $H$. Crucially, unlike standard settings that assume time-invariant marginal distributions, we explicitly account for non-stationary regime shifts where $\mathbb{E}[x_t]$ and $\text{Cov}(x_t)$ evolve continuously between the lookback and forecasting windows.
\end{ExampleBox}
\begin{TipBox}{Rhetorical Move Breakdown \& Constituent Analysis}
\begin{itemize}[leftmargin=1.5em]
  \item \textbf{S1 [任务总体形式化 Task Definition]:} \textit{Consider a multivariate time series $\mathbf{X}_{1:L} = [x_1, \dots, x_L]^\top \in \mathbb{R}^{L \times C}$ recording $C$ correlated sensors over lookback horizon $L$.}\par\noindent{\small 对应步骤：S1 [任务总体形式化 Task Definition]：Consider a [System/Dataset] observed over a lookback window of length $L$, denoted as $\mathbf{X} = [x_1, \dots, x_L] \in \mathbb{R}^{L \times C}$, where $C$ represents [Meaning]. | 句首锚定旧信息/逻辑视角，主干动词紧跟主语，复杂修饰置于宾语之后（符合尾重原则 L011 与旧新信息流 L001）。}
  \item \textbf{S2 [补充条件或模态输入 Auxiliary/Condition Inputs]:} \textit{Given an binary observation mask $\mathbf{M} \in \{0, 1\}^{L \times C}$ indicating irregular sensor dropouts, our objective is to learn a generative mapping $p_\theta(\mathbf{Y}_{L+1:L+H} \mid \mathbf{X}_{1:L} \odot \mathbf{M})$ over future horizon $H$.}\par\noindent{\small 对应步骤：S2 [补充条件或模态输入 Auxiliary/Condition Inputs]：At each step $t$, the system also receives [Condition/Observation] $z_t \in \mathcal{Z}$ under [Noise/Missingness/Dynamic Regime]. | 句首锚定旧信息/逻辑视角，主干动词紧跟主语，复杂修饰置于宾语之后（符合尾重原则 L011 与旧新信息流 L001）。}
  \item \textbf{S3 [预测/控制目标 Objective]:} \textit{Crucially, unlike standard settings that assume time-invariant marginal distributions, we explicitly account for non-stationary regime shifts where $\mathbb{E}[x_t]$ and $\text{Cov}(x_t)$ evolve continuously between the lookback and forecasting windows.}\par\noindent{\small 对应步骤：S3 [预测/控制目标 Objective]：Our goal is to learn a parameterized mapping $f_\theta(\cdot)$ that predicts future states $\hat{\mathbf{Y}} \in \mathbb{R}^{H \times C}$ over horizon $H$ (or action trajectory $a_{t:t+K}$). | 句首锚定旧信息/逻辑视角，主干动词紧跟主语，复杂修饰置于宾语之后（符合尾重原则 L011 与旧新信息流 L001）。}
\end{itemize}
\end{TipBox}
\medskip
\subsection*{Paragraph 5: 整体模型架构数据流俯瞰段（Framework Overview \& Pipeline Walkthrough） (\texttt{BP12})}
\noindent\textbf{Expressive Intent:}\enspace 我这一段想配合论文里的 Figure 2（模型总架构图），用一段话把输入张量从左到右流经各个模块直到输出预测结果的完整全貌讲清楚。\par
\noindent\textbf{Cohesion Formula:}\enspace {\small\color{AcadSlate}严格的链式数据流（$\mathbf{X} \to \mathbf{Z}^{(1)} \to \mathbf{Z}^{(2)} \to \hat{\mathbf{Y}}$），每一句句首承接上一句的输出张量。绑定逻辑规则：L001, L003, L004.}\par\vspace{4pt}
\begin{ExampleBox}{Publication-Grade Authentic Text}
Figure 2 illustrates the overall architecture of DyTS-Net, which maps raw sensor streams to multi-horizon forecasts through three sequential stages. Given an input sequence $\mathbf{X}$, the Adaptive De-stationary Tokenizer first separates local trend-scale statistics from normalized sub-series patches $\mathbf{Z}_0$. These normalized patches then pass through stacked Spectral-Spatial Mixing blocks, which jointly filter high-frequency channel noise in the Fourier domain and propagate cross-sensor phase dependencies via a sparse directed graph. Finally, the Scale-Guided Decoder fuses the refined structural representations with the projected neural-ODE scale trajectory to reconstruct the multi-horizon output $\hat{\mathbf{Y}}$.
\end{ExampleBox}
\begin{TipBox}{Rhetorical Move Breakdown \& Constituent Analysis}
\begin{itemize}[leftmargin=1.5em]
  \item \textbf{S1 [总图锚定与设计哲学 Visual Anchor \& Philosophy]:} \textit{Figure 2 illustrates the overall architecture of DyTS-Net, which maps raw sensor streams to multi-horizon forecasts through three sequential stages.}\par\noindent{\small 对应步骤：S1 [总图锚定与设计哲学 Visual Anchor \& Philosophy]：Figure 2 illustrates the overall architecture of [Model], which processes [Input] through three sequential stages: [Stage 1], [Stage 2], and [Stage 3]. | 句首锚定旧信息/逻辑视角，主干动词紧跟主语，复杂修饰置于宾语之后（符合尾重原则 L011 与旧新信息流 L001）。}
  \item \textbf{S2 [阶段 1 数据流 Stage 1 Flow]:} \textit{Given an input sequence $\mathbf{X}$, the Adaptive De-stationary Tokenizer first separates local trend-scale statistics from normalized sub-series patches $\mathbf{Z}_0$.}\par\noindent{\small 对应步骤：S2 [阶段 1 数据流 Stage 1 Flow]：Given raw input $\mathbf{X}$, the [Module 1] first transforms [Raw Input] into [Intermediate Representation $\mathbf{Z}^{(1)}$] to [Functional Purpose 1]. | 句首锚定旧信息/逻辑视角，主干动词紧跟主语，复杂修饰置于宾语之后（符合尾重原则 L011 与旧新信息流 L001）。}
  \item \textbf{S3 [阶段 2 数据流 Stage 2 Flow]:} \textit{These normalized patches then pass through stacked Spectral-Spatial Mixing blocks, which jointly filter high-frequency channel noise in the Fourier domain and propagate cross-sensor phase dependencies via a sparse directed graph.}\par\noindent{\small 对应步骤：S3 [阶段 2 数据流 Stage 2 Flow]：Next, these [Summary Noun: representations $\mathbf{Z}^{(1)}$] are fed into [Module 2], which [Mechanism 2] to produce [Refined State $\mathbf{Z}^{(2)}$]. | 句首锚定旧信息/逻辑视角，主干动词紧跟主语，复杂修饰置于宾语之后（符合尾重原则 L011 与旧新信息流 L001）。}
  \item \textbf{S4 [阶段 3 解码与闭环输出 Stage 3 Decoding]:} \textit{Finally, the Scale-Guided Decoder fuses the refined structural representations with the projected neural-ODE scale trajectory to reconstruct the multi-horizon output $\hat{\mathbf{Y}}$.}\par\noindent{\small 对应步骤：S4 [阶段 3 解码与闭环输出 Stage 3 Decoding]：Finally, the [Module 3] integrates $\mathbf{Z}^{(2)}$ with [Prior/Scale Cue] to output the final [Prediction/Action Trajectory]. | 句首锚定旧信息/逻辑视角，主干动词紧跟主语，复杂修饰置于宾语之后（符合尾重原则 L011 与旧新信息流 L001）。}
\end{itemize}
\end{TipBox}
\medskip
\subsection*{Paragraph 6: 核心子模块：直觉动机与数学设计无缝衔接段（Intuition-before-Math） (\texttt{BP13})}
\noindent\textbf{Expressive Intent:}\enspace 我这一段想详细介绍我们提出的某个核心创新模块：先用两句话讲清楚为什么直接用传统算子会出问题，再顺理成章推导出我们的新公式。\par
\noindent\textbf{Cohesion Formula:}\enspace {\small\color{AcadSlate}Pathology of Standard Operator $\to$ Physical Intuition $\to$ Equation Integration $\to$ Theoretical Guarantee. 绑定逻辑规则：L014, L003.}\par\vspace{4pt}
\begin{ExampleBox}{Publication-Grade Authentic Text}
Standard dot-product attention measures token similarity purely by instantaneous amplitude alignment, often matching noisy spikes across unrelated time steps. To enforce temporal coherence, we observe that genuine physical interactions between industrial sensors exhibit consistent phase lag across dominant frequencies. Formally, given query and key projections $Q, K \in \mathbb{R}^{N \times d}$, we modulate the attention logits via a phase-alignment mask: $\mathcal{A}(Q, K, V) = \text{softmax}\big(\frac{QK^\top}{\sqrt{d}} - \lambda \|\Delta \Phi_{Q,K}\|_1\big)V$, where $\Delta \Phi_{Q,K} \in [-\pi, \pi]^{N \times N}$ denotes the instantaneous Hilbert phase difference matrix and $\lambda > 0$ controls phase selectivity. Penalizing large phase discrepancies prevents spurious attention weights between out-of-phase noise bursts.
\end{ExampleBox}
\begin{TipBox}{Rhetorical Move Breakdown \& Constituent Analysis}
\begin{itemize}[leftmargin=1.5em]
  \item \textbf{S1 [传统算子缺陷 Motivation]:} \textit{Standard dot-product attention measures token similarity purely by instantaneous amplitude alignment, often matching noisy spikes across unrelated time steps.}\par\noindent{\small 对应步骤：S1 [传统算子缺陷 Motivation]：Standard [Conventional Operator, e.g., softmax self-attention] computes [Operation], which suffers from [Pathology] when applied to [Domain Condition]. | 句首锚定旧信息/逻辑视角，主干动词紧跟主语，复杂修饰置于宾语之后（符合尾重原则 L011 与旧新信息流 L001）。}
  \item \textbf{S2 [物理/统计直觉 Physical Intuition]:} \textit{To enforce temporal coherence, we observe that genuine physical interactions between industrial sensors exhibit consistent phase lag across dominant frequencies.}\par\noindent{\small 对应步骤：S2 [物理/统计直觉 Physical Intuition]：To overcome this pathology, we observe that [Key Domain Property] can be effectively regularized by [Intuitive Operation]. | 句首锚定旧信息/逻辑视角，主干动词紧跟主语，复杂修饰置于宾语之后（符合尾重原则 L011 与旧新信息流 L001）。}
  \item \textbf{S3 [数学公式定义 Mathematical Formulation]:} \textit{Formally, given query and key projections $Q, K \in \mathbb{R}^{N \times d}$, we modulate the attention logits via a phase-alignment mask: $\mathcal{A}(Q, K, V) = \text{softmax}\big(\frac{QK^\top}{\sqrt{d}} - \lambda \|\Delta \Phi_{Q,K}\|_1\big)V$, where $\Delta \Phi_{Q,K} \in [-\pi, \pi]^{N \times N}$ denotes the instantaneous Hilbert phase difference matrix and $\lambda > 0$ controls phase selectivity.}\par\noindent{\small 对应步骤：S3 [数学公式定义 Mathematical Formulation]：Formally, given input representation $H \in \mathbb{R}^{N \times d}$, we formulate the [Module Name] as: [Equation], where [Symbol Explanations]. | 句首锚定旧信息/逻辑视角，主干动词紧跟主语，复杂修饰置于宾语之后（符合尾重原则 L011 与旧新信息流 L001）。}
  \item \textbf{S4 [理论/计算性质 Guarantee]:} \textit{Penalizing large phase discrepancies prevents spurious attention weights between out-of-phase noise bursts.}\par\noindent{\small 对应步骤：S4 [理论/计算性质 Guarantee]：By construction, this formulation guarantees [Desirable Property, e.g., $O(N)$ complexity / translation invariance]. | 句首锚定旧信息/逻辑视角，主干动词紧跟主语，复杂修饰置于宾语之后（符合尾重原则 L011 与旧新信息流 L001）。}
\end{itemize}
\end{TipBox}
\medskip
\subsection*{Paragraph 7: 实验设置、数据集特性与强基线选型说明段（Experimental Setup \& Baselines） (\texttt{BP17})}
\noindent\textbf{Expressive Intent:}\enspace 我这一段想在 Experiments 开篇说明：我们为什么选这几类数据集和这几类基线模型，证明我们的实验设计是全面、公平且直击痛点的。\par
\noindent\textbf{Cohesion Formula:}\enspace {\small\color{AcadSlate}Dataset Diversity $\to$ Challenging Regime Highlights $\to$ Categorized Strong Baselines $\to$ Strict Fairness Protocol. 绑定逻辑规则：L009, L010.}\par\vspace{4pt}
\begin{ExampleBox}{Publication-Grade Authentic Text}
To evaluate robustness across diverse non-stationary regimes, we conduct experiments on seven public benchmarks spanning energy grid load (ETT, ECL), meteorological telemetry (Weather), and urban traffic flow (Traffic). Notably, Weather contains abrupt atmospheric pressure transitions that test non-stationary adaptability, whereas Traffic spans 862 synchronized sensors that challenge high-dimensional scalability. We benchmark our approach against 11 state-of-the-art methods across three paradigms: channel-independent patch Transformers (PatchTST), variate-mixing architectures (iTransformer, Crossformer), and frequency-domain forecasters (TimesNet, FEDformer). All baselines are reproduced under identical lookback windows ($L=96$) and prediction horizons ($H \in \{96, 192, 336, 720\}$) using standard MSE and MAE metrics.
\end{ExampleBox}
\begin{TipBox}{Rhetorical Move Breakdown \& Constituent Analysis}
\begin{itemize}[leftmargin=1.5em]
  \item \textbf{S1 [数据集选型逻辑 Dataset Selection Rationale]:} \textit{To evaluate robustness across diverse non-stationary regimes, we conduct experiments on seven public benchmarks spanning energy grid load (ETT, ECL), meteorological telemetry (Weather), and urban traffic flow (Traffic).}\par\noindent{\small 对应步骤：S1 [数据集选型逻辑 Dataset Selection Rationale]：To evaluate performance across diverse non-stationary regimes, we conduct experiments on [N] benchmarks spanning [Domain 1], [Domain 2], and [Domain 3]. | 句首锚定旧信息/逻辑视角，主干动词紧跟主语，复杂修饰置于宾语之后（符合尾重原则 L011 与旧新信息流 L001）。}
  \item \textbf{S2 [极端/困难特性突出 Challenging Dataset Properties]:} \textit{Notably, Weather contains abrupt atmospheric pressure transitions that test non-stationary adaptability, whereas Traffic spans 862 synchronized sensors that challenge high-dimensional scalability.}\par\noindent{\small 对应步骤：S2 [极端/困难特性突出 Challenging Dataset Properties]：Notably, [Dataset A] exhibits [Hard Property 1, e.g., strong covariate shift], whereas [Dataset B] tests scalability across [Hard Property 2, e.g., 862 correlated channels]. | 句首锚定旧信息/逻辑视角，主干动词紧跟主语，复杂修饰置于宾语之后（符合尾重原则 L011 与旧新信息流 L001）。}
  \item \textbf{S3 [基线流派全覆盖 Baseline Family Coverage]:} \textit{We benchmark our approach against 11 state-of-the-art methods across three paradigms: channel-independent patch Transformers (PatchTST), variate-mixing architectures (iTransformer, Crossformer), and frequency-domain forecasters (TimesNet, FEDformer).}\par\noindent{\small 对应步骤：S3 [基线流派全覆盖 Baseline Family Coverage]：We compare our method against [K] competitive baselines grouped into three families: [Family 1], [Family 2], and [Family 3]. | 句首锚定旧信息/逻辑视角，主干动词紧跟主语，复杂修饰置于宾语之后（符合尾重原则 L011 与旧新信息流 L001）。}
  \item \textbf{S4 [公平评测协议 Fair Evaluation Protocol]:} \textit{All baselines are reproduced under identical lookback windows ($L=96$) and prediction horizons ($H \in \{96, 192, 336, 720\}$) using standard MSE and MAE metrics.}\par\noindent{\small 对应步骤：S4 [公平评测协议 Fair Evaluation Protocol]：All models are evaluated under identical lookback windows ($L=96$), prediction horizons ($H \in \{96, 192, 336, 720\}$), and hyperparameter search budgets. | 句首锚定旧信息/逻辑视角，主干动词紧跟主语，复杂修饰置于宾语之后（符合尾重原则 L011 与旧新信息流 L001）。}
\end{itemize}
\end{TipBox}
\medskip
\subsection*{Paragraph 8: 主实验结果整体对比与核心胜出原因剖析段（Main Results Analysis） (\texttt{BP18})}
\noindent\textbf{Expressive Intent:}\enspace 我这一段想解读主实验大表（Table 1）：既报告全局胜率，又挑出最有代表性的对比案例，解释为什么竞争对手在那里跌倒而我们稳赢。\par
\noindent\textbf{Cohesion Formula:}\enspace {\small\color{AcadSlate}Global Win Rate $\to$ Advantage over Family A $\to$ Resilience at Long Horizon vs. Family B $\to$ Mechanistic Takeaway. 绑定逻辑规则：L015, L013, L010.}\par\vspace{4pt}
\begin{ExampleBox}{Publication-Grade Authentic Text}
Table 1 reports the multivariate forecasting results, where DyTS-Net secures the lowest error in 25 out of 28 dataset-horizon configurations. Compared with channel-independent PatchTST, DyTS-Net reduces average MSE by 14.2\% on the 862-channel Traffic benchmark, where modeling spatial-temporal congestion propagation across adjacent road sensors is critical. Furthermore, while variate-mixing iTransformer performs competitively at $H=96$, its MSE escalates by 68\% at $H=720$ on volatile Weather series as unregularized cross-channel attention overfits transient wind gusts. By decoupling stationary phase coupling from neural-ODE scale projection, DyTS-Net suppresses long-horizon drift and maintains an 11.6\% margin over the strongest baseline at $H=720$.
\end{ExampleBox}
\begin{TipBox}{Rhetorical Move Breakdown \& Constituent Analysis}
\begin{itemize}[leftmargin=1.5em]
  \item \textbf{S1 [全局统治力概述 Overall Dominance]:} \textit{Table 1 reports the multivariate forecasting results, where DyTS-Net secures the lowest error in 25 out of 28 dataset-horizon configurations.}\par\noindent{\small 对应步骤：S1 [全局统治力概述 Overall Dominance]：Table 1 summarizes the multivariate forecasting results, where [Our Model] achieves the lowest error in [X] out of [Y] dataset-horizon configurations. | 句首锚定旧信息/逻辑视角，主干动词紧跟主语，复杂修饰置于宾语之后（符合尾重原则 L011 与旧新信息流 L001）。}
  \item \textbf{S2 [对比流派 1 的差距与原因 Contrast vs. Family 1]:} \textit{Compared with channel-independent PatchTST, DyTS-Net reduces average MSE by 14.2\% on the 862-channel Traffic benchmark, where modeling spatial-temporal congestion propagation across adjacent road sensors is critical.}\par\noindent{\small 对应步骤：S2 [对比流派 1 的差距与原因 Contrast vs. Family 1]：Compared with channel-independent baselines such as [Baseline A], our method reduces average MSE by [X\%] on [Dataset], where strong cross-sensor coupling is critical. | 句首锚定旧信息/逻辑视角，主干动词紧跟主语，复杂修饰置于宾语之后（符合尾重原则 L011 与旧新信息流 L001）。}
  \item \textbf{S3 [对比流派 2 在困难场景下的崩溃 Contrast vs. Family 2 under Hard Regime]:} \textit{Furthermore, while variate-mixing iTransformer performs competitively at $H=96$, its MSE escalates by 68\% at $H=720$ on volatile Weather series as unregularized cross-channel attention overfits transient wind gusts.}\par\noindent{\small 对应步骤：S3 [对比流派 2 在困难场景下的崩溃 Contrast vs. Family 2 under Hard Regime]：Furthermore, while [Baseline B] performs competitively at short horizons ($H=96$), its error escalates by [Y\%] at $H=720$ due to [Failure Mechanism]. | 句首锚定旧信息/逻辑视角，主干动词紧跟主语，复杂修饰置于宾语之后（符合尾重原则 L011 与旧新信息流 L001）。}
  \item \textbf{S4 [回扣本文核心机制 Mechanistic Confirmation]:} \textit{By decoupling stationary phase coupling from neural-ODE scale projection, DyTS-Net suppresses long-horizon drift and maintains an 11.6\% margin over the strongest baseline at $H=720$.}\par\noindent{\small 对应步骤：S4 [回扣本文核心机制 Mechanistic Confirmation]：By contrast, our [Core Module] bounds long-range drift, maintaining a [Z\%] margin over the second-best baseline at $H=720$. | 句首锚定旧信息/逻辑视角，主干动词紧跟主语，复杂修饰置于宾语之后（符合尾重原则 L011 与旧新信息流 L001）。}
\end{itemize}
\end{TipBox}
\medskip
\subsection*{Paragraph 9: 消融实验逐组件归因论证段（Counterfactual Ablation Study） (\texttt{BP19})}
\noindent\textbf{Expressive Intent:}\enspace 我这一段想汇报消融实验（Table 2）：证明我们的模块 A、模块 B 和正则项 C 每一个都不可或缺，且各自负责解决不同的特定子问题。\par
\noindent\textbf{Cohesion Formula:}\enspace {\small\color{AcadSlate}Ablation Overview $\to$ w/o Module A (Specific Failure) $\to$ w/o Module B (Specific Failure) $\to$ Cross-Module Synergy. 绑定逻辑规则：L016, L007.}\par\vspace{4pt}
\begin{ExampleBox}{Publication-Grade Authentic Text}
To isolate the contribution of each design choice, Table 2 compares full DyTS-Net against three ablated variants across ETTm2 and ECL. Removing the neural-ODE scale tracker (w/o Scale-ODE) triggers the sharpest degradation during seasonal transition months on ETTm2, raising 720-step MSE by 19.4\% as the decoder fails to track shifting baseline temperatures. Conversely, replacing phase-aligned graph attention with standard dot-product attention (w/o Phase-Attn) degrades performance most severely on the 321-channel ECL grid (+13.8\% MSE), where unaligned high-frequency load spikes corrupt cross-substation aggregation. Jointly deploying both modules yields complementary gains, as accurate scale tracking prevents amplitude distortion from biasing phase estimation.
\end{ExampleBox}
\begin{TipBox}{Rhetorical Move Breakdown \& Constituent Analysis}
\begin{itemize}[leftmargin=1.5em]
  \item \textbf{S1 [消融实验总括 Ablation Setup]:} \textit{To isolate the contribution of each design choice, Table 2 compares full DyTS-Net against three ablated variants across ETTm2 and ECL.}\par\noindent{\small 对应步骤：S1 [消融实验总括 Ablation Setup]：To isolate the contribution of each architectural component, we evaluate four ablated variants in Table 2 by progressively removing [Mod A], [Mod B], and [Loss C]. | 句首锚定旧信息/逻辑视角，主干动词紧跟主语，复杂修饰置于宾语之后（符合尾重原则 L011 与旧新信息流 L001）。}
  \item \textbf{S2 [移除模块 A 的冲击与归因 Impact of Removing Mod A]:} \textit{Removing the neural-ODE scale tracker (w/o Scale-ODE) triggers the sharpest degradation during seasonal transition months on ETTm2, raising 720-step MSE by 19.4\% as the decoder fails to track shifting baseline temperatures.}\par\noindent{\small 对应步骤：S2 [移除模块 A 的冲击与归因 Impact of Removing Mod A]：Removing [Mod A] (w/o A) causes the largest degradation on [Non-stationary/Volatile Subset], increasing MSE by [X\%] because the model loses [Specific Capability A]. | 句首锚定旧信息/逻辑视角，主干动词紧跟主语，复杂修饰置于宾语之后（符合尾重原则 L011 与旧新信息流 L001）。}
  \item \textbf{S3 [移除模块 B 的冲击与归因 Impact of Removing Mod B]:} \textit{Conversely, replacing phase-aligned graph attention with standard dot-product attention (w/o Phase-Attn) degrades performance most severely on the 321-channel ECL grid (+13.8\% MSE), where unaligned high-frequency load spikes corrupt cross-substation aggregation.}\par\noindent{\small 对应步骤：S3 [移除模块 B 的冲击与归因 Impact of Removing Mod B]：Similarly, replacing [Mod B] with standard [Naive Operator] degrades performance by [Y\%] on high-dimensional benchmarks, confirming its role in [Specific Capability B]. | 句首锚定旧信息/逻辑视角，主干动词紧跟主语，复杂修饰置于宾语之后（符合尾重原则 L011 与旧新信息流 L001）。}
  \item \textbf{S4 [协同效应总结 Synergistic Conclusion]:} \textit{Jointly deploying both modules yields complementary gains, as accurate scale tracking prevents amplitude distortion from biasing phase estimation.}\par\noindent{\small 对应步骤：S4 [协同效应总结 Synergistic Conclusion]：Combining all components yields super-additive gains, as [Mod A] stabilizes input statistics for [Mod B]'s relational reasoning. | 句首锚定旧信息/逻辑视角，主干动词紧跟主语，复杂修饰置于宾语之后（符合尾重原则 L011 与旧新信息流 L001）。}
\end{itemize}
\end{TipBox}
\medskip
\subsection*{Paragraph 10: 超参数敏感性与极端噪声鲁棒性分析段（Sensitivity \& Robustness Analysis） (\texttt{BP20})}
\noindent\textbf{Expressive Intent:}\enspace 我这一段想回应审稿人最常问的：你的模型是不是对超参数很敏感？在缺失值或传感器强噪声下会不会突然崩盘？\par
\noindent\textbf{Cohesion Formula:}\enspace {\small\color{AcadSlate}Goal $\to$ Wide Hyperparameter Plateau $\to$ Stress-Test Comparison $\to$ Mechanistic Explanation. 绑定逻辑规则：L010, L019.}\par\vspace{4pt}
\begin{ExampleBox}{Publication-Grade Authentic Text}
Figure 4 evaluates hyperparameter sensitivity and resilience against severe observation corruption. Varying the phase regularization weight $\lambda$ across $[0.05, 0.50]$ alters validation MSE by less than 1.4\%, indicating a wide, tuning-friendly operating plateau around $\lambda=0.15$. Under stress tests where 30\% of sensor readings are randomly masked or corrupted by impulse noise ($\text{SNR} = 5\text{ dB}$), point-wise Transformer baselines suffer a 34.2\% error surge, whereas DyTS-Net degrades by only 4.8\%. This robustness stems from our spectral energy gate, which automatically attenuates high-entropy corrupted channels prior to cross-variate message passing.
\end{ExampleBox}
\begin{TipBox}{Rhetorical Move Breakdown \& Constituent Analysis}
\begin{itemize}[leftmargin=1.5em]
  \item \textbf{S1 [分析目的 Objective]:} \textit{Figure 4 evaluates hyperparameter sensitivity and resilience against severe observation corruption.}\par\noindent{\small 对应步骤：S1 [分析目的 Objective]：We further examine sensitivity to key hyperparameters and resilience against severe sensor corruption in Figure 4. | 句首锚定旧信息/逻辑视角，主干动词紧跟主语，复杂修饰置于宾语之后（符合尾重原则 L011 与旧新信息流 L001）。}
  \item \textbf{S2 [超参数平坦区间 Plateau Stability]:} \textit{Varying the phase regularization weight $\lambda$ across $[0.05, 0.50]$ alters validation MSE by less than 1.4\%, indicating a wide, tuning-friendly operating plateau around $\lambda=0.15$.}\par\noindent{\small 对应步骤：S2 [超参数平坦区间 Plateau Stability]：Varying [Hyperparameter $\lambda$] across two orders of magnitude ($[0.01, 1.0]$) alters validation MSE by less than [1.8\%], demonstrating a wide stable operating plateau. | 句首锚定旧信息/逻辑视角，主干动词紧跟主语，复杂修饰置于宾语之后（符合尾重原则 L011 与旧新信息流 L001）。}
  \item \textbf{S3 [极端边缘工况压力测试 Stress Test under Corruption]:} \textit{Under stress tests where 30\% of sensor readings are randomly masked or corrupted by impulse noise ($\text{SNR} = 5\text{ dB}$), point-wise Transformer baselines suffer a 34.2\% error surge, whereas DyTS-Net degrades by only 4.8\%.}\par\noindent{\small 对应步骤：S3 [极端边缘工况压力测试 Stress Test under Corruption]：Under synthetic sensor dropouts up to [40\%] and Gaussian noise ($\sigma=0.3$), baseline accuracy plummets by over [35\%], whereas our model degrades by less than [6.5\%]. | 句首锚定旧信息/逻辑视角，主干动词紧跟主语，复杂修饰置于宾语之后（符合尾重原则 L011 与旧新信息流 L001）。}
  \item \textbf{S4 [鲁棒性物理归因 Robustness Mechanism]:} \textit{This robustness stems from our spectral energy gate, which automatically attenuates high-entropy corrupted channels prior to cross-variate message passing.}\par\noindent{\small 对应步骤：S4 [鲁棒性物理归因 Robustness Mechanism]：This resilience stems from [Specific Mechanism], which filters out corrupted channels before global aggregation. | 句首锚定旧信息/逻辑视角，主干动词紧跟主语，复杂修饰置于宾语之后（符合尾重原则 L011 与旧新信息流 L001）。}
\end{itemize}
\end{TipBox}
\medskip
\subsection*{Paragraph 11: 全文总结与学术影响收尾段（Conclusion \& Broader Outlook） (\texttt{BP24})}
\noindent\textbf{Expressive Intent:}\enspace 我这一段想写整篇小论文的 Conclusion：不简单复制粘贴 Abstract，而是从更高维度总结本文解决的科学矛盾、方法精髓、实证高度与领域启示。\par
\noindent\textbf{Cohesion Formula:}\enspace {\small\color{AcadSlate}Tension Resolved $\to$ Methodological Insight $\to$ Empirical Confirmation $\to$ Generalizable Blueprint. 绑定逻辑规则：L013, L023, L029.}\par\vspace{4pt}
\begin{ExampleBox}{Publication-Grade Authentic Text}
In this work, we resolved the long-standing tension between shift-invariant normalization and physical amplitude preservation in non-stationary time series forecasting. By decoupling frequency-domain phase attention from continuous neural-ODE scale tracking, DyTS-Net captures long-range seasonal dependencies while accurately reconstructing extreme peak magnitudes. Empirical evaluations across eight benchmarks and industrial wind-power telemetry demonstrated MSE reductions of up to 17.4\% alongside a 45\% reduction in computational overhead. Beyond forecasting, this structural-scale disentanglement principle provides a generalizable foundation for anomaly detection and closed-loop predictive control in non-stationary cyber-physical systems.
\end{ExampleBox}
\begin{TipBox}{Rhetorical Move Breakdown \& Constituent Analysis}
\begin{itemize}[leftmargin=1.5em]
  \item \textbf{S1 [核心科学问题闭合 Scientific Closure]:} \textit{In this work, we resolved the long-standing tension between shift-invariant normalization and physical amplitude preservation in non-stationary time series forecasting.}\par\noindent{\small 对应步骤：S1 [核心科学问题闭合 Scientific Closure]：In this paper, we addressed the fundamental tension between [Goal A] and [Goal B] in [Domain]. | 句首锚定旧信息/逻辑视角，主干动词紧跟主语，复杂修饰置于宾语之后（符合尾重原则 L011 与旧新信息流 L001）。}
  \item \textbf{S2 [方法论精髓回顾 Methodological Synthesis]:} \textit{By decoupling frequency-domain phase attention from continuous neural-ODE scale tracking, DyTS-Net captures long-range seasonal dependencies while accurately reconstructing extreme peak magnitudes.}\par\noindent{\small 对应步骤：S2 [方法论精髓回顾 Methodological Synthesis]：By introducing [Model/Mechanism], we demonstrated that [Core Insight] enables [Capability] without incurring [Traditional Penalty]. | 句首锚定旧信息/逻辑视角，主干动词紧跟主语，复杂修饰置于宾语之后（符合尾重原则 L011 与旧新信息流 L001）。}
  \item \textbf{S3 [实证结果高度概括 Empirical Recap]:} \textit{Empirical evaluations across eight benchmarks and industrial wind-power telemetry demonstrated MSE reductions of up to 17.4\% alongside a 45\% reduction in computational overhead.}\par\noindent{\small 对应步骤：S3 [实证结果高度概括 Empirical Recap]：Comprehensive evaluations across [Benchmarks/Hardware] confirmed consistent error reductions of up to [X\%] alongside substantial inference speedups. | 句首锚定旧信息/逻辑视角，主干动词紧跟主语，复杂修饰置于宾语之后（符合尾重原则 L011 与旧新信息流 L001）。}
  \item \textbf{S4 [跨领域迁移展望 Broader Outlook]:} \textit{Beyond forecasting, this structural-scale disentanglement principle provides a generalizable foundation for anomaly detection and closed-loop predictive control in non-stationary cyber-physical systems.}\par\noindent{\small 对应步骤：S4 [跨领域迁移展望 Broader Outlook]：Beyond [Current Task], our [Core Principle] offers a generalizable blueprint for broader sequential decision-making tasks under non-stationary environments. | 句首锚定旧信息/逻辑视角，主干动词紧跟主语，复杂修饰置于宾语之后（符合尾重原则 L011 与旧新信息流 L001）。}
\end{itemize}
\end{TipBox}
\medskip
\clearpage
\section{FT07: 顶会 5 分钟 Oral 演讲完整逐字稿：非平稳时间序列预测（含开场 Hook、图表引导、公式直觉与 Q\&A）}\label{comp:FT07}
\begin{ConceptBox}{Composition Architecture \& Structural Blueprint Chain}
\textbf{Genre:} \texttt{Academic\_Speech} \quad \textbf{Paragraph Count:} 6\par
\textbf{Blueprint Chain:} \texttt{BP41 $\to$ BP03 $\to$ BP04 $\to$ BP43 $\to$ BP42 $\to$ BP44}
\end{ConceptBox}
\subsection*{Paragraph 1: 学术演讲开场 60 秒 Hook：用直观痛点或反直觉现象抓住全场（Oral Opening Hook） (\texttt{BP41})}
\noindent\textbf{Expressive Intent:}\enspace 我这一段想作为顶会 Oral Presentation、组会汇报或博士面试 PPT 第一页的开场白：在 45--60 秒内不用任何晦涩公式，让全场听众立刻听懂“我们为什么要做这件事、现有方法到底卡在哪”。\par
\noindent\textbf{Cohesion Formula:}\enspace {\small\color{AcadSlate}Concrete Visual Scenario $\to$ Rhetorical Question $\to$ Crisp Culprit Identification $\to$ Talk Promise. 绑定逻辑规则：L018, L029.}\par\vspace{4pt}
\begin{ExampleBox}{Publication-Grade Authentic Text}
Good morning everyone. Suppose you are operating a regional wind-power grid when a sudden cold front sweeps across the turbines. Standard Transformer forecasters look accurate during calm weather, yet why do their predictions lag by hours right when the storm hits? The culprit is instance normalization: in trying to stabilize attention weights, it strips away the exact amplitude surges needed to predict peak loads. In the next five minutes, I will show how our model, DyTS-Net, resolves this trade-off---cutting extreme-weather forecasting error by 17\% while running twice as fast.
\end{ExampleBox}
\begin{TipBox}{Rhetorical Move Breakdown \& Constituent Analysis}
\begin{itemize}[leftmargin=1.5em]
  \item \textbf{S1 [场景直觉切入 Concrete Scenario Hook]:} \textit{Good morning everyone.}\par\noindent{\small 对应步骤：S1 [场景直觉切入 Concrete Scenario Hook]：Good morning everyone. Imagine deploying a humanoid robot to assemble flexible cables on a moving factory line---or forecasting grid load when a storm hits. | 句首锚定旧信息/逻辑视角，主干动词紧跟主语，复杂修饰置于宾语之后（符合尾重原则 L011 与旧新信息流 L001）。}
  \item \textbf{S2 [反直觉痛点设问 Counter-intuitive Breakdown]:} \textit{Suppose you are operating a regional wind-power grid when a sudden cold front sweeps across the turbines.}\par\noindent{\small 对应步骤：S2 [反直觉痛点设问 Counter-intuitive Breakdown]：Standard deep models perform remarkably well when conditions stay calm, yet why do they fail abruptly the moment contact dynamics or weather patterns shift? | 句首锚定旧信息/逻辑视角，主干动词紧跟主语，复杂修饰置于宾语之后（符合尾重原则 L011 与旧新信息流 L001）。}
  \item \textbf{S3 [一针见血揭晓答案 Crisp Root Cause]:} \textit{Standard Transformer forecasters look accurate during calm weather, yet why do their predictions lag by hours right when the storm hits? The culprit is instance normalization: in trying to stabilize attention weights, it strips away the exact amplitude surges needed to predict peak loads.}\par\noindent{\small 对应步骤：S3 [一针见血揭晓答案 Crisp Root Cause]：The culprit is a hidden stationarity assumption baked into standard attention layers. | 句首锚定旧信息/逻辑视角，主干动词紧跟主语，复杂修饰置于宾语之后（符合尾重原则 L011 与旧新信息流 L001）。}
  \item \textbf{S4 [亮出今日主旨 Talk Roadmap]:} \textit{In the next five minutes, I will show how our model, DyTS-Net, resolves this trade-off---cutting extreme-weather forecasting error by 17\% while running twice as fast.}\par\noindent{\small 对应步骤：S4 [亮出今日主旨 Talk Roadmap]：In this talk, I will show how our framework, [Model Name], breaks this bottleneck and achieves [Key Metric Gain] in real-time deployment. | 句首锚定旧信息/逻辑视角，主干动词紧跟主语，复杂修饰置于宾语之后（符合尾重原则 L011 与旧新信息流 L001）。}
\end{itemize}
\end{TipBox}
\medskip
\subsection*{Paragraph 2: 尖锐指出核心技术瓶颈与失效机理（Gap Identification） (\texttt{BP03})}
\noindent\textbf{Expressive Intent:}\enspace 我这一段想深入剖析现有方法在特定难题（如非平稳突变、长程自回归误差累积、Sim-to-Real 接触失配）上为什么会失效，把机理讲透。\par
\noindent\textbf{Cohesion Formula:}\enspace {\small\color{AcadSlate}因果链式推进（A induces B $\to$ This B compounds into C $\to$ Naive fix D sacrifices E）。绑定逻辑规则：L003, L004, L024/L025.}\par\vspace{4pt}
\begin{ExampleBox}{Publication-Grade Authentic Text}
Bridging non-stationary robustness and amplitude fidelity remains a core bottleneck in long-horizon forecasting. While instance normalization strips time-varying mean and variance shifts to stabilize attention weights, it simultaneously erases physical scale cues required to predict extreme load peaks. Conversely, retaining raw unnormalized inputs exposes self-attention layers to severe covariate shift, causing validation error to diverge over extended horizons. Resolving this tension requires a mechanism that disentangles intrinsic temporal dynamics from transient scale drift without discarding magnitude information.
\end{ExampleBox}
\begin{TipBox}{Rhetorical Move Breakdown \& Constituent Analysis}
\begin{itemize}[leftmargin=1.5em]
  \item \textbf{S1 [核心痛点立论 Bottleneck Claim]:} \textit{Bridging non-stationary robustness and amplitude fidelity remains a core bottleneck in long-horizon forecasting.}\par\noindent{\small 对应步骤：S1 [核心痛点立论 Bottleneck Claim]：一句话直接命名该核心瓶颈（如 Compound error accumulation / Sim-to-real contact mismatch）。 | 句首锚定旧信息/逻辑视角，主干动词紧跟主语，复杂修饰置于宾语之后（符合尾重原则 L011 与旧新信息流 L001）。}
  \item \textbf{S2 [第一层失效机理 Primary Failure Mechanism]:} \textit{While instance normalization strips time-varying mean and variance shifts to stabilize attention weights, it simultaneously erases physical scale cues required to predict extreme load peaks.}\par\noindent{\small 对应步骤：S2 [第一层失效机理 Primary Failure Mechanism]：从数学或物理传导链条解释误差是如何产生并放大的。 | 句首锚定旧信息/逻辑视角，主干动词紧跟主语，复杂修饰置于宾语之后（符合尾重原则 L011 与旧新信息流 L001）。}
  \item \textbf{S3 [第二层耦合困境 Secondary Dilemma]:} \textit{Conversely, retaining raw unnormalized inputs exposes self-attention layers to severe covariate shift, causing validation error to diverge over extended horizons.}\par\noindent{\small 对应步骤：S3 [第二层耦合困境 Secondary Dilemma]：指出为什么简单的修补手段（如单纯加大模型、简单归一化或域随机化）不仅无效还会带来副作用。 | 句首锚定旧信息/逻辑视角，主干动词紧跟主语，复杂修饰置于宾语之后（符合尾重原则 L011 与旧新信息流 L001）。}
  \item \textbf{S4 [凝练核心研究问题 Research Question Formulation]:} \textit{Resolving this tension requires a mechanism that disentangles intrinsic temporal dynamics from transient scale drift without discarding magnitude information.}\par\noindent{\small 对应步骤：S4 [凝练核心研究问题 Research Question Formulation]：用 Therefore, a fundamental question arises: how can we... 引出本文目标。 | 句首锚定旧信息/逻辑视角，主干动词紧跟主语，复杂修饰置于宾语之后（符合尾重原则 L011 与旧新信息流 L001）。}
\end{itemize}
\end{TipBox}
\medskip
\subsection*{Paragraph 3: 亮出本文核心洞察与破局方法总览（Key Insight \& Method Overview） (\texttt{BP04})}
\noindent\textbf{Expressive Intent:}\enspace 我这一段想在指出前人痛点后，正式亮出我们的方法（To address this dilemma, we propose...），并用最通俗清晰的直觉告诉审稿人为什么我们的设计能破局。\par
\noindent\textbf{Cohesion Formula:}\enspace {\small\color{AcadSlate}Problem Anchor $\to$ Key Insight $\to$ Dual-Module Mechanism $\to$ Pareto Resolution. 绑定逻辑规则：L002, L023.}\par\vspace{4pt}
\begin{ExampleBox}{Publication-Grade Authentic Text}
To resolve this dilemma, we propose DyTS-Net, a dual-stream forecasting architecture that decouples stationary structural modeling from dynamic scale trajectory estimation. Our key insight is that while high-frequency phase relationships benefit from normalized shift-invariant attention, macro-scale amplitude evolution follows a low-dimensional continuous dynamical system. Specifically, DyTS-Net first extracts scale-invariant temporal tokens via spectral patch attention, and concurrently models non-stationary mean-variance evolution through a lightweight neural ODE branch. Recombining both streams at the decoding stage preserves sharp peak magnitudes without sacrificing training stability.
\end{ExampleBox}
\begin{TipBox}{Rhetorical Move Breakdown \& Constituent Analysis}
\begin{itemize}[leftmargin=1.5em]
  \item \textbf{S1 [方法登场与定位 Method Introduction]:} \textit{To resolve this dilemma, we propose DyTS-Net, a dual-stream forecasting architecture that decouples stationary structural modeling from dynamic scale trajectory estimation.}\par\noindent{\small 对应步骤：S1 [方法登场与定位 Method Introduction]：To address this dilemma, we propose [Model Name], a [Core Paradigm] framework that [Primary Goal]. | 句首锚定旧信息/逻辑视角，主干动词紧跟主语，复杂修饰置于宾语之后（符合尾重原则 L011 与旧新信息流 L001）。}
  \item \textbf{S2 [核心直觉洞察 Key Insight]:} \textit{Our key insight is that while high-frequency phase relationships benefit from normalized shift-invariant attention, macro-scale amplitude evolution follows a low-dimensional continuous dynamical system.}\par\noindent{\small 对应步骤：S2 [核心直觉洞察 Key Insight]：Our key insight is that [Physical/Statistical Principle]，点破破局的底层支点。 | 句首锚定旧信息/逻辑视角，主干动词紧跟主语，复杂修饰置于宾语之后（符合尾重原则 L011 与旧新信息流 L001）。}
  \item \textbf{S3 [模块协同主干 Module Synergy]:} \textit{Specifically, DyTS-Net first extracts scale-invariant temporal tokens via spectral patch attention, and concurrently models non-stationary mean-variance evolution through a lightweight neural ODE branch.}\par\noindent{\small 对应步骤：S3 [模块协同主干 Module Synergy]：用 First / Specifically 与 Concurrently / Furthermore 简述两大核心模块如何分工咬合。 | 句首锚定旧信息/逻辑视角，主干动词紧跟主语，复杂修饰置于宾语之后（符合尾重原则 L011 与旧新信息流 L001）。}
  \item \textbf{S4 [闭环收益总结 Closed-Loop Benefit]:} \textit{Recombining both streams at the decoding stage preserves sharp peak magnitudes without sacrificing training stability.}\par\noindent{\small 对应步骤：S4 [闭环收益总结 Closed-Loop Benefit]：说明这种协同设计如何同时兼顾原先冲突的两个指标（如兼顾平稳表征与幅值保真，或兼顾多模态多样性与单步实时控制）。 | 句首锚定旧信息/逻辑视角，主干动词紧跟主语，复杂修饰置于宾语之后（符合尾重原则 L011 与旧新信息流 L001）。}
\end{itemize}
\end{TipBox}
\medskip
\subsection*{Paragraph 4: 演讲中化繁为简讲解核心公式与算法直觉（Explaining Math Orally） (\texttt{BP43})}
\noindent\textbf{Expressive Intent:}\enspace 我这一段想在 PPT 上出现核心损失函数或注意力公式时，用 4 句通俗、精准的口语把公式背后的物理/几何直觉讲透，而不是在台上干巴巴念希腊字母。\par
\noindent\textbf{Cohesion Formula:}\enspace {\small\color{AcadSlate}Intuition Before Symbols $\to$ Physical Reality $\to$ Pointer to Highlighted Math Term $\to$ Operational Effect. 绑定逻辑规则：L014, L018.}\par\vspace{4pt}
\begin{ExampleBox}{Publication-Grade Authentic Text}
Before looking at the equation in the middle box, here is the physical intuition behind it. When two sensors sit on the same turbine shaft, genuine mechanical vibrations travel with a consistent phase lag, whereas random sensor glitches fire completely out of phase. That physical rule is captured by the highlighted red term, $\|\Delta \Phi_{Q,K}\|_1$, right inside the softmax. Whenever two channels show inconsistent phase lag, this penalty automatically zeroes out their attention weight.
\end{ExampleBox}
\begin{TipBox}{Rhetorical Move Breakdown \& Constituent Analysis}
\begin{itemize}[leftmargin=1.5em]
  \item \textbf{S1 [先讲直觉目标 Intuition First]:} \textit{Before looking at the equation in the middle box, here is the physical intuition behind it.}\par\noindent{\small 对应步骤：S1 [先讲直觉目标 Intuition First]：Before looking at the math at the bottom of the slide, here is the simple physical intuition we want to capture. | 句首锚定旧信息/逻辑视角，主干动词紧跟主语，复杂修饰置于宾语之后（符合尾重原则 L011 与旧新信息流 L001）。}
  \item \textbf{S2 [日常/物理类比 Physical Analogy]:} \textit{When two sensors sit on the same turbine shaft, genuine mechanical vibrations travel with a consistent phase lag, whereas random sensor glitches fire completely out of phase.}\par\noindent{\small 对应步骤：S2 [日常/物理类比 Physical Analogy]：When two sensors monitor the same mechanical shaft, genuine vibrations share a consistent phase lag, whereas random electrical glitches fire out of phase. | 句首锚定旧信息/逻辑视角，主干动词紧跟主语，复杂修饰置于宾语之后（符合尾重原则 L011 与旧新信息流 L001）。}
  \item \textbf{S3 [指向公式核心项 Point to Key Term]:} \textit{That physical rule is captured by the highlighted red term, $\|\Delta \Phi_{Q,K}\|_1$, right inside the softmax.}\par\noindent{\small 对应步骤：S3 [指向公式核心项 Point to Key Term]：That is exactly what the red penalty term $\|\Delta \Phi\|_1$ inside the softmax does: it measures phase mismatch between channels. | 句首锚定旧信息/逻辑视角，主干动词紧跟主语，复杂修饰置于宾语之后（符合尾重原则 L011 与旧新信息流 L001）。}
  \item \textbf{S4 [一句话总结效果 Operational Effect]:} \textit{Whenever two channels show inconsistent phase lag, this penalty automatically zeroes out their attention weight.}\par\noindent{\small 对应步骤：S4 [一句话总结效果 Operational Effect]：Whenever two channels fall out of phase, this term suppresses their attention weight automatically. | 句首锚定旧信息/逻辑视角，主干动词紧跟主语，复杂修饰置于宾语之后（符合尾重原则 L011 与旧新信息流 L001）。}
\end{itemize}
\end{TipBox}
\medskip
\subsection*{Paragraph 5: 演讲幻灯片与复杂图表视觉引导段（Visual Signposting in Tech Talks） (\texttt{BP42})}
\noindent\textbf{Expressive Intent:}\enspace 我这一段想在演讲到架构图或核心实验曲线（Slide 4 / Slide 7）时，用最自然的口语引导听众的眼睛先看哪、再看哪，瞬间读懂图表背后的胜出逻辑。\par
\noindent\textbf{Cohesion Formula:}\enspace {\small\color{AcadSlate}Spatial Coordinate $\to$ Notice Baseline Failure $\to$ Compare Our Curve $\to$ One-Sentence Mechanistic Takeaway. 绑定逻辑规则：L018, L015.}\par\vspace{4pt}
\begin{ExampleBox}{Publication-Grade Authentic Text}
Let's focus on the right panel of Slide 6, which plots forecasting error as the prediction window stretches from 96 to 720 hours. Follow the dashed orange curve for iTransformer---notice how rapidly the error climbs once we pass 336 hours. Now compare that against our solid blue curve along the bottom, which remains nearly flat all the way to 720 hours. That 17\% gap at the far right comes directly from decoupling phase attention from continuous scale tracking.
\end{ExampleBox}
\begin{TipBox}{Rhetorical Move Breakdown \& Constituent Analysis}
\begin{itemize}[leftmargin=1.5em]
  \item \textbf{S1 [视觉坐标锁定 Spatial Pointer]:} \textit{Let's focus on the right panel of Slide 6, which plots forecasting error as the prediction window stretches from 96 to 720 hours.}\par\noindent{\small 对应步骤：S1 [视觉坐标锁定 Spatial Pointer]：Let's look at the comparison in the right panel of Slide 5, where the horizontal axis tracks prediction horizon and the vertical axis shows MSE. | 句首锚定旧信息/逻辑视角，主干动词紧跟主语，复杂修饰置于宾语之后（符合尾重原则 L011 与旧新信息流 L001）。}
  \item \textbf{S2 [引导观察反方崩溃曲线 Guide to Baseline Curve]:} \textit{Follow the dashed orange curve for iTransformer---notice how rapidly the error climbs once we pass 336 hours.}\par\noindent{\small 对应步骤：S2 [引导观察反方崩溃曲线 Guide to Baseline Curve]：Notice the dashed orange curve representing the standard Transformer baseline: past 336 steps, error shoots up almost vertically. | 句首锚定旧信息/逻辑视角，主干动词紧跟主语，复杂修饰置于宾语之后（符合尾重原则 L011 与旧新信息流 L001）。}
  \item \textbf{S3 [引导观察我方稳定曲线 Guide to Our Curve]:} \textit{Now compare that against our solid blue curve along the bottom, which remains nearly flat all the way to 720 hours.}\par\noindent{\small 对应步骤：S3 [引导观察我方稳定曲线 Guide to Our Curve]：Now compare that with our solid blue curve at the bottom, which stays flat even out to 720 steps. | 句首锚定旧信息/逻辑视角，主干动词紧跟主语，复杂修饰置于宾语之后（符合尾重原则 L011 与旧新信息流 L001）。}
  \item \textbf{S4 [一句话点破视觉差距原因 One-Sentence Takeaway]:} \textit{That 17\% gap at the far right comes directly from decoupling phase attention from continuous scale tracking.}\par\noindent{\small 对应步骤：S4 [一句话点破视觉差距原因 One-Sentence Takeaway]：That gap comes directly from our neural-ODE scale tracker preventing long-range drift. | 句首锚定旧信息/逻辑视角，主干动词紧跟主语，复杂修饰置于宾语之后（符合尾重原则 L011 与旧新信息流 L001）。}
\end{itemize}
\end{TipBox}
\medskip
\subsection*{Paragraph 6: 演讲总结页（Take-home Message）与问答（Q\&A）从容应答段 (\texttt{BP44})}
\noindent\textbf{Expressive Intent:}\enspace 我这一段想用于演讲最后一页总结（Take-home Messages）以及回答台下教授/审稿人犀利提问时的标准应答模板。\par
\noindent\textbf{Cohesion Formula:}\enspace {\small\color{AcadSlate}Single Takeaway $\to$ Quantitative Triad $\to$ Q\&A Direct Answer $\to$ Empirical Grounding. 绑定逻辑规则：L018, L019.}\par\vspace{4pt}
\begin{ExampleBox}{Publication-Grade Authentic Text}
To wrap up, the key takeaway from our work is that non-stationary forecasting requires separating shift-invariant phase structure from continuous amplitude dynamics. Doing so lets DyTS-Net cut long-horizon error by up to 17\% while running at half the computational cost, and our full code and weights are live at the QR code on screen. --- [Q\&A Response]: Thank you for that sharp question on irregular missing timestamps. The short answer is that our neural-ODE branch handles arbitrary time intervals natively: in our 30\% random-dropout stress test, MSE degraded by less than 5\%, and we can share the backup slide on irregular sampling right here.
\end{ExampleBox}
\begin{TipBox}{Rhetorical Move Breakdown \& Constituent Analysis}
\begin{itemize}[leftmargin=1.5em]
  \item \textbf{S1 [凝练核心带走信息 Core Take-Home Message]:} \textit{To wrap up, the key takeaway from our work is that non-stationary forecasting requires separating shift-invariant phase structure from continuous amplitude dynamics.}\par\noindent{\small 对应步骤：S1 [凝练核心带走信息 Core Take-Home Message]：To wrap up, if there is one takeaway from our work today, it is that [One-Sentence Core Principle]. | 句首锚定旧信息/逻辑视角，主干动词紧跟主语，复杂修饰置于宾语之后（符合尾重原则 L011 与旧新信息流 L001）。}
  \item \textbf{S2 [成果三点速览 Rapid Recap]:} \textit{Doing so lets DyTS-Net cut long-horizon error by up to 17\% while running at half the computational cost, and our full code and weights are live at the QR code on screen.}\par\noindent{\small 对应步骤：S2 [成果三点速览 Rapid Recap]：By combining [A] and [B], [Model] delivers [Metric 1], [Metric 2], and [Open-Source/Practical Asset]. | 句首锚定旧信息/逻辑视角，主干动词紧跟主语，复杂修饰置于宾语之后（符合尾重原则 L011 与旧新信息流 L001）。}
  \item \textbf{S3 [Q\&A 礼貌接招与结论前置 Q\&A Direct Answer]:} \textit{Thank you for that sharp question on irregular missing timestamps.}\par\noindent{\small 对应步骤：S3 [Q\&A 礼貌接招与结论前置 Q\&A Direct Answer]：(For Q\&A) That is a great question regarding [Topic]---the short answer is [Direct Conclusion]. | 句首锚定旧信息/逻辑视角，主干动词紧跟主语，复杂修饰置于宾语之后（符合尾重原则 L011 与旧新信息流 L001）。}
  \item \textbf{S4 [Q\&A 展开与诚实边界 Q\&A Elaboration]:} \textit{The short answer is that our neural-ODE branch handles arbitrary time intervals natively: in our 30\% random-dropout stress test, MSE degraded by less than 5\%, and we can share the backup slide on irregular sampling right here.}\par\noindent{\small 对应步骤：S4 [Q\&A 展开与诚实边界 Q\&A Elaboration]：In our current experiments, we observed [Empirical Fact], and extending this to [Questioner's Scenario] can be handled by [Concrete Mechanism]. | 句首锚定旧信息/逻辑视角，主干动词紧跟主语，复杂修饰置于宾语之后（符合尾重原则 L011 与旧新信息流 L001）。}
\end{itemize}
\end{TipBox}
\medskip
\clearpage
\section{FT08: 顶会 5 分钟 Spotlight 演讲完整逐字稿：具身智能实时视觉运动控制（含真机视频解说与延迟对比）}\label{comp:FT08}
\begin{ConceptBox}{Composition Architecture \& Structural Blueprint Chain}
\textbf{Genre:} \texttt{Academic\_Speech} \quad \textbf{Paragraph Count:} 6\par
\textbf{Blueprint Chain:} \texttt{BP41 $\to$ BP05 $\to$ BP16 $\to$ BP42 $\to$ BP22 $\to$ BP44}
\end{ConceptBox}
\subsection*{Paragraph 1: 学术演讲开场 60 秒 Hook：用直观痛点或反直觉现象抓住全场（Oral Opening Hook） (\texttt{BP41})}
\noindent\textbf{Expressive Intent:}\enspace 我这一段想作为顶会 Oral Presentation、组会汇报或博士面试 PPT 第一页的开场白：在 45--60 秒内不用任何晦涩公式，让全场听众立刻听懂“我们为什么要做这件事、现有方法到底卡在哪”。\par
\noindent\textbf{Cohesion Formula:}\enspace {\small\color{AcadSlate}Concrete Visual Scenario $\to$ Rhetorical Question $\to$ Crisp Culprit Identification $\to$ Talk Promise. 绑定逻辑规则：L018, L029.}\par\vspace{4pt}
\begin{ExampleBox}{Publication-Grade Authentic Text}
Good morning everyone. Suppose you are operating a regional wind-power grid when a sudden cold front sweeps across the turbines. Standard Transformer forecasters look accurate during calm weather, yet why do their predictions lag by hours right when the storm hits? The culprit is instance normalization: in trying to stabilize attention weights, it strips away the exact amplitude surges needed to predict peak loads. In the next five minutes, I will show how our model, DyTS-Net, resolves this trade-off---cutting extreme-weather forecasting error by 17\% while running twice as fast.
\end{ExampleBox}
\begin{TipBox}{Rhetorical Move Breakdown \& Constituent Analysis}
\begin{itemize}[leftmargin=1.5em]
  \item \textbf{S1 [场景直觉切入 Concrete Scenario Hook]:} \textit{Good morning everyone.}\par\noindent{\small 对应步骤：S1 [场景直觉切入 Concrete Scenario Hook]：Good morning everyone. Imagine deploying a humanoid robot to assemble flexible cables on a moving factory line---or forecasting grid load when a storm hits. | 句首锚定旧信息/逻辑视角，主干动词紧跟主语，复杂修饰置于宾语之后（符合尾重原则 L011 与旧新信息流 L001）。}
  \item \textbf{S2 [反直觉痛点设问 Counter-intuitive Breakdown]:} \textit{Suppose you are operating a regional wind-power grid when a sudden cold front sweeps across the turbines.}\par\noindent{\small 对应步骤：S2 [反直觉痛点设问 Counter-intuitive Breakdown]：Standard deep models perform remarkably well when conditions stay calm, yet why do they fail abruptly the moment contact dynamics or weather patterns shift? | 句首锚定旧信息/逻辑视角，主干动词紧跟主语，复杂修饰置于宾语之后（符合尾重原则 L011 与旧新信息流 L001）。}
  \item \textbf{S3 [一针见血揭晓答案 Crisp Root Cause]:} \textit{Standard Transformer forecasters look accurate during calm weather, yet why do their predictions lag by hours right when the storm hits? The culprit is instance normalization: in trying to stabilize attention weights, it strips away the exact amplitude surges needed to predict peak loads.}\par\noindent{\small 对应步骤：S3 [一针见血揭晓答案 Crisp Root Cause]：The culprit is a hidden stationarity assumption baked into standard attention layers. | 句首锚定旧信息/逻辑视角，主干动词紧跟主语，复杂修饰置于宾语之后（符合尾重原则 L011 与旧新信息流 L001）。}
  \item \textbf{S4 [亮出今日主旨 Talk Roadmap]:} \textit{In the next five minutes, I will show how our model, DyTS-Net, resolves this trade-off---cutting extreme-weather forecasting error by 17\% while running twice as fast.}\par\noindent{\small 对应步骤：S4 [亮出今日主旨 Talk Roadmap]：In this talk, I will show how our framework, [Model Name], breaks this bottleneck and achieves [Key Metric Gain] in real-time deployment. | 句首锚定旧信息/逻辑视角，主干动词紧跟主语，复杂修饰置于宾语之后（符合尾重原则 L011 与旧新信息流 L001）。}
\end{itemize}
\end{TipBox}
\medskip
\subsection*{Paragraph 2: 挑战与机制一对一映射段（Challenge-to-Mechanism Mapping） (\texttt{BP05})}
\noindent\textbf{Expressive Intent:}\enspace 我这一段想说明：把我们的核心思想落地时遇到了两个具体工程/数学难点，我们分别设计了什么精巧机制去攻克它们。\par
\noindent\textbf{Cohesion Formula:}\enspace {\small\color{AcadSlate}Challenge 1 $\to$ Mechanism 1; Challenge 2 $\to$ Mechanism 2 $\to$ Joint System Property. 绑定逻辑规则：L003, L011.}\par\vspace{4pt}
\begin{ExampleBox}{Publication-Grade Authentic Text}
Instantiating closed-loop visuomotor diffusion on physical robots entails two technical challenges. First, because raw depth camera streams contain occlusion artifacts around reflective object edges, we introduce a tactile-guided cross-attention gate that reweights visual tokens upon contact onset. Second, to eliminate the 200 ms sampling bottleneck of iterative denoising, we formulate a trajectory consistency distillation objective that maps Gaussian noise directly to collision-free action chunks in two forward steps. Together, these designs sustain 50 Hz reactive control under partial visual occlusion.
\end{ExampleBox}
\begin{TipBox}{Rhetorical Move Breakdown \& Constituent Analysis}
\begin{itemize}[leftmargin=1.5em]
  \item \textbf{S1 [落地难点总起 Implementation Hurdle]:} \textit{Instantiating closed-loop visuomotor diffusion on physical robots entails two technical challenges.}\par\noindent{\small 对应步骤：S1 [落地难点总起 Implementation Hurdle]：Translating this principle into a scalable architecture entails two technical challenges. | 句首锚定旧信息/逻辑视角，主干动词紧跟主语，复杂修饰置于宾语之后（符合尾重原则 L011 与旧新信息流 L001）。}
  \item \textbf{S2 [难点 1 + 靶向设计 1]:} \textit{First, because raw depth camera streams contain occlusion artifacts around reflective object edges, we introduce a tactile-guided cross-attention gate that reweights visual tokens upon contact onset.}\par\noindent{\small 对应步骤：S2 [难点 1 + 靶向设计 1]：First, because [Challenge 1], we design [Mechanism 1] to [Specific Function]. | 句首锚定旧信息/逻辑视角，主干动词紧跟主语，复杂修饰置于宾语之后（符合尾重原则 L011 与旧新信息流 L001）。}
  \item \textbf{S3 [难点 2 + 靶向设计 2]:} \textit{Second, to eliminate the 200 ms sampling bottleneck of iterative denoising, we formulate a trajectory consistency distillation objective that maps Gaussian noise directly to collision-free action chunks in two forward steps.}\par\noindent{\small 对应步骤：S3 [难点 2 + 靶向设计 2]：Second, to prevent [Challenge 2] under [Constraint], we introduce [Mechanism 2] that [Specific Function]. | 句首锚定旧信息/逻辑视角，主干动词紧跟主语，复杂修饰置于宾语之后（符合尾重原则 L011 与旧新信息流 L001）。}
  \item \textbf{S4 [整体计算/泛化特性 System Property]:} \textit{Together, these designs sustain 50 Hz reactive control under partial visual occlusion.}\par\noindent{\small 对应步骤：S4 [整体计算/泛化特性 System Property]：Together, these components enable [Model] to operate at [Efficiency/Robustness Metric]. | 句首锚定旧信息/逻辑视角，主干动词紧跟主语，复杂修饰置于宾语之后（符合尾重原则 L011 与旧新信息流 L001）。}
\end{itemize}
\end{TipBox}
\medskip
\subsection*{Paragraph 3: 具身智能专属：多模态感知-动作闭环与 Sim-to-Real 阐述段（Embodied AI Policy） (\texttt{BP16})}
\noindent\textbf{Expressive Intent:}\enspace 我这一段想阐述具身智能机器人策略（Embodied Policy）如何融合视觉、本体感觉（Proprioception）或触觉，并克服仿真到真实硬件的动力学差距（Sim-to-Real Gap）。\par
\noindent\textbf{Cohesion Formula:}\enspace {\small\color{AcadSlate}Visual Occlusion Bottleneck $\to$ Asynchronous Multimodal Alignment $\to$ Action Chunking $\to$ Online Sim-to-Real System Identification. 绑定逻辑规则：L025, L003.}\par\vspace{4pt}
\begin{ExampleBox}{Publication-Grade Authentic Text}
Relying solely on third-person RGB-D observations leaves visuomotor policies vulnerable to severe self-occlusion once the gripper approaches the target object. To preserve observability during contact-rich manipulation, our encoder aligns 30 Hz visual spatial tokens with 500 Hz joint proprioception and wrist force-torque histories via cross-modal rotary embeddings. Rather than issuing single-step commands prone to high-frequency jitter, the policy predicts 16-step SE(3) end-effector action chunks smoothed through exponential temporal ensembling. At deployment, a recurrent history adapter estimates latent surface friction and payload mass online, bridging the sim-to-real dynamics gap in zero shot.
\end{ExampleBox}
\begin{TipBox}{Rhetorical Move Breakdown \& Constituent Analysis}
\begin{itemize}[leftmargin=1.5em]
  \item \textbf{S1 [单模态开环局限 Open-Loop Visual Limitation]:} \textit{Relying solely on third-person RGB-D observations leaves visuomotor policies vulnerable to severe self-occlusion once the gripper approaches the target object.}\par\noindent{\small 对应步骤：S1 [单模态开环局限 Open-Loop Visual Limitation]：Relying solely on RGB-D point clouds leaves manipulation policies vulnerable to self-occlusion by the robot gripper during fine-grained contact. | 句首锚定旧信息/逻辑视角，主干动词紧跟主语，复杂修饰置于宾语之后（符合尾重原则 L011 与旧新信息流 L001）。}
  \item \textbf{S2 [多模态时空对齐 Multi-Modal Alignment]:} \textit{To preserve observability during contact-rich manipulation, our encoder aligns 30 Hz visual spatial tokens with 500 Hz joint proprioception and wrist force-torque histories via cross-modal rotary embeddings.}\par\noindent{\small 对应步骤：S2 [多模态时空对齐 Multi-Modal Alignment]：To maintain state observability during contact, our encoder aligns asynchronous 30 Hz visual tokens with 500 Hz joint proprioception and wrist force-torque histories. | 句首锚定旧信息/逻辑视角，主干动词紧跟主语，复杂修饰置于宾语之后（符合尾重原则 L011 与旧新信息流 L001）。}
  \item \textbf{S3 [动作块预测与平滑执行 Action Chunking \& Temporal Ensembling]:} \textit{Rather than issuing single-step commands prone to high-frequency jitter, the policy predicts 16-step SE(3) end-effector action chunks smoothed through exponential temporal ensembling.}\par\noindent{\small 对应步骤：S3 [动作块预测与平滑执行 Action Chunking \& Temporal Ensembling]：Rather than predicting single-step torques prone to jitter, the policy outputs overlapping $K$-step SE(3) end-effector action chunks smoothed via exponential temporal ensembling. | 句首锚定旧信息/逻辑视角，主干动词紧跟主语，复杂修饰置于宾语之后（符合尾重原则 L011 与旧新信息流 L001）。}
  \item \textbf{S4 [Sim-to-Real 自适应机制 Sim-to-Real Adaptation]:} \textit{At deployment, a recurrent history adapter estimates latent surface friction and payload mass online, bridging the sim-to-real dynamics gap in zero shot.}\par\noindent{\small 对应步骤：S4 [Sim-to-Real 自适应机制 Sim-to-Real Adaptation]：During physical deployment, a lightweight history adapter infers latent friction and payload parameters online, bridging the sim-to-real dynamics gap without policy fine-tuning. | 句首锚定旧信息/逻辑视角，主干动词紧跟主语，复杂修饰置于宾语之后（符合尾重原则 L011 与旧新信息流 L001）。}
\end{itemize}
\end{TipBox}
\medskip
\subsection*{Paragraph 4: 演讲幻灯片与复杂图表视觉引导段（Visual Signposting in Tech Talks） (\texttt{BP42})}
\noindent\textbf{Expressive Intent:}\enspace 我这一段想在演讲到架构图或核心实验曲线（Slide 4 / Slide 7）时，用最自然的口语引导听众的眼睛先看哪、再看哪，瞬间读懂图表背后的胜出逻辑。\par
\noindent\textbf{Cohesion Formula:}\enspace {\small\color{AcadSlate}Spatial Coordinate $\to$ Notice Baseline Failure $\to$ Compare Our Curve $\to$ One-Sentence Mechanistic Takeaway. 绑定逻辑规则：L018, L015.}\par\vspace{4pt}
\begin{ExampleBox}{Publication-Grade Authentic Text}
Let's focus on the right panel of Slide 6, which plots forecasting error as the prediction window stretches from 96 to 720 hours. Follow the dashed orange curve for iTransformer---notice how rapidly the error climbs once we pass 336 hours. Now compare that against our solid blue curve along the bottom, which remains nearly flat all the way to 720 hours. That 17\% gap at the far right comes directly from decoupling phase attention from continuous scale tracking.
\end{ExampleBox}
\begin{TipBox}{Rhetorical Move Breakdown \& Constituent Analysis}
\begin{itemize}[leftmargin=1.5em]
  \item \textbf{S1 [视觉坐标锁定 Spatial Pointer]:} \textit{Let's focus on the right panel of Slide 6, which plots forecasting error as the prediction window stretches from 96 to 720 hours.}\par\noindent{\small 对应步骤：S1 [视觉坐标锁定 Spatial Pointer]：Let's look at the comparison in the right panel of Slide 5, where the horizontal axis tracks prediction horizon and the vertical axis shows MSE. | 句首锚定旧信息/逻辑视角，主干动词紧跟主语，复杂修饰置于宾语之后（符合尾重原则 L011 与旧新信息流 L001）。}
  \item \textbf{S2 [引导观察反方崩溃曲线 Guide to Baseline Curve]:} \textit{Follow the dashed orange curve for iTransformer---notice how rapidly the error climbs once we pass 336 hours.}\par\noindent{\small 对应步骤：S2 [引导观察反方崩溃曲线 Guide to Baseline Curve]：Notice the dashed orange curve representing the standard Transformer baseline: past 336 steps, error shoots up almost vertically. | 句首锚定旧信息/逻辑视角，主干动词紧跟主语，复杂修饰置于宾语之后（符合尾重原则 L011 与旧新信息流 L001）。}
  \item \textbf{S3 [引导观察我方稳定曲线 Guide to Our Curve]:} \textit{Now compare that against our solid blue curve along the bottom, which remains nearly flat all the way to 720 hours.}\par\noindent{\small 对应步骤：S3 [引导观察我方稳定曲线 Guide to Our Curve]：Now compare that with our solid blue curve at the bottom, which stays flat even out to 720 steps. | 句首锚定旧信息/逻辑视角，主干动词紧跟主语，复杂修饰置于宾语之后（符合尾重原则 L011 与旧新信息流 L001）。}
  \item \textbf{S4 [一句话点破视觉差距原因 One-Sentence Takeaway]:} \textit{That 17\% gap at the far right comes directly from decoupling phase attention from continuous scale tracking.}\par\noindent{\small 对应步骤：S4 [一句话点破视觉差距原因 One-Sentence Takeaway]：That gap comes directly from our neural-ODE scale tracker preventing long-range drift. | 句首锚定旧信息/逻辑视角，主干动词紧跟主语，复杂修饰置于宾语之后（符合尾重原则 L011 与旧新信息流 L001）。}
\end{itemize}
\end{TipBox}
\medskip
\subsection*{Paragraph 5: 计算开销、推理延迟与精度帕累托权衡段（Efficiency-Accuracy Pareto Analysis） (\texttt{BP22})}
\noindent\textbf{Expressive Intent:}\enspace 我这一段想说明我们的模型不仅准，而且算得快、显存省，在“精度-延迟/显存”帕累托前沿（Pareto Frontier）上显著优于现有大模型。\par
\noindent\textbf{Cohesion Formula:}\enspace {\small\color{AcadSlate}Hardware Benchmark $\to$ Complexity Reduction \& Memory Gain $\to$ Real-Time Latency Margin $\to$ Pareto Conclusion. 绑定逻辑规则：L010, L023.}\par\vspace{4pt}
\begin{ExampleBox}{Publication-Grade Authentic Text}
Figure 6 plots the MSE-latency Pareto frontier measured on an onboard NVIDIA Orin AGX module. By replacing $O(L^2)$ point-wise self-attention with $O(L \log L)$ spectral filtering and compact patch tokens, DyTS-Net reduces peak GPU memory by 62\% and training wall-clock time by 3.4x compared with iTransformer. At inference time, generating a 16-step visuomotor trajectory requires only 6.8 ms---well below the 20 ms cycle budget of a 50 Hz industrial robot controller---while outperforming a 200 ms DDPM baseline by 11.4\% in task success rate.
\end{ExampleBox}
\begin{TipBox}{Rhetorical Move Breakdown \& Constituent Analysis}
\begin{itemize}[leftmargin=1.5em]
  \item \textbf{S1 [硬件与测试基准 Hardware Setup]:} \textit{Figure 6 plots the MSE-latency Pareto frontier measured on an onboard NVIDIA Orin AGX module.}\par\noindent{\small 对应步骤：S1 [硬件与测试基准 Hardware Setup]：Figure 6 compares the accuracy-efficiency Pareto frontier across all baselines measured on a single NVIDIA RTX 4090 GPU. | 句首锚定旧信息/逻辑视角，主干动词紧跟主语，复杂修饰置于宾语之后（符合尾重原则 L011 与旧新信息流 L001）。}
  \item \textbf{S2 [复杂度理论与实测对比 Complexity \& Runtime]:} \textit{By replacing $O(L^2)$ point-wise self-attention with $O(L \log L)$ spectral filtering and compact patch tokens, DyTS-Net reduces peak GPU memory by 62\% and training wall-clock time by 3.4x compared with iTransformer.}\par\noindent{\small 对应步骤：S2 [复杂度理论与实测对比 Complexity \& Runtime]：By replacing $O(L^2)$ point-wise attention with $O(L \log L)$ spectral filtering and patch aggregation, our model cuts training time per epoch by [X\%] and peak VRAM by [Y\%]. | 句首锚定旧信息/逻辑视角，主干动词紧跟主语，复杂修饰置于宾语之后（符合尾重原则 L011 与旧新信息流 L001）。}
  \item \textbf{S3 [实时闭环部署可行性 Real-Time Viability]:} \textit{At inference time, generating a 16-step visuomotor trajectory requires only 6.8 ms---well below the 20 ms cycle budget of a 50 Hz industrial robot controller---while outperforming a 200 ms DDPM baseline by 11.4\% in task success rate.}\par\noindent{\small 对应步骤：S3 [实时闭环部署可行性 Real-Time Viability]：At inference time, generating a 720-step forecast (or 16-step robot action chunk) requires only [4.2 ms], well below the [20 ms] real-time control deadline. | 句首锚定旧信息/逻辑视角，主干动词紧跟主语，复杂修饰置于宾语之后（符合尾重原则 L011 与旧新信息流 L001）。}
\end{itemize}
\end{TipBox}
\medskip
\subsection*{Paragraph 6: 演讲总结页（Take-home Message）与问答（Q\&A）从容应答段 (\texttt{BP44})}
\noindent\textbf{Expressive Intent:}\enspace 我这一段想用于演讲最后一页总结（Take-home Messages）以及回答台下教授/审稿人犀利提问时的标准应答模板。\par
\noindent\textbf{Cohesion Formula:}\enspace {\small\color{AcadSlate}Single Takeaway $\to$ Quantitative Triad $\to$ Q\&A Direct Answer $\to$ Empirical Grounding. 绑定逻辑规则：L018, L019.}\par\vspace{4pt}
\begin{ExampleBox}{Publication-Grade Authentic Text}
To wrap up, the key takeaway from our work is that non-stationary forecasting requires separating shift-invariant phase structure from continuous amplitude dynamics. Doing so lets DyTS-Net cut long-horizon error by up to 17\% while running at half the computational cost, and our full code and weights are live at the QR code on screen. --- [Q\&A Response]: Thank you for that sharp question on irregular missing timestamps. The short answer is that our neural-ODE branch handles arbitrary time intervals natively: in our 30\% random-dropout stress test, MSE degraded by less than 5\%, and we can share the backup slide on irregular sampling right here.
\end{ExampleBox}
\begin{TipBox}{Rhetorical Move Breakdown \& Constituent Analysis}
\begin{itemize}[leftmargin=1.5em]
  \item \textbf{S1 [凝练核心带走信息 Core Take-Home Message]:} \textit{To wrap up, the key takeaway from our work is that non-stationary forecasting requires separating shift-invariant phase structure from continuous amplitude dynamics.}\par\noindent{\small 对应步骤：S1 [凝练核心带走信息 Core Take-Home Message]：To wrap up, if there is one takeaway from our work today, it is that [One-Sentence Core Principle]. | 句首锚定旧信息/逻辑视角，主干动词紧跟主语，复杂修饰置于宾语之后（符合尾重原则 L011 与旧新信息流 L001）。}
  \item \textbf{S2 [成果三点速览 Rapid Recap]:} \textit{Doing so lets DyTS-Net cut long-horizon error by up to 17\% while running at half the computational cost, and our full code and weights are live at the QR code on screen.}\par\noindent{\small 对应步骤：S2 [成果三点速览 Rapid Recap]：By combining [A] and [B], [Model] delivers [Metric 1], [Metric 2], and [Open-Source/Practical Asset]. | 句首锚定旧信息/逻辑视角，主干动词紧跟主语，复杂修饰置于宾语之后（符合尾重原则 L011 与旧新信息流 L001）。}
  \item \textbf{S3 [Q\&A 礼貌接招与结论前置 Q\&A Direct Answer]:} \textit{Thank you for that sharp question on irregular missing timestamps.}\par\noindent{\small 对应步骤：S3 [Q\&A 礼貌接招与结论前置 Q\&A Direct Answer]：(For Q\&A) That is a great question regarding [Topic]---the short answer is [Direct Conclusion]. | 句首锚定旧信息/逻辑视角，主干动词紧跟主语，复杂修饰置于宾语之后（符合尾重原则 L011 与旧新信息流 L001）。}
  \item \textbf{S4 [Q\&A 展开与诚实边界 Q\&A Elaboration]:} \textit{The short answer is that our neural-ODE branch handles arbitrary time intervals natively: in our 30\% random-dropout stress test, MSE degraded by less than 5\%, and we can share the backup slide on irregular sampling right here.}\par\noindent{\small 对应步骤：S4 [Q\&A 展开与诚实边界 Q\&A Elaboration]：In our current experiments, we observed [Empirical Fact], and extending this to [Questioner's Scenario] can be handled by [Concrete Mechanism]. | 句首锚定旧信息/逻辑视角，主干动词紧跟主语，复杂修饰置于宾语之后（符合尾重原则 L011 与旧新信息流 L001）。}
\end{itemize}
\end{TipBox}
\medskip
\clearpage
\section{FT09: ICLR/NeurIPS 满分审稿回复全文（Author Rebuttal：含置顶全局回复 + 创新性辩护 + 新增基线与消融表）}\label{comp:FT09}
\begin{ConceptBox}{Composition Architecture \& Structural Blueprint Chain}
\textbf{Genre:} \texttt{Conference\_Rebuttal} \quad \textbf{Paragraph Count:} 6\par
\textbf{Blueprint Chain:} \texttt{BP45 $\to$ BP46 $\to$ BP47 $\to$ BP19 $\to$ BP20 $\to$ BP23}
\end{ConceptBox}
\subsection*{Paragraph 1: 顶会 Rebuttal 全局总括与新增实验汇总段（Author Rebuttal General Response） (\texttt{BP45})}
\noindent\textbf{Expressive Intent:}\enspace 我这一段想写 ICLR / NeurIPS / ICML / CVPR 审稿回复（Rebuttal）最顶部的致全体审稿人总括段（Global Response）：在 150 词内用最提气的语气感谢审稿人肯定，并一眼亮出我们在 Rebuttal 期间赶做出的所有硬核新增实验。\par
\noindent\textbf{Cohesion Formula:}\enspace {\small\color{AcadSlate}Synthesize Reviewer Praise $\to$ Announce Concrete Revisions $\to$ Numbered Evidence Bullets mapped to Reviewers $\to$ Courteous Hook. 绑定逻辑规则：L019, L010.}\par\vspace{4pt}
\begin{ExampleBox}{Publication-Grade Authentic Text}
We thank all reviewers for their constructive feedback and for recognizing our framework's 'principled formulation of the stationarity-amplitude trade-off' (R1, R3), 'strong empirical gains across eight benchmarks' (R2), and 'clear ablation design' (R4). To address your valuable suggestions, we conducted four sets of new experiments during the rebuttal period (incorporated in blue in the revised manuscript): (1) Benchmarked against two latest 2025 foundation models, Chronos and MOIRAI, across all horizons (Table R1, addressing R1, R3); (2) Added real-robot grasping evaluations under dynamic lighting and 25\% camera occlusion (Section 4.5, R2); and (3) Included full GPU memory and wall-clock latency scaling curves up to $L=4,096$ (Appendix E, R4). Below we address each reviewer's specific questions in detail.
\end{ExampleBox}
\begin{TipBox}{Rhetorical Move Breakdown \& Constituent Analysis}
\begin{itemize}[leftmargin=1.5em]
  \item \textbf{S1 [致谢并提炼审稿人共同好评 Gratitude \& Shared Praise]:} \textit{We thank all reviewers for their constructive feedback and for recognizing our framework's 'principled formulation of the stationarity-amplitude trade-off' (R1, R3), 'strong empirical gains across eight benchmarks' (R2), and 'clear ablation design' (R4).}\par\noindent{\small 对应步骤：S1 [致谢并提炼审稿人共同好评 Gratitude \& Shared Praise]：We thank all reviewers for their constructive feedback and for recognizing our work's [Quote Reviewer 1 Praise] (R1), [Quote Reviewer 2 Praise] (R2), and [Quote Reviewer 3 Praise] (R3). | 句首锚定旧信息/逻辑视角，主干动词紧跟主语，复杂修饰置于宾语之后（符合尾重原则 L011 与旧新信息流 L001）。}
  \item \textbf{S2 [概览新增实证工作 Overview of New Experiments]:} \textit{To address your valuable suggestions, we conducted four sets of new experiments during the rebuttal period (incorporated in blue in the revised manuscript): (1) Benchmarked against two latest 2025 foundation models, Chronos and MOIRAI, across all horizons (Table R1, addressing R1, R3); (2) Added real-robot grasping evaluations under dynamic lighting and 25\% camera occlusion (Section 4.5, R2); and (3) Included full GPU memory and wall-clock latency scaling curves up to $L=4,096$ (Appendix E, R4).}\par\noindent{\small 对应步骤：S2 [概览新增实证工作 Overview of New Experiments]：In response to your valuable suggestions, we have conducted [N] sets of new experiments and theoretical analyses, summarized below and highlighted in blue in the revised PDF: | 句首锚定旧信息/逻辑视角，主干动词紧跟主语，复杂修饰置于宾语之后（符合尾重原则 L011 与旧新信息流 L001）。}
  \item \textbf{S3 [分条列出硬核新增实验 Bullet Summary of Additions]:} \textit{Below we address each reviewer's specific questions in detail.}\par\noindent{\small 对应步骤：S3 [分条列出硬核新增实验 Bullet Summary of Additions]：(1) Added comparisons against recent baselines [A, B] on [Dataset] (Table R1, addressing R1, R3); (2) Evaluated real-world robustness under [Harsh Condition] (Appendix D, R2); (3) Provided wall-clock latency and memory profiling across sequence lengths up to $L=4,096$ (Appendix E, R3). | 句首锚定旧信息/逻辑视角，主干动词紧跟主语，复杂修饰置于宾语之后（符合尾重原则 L011 与旧新信息流 L001）。}
\end{itemize}
\end{TipBox}
\medskip
\subsection*{Paragraph 2: 顶会 Rebuttal 针对“创新性存疑（Limited Novelty / Incremental）”的有力澄清段 (\texttt{BP46})}
\noindent\textbf{Expressive Intent:}\enspace 我这一段想回复审稿人最致命的质疑：'The method seems like a simple combination of existing Module A and Module B (limited novelty).' 既不顶撞审稿人，又能用理论+消融数据证明“绝对不是简单拼凑”。\par
\noindent\textbf{Cohesion Formula:}\enspace {\small\color{AcadSlate}Clarify Distinction $\to$ Explain Why Naive Stacking Fails $\to$ Highlight Core Coupling Mechanism $\to$ Prove with 'Naive A+B' Ablation Data. 绑定逻辑规则：L019, L016.}\par\vspace{4pt}
\begin{ExampleBox}{Publication-Grade Authentic Text}
We appreciate the opportunity to clarify how our architecture differs fundamentally from a naive combination of spectral filtering and neural ODEs. Directly stacking an FFT block with a standard neural ODE fails under abrupt regime shifts because spectral normalization strips the exact boundary values required to initialize the ODE trajectory. Our core technical innovation lies in the bidirectional phase-scale coupling operator (Eq. 7), which uses instantaneous spectral entropy to gate ODE stiffness while feeding predicted scale bounds back to regularize frequency attention. To verify this empirically, we implemented a direct `Naive FFT + ODE` concatenation baseline in Table R2 below: our coupled formulation reduces 720-step MSE by 15.2\% over `Naive FFT + ODE`, confirming that the gain stems from our mathematical coupling rather than module stacking.
\end{ExampleBox}
\begin{TipBox}{Rhetorical Move Breakdown \& Constituent Analysis}
\begin{itemize}[leftmargin=1.5em]
  \item \textbf{S1 [礼貌接招并定位误解核心 Courteous Reframing]:} \textit{We appreciate the opportunity to clarify how our architecture differs fundamentally from a naive combination of spectral filtering and neural ODEs.}\par\noindent{\small 对应步骤：S1 [礼貌接招并定位误解核心 Courteous Reframing]：We appreciate Reviewer X's comment on architectural design and would like to clarify the fundamental distinction between our coupled formulation and a naive combination of [A] and [B]. | 句首锚定旧信息/逻辑视角，主干动词紧跟主语，复杂修饰置于宾语之后（符合尾重原则 L011 与旧新信息流 L001）。}
  \item \textbf{S2 [点破简单拼接为何会失败 Why Naive Combination Fails]:} \textit{Directly stacking an FFT block with a standard neural ODE fails under abrupt regime shifts because spectral normalization strips the exact boundary values required to initialize the ODE trajectory.}\par\noindent{\small 对应步骤：S2 [点破简单拼接为何会失败 Why Naive Combination Fails]：Directly stacking standard [A] and [B] fails on [Problem] because [Theoretical Conflict / Gradient Interference]. | 句首锚定旧信息/逻辑视角，主干动词紧跟主语，复杂修饰置于宾语之后（符合尾重原则 L011 与旧新信息流 L001）。}
  \item \textbf{S3 [阐明本文核心数学/机制桥梁 Our Core Coupling Innovation]:} \textit{Our core technical innovation lies in the bidirectional phase-scale coupling operator (Eq.}\par\noindent{\small 对应步骤：S3 [阐明本文核心数学/机制桥梁 Our Core Coupling Innovation]：Our core contribution lies not in using [A] or [B] in isolation, but in deriving the [Specific Coupling Mechanism C], which resolves [Conflict] by [Mathematical Guarantee]. | 句首锚定旧信息/逻辑视角，主干动词紧跟主语，复杂修饰置于宾语之后（符合尾重原则 L011 与旧新信息流 L001）。}
  \item \textbf{S4 [用直接对比实验一锤定音 Empirical Proof of Non-Triviality]:} \textit{7), which uses instantaneous spectral entropy to gate ODE stiffness while feeding predicted scale bounds back to regularize frequency attention.}\par\noindent{\small 对应步骤：S4 [用直接对比实验一锤定音 Empirical Proof of Non-Triviality]：To verify this empirically, we implemented the exact naive concatenation baseline (`Naive A+B`) in Table R2: our coupled formulation outperforms `Naive A+B` by [14.6\%] in MSE, confirming that the performance gain stems from our coupling mechanism rather than mere module stacking. | 句首锚定旧信息/逻辑视角，主干动词紧跟主语，复杂修饰置于宾语之后（符合尾重原则 L011 与旧新信息流 L001）。}
  \item \textbf{S4 [用直接对比实验一锤定音 Empirical Proof of Non-Triviality]:} \textit{To verify this empirically, we implemented a direct `Naive FFT + ODE` concatenation baseline in Table R2 below: our coupled formulation reduces 720-step MSE by 15.2\% over `Naive FFT + ODE`, confirming that the gain stems from our mathematical coupling rather than module stacking.}\par\noindent{\small 对应步骤：S4 [用直接对比实验一锤定音 Empirical Proof of Non-Triviality]：To verify this empirically, we implemented the exact naive concatenation baseline (`Naive A+B`) in Table R2: our coupled formulation outperforms `Naive A+B` by [14.6\%] in MSE, confirming that the performance gain stems from our coupling mechanism rather than mere module stacking. | 句首锚定旧信息/逻辑视角，主干动词紧跟主语，复杂修饰置于宾语之后（符合尾重原则 L011 与旧新信息流 L001）。}
\end{itemize}
\end{TipBox}
\medskip
\subsection*{Paragraph 3: 顶会 Rebuttal 针对“缺少最新基线或算力不公平（Missing Baselines / Fairness）”的补全段 (\texttt{BP47})}
\noindent\textbf{Expressive Intent:}\enspace 我这一段想回复审稿人要求补测最新基线模型，或者质疑我们的提升是不是因为参数量更大/训练更久。\par
\noindent\textbf{Cohesion Formula:}\enspace {\small\color{AcadSlate}Direct Fulfilment $\to$ Margin over New Baselines $\to$ Parameter-Reduced Variant (`Ours-Small`) Proof $\to$ Inductive Bias Takeaway. 绑定逻辑规则：L019, L010.}\par\vspace{4pt}
\begin{ExampleBox}{Publication-Grade Authentic Text}
Following your suggestion, we benchmarked against the two recent baselines (TimeMixer and FreTS) and evaluated a capacity-reduced variant of our architecture (`DyTS-Lite`) in Table R3. Across all horizons on Weather and ECL, standard DyTS-Net reduces MSE by 9.6\% over TimeMixer and 13.4\% over FreTS. Furthermore, to ensure strict capacity fairness, `DyTS-Lite` shrinks the hidden dimension to $d=64$---requiring 42\% fewer parameters and 55\% fewer FLOPs than TimeMixer---yet still outperforms TimeMixer by 6.8\% in average MSE. This confirms that our gains stem from structural-scale disentanglement rather than model size.
\end{ExampleBox}
\begin{TipBox}{Rhetorical Move Breakdown \& Constituent Analysis}
\begin{itemize}[leftmargin=1.5em]
  \item \textbf{S1 [直接响应并呈上新表 Direct Fulfilment]:} \textit{Following your suggestion, we benchmarked against the two recent baselines (TimeMixer and FreTS) and evaluated a capacity-reduced variant of our architecture (`DyTS-Lite`) in Table R3.}\par\noindent{\small 对应步骤：S1 [直接响应并呈上新表 Direct Fulfilment]：Following your suggestion, we have evaluated [Requested Baselines X and Y] under identical splits and added a parameter-matched variant of our model in Table R3. | 句首锚定旧信息/逻辑视角，主干动词紧跟主语，复杂修饰置于宾语之后（符合尾重原则 L011 与旧新信息流 L001）。}
  \item \textbf{S2 [对比最新基线结果 Performance vs. New Baselines]:} \textit{Across all horizons on Weather and ECL, standard DyTS-Net reduces MSE by 9.6\% over TimeMixer and 13.4\% over FreTS.}\par\noindent{\small 对应步骤：S2 [对比最新基线结果 Performance vs. New Baselines]：Across all four horizons on [Datasets], our method outperforms [New Baseline X] by an average of [9.8\%] in MSE and [New Baseline Y] by [12.1\%]. | 句首锚定旧信息/逻辑视角，主干动词紧跟主语，复杂修饰置于宾语之后（符合尾重原则 L011 与旧新信息流 L001）。}
  \item \textbf{S3 [打消参数量/算力质疑 Parameter \& FLOPs Parity Proof]:} \textit{Furthermore, to ensure strict capacity fairness, `DyTS-Lite` shrinks the hidden dimension to $d=64$---requiring 42\% fewer parameters and 55\% fewer FLOPs than TimeMixer---yet still outperforms TimeMixer by 6.8\% in average MSE.}\par\noindent{\small 对应步骤：S3 [打消参数量/算力质疑 Parameter \& FLOPs Parity Proof]：Regarding model capacity, even when we scale down our hidden dimension from $d=256$ to $d=128$ (`Ours-Small`, using 38\% fewer parameters and 52\% fewer FLOPs than [Baseline X]), `Ours-Small` still outperforms [Baseline X] by [7.4\%]. | 句首锚定旧信息/逻辑视角，主干动词紧跟主语，复杂修饰置于宾语之后（符合尾重原则 L011 与旧新信息流 L001）。}
  \item \textbf{S4 [结论锁定 Conclusion]:} \textit{This confirms that our gains stem from structural-scale disentanglement rather than model size.}\par\noindent{\small 对应步骤：S4 [结论锁定 Conclusion]：These results confirm that our improvements derive from inductive bias rather than increased parameter count. | 句首锚定旧信息/逻辑视角，主干动词紧跟主语，复杂修饰置于宾语之后（符合尾重原则 L011 与旧新信息流 L001）。}
\end{itemize}
\end{TipBox}
\medskip
\subsection*{Paragraph 4: 消融实验逐组件归因论证段（Counterfactual Ablation Study） (\texttt{BP19})}
\noindent\textbf{Expressive Intent:}\enspace 我这一段想汇报消融实验（Table 2）：证明我们的模块 A、模块 B 和正则项 C 每一个都不可或缺，且各自负责解决不同的特定子问题。\par
\noindent\textbf{Cohesion Formula:}\enspace {\small\color{AcadSlate}Ablation Overview $\to$ w/o Module A (Specific Failure) $\to$ w/o Module B (Specific Failure) $\to$ Cross-Module Synergy. 绑定逻辑规则：L016, L007.}\par\vspace{4pt}
\begin{ExampleBox}{Publication-Grade Authentic Text}
To isolate the contribution of each design choice, Table 2 compares full DyTS-Net against three ablated variants across ETTm2 and ECL. Removing the neural-ODE scale tracker (w/o Scale-ODE) triggers the sharpest degradation during seasonal transition months on ETTm2, raising 720-step MSE by 19.4\% as the decoder fails to track shifting baseline temperatures. Conversely, replacing phase-aligned graph attention with standard dot-product attention (w/o Phase-Attn) degrades performance most severely on the 321-channel ECL grid (+13.8\% MSE), where unaligned high-frequency load spikes corrupt cross-substation aggregation. Jointly deploying both modules yields complementary gains, as accurate scale tracking prevents amplitude distortion from biasing phase estimation.
\end{ExampleBox}
\begin{TipBox}{Rhetorical Move Breakdown \& Constituent Analysis}
\begin{itemize}[leftmargin=1.5em]
  \item \textbf{S1 [消融实验总括 Ablation Setup]:} \textit{To isolate the contribution of each design choice, Table 2 compares full DyTS-Net against three ablated variants across ETTm2 and ECL.}\par\noindent{\small 对应步骤：S1 [消融实验总括 Ablation Setup]：To isolate the contribution of each architectural component, we evaluate four ablated variants in Table 2 by progressively removing [Mod A], [Mod B], and [Loss C]. | 句首锚定旧信息/逻辑视角，主干动词紧跟主语，复杂修饰置于宾语之后（符合尾重原则 L011 与旧新信息流 L001）。}
  \item \textbf{S2 [移除模块 A 的冲击与归因 Impact of Removing Mod A]:} \textit{Removing the neural-ODE scale tracker (w/o Scale-ODE) triggers the sharpest degradation during seasonal transition months on ETTm2, raising 720-step MSE by 19.4\% as the decoder fails to track shifting baseline temperatures.}\par\noindent{\small 对应步骤：S2 [移除模块 A 的冲击与归因 Impact of Removing Mod A]：Removing [Mod A] (w/o A) causes the largest degradation on [Non-stationary/Volatile Subset], increasing MSE by [X\%] because the model loses [Specific Capability A]. | 句首锚定旧信息/逻辑视角，主干动词紧跟主语，复杂修饰置于宾语之后（符合尾重原则 L011 与旧新信息流 L001）。}
  \item \textbf{S3 [移除模块 B 的冲击与归因 Impact of Removing Mod B]:} \textit{Conversely, replacing phase-aligned graph attention with standard dot-product attention (w/o Phase-Attn) degrades performance most severely on the 321-channel ECL grid (+13.8\% MSE), where unaligned high-frequency load spikes corrupt cross-substation aggregation.}\par\noindent{\small 对应步骤：S3 [移除模块 B 的冲击与归因 Impact of Removing Mod B]：Similarly, replacing [Mod B] with standard [Naive Operator] degrades performance by [Y\%] on high-dimensional benchmarks, confirming its role in [Specific Capability B]. | 句首锚定旧信息/逻辑视角，主干动词紧跟主语，复杂修饰置于宾语之后（符合尾重原则 L011 与旧新信息流 L001）。}
  \item \textbf{S4 [协同效应总结 Synergistic Conclusion]:} \textit{Jointly deploying both modules yields complementary gains, as accurate scale tracking prevents amplitude distortion from biasing phase estimation.}\par\noindent{\small 对应步骤：S4 [协同效应总结 Synergistic Conclusion]：Combining all components yields super-additive gains, as [Mod A] stabilizes input statistics for [Mod B]'s relational reasoning. | 句首锚定旧信息/逻辑视角，主干动词紧跟主语，复杂修饰置于宾语之后（符合尾重原则 L011 与旧新信息流 L001）。}
\end{itemize}
\end{TipBox}
\medskip
\subsection*{Paragraph 5: 超参数敏感性与极端噪声鲁棒性分析段（Sensitivity \& Robustness Analysis） (\texttt{BP20})}
\noindent\textbf{Expressive Intent:}\enspace 我这一段想回应审稿人最常问的：你的模型是不是对超参数很敏感？在缺失值或传感器强噪声下会不会突然崩盘？\par
\noindent\textbf{Cohesion Formula:}\enspace {\small\color{AcadSlate}Goal $\to$ Wide Hyperparameter Plateau $\to$ Stress-Test Comparison $\to$ Mechanistic Explanation. 绑定逻辑规则：L010, L019.}\par\vspace{4pt}
\begin{ExampleBox}{Publication-Grade Authentic Text}
Figure 4 evaluates hyperparameter sensitivity and resilience against severe observation corruption. Varying the phase regularization weight $\lambda$ across $[0.05, 0.50]$ alters validation MSE by less than 1.4\%, indicating a wide, tuning-friendly operating plateau around $\lambda=0.15$. Under stress tests where 30\% of sensor readings are randomly masked or corrupted by impulse noise ($\text{SNR} = 5\text{ dB}$), point-wise Transformer baselines suffer a 34.2\% error surge, whereas DyTS-Net degrades by only 4.8\%. This robustness stems from our spectral energy gate, which automatically attenuates high-entropy corrupted channels prior to cross-variate message passing.
\end{ExampleBox}
\begin{TipBox}{Rhetorical Move Breakdown \& Constituent Analysis}
\begin{itemize}[leftmargin=1.5em]
  \item \textbf{S1 [分析目的 Objective]:} \textit{Figure 4 evaluates hyperparameter sensitivity and resilience against severe observation corruption.}\par\noindent{\small 对应步骤：S1 [分析目的 Objective]：We further examine sensitivity to key hyperparameters and resilience against severe sensor corruption in Figure 4. | 句首锚定旧信息/逻辑视角，主干动词紧跟主语，复杂修饰置于宾语之后（符合尾重原则 L011 与旧新信息流 L001）。}
  \item \textbf{S2 [超参数平坦区间 Plateau Stability]:} \textit{Varying the phase regularization weight $\lambda$ across $[0.05, 0.50]$ alters validation MSE by less than 1.4\%, indicating a wide, tuning-friendly operating plateau around $\lambda=0.15$.}\par\noindent{\small 对应步骤：S2 [超参数平坦区间 Plateau Stability]：Varying [Hyperparameter $\lambda$] across two orders of magnitude ($[0.01, 1.0]$) alters validation MSE by less than [1.8\%], demonstrating a wide stable operating plateau. | 句首锚定旧信息/逻辑视角，主干动词紧跟主语，复杂修饰置于宾语之后（符合尾重原则 L011 与旧新信息流 L001）。}
  \item \textbf{S3 [极端边缘工况压力测试 Stress Test under Corruption]:} \textit{Under stress tests where 30\% of sensor readings are randomly masked or corrupted by impulse noise ($\text{SNR} = 5\text{ dB}$), point-wise Transformer baselines suffer a 34.2\% error surge, whereas DyTS-Net degrades by only 4.8\%.}\par\noindent{\small 对应步骤：S3 [极端边缘工况压力测试 Stress Test under Corruption]：Under synthetic sensor dropouts up to [40\%] and Gaussian noise ($\sigma=0.3$), baseline accuracy plummets by over [35\%], whereas our model degrades by less than [6.5\%]. | 句首锚定旧信息/逻辑视角，主干动词紧跟主语，复杂修饰置于宾语之后（符合尾重原则 L011 与旧新信息流 L001）。}
  \item \textbf{S4 [鲁棒性物理归因 Robustness Mechanism]:} \textit{This robustness stems from our spectral energy gate, which automatically attenuates high-entropy corrupted channels prior to cross-variate message passing.}\par\noindent{\small 对应步骤：S4 [鲁棒性物理归因 Robustness Mechanism]：This resilience stems from [Specific Mechanism], which filters out corrupted channels before global aggregation. | 句首锚定旧信息/逻辑视角，主干动词紧跟主语，复杂修饰置于宾语之后（符合尾重原则 L011 与旧新信息流 L001）。}
\end{itemize}
\end{TipBox}
\medskip
\subsection*{Paragraph 6: 局限性（Limitations）客观陈述与未来扩展段（Constructive Limitations） (\texttt{BP23})}
\noindent\textbf{Expressive Intent:}\enspace 我这一段想写论文末尾或 NeurIPS/ICLR Checklist 要求的 Limitations 段落：诚实指出当前方法的两个客观边界，但绝不贬损论文的核心贡献。\par
\noindent\textbf{Cohesion Formula:}\enspace {\small\color{AcadSlate}Two Objective Boundaries $\to$ [Boundary 1 + Technical Remedy] $\to$ [Boundary 2 + Stochastic Extension]. 绑定逻辑规则：L017.}\par\vspace{4pt}
\begin{ExampleBox}{Publication-Grade Authentic Text}
Despite its strong empirical performance, our framework exhibits two theoretical boundaries worth noting. First, while our spectral branch scales quasi-linearly with sequence length, computing pairwise Hilbert phase differences scales quadratically with channel count ($O(C^2)$), which increases memory footprint when $C > 2,000$; incorporating anchor-based low-rank graph sparsification provides a direct remedy. Second, our continuous neural-ODE scale tracker assumes piece-wise smooth drift and may smooth over instantaneous step discontinuities caused by sensor recalibration. Extending the continuous drift prior to jump-diffusion stochastic differential equations offers a principled direction for future investigation.
\end{ExampleBox}
\begin{TipBox}{Rhetorical Move Breakdown \& Constituent Analysis}
\begin{itemize}[leftmargin=1.5em]
  \item \textbf{S1 [边界 1:} \textit{Despite its strong empirical performance, our framework exhibits two theoretical boundaries worth noting.}\par\noindent{\small 对应步骤：S1 [边界 1：极端高维或极度稀疏场景 Boundary 1]：Despite its effectiveness across standard benchmarks, our framework currently exhibits two practical boundaries. | 句首锚定旧信息/逻辑视角，主干动词紧跟主语，复杂修饰置于宾语之后（符合尾重原则 L011 与旧新信息流 L001）。}
  \item \textbf{S2 [边界 1 机理分析与解法 Mechanism \& Solution 1]:} \textit{First, while our spectral branch scales quasi-linearly with sequence length, computing pairwise Hilbert phase differences scales quadratically with channel count ($O(C^2)$), which increases memory footprint when $C > 2,000$; incorporating anchor-based low-rank graph sparsification provides a direct remedy.}\par\noindent{\small 对应步骤：S2 [边界 1 机理分析与解法 Mechanism \& Solution 1]：First, when the variate dimension scales beyond $C > 5,000$ unclustered sensors, constructing the cross-channel phase graph incurs noticeable memory overhead; integrating hierarchical sensor clustering offers a natural remedy. | 句首锚定旧信息/逻辑视角，主干动词紧跟主语，复杂修饰置于宾语之后（符合尾重原则 L011 与旧新信息流 L001）。}
  \item \textbf{S3 [边界 2:} \textit{Second, our continuous neural-ODE scale tracker assumes piece-wise smooth drift and may smooth over instantaneous step discontinuities caused by sensor recalibration.}\par\noindent{\small 对应步骤：S3 [边界 2：特定假设限制 Boundary 2]：Second, our continuous scale tracker assumes piece-wise differentiable drift and may under-estimate instantaneous discontinuous jumps caused by hardware resets. | 句首锚定旧信息/逻辑视角，主干动词紧跟主语，复杂修饰置于宾语之后（符合尾重原则 L011 与旧新信息流 L001）。}
  \item \textbf{S4 [未来演进方向 Future Extension]:} \textit{Extending the continuous drift prior to jump-diffusion stochastic differential equations offers a principled direction for future investigation.}\par\noindent{\small 对应步骤：S4 [未来演进方向 Future Extension]：Coupling our neural ODE tracker with jump-diffusion stochastic processes represents a promising avenue for future work. | 句首锚定旧信息/逻辑视角，主干动词紧跟主语，复杂修饰置于宾语之后（符合尾重原则 L011 与旧新信息流 L001）。}
\end{itemize}
\end{TipBox}
\medskip
\clearpage
\section{FT10: 学术答辩与组会 6 大高频刁钻提问应答全案（Novelty、超参敏感性、极端噪声、计算复杂度、失败案例与未来扩展）}\label{comp:FT10}
\begin{ConceptBox}{Composition Architecture \& Structural Blueprint Chain}
\textbf{Genre:} \texttt{Academic\_QA\_Defense} \quad \textbf{Paragraph Count:} 6\par
\textbf{Blueprint Chain:} \texttt{BP44 $\to$ BP46 $\to$ BP20 $\to$ BP21 $\to$ BP22 $\to$ BP23}
\end{ConceptBox}
\subsection*{Paragraph 1: 演讲总结页（Take-home Message）与问答（Q\&A）从容应答段 (\texttt{BP44})}
\noindent\textbf{Expressive Intent:}\enspace 我这一段想用于演讲最后一页总结（Take-home Messages）以及回答台下教授/审稿人犀利提问时的标准应答模板。\par
\noindent\textbf{Cohesion Formula:}\enspace {\small\color{AcadSlate}Single Takeaway $\to$ Quantitative Triad $\to$ Q\&A Direct Answer $\to$ Empirical Grounding. 绑定逻辑规则：L018, L019.}\par\vspace{4pt}
\begin{ExampleBox}{Publication-Grade Authentic Text}
To wrap up, the key takeaway from our work is that non-stationary forecasting requires separating shift-invariant phase structure from continuous amplitude dynamics. Doing so lets DyTS-Net cut long-horizon error by up to 17\% while running at half the computational cost, and our full code and weights are live at the QR code on screen. --- [Q\&A Response]: Thank you for that sharp question on irregular missing timestamps. The short answer is that our neural-ODE branch handles arbitrary time intervals natively: in our 30\% random-dropout stress test, MSE degraded by less than 5\%, and we can share the backup slide on irregular sampling right here.
\end{ExampleBox}
\begin{TipBox}{Rhetorical Move Breakdown \& Constituent Analysis}
\begin{itemize}[leftmargin=1.5em]
  \item \textbf{S1 [凝练核心带走信息 Core Take-Home Message]:} \textit{To wrap up, the key takeaway from our work is that non-stationary forecasting requires separating shift-invariant phase structure from continuous amplitude dynamics.}\par\noindent{\small 对应步骤：S1 [凝练核心带走信息 Core Take-Home Message]：To wrap up, if there is one takeaway from our work today, it is that [One-Sentence Core Principle]. | 句首锚定旧信息/逻辑视角，主干动词紧跟主语，复杂修饰置于宾语之后（符合尾重原则 L011 与旧新信息流 L001）。}
  \item \textbf{S2 [成果三点速览 Rapid Recap]:} \textit{Doing so lets DyTS-Net cut long-horizon error by up to 17\% while running at half the computational cost, and our full code and weights are live at the QR code on screen.}\par\noindent{\small 对应步骤：S2 [成果三点速览 Rapid Recap]：By combining [A] and [B], [Model] delivers [Metric 1], [Metric 2], and [Open-Source/Practical Asset]. | 句首锚定旧信息/逻辑视角，主干动词紧跟主语，复杂修饰置于宾语之后（符合尾重原则 L011 与旧新信息流 L001）。}
  \item \textbf{S3 [Q\&A 礼貌接招与结论前置 Q\&A Direct Answer]:} \textit{Thank you for that sharp question on irregular missing timestamps.}\par\noindent{\small 对应步骤：S3 [Q\&A 礼貌接招与结论前置 Q\&A Direct Answer]：(For Q\&A) That is a great question regarding [Topic]---the short answer is [Direct Conclusion]. | 句首锚定旧信息/逻辑视角，主干动词紧跟主语，复杂修饰置于宾语之后（符合尾重原则 L011 与旧新信息流 L001）。}
  \item \textbf{S4 [Q\&A 展开与诚实边界 Q\&A Elaboration]:} \textit{The short answer is that our neural-ODE branch handles arbitrary time intervals natively: in our 30\% random-dropout stress test, MSE degraded by less than 5\%, and we can share the backup slide on irregular sampling right here.}\par\noindent{\small 对应步骤：S4 [Q\&A 展开与诚实边界 Q\&A Elaboration]：In our current experiments, we observed [Empirical Fact], and extending this to [Questioner's Scenario] can be handled by [Concrete Mechanism]. | 句首锚定旧信息/逻辑视角，主干动词紧跟主语，复杂修饰置于宾语之后（符合尾重原则 L011 与旧新信息流 L001）。}
\end{itemize}
\end{TipBox}
\medskip
\subsection*{Paragraph 2: 顶会 Rebuttal 针对“创新性存疑（Limited Novelty / Incremental）”的有力澄清段 (\texttt{BP46})}
\noindent\textbf{Expressive Intent:}\enspace 我这一段想回复审稿人最致命的质疑：'The method seems like a simple combination of existing Module A and Module B (limited novelty).' 既不顶撞审稿人，又能用理论+消融数据证明“绝对不是简单拼凑”。\par
\noindent\textbf{Cohesion Formula:}\enspace {\small\color{AcadSlate}Clarify Distinction $\to$ Explain Why Naive Stacking Fails $\to$ Highlight Core Coupling Mechanism $\to$ Prove with 'Naive A+B' Ablation Data. 绑定逻辑规则：L019, L016.}\par\vspace{4pt}
\begin{ExampleBox}{Publication-Grade Authentic Text}
We appreciate the opportunity to clarify how our architecture differs fundamentally from a naive combination of spectral filtering and neural ODEs. Directly stacking an FFT block with a standard neural ODE fails under abrupt regime shifts because spectral normalization strips the exact boundary values required to initialize the ODE trajectory. Our core technical innovation lies in the bidirectional phase-scale coupling operator (Eq. 7), which uses instantaneous spectral entropy to gate ODE stiffness while feeding predicted scale bounds back to regularize frequency attention. To verify this empirically, we implemented a direct `Naive FFT + ODE` concatenation baseline in Table R2 below: our coupled formulation reduces 720-step MSE by 15.2\% over `Naive FFT + ODE`, confirming that the gain stems from our mathematical coupling rather than module stacking.
\end{ExampleBox}
\begin{TipBox}{Rhetorical Move Breakdown \& Constituent Analysis}
\begin{itemize}[leftmargin=1.5em]
  \item \textbf{S1 [礼貌接招并定位误解核心 Courteous Reframing]:} \textit{We appreciate the opportunity to clarify how our architecture differs fundamentally from a naive combination of spectral filtering and neural ODEs.}\par\noindent{\small 对应步骤：S1 [礼貌接招并定位误解核心 Courteous Reframing]：We appreciate Reviewer X's comment on architectural design and would like to clarify the fundamental distinction between our coupled formulation and a naive combination of [A] and [B]. | 句首锚定旧信息/逻辑视角，主干动词紧跟主语，复杂修饰置于宾语之后（符合尾重原则 L011 与旧新信息流 L001）。}
  \item \textbf{S2 [点破简单拼接为何会失败 Why Naive Combination Fails]:} \textit{Directly stacking an FFT block with a standard neural ODE fails under abrupt regime shifts because spectral normalization strips the exact boundary values required to initialize the ODE trajectory.}\par\noindent{\small 对应步骤：S2 [点破简单拼接为何会失败 Why Naive Combination Fails]：Directly stacking standard [A] and [B] fails on [Problem] because [Theoretical Conflict / Gradient Interference]. | 句首锚定旧信息/逻辑视角，主干动词紧跟主语，复杂修饰置于宾语之后（符合尾重原则 L011 与旧新信息流 L001）。}
  \item \textbf{S3 [阐明本文核心数学/机制桥梁 Our Core Coupling Innovation]:} \textit{Our core technical innovation lies in the bidirectional phase-scale coupling operator (Eq.}\par\noindent{\small 对应步骤：S3 [阐明本文核心数学/机制桥梁 Our Core Coupling Innovation]：Our core contribution lies not in using [A] or [B] in isolation, but in deriving the [Specific Coupling Mechanism C], which resolves [Conflict] by [Mathematical Guarantee]. | 句首锚定旧信息/逻辑视角，主干动词紧跟主语，复杂修饰置于宾语之后（符合尾重原则 L011 与旧新信息流 L001）。}
  \item \textbf{S4 [用直接对比实验一锤定音 Empirical Proof of Non-Triviality]:} \textit{7), which uses instantaneous spectral entropy to gate ODE stiffness while feeding predicted scale bounds back to regularize frequency attention.}\par\noindent{\small 对应步骤：S4 [用直接对比实验一锤定音 Empirical Proof of Non-Triviality]：To verify this empirically, we implemented the exact naive concatenation baseline (`Naive A+B`) in Table R2: our coupled formulation outperforms `Naive A+B` by [14.6\%] in MSE, confirming that the performance gain stems from our coupling mechanism rather than mere module stacking. | 句首锚定旧信息/逻辑视角，主干动词紧跟主语，复杂修饰置于宾语之后（符合尾重原则 L011 与旧新信息流 L001）。}
  \item \textbf{S4 [用直接对比实验一锤定音 Empirical Proof of Non-Triviality]:} \textit{To verify this empirically, we implemented a direct `Naive FFT + ODE` concatenation baseline in Table R2 below: our coupled formulation reduces 720-step MSE by 15.2\% over `Naive FFT + ODE`, confirming that the gain stems from our mathematical coupling rather than module stacking.}\par\noindent{\small 对应步骤：S4 [用直接对比实验一锤定音 Empirical Proof of Non-Triviality]：To verify this empirically, we implemented the exact naive concatenation baseline (`Naive A+B`) in Table R2: our coupled formulation outperforms `Naive A+B` by [14.6\%] in MSE, confirming that the performance gain stems from our coupling mechanism rather than mere module stacking. | 句首锚定旧信息/逻辑视角，主干动词紧跟主语，复杂修饰置于宾语之后（符合尾重原则 L011 与旧新信息流 L001）。}
\end{itemize}
\end{TipBox}
\medskip
\subsection*{Paragraph 3: 超参数敏感性与极端噪声鲁棒性分析段（Sensitivity \& Robustness Analysis） (\texttt{BP20})}
\noindent\textbf{Expressive Intent:}\enspace 我这一段想回应审稿人最常问的：你的模型是不是对超参数很敏感？在缺失值或传感器强噪声下会不会突然崩盘？\par
\noindent\textbf{Cohesion Formula:}\enspace {\small\color{AcadSlate}Goal $\to$ Wide Hyperparameter Plateau $\to$ Stress-Test Comparison $\to$ Mechanistic Explanation. 绑定逻辑规则：L010, L019.}\par\vspace{4pt}
\begin{ExampleBox}{Publication-Grade Authentic Text}
Figure 4 evaluates hyperparameter sensitivity and resilience against severe observation corruption. Varying the phase regularization weight $\lambda$ across $[0.05, 0.50]$ alters validation MSE by less than 1.4\%, indicating a wide, tuning-friendly operating plateau around $\lambda=0.15$. Under stress tests where 30\% of sensor readings are randomly masked or corrupted by impulse noise ($\text{SNR} = 5\text{ dB}$), point-wise Transformer baselines suffer a 34.2\% error surge, whereas DyTS-Net degrades by only 4.8\%. This robustness stems from our spectral energy gate, which automatically attenuates high-entropy corrupted channels prior to cross-variate message passing.
\end{ExampleBox}
\begin{TipBox}{Rhetorical Move Breakdown \& Constituent Analysis}
\begin{itemize}[leftmargin=1.5em]
  \item \textbf{S1 [分析目的 Objective]:} \textit{Figure 4 evaluates hyperparameter sensitivity and resilience against severe observation corruption.}\par\noindent{\small 对应步骤：S1 [分析目的 Objective]：We further examine sensitivity to key hyperparameters and resilience against severe sensor corruption in Figure 4. | 句首锚定旧信息/逻辑视角，主干动词紧跟主语，复杂修饰置于宾语之后（符合尾重原则 L011 与旧新信息流 L001）。}
  \item \textbf{S2 [超参数平坦区间 Plateau Stability]:} \textit{Varying the phase regularization weight $\lambda$ across $[0.05, 0.50]$ alters validation MSE by less than 1.4\%, indicating a wide, tuning-friendly operating plateau around $\lambda=0.15$.}\par\noindent{\small 对应步骤：S2 [超参数平坦区间 Plateau Stability]：Varying [Hyperparameter $\lambda$] across two orders of magnitude ($[0.01, 1.0]$) alters validation MSE by less than [1.8\%], demonstrating a wide stable operating plateau. | 句首锚定旧信息/逻辑视角，主干动词紧跟主语，复杂修饰置于宾语之后（符合尾重原则 L011 与旧新信息流 L001）。}
  \item \textbf{S3 [极端边缘工况压力测试 Stress Test under Corruption]:} \textit{Under stress tests where 30\% of sensor readings are randomly masked or corrupted by impulse noise ($\text{SNR} = 5\text{ dB}$), point-wise Transformer baselines suffer a 34.2\% error surge, whereas DyTS-Net degrades by only 4.8\%.}\par\noindent{\small 对应步骤：S3 [极端边缘工况压力测试 Stress Test under Corruption]：Under synthetic sensor dropouts up to [40\%] and Gaussian noise ($\sigma=0.3$), baseline accuracy plummets by over [35\%], whereas our model degrades by less than [6.5\%]. | 句首锚定旧信息/逻辑视角，主干动词紧跟主语，复杂修饰置于宾语之后（符合尾重原则 L011 与旧新信息流 L001）。}
  \item \textbf{S4 [鲁棒性物理归因 Robustness Mechanism]:} \textit{This robustness stems from our spectral energy gate, which automatically attenuates high-entropy corrupted channels prior to cross-variate message passing.}\par\noindent{\small 对应步骤：S4 [鲁棒性物理归因 Robustness Mechanism]：This resilience stems from [Specific Mechanism], which filters out corrupted channels before global aggregation. | 句首锚定旧信息/逻辑视角，主干动词紧跟主语，复杂修饰置于宾语之后（符合尾重原则 L011 与旧新信息流 L001）。}
\end{itemize}
\end{TipBox}
\medskip
\subsection*{Paragraph 4: 极端/反常情况与失败案例坦诚剖析段（Failure Case \& Outlier Analysis） (\texttt{BP21})}
\noindent\textbf{Expressive Intent:}\enspace 我这一段想通过可视化案例（Case Study / Qualitative Analysis），既展示模型在复杂跳变下的精准跟踪，又坦诚剖析一个模型预测失准的极端案例。\par
\noindent\textbf{Cohesion Formula:}\enspace {\small\color{AcadSlate}Visual Setup $\to$ Success under Shift $\to$ Honest Failure Case under Unobserved Exogenous Shock $\to$ Actionable Insight. 绑定逻辑规则：L017, L018.}\par\vspace{4pt}
\begin{ExampleBox}{Publication-Grade Authentic Text}
Figure 5 visualizes 336-step forecasts during an abrupt voltage drop on the ETTm1 transformer load benchmark. Whereas PatchTST and iTransformer exhibit a 14-step phase lag and severe post-recovery overshoot, DyTS-Net tracks both the sharp downward ramp and subsequent stabilization within 5\% amplitude error. Nevertheless, the failure case in Panel (d) reveals that when an unscheduled manual grid curtailment occurs without preceding spectral precursors, our endogenous model experiences a 6-step delayed response. Anticipating such zero-precursor manual interventions requires integrating exogenous dispatch logs alongside historical telemetry.
\end{ExampleBox}
\begin{TipBox}{Rhetorical Move Breakdown \& Constituent Analysis}
\begin{itemize}[leftmargin=1.5em]
  \item \textbf{S1 [定性对比切入 Qualitative Setup]:} \textit{Figure 5 visualizes 336-step forecasts during an abrupt voltage drop on the ETTm1 transformer load benchmark.}\par\noindent{\small 对应步骤：S1 [定性对比切入 Qualitative Setup]：Figure 5 visualizes predicted trajectories during an abrupt regime shift on [Dataset/Task]. | 句首锚定旧信息/逻辑视角，主干动词紧跟主语，复杂修饰置于宾语之后（符合尾重原则 L011 与旧新信息流 L001）。}
  \item \textbf{S2 [成功案例亮点 Success Highlight]:} \textit{Whereas PatchTST and iTransformer exhibit a 14-step phase lag and severe post-recovery overshoot, DyTS-Net tracks both the sharp downward ramp and subsequent stabilization within 5\% amplitude error.}\par\noindent{\small 对应步骤：S2 [成功案例亮点 Success Highlight]：While [Baseline] lags behind the sudden step change by [K steps] and overshoots, our model rapidly locks onto the new amplitude regime within [M steps]. | 句首锚定旧信息/逻辑视角，主干动词紧跟主语，复杂修饰置于宾语之后（符合尾重原则 L011 与旧新信息流 L001）。}
  \item \textbf{S3 [失败案例坦诚剖析 Honest Failure Mode]:} \textit{Nevertheless, the failure case in Panel (d) reveals that when an unscheduled manual grid curtailment occurs without preceding spectral precursors, our endogenous model experiences a 6-step delayed response.}\par\noindent{\small 对应步骤：S3 [失败案例坦诚剖析 Honest Failure Mode]：Nevertheless, inspecting the bottom-right failure case reveals that when an external manual override occurs without historical precursors, all purely endogenous forecasters---including ours---exhibit a transient lag. | 句首锚定旧信息/逻辑视角，主干动词紧跟主语，复杂修饰置于宾语之后（符合尾重原则 L011 与旧新信息流 L001）。}
  \item \textbf{S4 [改进启示 Constructive Takeaway]:} \textit{Anticipating such zero-precursor manual interventions requires integrating exogenous dispatch logs alongside historical telemetry.}\par\noindent{\small 对应步骤：S4 [改进启示 Constructive Takeaway]：Incorporating exogenous event logs or textual maintenance records is necessary to anticipate such unscheduled interventions. | 句首锚定旧信息/逻辑视角，主干动词紧跟主语，复杂修饰置于宾语之后（符合尾重原则 L011 与旧新信息流 L001）。}
\end{itemize}
\end{TipBox}
\medskip
\subsection*{Paragraph 5: 计算开销、推理延迟与精度帕累托权衡段（Efficiency-Accuracy Pareto Analysis） (\texttt{BP22})}
\noindent\textbf{Expressive Intent:}\enspace 我这一段想说明我们的模型不仅准，而且算得快、显存省，在“精度-延迟/显存”帕累托前沿（Pareto Frontier）上显著优于现有大模型。\par
\noindent\textbf{Cohesion Formula:}\enspace {\small\color{AcadSlate}Hardware Benchmark $\to$ Complexity Reduction \& Memory Gain $\to$ Real-Time Latency Margin $\to$ Pareto Conclusion. 绑定逻辑规则：L010, L023.}\par\vspace{4pt}
\begin{ExampleBox}{Publication-Grade Authentic Text}
Figure 6 plots the MSE-latency Pareto frontier measured on an onboard NVIDIA Orin AGX module. By replacing $O(L^2)$ point-wise self-attention with $O(L \log L)$ spectral filtering and compact patch tokens, DyTS-Net reduces peak GPU memory by 62\% and training wall-clock time by 3.4x compared with iTransformer. At inference time, generating a 16-step visuomotor trajectory requires only 6.8 ms---well below the 20 ms cycle budget of a 50 Hz industrial robot controller---while outperforming a 200 ms DDPM baseline by 11.4\% in task success rate.
\end{ExampleBox}
\begin{TipBox}{Rhetorical Move Breakdown \& Constituent Analysis}
\begin{itemize}[leftmargin=1.5em]
  \item \textbf{S1 [硬件与测试基准 Hardware Setup]:} \textit{Figure 6 plots the MSE-latency Pareto frontier measured on an onboard NVIDIA Orin AGX module.}\par\noindent{\small 对应步骤：S1 [硬件与测试基准 Hardware Setup]：Figure 6 compares the accuracy-efficiency Pareto frontier across all baselines measured on a single NVIDIA RTX 4090 GPU. | 句首锚定旧信息/逻辑视角，主干动词紧跟主语，复杂修饰置于宾语之后（符合尾重原则 L011 与旧新信息流 L001）。}
  \item \textbf{S2 [复杂度理论与实测对比 Complexity \& Runtime]:} \textit{By replacing $O(L^2)$ point-wise self-attention with $O(L \log L)$ spectral filtering and compact patch tokens, DyTS-Net reduces peak GPU memory by 62\% and training wall-clock time by 3.4x compared with iTransformer.}\par\noindent{\small 对应步骤：S2 [复杂度理论与实测对比 Complexity \& Runtime]：By replacing $O(L^2)$ point-wise attention with $O(L \log L)$ spectral filtering and patch aggregation, our model cuts training time per epoch by [X\%] and peak VRAM by [Y\%]. | 句首锚定旧信息/逻辑视角，主干动词紧跟主语，复杂修饰置于宾语之后（符合尾重原则 L011 与旧新信息流 L001）。}
  \item \textbf{S3 [实时闭环部署可行性 Real-Time Viability]:} \textit{At inference time, generating a 16-step visuomotor trajectory requires only 6.8 ms---well below the 20 ms cycle budget of a 50 Hz industrial robot controller---while outperforming a 200 ms DDPM baseline by 11.4\% in task success rate.}\par\noindent{\small 对应步骤：S3 [实时闭环部署可行性 Real-Time Viability]：At inference time, generating a 720-step forecast (or 16-step robot action chunk) requires only [4.2 ms], well below the [20 ms] real-time control deadline. | 句首锚定旧信息/逻辑视角，主干动词紧跟主语，复杂修饰置于宾语之后（符合尾重原则 L011 与旧新信息流 L001）。}
\end{itemize}
\end{TipBox}
\medskip
\subsection*{Paragraph 6: 局限性（Limitations）客观陈述与未来扩展段（Constructive Limitations） (\texttt{BP23})}
\noindent\textbf{Expressive Intent:}\enspace 我这一段想写论文末尾或 NeurIPS/ICLR Checklist 要求的 Limitations 段落：诚实指出当前方法的两个客观边界，但绝不贬损论文的核心贡献。\par
\noindent\textbf{Cohesion Formula:}\enspace {\small\color{AcadSlate}Two Objective Boundaries $\to$ [Boundary 1 + Technical Remedy] $\to$ [Boundary 2 + Stochastic Extension]. 绑定逻辑规则：L017.}\par\vspace{4pt}
\begin{ExampleBox}{Publication-Grade Authentic Text}
Despite its strong empirical performance, our framework exhibits two theoretical boundaries worth noting. First, while our spectral branch scales quasi-linearly with sequence length, computing pairwise Hilbert phase differences scales quadratically with channel count ($O(C^2)$), which increases memory footprint when $C > 2,000$; incorporating anchor-based low-rank graph sparsification provides a direct remedy. Second, our continuous neural-ODE scale tracker assumes piece-wise smooth drift and may smooth over instantaneous step discontinuities caused by sensor recalibration. Extending the continuous drift prior to jump-diffusion stochastic differential equations offers a principled direction for future investigation.
\end{ExampleBox}
\begin{TipBox}{Rhetorical Move Breakdown \& Constituent Analysis}
\begin{itemize}[leftmargin=1.5em]
  \item \textbf{S1 [边界 1:} \textit{Despite its strong empirical performance, our framework exhibits two theoretical boundaries worth noting.}\par\noindent{\small 对应步骤：S1 [边界 1：极端高维或极度稀疏场景 Boundary 1]：Despite its effectiveness across standard benchmarks, our framework currently exhibits two practical boundaries. | 句首锚定旧信息/逻辑视角，主干动词紧跟主语，复杂修饰置于宾语之后（符合尾重原则 L011 与旧新信息流 L001）。}
  \item \textbf{S2 [边界 1 机理分析与解法 Mechanism \& Solution 1]:} \textit{First, while our spectral branch scales quasi-linearly with sequence length, computing pairwise Hilbert phase differences scales quadratically with channel count ($O(C^2)$), which increases memory footprint when $C > 2,000$; incorporating anchor-based low-rank graph sparsification provides a direct remedy.}\par\noindent{\small 对应步骤：S2 [边界 1 机理分析与解法 Mechanism \& Solution 1]：First, when the variate dimension scales beyond $C > 5,000$ unclustered sensors, constructing the cross-channel phase graph incurs noticeable memory overhead; integrating hierarchical sensor clustering offers a natural remedy. | 句首锚定旧信息/逻辑视角，主干动词紧跟主语，复杂修饰置于宾语之后（符合尾重原则 L011 与旧新信息流 L001）。}
  \item \textbf{S3 [边界 2:} \textit{Second, our continuous neural-ODE scale tracker assumes piece-wise smooth drift and may smooth over instantaneous step discontinuities caused by sensor recalibration.}\par\noindent{\small 对应步骤：S3 [边界 2：特定假设限制 Boundary 2]：Second, our continuous scale tracker assumes piece-wise differentiable drift and may under-estimate instantaneous discontinuous jumps caused by hardware resets. | 句首锚定旧信息/逻辑视角，主干动词紧跟主语，复杂修饰置于宾语之后（符合尾重原则 L011 与旧新信息流 L001）。}
  \item \textbf{S4 [未来演进方向 Future Extension]:} \textit{Extending the continuous drift prior to jump-diffusion stochastic differential equations offers a principled direction for future investigation.}\par\noindent{\small 对应步骤：S4 [未来演进方向 Future Extension]：Coupling our neural ODE tracker with jump-diffusion stochastic processes represents a promising avenue for future work. | 句首锚定旧信息/逻辑视角，主干动词紧跟主语，复杂修饰置于宾语之后（符合尾重原则 L011 与旧新信息流 L001）。}
\end{itemize}
\end{TipBox}
\medskip
\clearpage
\chapter{Economics, Finance \& Management / 经济金融与管理科学}\label{ch:anthology_20}
\Needspace{6\baselineskip}
本章精选该学科领域及对应文体的完整学术范文。每一篇范文不仅展示母语级发表标准的正文全貌，更清晰拆解全篇的“段落蓝图链条（Structural Blueprint Chain）”，并对每个自然段逐一剖析其核心交际意图、句际衔接逻辑公式（Cohesion Formula）以及每句话的句法成分与推进机制。

\section{FT11: 人工智能方向博士/访问学者申请个人陈述（SoP 全文：从时间序列与生成模型迈向具身智能世界模型）}\label{comp:FT11}
\begin{ConceptBox}{Composition Architecture \& Structural Blueprint Chain}
\textbf{Genre:} \texttt{Statement\_of\_Purpose} \quad \textbf{Paragraph Count:} 6\par
\textbf{Blueprint Chain:} \texttt{BP48 $\to$ BP03 $\to$ BP04 $\to$ BP15 $\to$ BP16 $\to$ BP24}
\end{ConceptBox}
\subsection*{Paragraph 1: 科研个人陈述（SoP / Research Narrative）：从时间序列/生成建模迈向具身智能的闭环叙事段 (\texttt{BP48})}
\noindent\textbf{Expressive Intent:}\enspace 我这一段想写申请博士/博后/访问学者或个人学术主页的 Research Statement / SoP 核心段：把我已经做过的深度学习/时间序列/生成模型工作，与我未来想探索的具身智能（Embodied AI）无缝串联成一条充满内驱力的科研主线。\par
\noindent\textbf{Cohesion Formula:}\enspace {\small\color{AcadSlate}Past Core Thread $\to$ Concrete Method \& Metric $\to$ Intellectual Bridge (Passive Forecasting $\to$ Active Closed-Loop Control) $\to$ Future Embodied AI Agenda. 绑定逻辑规则：L003, L025, L029.}\par\vspace{4pt}
\begin{ExampleBox}{Publication-Grade Authentic Text}
My research focuses on learning reliable representations and generative trajectories from non-stationary sequential data. While investigating multivariate industrial telemetry, I found that conventional attention mechanisms collapse when underlying system dynamics shift---motivating me to design a neural-ODE coupled spectral forecaster that reduced long-range drift by 17.4\%. Yet passive sequence forecasting only observes how a system evolves; intelligent physical agents must actively intervene within that loop, where every motor command alters future state distributions. Bridging my background in non-stationary time-series modeling and fast flow-matching samplers, my future research aims to build adaptive visuomotor world models that allow embodied robots to generalize across unfamiliar contact regimes in real time.
\end{ExampleBox}
\begin{TipBox}{Rhetorical Move Breakdown \& Constituent Analysis}
\begin{itemize}[leftmargin=1.5em]
  \item \textbf{S1 [过往研究主线凝练 Past Research Thread]:} \textit{My research focuses on learning reliable representations and generative trajectories from non-stationary sequential data.}\par\noindent{\small 对应步骤：S1 [过往研究主线凝练 Past Research Thread]：My prior research has centered on modeling non-stationary dynamics and uncertainty in high-dimensional sequential data. | 句首锚定旧信息/逻辑视角，主干动词紧跟主语，复杂修饰置于宾语之后（符合尾重原则 L011 与旧新信息流 L001）。}
  \item \textbf{S2 [代表作攻克过程与硬核成果 Concrete Achievement]:} \textit{While investigating multivariate industrial telemetry, I found that conventional attention mechanisms collapse when underlying system dynamics shift---motivating me to design a neural-ODE coupled spectral forecaster that reduced long-range drift by 17.4\%.}\par\noindent{\small 对应步骤：S2 [代表作攻克过程与硬核成果 Concrete Achievement]：In my recent work on [Time Series / Generative Modeling], I observed that [Core Bottleneck], which led me to formulate [Method]---reducing long-horizon drift by [X\%] across industrial telemetry benchmarks. | 句首锚定旧信息/逻辑视角，主干动词紧跟主语，复杂修饰置于宾语之后（符合尾重原则 L011 与旧新信息流 L001）。}
  \item \textbf{S3 [向具身智能跃迁的科学顿悟 Scientific Pivot to Embodied AI]:} \textit{Yet passive sequence forecasting only observes how a system evolves; intelligent physical agents must actively intervene within that loop, where every motor command alters future state distributions.}\par\noindent{\small 对应步骤：S3 [向具身智能跃迁的科学顿悟 Scientific Pivot to Embodied AI]：While passive sequence forecasting predicts how an environment evolves, physical decision-making requires closing the loop: actions actively alter future state distributions. | 句首锚定旧信息/逻辑视角，主干动词紧跟主语，复杂修饰置于宾语之后（符合尾重原则 L011 与旧新信息流 L001）。}
  \item \textbf{S4 [未来研究愿景 Future Research Agenda]:} \textit{Bridging my background in non-stationary time-series modeling and fast flow-matching samplers, my future research aims to build adaptive visuomotor world models that allow embodied robots to generalize across unfamiliar contact regimes in real time.}\par\noindent{\small 对应步骤：S4 [未来研究愿景 Future Research Agenda]：Building on my foundation in continuous-time generative modeling, my doctoral/future research aims to develop sample-efficient visuomotor world models that enable robots to adapt online to non-stationary contact dynamics. | 句首锚定旧信息/逻辑视角，主干动词紧跟主语，复杂修饰置于宾语之后（符合尾重原则 L011 与旧新信息流 L001）。}
\end{itemize}
\end{TipBox}
\medskip
\subsection*{Paragraph 2: 尖锐指出核心技术瓶颈与失效机理（Gap Identification） (\texttt{BP03})}
\noindent\textbf{Expressive Intent:}\enspace 我这一段想深入剖析现有方法在特定难题（如非平稳突变、长程自回归误差累积、Sim-to-Real 接触失配）上为什么会失效，把机理讲透。\par
\noindent\textbf{Cohesion Formula:}\enspace {\small\color{AcadSlate}因果链式推进（A induces B $\to$ This B compounds into C $\to$ Naive fix D sacrifices E）。绑定逻辑规则：L003, L004, L024/L025.}\par\vspace{4pt}
\begin{ExampleBox}{Publication-Grade Authentic Text}
Bridging non-stationary robustness and amplitude fidelity remains a core bottleneck in long-horizon forecasting. While instance normalization strips time-varying mean and variance shifts to stabilize attention weights, it simultaneously erases physical scale cues required to predict extreme load peaks. Conversely, retaining raw unnormalized inputs exposes self-attention layers to severe covariate shift, causing validation error to diverge over extended horizons. Resolving this tension requires a mechanism that disentangles intrinsic temporal dynamics from transient scale drift without discarding magnitude information.
\end{ExampleBox}
\begin{TipBox}{Rhetorical Move Breakdown \& Constituent Analysis}
\begin{itemize}[leftmargin=1.5em]
  \item \textbf{S1 [核心痛点立论 Bottleneck Claim]:} \textit{Bridging non-stationary robustness and amplitude fidelity remains a core bottleneck in long-horizon forecasting.}\par\noindent{\small 对应步骤：S1 [核心痛点立论 Bottleneck Claim]：一句话直接命名该核心瓶颈（如 Compound error accumulation / Sim-to-real contact mismatch）。 | 句首锚定旧信息/逻辑视角，主干动词紧跟主语，复杂修饰置于宾语之后（符合尾重原则 L011 与旧新信息流 L001）。}
  \item \textbf{S2 [第一层失效机理 Primary Failure Mechanism]:} \textit{While instance normalization strips time-varying mean and variance shifts to stabilize attention weights, it simultaneously erases physical scale cues required to predict extreme load peaks.}\par\noindent{\small 对应步骤：S2 [第一层失效机理 Primary Failure Mechanism]：从数学或物理传导链条解释误差是如何产生并放大的。 | 句首锚定旧信息/逻辑视角，主干动词紧跟主语，复杂修饰置于宾语之后（符合尾重原则 L011 与旧新信息流 L001）。}
  \item \textbf{S3 [第二层耦合困境 Secondary Dilemma]:} \textit{Conversely, retaining raw unnormalized inputs exposes self-attention layers to severe covariate shift, causing validation error to diverge over extended horizons.}\par\noindent{\small 对应步骤：S3 [第二层耦合困境 Secondary Dilemma]：指出为什么简单的修补手段（如单纯加大模型、简单归一化或域随机化）不仅无效还会带来副作用。 | 句首锚定旧信息/逻辑视角，主干动词紧跟主语，复杂修饰置于宾语之后（符合尾重原则 L011 与旧新信息流 L001）。}
  \item \textbf{S4 [凝练核心研究问题 Research Question Formulation]:} \textit{Resolving this tension requires a mechanism that disentangles intrinsic temporal dynamics from transient scale drift without discarding magnitude information.}\par\noindent{\small 对应步骤：S4 [凝练核心研究问题 Research Question Formulation]：用 Therefore, a fundamental question arises: how can we... 引出本文目标。 | 句首锚定旧信息/逻辑视角，主干动词紧跟主语，复杂修饰置于宾语之后（符合尾重原则 L011 与旧新信息流 L001）。}
\end{itemize}
\end{TipBox}
\medskip
\subsection*{Paragraph 3: 亮出本文核心洞察与破局方法总览（Key Insight \& Method Overview） (\texttt{BP04})}
\noindent\textbf{Expressive Intent:}\enspace 我这一段想在指出前人痛点后，正式亮出我们的方法（To address this dilemma, we propose...），并用最通俗清晰的直觉告诉审稿人为什么我们的设计能破局。\par
\noindent\textbf{Cohesion Formula:}\enspace {\small\color{AcadSlate}Problem Anchor $\to$ Key Insight $\to$ Dual-Module Mechanism $\to$ Pareto Resolution. 绑定逻辑规则：L002, L023.}\par\vspace{4pt}
\begin{ExampleBox}{Publication-Grade Authentic Text}
To resolve this dilemma, we propose DyTS-Net, a dual-stream forecasting architecture that decouples stationary structural modeling from dynamic scale trajectory estimation. Our key insight is that while high-frequency phase relationships benefit from normalized shift-invariant attention, macro-scale amplitude evolution follows a low-dimensional continuous dynamical system. Specifically, DyTS-Net first extracts scale-invariant temporal tokens via spectral patch attention, and concurrently models non-stationary mean-variance evolution through a lightweight neural ODE branch. Recombining both streams at the decoding stage preserves sharp peak magnitudes without sacrificing training stability.
\end{ExampleBox}
\begin{TipBox}{Rhetorical Move Breakdown \& Constituent Analysis}
\begin{itemize}[leftmargin=1.5em]
  \item \textbf{S1 [方法登场与定位 Method Introduction]:} \textit{To resolve this dilemma, we propose DyTS-Net, a dual-stream forecasting architecture that decouples stationary structural modeling from dynamic scale trajectory estimation.}\par\noindent{\small 对应步骤：S1 [方法登场与定位 Method Introduction]：To address this dilemma, we propose [Model Name], a [Core Paradigm] framework that [Primary Goal]. | 句首锚定旧信息/逻辑视角，主干动词紧跟主语，复杂修饰置于宾语之后（符合尾重原则 L011 与旧新信息流 L001）。}
  \item \textbf{S2 [核心直觉洞察 Key Insight]:} \textit{Our key insight is that while high-frequency phase relationships benefit from normalized shift-invariant attention, macro-scale amplitude evolution follows a low-dimensional continuous dynamical system.}\par\noindent{\small 对应步骤：S2 [核心直觉洞察 Key Insight]：Our key insight is that [Physical/Statistical Principle]，点破破局的底层支点。 | 句首锚定旧信息/逻辑视角，主干动词紧跟主语，复杂修饰置于宾语之后（符合尾重原则 L011 与旧新信息流 L001）。}
  \item \textbf{S3 [模块协同主干 Module Synergy]:} \textit{Specifically, DyTS-Net first extracts scale-invariant temporal tokens via spectral patch attention, and concurrently models non-stationary mean-variance evolution through a lightweight neural ODE branch.}\par\noindent{\small 对应步骤：S3 [模块协同主干 Module Synergy]：用 First / Specifically 与 Concurrently / Furthermore 简述两大核心模块如何分工咬合。 | 句首锚定旧信息/逻辑视角，主干动词紧跟主语，复杂修饰置于宾语之后（符合尾重原则 L011 与旧新信息流 L001）。}
  \item \textbf{S4 [闭环收益总结 Closed-Loop Benefit]:} \textit{Recombining both streams at the decoding stage preserves sharp peak magnitudes without sacrificing training stability.}\par\noindent{\small 对应步骤：S4 [闭环收益总结 Closed-Loop Benefit]：说明这种协同设计如何同时兼顾原先冲突的两个指标（如兼顾平稳表征与幅值保真，或兼顾多模态多样性与单步实时控制）。 | 句首锚定旧信息/逻辑视角，主干动词紧跟主语，复杂修饰置于宾语之后（符合尾重原则 L011 与旧新信息流 L001）。}
\end{itemize}
\end{TipBox}
\medskip
\subsection*{Paragraph 4: 生成模型专属：条件控制与采样去噪轨迹阐述段（Generative Modeling Mechanism） (\texttt{BP15})}
\noindent\textbf{Expressive Intent:}\enspace 我这一段想阐述生成模型（Diffusion / Flow Matching / Autoregressive）如何在满足外部物理约束或条件观测的前提下，生成高质量且多样化的序列/轨迹。\par
\noindent\textbf{Cohesion Formula:}\enspace {\small\color{AcadSlate}Guidance Artifact Problem $\to$ Intrinsic Velocity Conditioning $\to$ Tangent Projection Step $\to$ Fast Exact Convergence. 绑定逻辑规则：L026, L014.}\par\vspace{4pt}
\begin{ExampleBox}{Publication-Grade Authentic Text}
Applying post-hoc gradient guidance during diffusion sampling often pushes intermediate noisy states off the data manifold, introducing unphysical high-frequency oscillations. To enforce physical boundary constraints without manifold deviation, we incorporate observation consistency directly into the conditional flow-matching velocity field. At flow time $\tau \in [0, 1]$, our network predicts a residual velocity $v_\theta(x_\tau, \tau \mid c)$ that projects the transport trajectory onto the null space of the linear constraint operator. This projection guarantees exact boundary satisfaction at terminal time $\tau=1$ while maintaining straight transport paths amenable to 4-step Euler integration.
\end{ExampleBox}
\begin{TipBox}{Rhetorical Move Breakdown \& Constituent Analysis}
\begin{itemize}[leftmargin=1.5em]
  \item \textbf{S1 [无条件/朴素引导缺陷 Naive Guidance Bottleneck]:} \textit{Applying post-hoc gradient guidance during diffusion sampling often pushes intermediate noisy states off the data manifold, introducing unphysical high-frequency oscillations.}\par\noindent{\small 对应步骤：S1 [无条件/朴素引导缺陷 Naive Guidance Bottleneck]：Applying post-hoc classifier guidance to pre-trained diffusion samplers often pushes intermediate noisy states off the data manifold, inducing high-frequency artifacts. | 句首锚定旧信息/逻辑视角，主干动词紧跟主语，复杂修饰置于宾语之后（符合尾重原则 L011 与旧新信息流 L001）。}
  \item \textbf{S2 [本文流形内条件注入机制 Manifold-Consistent Conditioning]:} \textit{To enforce physical boundary constraints without manifold deviation, we incorporate observation consistency directly into the conditional flow-matching velocity field.}\par\noindent{\small 对应步骤：S2 [本文流形内条件注入机制 Manifold-Consistent Conditioning]：To enforce hard boundary constraints without manifold deviation, we embed [Condition/Constraint] directly into the velocity field / denoising kernel. | 句首锚定旧信息/逻辑视角，主干动词紧跟主语，复杂修饰置于宾语之后（符合尾重原则 L011 与旧新信息流 L001）。}
  \item \textbf{S3 [具体轨迹更新方程 Trajectory Update Formulation]:} \textit{At flow time $\tau \in [0, 1]$, our network predicts a residual velocity $v_\theta(x_\tau, \tau \mid c)$ that projects the transport trajectory onto the null space of the linear constraint operator.}\par\noindent{\small 对应步骤：S3 [具体轨迹更新方程 Trajectory Update Formulation]：At solver step $\tau \in [0, 1]$, the conditional velocity field $v_\theta(x_\tau, \tau \mid c)$ projects the unconstrained flow onto the tangent space of feasible trajectories. | 句首锚定旧信息/逻辑视角，主干动词紧跟主语，复杂修饰置于宾语之后（符合尾重原则 L011 与旧新信息流 L001）。}
  \item \textbf{S4 [采样效率与分布保真收益 Sampling \& Fidelity Guarantee]:} \textit{This projection guarantees exact boundary satisfaction at terminal time $\tau=1$ while maintaining straight transport paths amenable to 4-step Euler integration.}\par\noindent{\small 对应步骤：S4 [采样效率与分布保真收益 Sampling \& Fidelity Guarantee]：This tangent projection guarantees exact boundary satisfaction at $\tau=1$ while permitting straight 4-step Euler integration. | 句首锚定旧信息/逻辑视角，主干动词紧跟主语，复杂修饰置于宾语之后（符合尾重原则 L011 与旧新信息流 L001）。}
\end{itemize}
\end{TipBox}
\medskip
\subsection*{Paragraph 5: 具身智能专属：多模态感知-动作闭环与 Sim-to-Real 阐述段（Embodied AI Policy） (\texttt{BP16})}
\noindent\textbf{Expressive Intent:}\enspace 我这一段想阐述具身智能机器人策略（Embodied Policy）如何融合视觉、本体感觉（Proprioception）或触觉，并克服仿真到真实硬件的动力学差距（Sim-to-Real Gap）。\par
\noindent\textbf{Cohesion Formula:}\enspace {\small\color{AcadSlate}Visual Occlusion Bottleneck $\to$ Asynchronous Multimodal Alignment $\to$ Action Chunking $\to$ Online Sim-to-Real System Identification. 绑定逻辑规则：L025, L003.}\par\vspace{4pt}
\begin{ExampleBox}{Publication-Grade Authentic Text}
Relying solely on third-person RGB-D observations leaves visuomotor policies vulnerable to severe self-occlusion once the gripper approaches the target object. To preserve observability during contact-rich manipulation, our encoder aligns 30 Hz visual spatial tokens with 500 Hz joint proprioception and wrist force-torque histories via cross-modal rotary embeddings. Rather than issuing single-step commands prone to high-frequency jitter, the policy predicts 16-step SE(3) end-effector action chunks smoothed through exponential temporal ensembling. At deployment, a recurrent history adapter estimates latent surface friction and payload mass online, bridging the sim-to-real dynamics gap in zero shot.
\end{ExampleBox}
\begin{TipBox}{Rhetorical Move Breakdown \& Constituent Analysis}
\begin{itemize}[leftmargin=1.5em]
  \item \textbf{S1 [单模态开环局限 Open-Loop Visual Limitation]:} \textit{Relying solely on third-person RGB-D observations leaves visuomotor policies vulnerable to severe self-occlusion once the gripper approaches the target object.}\par\noindent{\small 对应步骤：S1 [单模态开环局限 Open-Loop Visual Limitation]：Relying solely on RGB-D point clouds leaves manipulation policies vulnerable to self-occlusion by the robot gripper during fine-grained contact. | 句首锚定旧信息/逻辑视角，主干动词紧跟主语，复杂修饰置于宾语之后（符合尾重原则 L011 与旧新信息流 L001）。}
  \item \textbf{S2 [多模态时空对齐 Multi-Modal Alignment]:} \textit{To preserve observability during contact-rich manipulation, our encoder aligns 30 Hz visual spatial tokens with 500 Hz joint proprioception and wrist force-torque histories via cross-modal rotary embeddings.}\par\noindent{\small 对应步骤：S2 [多模态时空对齐 Multi-Modal Alignment]：To maintain state observability during contact, our encoder aligns asynchronous 30 Hz visual tokens with 500 Hz joint proprioception and wrist force-torque histories. | 句首锚定旧信息/逻辑视角，主干动词紧跟主语，复杂修饰置于宾语之后（符合尾重原则 L011 与旧新信息流 L001）。}
  \item \textbf{S3 [动作块预测与平滑执行 Action Chunking \& Temporal Ensembling]:} \textit{Rather than issuing single-step commands prone to high-frequency jitter, the policy predicts 16-step SE(3) end-effector action chunks smoothed through exponential temporal ensembling.}\par\noindent{\small 对应步骤：S3 [动作块预测与平滑执行 Action Chunking \& Temporal Ensembling]：Rather than predicting single-step torques prone to jitter, the policy outputs overlapping $K$-step SE(3) end-effector action chunks smoothed via exponential temporal ensembling. | 句首锚定旧信息/逻辑视角，主干动词紧跟主语，复杂修饰置于宾语之后（符合尾重原则 L011 与旧新信息流 L001）。}
  \item \textbf{S4 [Sim-to-Real 自适应机制 Sim-to-Real Adaptation]:} \textit{At deployment, a recurrent history adapter estimates latent surface friction and payload mass online, bridging the sim-to-real dynamics gap in zero shot.}\par\noindent{\small 对应步骤：S4 [Sim-to-Real 自适应机制 Sim-to-Real Adaptation]：During physical deployment, a lightweight history adapter infers latent friction and payload parameters online, bridging the sim-to-real dynamics gap without policy fine-tuning. | 句首锚定旧信息/逻辑视角，主干动词紧跟主语，复杂修饰置于宾语之后（符合尾重原则 L011 与旧新信息流 L001）。}
\end{itemize}
\end{TipBox}
\medskip
\subsection*{Paragraph 6: 全文总结与学术影响收尾段（Conclusion \& Broader Outlook） (\texttt{BP24})}
\noindent\textbf{Expressive Intent:}\enspace 我这一段想写整篇小论文的 Conclusion：不简单复制粘贴 Abstract，而是从更高维度总结本文解决的科学矛盾、方法精髓、实证高度与领域启示。\par
\noindent\textbf{Cohesion Formula:}\enspace {\small\color{AcadSlate}Tension Resolved $\to$ Methodological Insight $\to$ Empirical Confirmation $\to$ Generalizable Blueprint. 绑定逻辑规则：L013, L023, L029.}\par\vspace{4pt}
\begin{ExampleBox}{Publication-Grade Authentic Text}
In this work, we resolved the long-standing tension between shift-invariant normalization and physical amplitude preservation in non-stationary time series forecasting. By decoupling frequency-domain phase attention from continuous neural-ODE scale tracking, DyTS-Net captures long-range seasonal dependencies while accurately reconstructing extreme peak magnitudes. Empirical evaluations across eight benchmarks and industrial wind-power telemetry demonstrated MSE reductions of up to 17.4\% alongside a 45\% reduction in computational overhead. Beyond forecasting, this structural-scale disentanglement principle provides a generalizable foundation for anomaly detection and closed-loop predictive control in non-stationary cyber-physical systems.
\end{ExampleBox}
\begin{TipBox}{Rhetorical Move Breakdown \& Constituent Analysis}
\begin{itemize}[leftmargin=1.5em]
  \item \textbf{S1 [核心科学问题闭合 Scientific Closure]:} \textit{In this work, we resolved the long-standing tension between shift-invariant normalization and physical amplitude preservation in non-stationary time series forecasting.}\par\noindent{\small 对应步骤：S1 [核心科学问题闭合 Scientific Closure]：In this paper, we addressed the fundamental tension between [Goal A] and [Goal B] in [Domain]. | 句首锚定旧信息/逻辑视角，主干动词紧跟主语，复杂修饰置于宾语之后（符合尾重原则 L011 与旧新信息流 L001）。}
  \item \textbf{S2 [方法论精髓回顾 Methodological Synthesis]:} \textit{By decoupling frequency-domain phase attention from continuous neural-ODE scale tracking, DyTS-Net captures long-range seasonal dependencies while accurately reconstructing extreme peak magnitudes.}\par\noindent{\small 对应步骤：S2 [方法论精髓回顾 Methodological Synthesis]：By introducing [Model/Mechanism], we demonstrated that [Core Insight] enables [Capability] without incurring [Traditional Penalty]. | 句首锚定旧信息/逻辑视角，主干动词紧跟主语，复杂修饰置于宾语之后（符合尾重原则 L011 与旧新信息流 L001）。}
  \item \textbf{S3 [实证结果高度概括 Empirical Recap]:} \textit{Empirical evaluations across eight benchmarks and industrial wind-power telemetry demonstrated MSE reductions of up to 17.4\% alongside a 45\% reduction in computational overhead.}\par\noindent{\small 对应步骤：S3 [实证结果高度概括 Empirical Recap]：Comprehensive evaluations across [Benchmarks/Hardware] confirmed consistent error reductions of up to [X\%] alongside substantial inference speedups. | 句首锚定旧信息/逻辑视角，主干动词紧跟主语，复杂修饰置于宾语之后（符合尾重原则 L011 与旧新信息流 L001）。}
  \item \textbf{S4 [跨领域迁移展望 Broader Outlook]:} \textit{Beyond forecasting, this structural-scale disentanglement principle provides a generalizable foundation for anomaly detection and closed-loop predictive control in non-stationary cyber-physical systems.}\par\noindent{\small 对应步骤：S4 [跨领域迁移展望 Broader Outlook]：Beyond [Current Task], our [Core Principle] offers a generalizable blueprint for broader sequential decision-making tasks under non-stationary environments. | 句首锚定旧信息/逻辑视角，主干动词紧跟主语，复杂修饰置于宾语之后（符合尾重原则 L011 与旧新信息流 L001）。}
\end{itemize}
\end{TipBox}
\medskip
\clearpage
\section{FT12: 具身智能与非平稳时序建模交叉课题立项申请书（Research Proposal 全文：研究背景、核心科学问题、技术路线与预期成果）}\label{comp:FT12}
\begin{ConceptBox}{Composition Architecture \& Structural Blueprint Chain}
\textbf{Genre:} \texttt{Research\_Proposal} \quad \textbf{Paragraph Count:} 7\par
\textbf{Blueprint Chain:} \texttt{BP01 $\to$ BP02 $\to$ BP03 $\to$ BP04 $\to$ BP05 $\to$ BP06 $\to$ BP07}
\end{ConceptBox}
\subsection*{Paragraph 1: 确立研究背景与真实工程价值（拒绝空喊口号） (\texttt{BP01})}
\noindent\textbf{Expressive Intent:}\enspace 我这一段想作为 Introduction 第一段：说明多变量时间序列预测或具身智能控制在实际工业系统中的关键作用，以及高可靠性的刚需。\par
\noindent\textbf{Cohesion Formula:}\enspace {\small\color{AcadSlate}Given-to-New 链：Core Task (S1) $\to$ Practical Deployment (S2) $\to$ Real-world Disturbances (S3) $\to$ Consequence of Failure (S4). 绑定逻辑规则：L001, L011, L029.}\par\vspace{4pt}
\begin{ExampleBox}{Publication-Grade Authentic Text}
Multivariate time series forecasting underpins safety-critical scheduling across smart power grids, cloud clusters, and automated manufacturing. In these operational settings, control systems rely on multi-horizon predictions to balance supply fluctuations and prevent cascading overload. However, field sensors routinely operate under abrupt regime shifts, irregular sampling intervals, and high-frequency measurement noise. Failing to accommodate such non-stationary disturbances degrades long-range forecast reliability and triggers costly control oscillations.
\end{ExampleBox}
\begin{TipBox}{Rhetorical Move Breakdown \& Constituent Analysis}
\begin{itemize}[leftmargin=1.5em]
  \item \textbf{S1 [任务锚定 Topic Anchor]:} \textit{Multivariate time series forecasting underpins safety-critical scheduling across smart power grids, cloud clusters, and automated manufacturing.}\par\noindent{\small 对应步骤：S1 [任务锚定 Topic Anchor]：直接用紧凑主谓宾点明研究对象及其在真实系统中的核心功能（≤16词，禁用 In the era of AI）。 | 句首锚定旧信息/逻辑视角，主干动词紧跟主语，复杂修饰置于宾语之后（符合尾重原则 L011 与旧新信息流 L001）。}
  \item \textbf{S2 [应用场景展开 Concrete Scenarios]:} \textit{In these operational settings, control systems rely on multi-horizon predictions to balance supply fluctuations and prevent cascading overload.}\par\noindent{\small 对应步骤：S2 [应用场景展开 Concrete Scenarios]：列举 2--3 个具体工程场景（如电网负荷调度、灵巧机械手装配）及其决策依赖。 | 句首锚定旧信息/逻辑视角，主干动词紧跟主语，复杂修饰置于宾语之后（符合尾重原则 L011 与旧新信息流 L001）。}
  \item \textbf{S3 [苛刻工况约束 Operational Bottleneck]:} \textit{However, field sensors routinely operate under abrupt regime shifts, irregular sampling intervals, and high-frequency measurement noise.}\par\noindent{\small 对应步骤：S3 [苛刻工况约束 Operational Bottleneck]：指出真实部署环境下不可避免的物理/统计挑战（如非平稳漂移、传感器缺失、毫秒级延迟约束）。 | 句首锚定旧信息/逻辑视角，主干动词紧跟主语，复杂修饰置于宾语之后（符合尾重原则 L011 与旧新信息流 L001）。}
  \item \textbf{S4 [研究必要性收束 Imperative Hook]:} \textit{Failing to accommodate such non-stationary disturbances degrades long-range forecast reliability and triggers costly control oscillations.}\par\noindent{\small 对应步骤：S4 [研究必要性收束 Imperative Hook]：点明若无法克服上述约束会造成什么具体系统后果，自然引出下一段对现有方法的审视。 | 句首锚定旧信息/逻辑视角，主干动词紧跟主语，复杂修饰置于宾语之后（符合尾重原则 L011 与旧新信息流 L001）。}
\end{itemize}
\end{TipBox}
\medskip
\subsection*{Paragraph 2: 总结现有主流范式及其隐含假设（欲抑先扬铺垫段） (\texttt{BP02})}
\noindent\textbf{Expressive Intent:}\enspace 我这一段想客观概括前人主流模型（如 Transformer / Diffusion / MLP）取得了什么进展，但点出它们共同依赖的一个脆弱前提假设。\par
\noindent\textbf{Cohesion Formula:}\enspace {\small\color{AcadSlate}先让步肯定（While/Recent paradigms...）$\to$ 后釜底抽薪揭示隐含前提（Yet these architectures implicitly assume...）。绑定逻辑规则：L008, L009.}\par\vspace{4pt}
\begin{ExampleBox}{Publication-Grade Authentic Text}
Recent deep forecasters have substantially reduced long-horizon error by replacing point-wise attention with sub-series patch tokenization and frequency-domain decomposition. These architectures effectively capture recurring seasonal patterns when training and test sequences share stationary statistical moments. Yet they implicitly assume that inter-channel correlations and spectral energy distributions remain time-invariant across the prediction window. When external shocks alter underlying system dynamics, this stationarity assumption collapses, causing frozen spectral filters to project outdated trends.
\end{ExampleBox}
\begin{TipBox}{Rhetorical Move Breakdown \& Constituent Analysis}
\begin{itemize}[leftmargin=1.5em]
  \item \textbf{S1 [流派归类 Taxonomic Opening]:} \textit{Recent deep forecasters have substantially reduced long-horizon error by replacing point-wise attention with sub-series patch tokenization and frequency-domain decomposition.}\par\noindent{\small 对应步骤：S1 [流派归类 Taxonomic Opening]：按核心机制归纳现有主流范式及其已解决的子问题。 | 句首锚定旧信息/逻辑视角，主干动词紧跟主语，复杂修饰置于宾语之后（符合尾重原则 L011 与旧新信息流 L001）。}
  \item \textbf{S2 [代表性机制展开 Mechanism Recap]:} \textit{These architectures effectively capture recurring seasonal patterns when training and test sequences share stationary statistical moments.}\par\noindent{\small 对应步骤：S2 [代表性机制展开 Mechanism Recap]：简述这些方法如何通过特定架构（如 Patch 分块、频域分解或扩散去噪）提升表征能力。 | 句首锚定旧信息/逻辑视角，主干动词紧跟主语，复杂修饰置于宾语之后（符合尾重原则 L011 与旧新信息流 L001）。}
  \item \textbf{S3 [转折揭示隐含假设 Pivot to Hidden Assumption]:} \textit{Yet they implicitly assume that inter-channel correlations and spectral energy distributions remain time-invariant across the prediction window.}\par\noindent{\small 对应步骤：S3 [转折揭示隐含假设 Pivot to Hidden Assumption]：用 Yet / However 锋利转折，指出它们暗中依赖的理想化假设（如平稳分布、静态通道相关性、仿真动力学无偏差）。 | 句首锚定旧信息/逻辑视角，主干动词紧跟主语，复杂修饰置于宾语之后（符合尾重原则 L011 与旧新信息流 L001）。}
  \item \textbf{S4 [假设破裂后果 Breakdown Under Stress]:} \textit{When external shocks alter underlying system dynamics, this stationarity assumption collapses, causing frozen spectral filters to project outdated trends.}\par\noindent{\small 对应步骤：S4 [假设破裂后果 Breakdown Under Stress]：说明当真实数据违背该假设时，现有架构为何会遭遇性能断崖。 | 句首锚定旧信息/逻辑视角，主干动词紧跟主语，复杂修饰置于宾语之后（符合尾重原则 L011 与旧新信息流 L001）。}
\end{itemize}
\end{TipBox}
\medskip
\subsection*{Paragraph 3: 尖锐指出核心技术瓶颈与失效机理（Gap Identification） (\texttt{BP03})}
\noindent\textbf{Expressive Intent:}\enspace 我这一段想深入剖析现有方法在特定难题（如非平稳突变、长程自回归误差累积、Sim-to-Real 接触失配）上为什么会失效，把机理讲透。\par
\noindent\textbf{Cohesion Formula:}\enspace {\small\color{AcadSlate}因果链式推进（A induces B $\to$ This B compounds into C $\to$ Naive fix D sacrifices E）。绑定逻辑规则：L003, L004, L024/L025.}\par\vspace{4pt}
\begin{ExampleBox}{Publication-Grade Authentic Text}
Bridging non-stationary robustness and amplitude fidelity remains a core bottleneck in long-horizon forecasting. While instance normalization strips time-varying mean and variance shifts to stabilize attention weights, it simultaneously erases physical scale cues required to predict extreme load peaks. Conversely, retaining raw unnormalized inputs exposes self-attention layers to severe covariate shift, causing validation error to diverge over extended horizons. Resolving this tension requires a mechanism that disentangles intrinsic temporal dynamics from transient scale drift without discarding magnitude information.
\end{ExampleBox}
\begin{TipBox}{Rhetorical Move Breakdown \& Constituent Analysis}
\begin{itemize}[leftmargin=1.5em]
  \item \textbf{S1 [核心痛点立论 Bottleneck Claim]:} \textit{Bridging non-stationary robustness and amplitude fidelity remains a core bottleneck in long-horizon forecasting.}\par\noindent{\small 对应步骤：S1 [核心痛点立论 Bottleneck Claim]：一句话直接命名该核心瓶颈（如 Compound error accumulation / Sim-to-real contact mismatch）。 | 句首锚定旧信息/逻辑视角，主干动词紧跟主语，复杂修饰置于宾语之后（符合尾重原则 L011 与旧新信息流 L001）。}
  \item \textbf{S2 [第一层失效机理 Primary Failure Mechanism]:} \textit{While instance normalization strips time-varying mean and variance shifts to stabilize attention weights, it simultaneously erases physical scale cues required to predict extreme load peaks.}\par\noindent{\small 对应步骤：S2 [第一层失效机理 Primary Failure Mechanism]：从数学或物理传导链条解释误差是如何产生并放大的。 | 句首锚定旧信息/逻辑视角，主干动词紧跟主语，复杂修饰置于宾语之后（符合尾重原则 L011 与旧新信息流 L001）。}
  \item \textbf{S3 [第二层耦合困境 Secondary Dilemma]:} \textit{Conversely, retaining raw unnormalized inputs exposes self-attention layers to severe covariate shift, causing validation error to diverge over extended horizons.}\par\noindent{\small 对应步骤：S3 [第二层耦合困境 Secondary Dilemma]：指出为什么简单的修补手段（如单纯加大模型、简单归一化或域随机化）不仅无效还会带来副作用。 | 句首锚定旧信息/逻辑视角，主干动词紧跟主语，复杂修饰置于宾语之后（符合尾重原则 L011 与旧新信息流 L001）。}
  \item \textbf{S4 [凝练核心研究问题 Research Question Formulation]:} \textit{Resolving this tension requires a mechanism that disentangles intrinsic temporal dynamics from transient scale drift without discarding magnitude information.}\par\noindent{\small 对应步骤：S4 [凝练核心研究问题 Research Question Formulation]：用 Therefore, a fundamental question arises: how can we... 引出本文目标。 | 句首锚定旧信息/逻辑视角，主干动词紧跟主语，复杂修饰置于宾语之后（符合尾重原则 L011 与旧新信息流 L001）。}
\end{itemize}
\end{TipBox}
\medskip
\subsection*{Paragraph 4: 亮出本文核心洞察与破局方法总览（Key Insight \& Method Overview） (\texttt{BP04})}
\noindent\textbf{Expressive Intent:}\enspace 我这一段想在指出前人痛点后，正式亮出我们的方法（To address this dilemma, we propose...），并用最通俗清晰的直觉告诉审稿人为什么我们的设计能破局。\par
\noindent\textbf{Cohesion Formula:}\enspace {\small\color{AcadSlate}Problem Anchor $\to$ Key Insight $\to$ Dual-Module Mechanism $\to$ Pareto Resolution. 绑定逻辑规则：L002, L023.}\par\vspace{4pt}
\begin{ExampleBox}{Publication-Grade Authentic Text}
To resolve this dilemma, we propose DyTS-Net, a dual-stream forecasting architecture that decouples stationary structural modeling from dynamic scale trajectory estimation. Our key insight is that while high-frequency phase relationships benefit from normalized shift-invariant attention, macro-scale amplitude evolution follows a low-dimensional continuous dynamical system. Specifically, DyTS-Net first extracts scale-invariant temporal tokens via spectral patch attention, and concurrently models non-stationary mean-variance evolution through a lightweight neural ODE branch. Recombining both streams at the decoding stage preserves sharp peak magnitudes without sacrificing training stability.
\end{ExampleBox}
\begin{TipBox}{Rhetorical Move Breakdown \& Constituent Analysis}
\begin{itemize}[leftmargin=1.5em]
  \item \textbf{S1 [方法登场与定位 Method Introduction]:} \textit{To resolve this dilemma, we propose DyTS-Net, a dual-stream forecasting architecture that decouples stationary structural modeling from dynamic scale trajectory estimation.}\par\noindent{\small 对应步骤：S1 [方法登场与定位 Method Introduction]：To address this dilemma, we propose [Model Name], a [Core Paradigm] framework that [Primary Goal]. | 句首锚定旧信息/逻辑视角，主干动词紧跟主语，复杂修饰置于宾语之后（符合尾重原则 L011 与旧新信息流 L001）。}
  \item \textbf{S2 [核心直觉洞察 Key Insight]:} \textit{Our key insight is that while high-frequency phase relationships benefit from normalized shift-invariant attention, macro-scale amplitude evolution follows a low-dimensional continuous dynamical system.}\par\noindent{\small 对应步骤：S2 [核心直觉洞察 Key Insight]：Our key insight is that [Physical/Statistical Principle]，点破破局的底层支点。 | 句首锚定旧信息/逻辑视角，主干动词紧跟主语，复杂修饰置于宾语之后（符合尾重原则 L011 与旧新信息流 L001）。}
  \item \textbf{S3 [模块协同主干 Module Synergy]:} \textit{Specifically, DyTS-Net first extracts scale-invariant temporal tokens via spectral patch attention, and concurrently models non-stationary mean-variance evolution through a lightweight neural ODE branch.}\par\noindent{\small 对应步骤：S3 [模块协同主干 Module Synergy]：用 First / Specifically 与 Concurrently / Furthermore 简述两大核心模块如何分工咬合。 | 句首锚定旧信息/逻辑视角，主干动词紧跟主语，复杂修饰置于宾语之后（符合尾重原则 L011 与旧新信息流 L001）。}
  \item \textbf{S4 [闭环收益总结 Closed-Loop Benefit]:} \textit{Recombining both streams at the decoding stage preserves sharp peak magnitudes without sacrificing training stability.}\par\noindent{\small 对应步骤：S4 [闭环收益总结 Closed-Loop Benefit]：说明这种协同设计如何同时兼顾原先冲突的两个指标（如兼顾平稳表征与幅值保真，或兼顾多模态多样性与单步实时控制）。 | 句首锚定旧信息/逻辑视角，主干动词紧跟主语，复杂修饰置于宾语之后（符合尾重原则 L011 与旧新信息流 L001）。}
\end{itemize}
\end{TipBox}
\medskip
\subsection*{Paragraph 5: 挑战与机制一对一映射段（Challenge-to-Mechanism Mapping） (\texttt{BP05})}
\noindent\textbf{Expressive Intent:}\enspace 我这一段想说明：把我们的核心思想落地时遇到了两个具体工程/数学难点，我们分别设计了什么精巧机制去攻克它们。\par
\noindent\textbf{Cohesion Formula:}\enspace {\small\color{AcadSlate}Challenge 1 $\to$ Mechanism 1; Challenge 2 $\to$ Mechanism 2 $\to$ Joint System Property. 绑定逻辑规则：L003, L011.}\par\vspace{4pt}
\begin{ExampleBox}{Publication-Grade Authentic Text}
Instantiating closed-loop visuomotor diffusion on physical robots entails two technical challenges. First, because raw depth camera streams contain occlusion artifacts around reflective object edges, we introduce a tactile-guided cross-attention gate that reweights visual tokens upon contact onset. Second, to eliminate the 200 ms sampling bottleneck of iterative denoising, we formulate a trajectory consistency distillation objective that maps Gaussian noise directly to collision-free action chunks in two forward steps. Together, these designs sustain 50 Hz reactive control under partial visual occlusion.
\end{ExampleBox}
\begin{TipBox}{Rhetorical Move Breakdown \& Constituent Analysis}
\begin{itemize}[leftmargin=1.5em]
  \item \textbf{S1 [落地难点总起 Implementation Hurdle]:} \textit{Instantiating closed-loop visuomotor diffusion on physical robots entails two technical challenges.}\par\noindent{\small 对应步骤：S1 [落地难点总起 Implementation Hurdle]：Translating this principle into a scalable architecture entails two technical challenges. | 句首锚定旧信息/逻辑视角，主干动词紧跟主语，复杂修饰置于宾语之后（符合尾重原则 L011 与旧新信息流 L001）。}
  \item \textbf{S2 [难点 1 + 靶向设计 1]:} \textit{First, because raw depth camera streams contain occlusion artifacts around reflective object edges, we introduce a tactile-guided cross-attention gate that reweights visual tokens upon contact onset.}\par\noindent{\small 对应步骤：S2 [难点 1 + 靶向设计 1]：First, because [Challenge 1], we design [Mechanism 1] to [Specific Function]. | 句首锚定旧信息/逻辑视角，主干动词紧跟主语，复杂修饰置于宾语之后（符合尾重原则 L011 与旧新信息流 L001）。}
  \item \textbf{S3 [难点 2 + 靶向设计 2]:} \textit{Second, to eliminate the 200 ms sampling bottleneck of iterative denoising, we formulate a trajectory consistency distillation objective that maps Gaussian noise directly to collision-free action chunks in two forward steps.}\par\noindent{\small 对应步骤：S3 [难点 2 + 靶向设计 2]：Second, to prevent [Challenge 2] under [Constraint], we introduce [Mechanism 2] that [Specific Function]. | 句首锚定旧信息/逻辑视角，主干动词紧跟主语，复杂修饰置于宾语之后（符合尾重原则 L011 与旧新信息流 L001）。}
  \item \textbf{S4 [整体计算/泛化特性 System Property]:} \textit{Together, these designs sustain 50 Hz reactive control under partial visual occlusion.}\par\noindent{\small 对应步骤：S4 [整体计算/泛化特性 System Property]：Together, these components enable [Model] to operate at [Efficiency/Robustness Metric]. | 句首锚定旧信息/逻辑视角，主干动词紧跟主语，复杂修饰置于宾语之后（符合尾重原则 L011 与旧新信息流 L001）。}
\end{itemize}
\end{TipBox}
\medskip
\subsection*{Paragraph 6: 实验证言与性能突破预览段（Quantitative Empirical Preview） (\texttt{BP06})}
\noindent\textbf{Expressive Intent:}\enspace 我这一段想在 Introduction 结尾前，用三句话硬核总结我们在哪些基准数据集/真实硬件上做了实验，取得了多大幅度的领先。\par
\noindent\textbf{Cohesion Formula:}\enspace {\small\color{AcadSlate}Evaluation Scope $\to$ Accuracy Gain over Strongest Baselines $\to$ Efficiency/Generalization Gain. 绑定逻辑规则：L010, L015.}\par\vspace{4pt}
\begin{ExampleBox}{Publication-Grade Authentic Text}
We evaluate DyTS-Net across eight long-horizon forecasting benchmarks and a real-world wind-turbine telemetry system spanning 1,400 sensors. Across all prediction horizons up to 720 steps, DyTS-Net reduces MSE by 12.8\% over PatchTST and 17.4\% over iTransformer, with the largest gains emerging during sudden weather-front transitions. Furthermore, our architecture requires 45\% fewer FLOPs than full channel-mixing Transformers and transfers zero-shot to newly deployed turbine units without fine-tuning.
\end{ExampleBox}
\begin{TipBox}{Rhetorical Move Breakdown \& Constituent Analysis}
\begin{itemize}[leftmargin=1.5em]
  \item \textbf{S1 [评测广度 Scope of Evaluation]:} \textit{We evaluate DyTS-Net across eight long-horizon forecasting benchmarks and a real-world wind-turbine telemetry system spanning 1,400 sensors.}\par\noindent{\small 对应步骤：S1 [评测广度 Scope of Evaluation]：We evaluate [Model] across [N] public benchmarks and [Real-world Hardware/Scenario], spanning [Tasks]. | 句首锚定旧信息/逻辑视角，主干动词紧跟主语，复杂修饰置于宾语之后（符合尾重原则 L011 与旧新信息流 L001）。}
  \item \textbf{S2 [核心精度突破 Primary Accuracy Gain]:} \textit{Across all prediction horizons up to 720 steps, DyTS-Net reduces MSE by 12.8\% over PatchTST and 17.4\% over iTransformer, with the largest gains emerging during sudden weather-front transitions.}\par\noindent{\small 对应步骤：S2 [核心精度突破 Primary Accuracy Gain]：实验主指标对比（必须包含具体百分比、最强基线名称及困难场景表现）。 | 句首锚定旧信息/逻辑视角，主干动词紧跟主语，复杂修饰置于宾语之后（符合尾重原则 L011 与旧新信息流 L001）。}
  \item \textbf{S3 [效率或泛化突破 Efficiency / Zero-Shot Gain]:} \textit{Furthermore, our architecture requires 45\% fewer FLOPs than full channel-mixing Transformers and transfers zero-shot to newly deployed turbine units without fine-tuning.}\par\noindent{\small 对应步骤：S3 [效率或泛化突破 Efficiency / Zero-Shot Gain]：报告计算速度、显存占用或零样本跨场景泛化收益。 | 句首锚定旧信息/逻辑视角，主干动词紧跟主语，复杂修饰置于宾语之后（符合尾重原则 L011 与旧新信息流 L001）。}
\end{itemize}
\end{TipBox}
\medskip
\subsection*{Paragraph 7: 凝练论文三点核心贡献（Contributions Summary） (\texttt{BP07})}
\noindent\textbf{Expressive Intent:}\enspace 我这一段想在 Introduction 最末尾清晰列出本文的三点核心贡献（理论/视角贡献 + 核心算法/架构贡献 + 实验与开源实证贡献）。\par
\noindent\textbf{Cohesion Formula:}\enspace {\small\color{AcadSlate}Parallel Grammatical Structure（严格平行结构：全部用 We + 强及物动词 identify / propose / demonstrate 开头）。绑定逻辑规则：L002, L010.}\par\vspace{4pt}
\begin{ExampleBox}{Publication-Grade Authentic Text}
In summary, our main contributions are threefold: (1) We formalize the amplitude-stationarity trade-off in long-horizon forecasting and prove that standard instance normalization bounds peak-load recoverability. (2) We propose DyTS-Net, a hybrid architecture that couples shift-invariant spectral attention with a continuous neural ODE scale tracker. (3) We demonstrate across eight public benchmarks and industrial wind-power telemetry that DyTS-Net lowers long-range MSE by up to 17.4\% while cutting inference latency nearly in half.
\end{ExampleBox}
\begin{TipBox}{Rhetorical Move Breakdown \& Constituent Analysis}
\begin{itemize}[leftmargin=1.5em]
  \item \textbf{S1 [总起引导句 Lead-in]:} \textit{In summary, our main contributions are threefold: (1) We formalize the amplitude-stationarity trade-off in long-horizon forecasting and prove that standard instance normalization bounds peak-load recoverability.}\par\noindent{\small 对应步骤：S1 [总起引导句 Lead-in]：In summary, our main contributions are threefold: | 句首锚定旧信息/逻辑视角，主干动词紧跟主语，复杂修饰置于宾语之后（符合尾重原则 L011 与旧新信息流 L001）。}
  \item \textbf{S2 [贡献 1:} \textit{(2) We propose DyTS-Net, a hybrid architecture that couples shift-invariant spectral attention with a continuous neural ODE scale tracker.}\par\noindent{\small 对应步骤：S2 [贡献 1：新视角/问题形式化 Conceptual/Formulation Contribution]：用强动词名物化或动词原形开头（We identify and formalize... / We establish...），点明打破了什么传统局限。 | 句首锚定旧信息/逻辑视角，主干动词紧跟主语，复杂修饰置于宾语之后（符合尾重原则 L011 与旧新信息流 L001）。}
  \item \textbf{S3 [贡献 2:} \textit{(3) We demonstrate across eight public benchmarks and industrial wind-power telemetry that DyTS-Net lowers long-range MSE by up to 17.4\% while cutting inference latency nearly in half.}\par\noindent{\small 对应步骤：S3 [贡献 2：核心方法与模块 Algorithmic/Architectural Contribution]：We propose [Model/Mechanism], which integrates [A] and [B] to achieve [Theoretical/Functional Guarantee]. | 句首锚定旧信息/逻辑视角，主干动词紧跟主语，复杂修饰置于宾语之后（符合尾重原则 L011 与旧新信息流 L001）。}
\end{itemize}
\end{TipBox}
\medskip
\clearpage
\section{FT13: 雅思小作文（Task 1）动态多折线图满分范文：四国清洁能源与煤炭发电占比 25 年演进及主次反转}\label{comp:FT13}
\begin{ConceptBox}{Composition Architecture \& Structural Blueprint Chain}
\textbf{Genre:} \texttt{IELTS\_Task1\_Dynamic} \quad \textbf{Paragraph Count:} 3\par
\textbf{Blueprint Chain:} \texttt{BP25 $\to$ BP26 $\to$ BP32}
\end{ConceptBox}
\subsection*{Paragraph 1: 图表小作文首段：题目改写与总体趋势概览（Paraphrase + Macro Overview） (\texttt{BP25})}
\noindent\textbf{Expressive Intent:}\enspace 我这一段想写雅思 Task 1 或数据报告的第一段：用同义句法改写图表标题，并用两句话高度凝练全图最核心的 2 个宏观趋势（极值、反超拐点或主次格局），绝不堆砌零碎数字。\par
\noindent\textbf{Cohesion Formula:}\enspace {\small\color{AcadSlate}Paraphrase $\to$ Overall Trajectory Contrast $\to$ Dominant Crossover/Extremum (Zero raw numbers in Overview!). 绑定逻辑规则：L013, L015.}\par\vspace{4pt}
\begin{ExampleBox}{Publication-Grade Authentic Text}
The line graph compares the proportion of electricity generated from four fuel sources across three industrial economies between 2000 and 2025. Overall, renewable sources---particularly solar and wind power---followed a pronounced upward trajectory over the 25-year span, whereas coal reliance contracted steadily. Most notably, clean energy generation overtook fossil-fuel consumption around 2018, shifting from a marginal supplement to the dominant power supply by the end of the period.
\end{ExampleBox}
\begin{TipBox}{Rhetorical Move Breakdown \& Constituent Analysis}
\begin{itemize}[leftmargin=1.5em]
  \item \textbf{S1 [图表命题同义改写 Prompt Paraphrase]:} \textit{The line graph compares the proportion of electricity generated from four fuel sources across three industrial economies between 2000 and 2025.}\par\noindent{\small 对应步骤：S1 [图表命题同义改写 Prompt Paraphrase]：The line graph / bar chart compares [Variable/Metric] across [Categories/Regions] over a [N]-year period from [Year A] to [Year B]. | 句首锚定旧信息/逻辑视角，主干动词紧跟主语，复杂修饰置于宾语之后（符合尾重原则 L011 与旧新信息流 L001）。}
  \item \textbf{S2 [宏观趋势 1:} \textit{Overall, renewable sources---particularly solar and wind power---followed a pronounced upward trajectory over the 25-year span, whereas coal reliance contracted steadily.}\par\noindent{\small 对应步骤：S2 [宏观趋势 1：总体走向与反超 Overview Trend 1]：Overall, while [Category A and B] experienced a steady upward trajectory, [Category C] declined markedly, leading to a reversal in dominance around [Mid-Year]. | 句首锚定旧信息/逻辑视角，主干动词紧跟主语，复杂修饰置于宾语之后（符合尾重原则 L011 与旧新信息流 L001）。}
  \item \textbf{S3 [宏观趋势 2:} \textit{Most notably, clean energy generation overtook fossil-fuel consumption around 2018, shifting from a marginal supplement to the dominant power supply by the end of the period.}\par\noindent{\small 对应步骤：S3 [宏观趋势 2：极值或差距特征 Overview Trend 2]：Notably, [Category A] registered the most dramatic growth throughout the period, ultimately overtaking all other sectors to become the primary contributor. | 句首锚定旧信息/逻辑视角，主干动词紧跟主语，复杂修饰置于宾语之后（符合尾重原则 L011 与旧新信息流 L001）。}
\end{itemize}
\end{TipBox}
\medskip
\subsection*{Paragraph 2: 动态时间序列趋势分段描述（暴涨、缓降、平台期与拐点） (\texttt{BP26})}
\noindent\textbf{Expressive Intent:}\enspace 我这一段想写动态折线图（或模型训练收敛曲线）的主体细节段：把一条经历“初期缓慢爬升 $\to$ 中期爆发增长 $\to$ 后期高位平台期”的曲线，以及另一条“冲高回落”的曲线对比写活。\par
\noindent\textbf{Cohesion Formula:}\enspace {\small\color{AcadSlate}Initial Disparity $\to$ Multi-Phase Surge \& Plateau (Entity A) $\to$ Peak-then-Decline (Entity B) $\to$ Gap Reversal. 绑定逻辑规则：L010, L013.}\par\vspace{4pt}
\begin{ExampleBox}{Publication-Grade Authentic Text}
In 2000, battery-electric vehicles accounted for a mere 4\% of municipal fleet purchases, trailing diesel buses at 62\%. Following five years of marginal fluctuation, electric adoption accelerated sharply after 2006, quadrupling to 38\% by 2018 before leveling off at approximately 41\% through 2025. By contrast, diesel procurement peaked briefly at 67\% in 2004 but subsequently underwent a sustained contraction, plunging to 19\% by the end of the survey. This divergent trajectory reversed the initial 58-percentage-point deficit into a more than twofold lead for electric fleets.
\end{ExampleBox}
\begin{TipBox}{Rhetorical Move Breakdown \& Constituent Analysis}
\begin{itemize}[leftmargin=1.5em]
  \item \textbf{S1 [起始点位对比 Initial Standing]:} \textit{In 2000, battery-electric vehicles accounted for a mere 4\% of municipal fleet purchases, trailing diesel buses at 62\%.}\par\noindent{\small 对应步骤：S1 [起始点位对比 Initial Standing]：At the beginning of the period in [2000], [Entity A] stood at a modest [12\%],less than a quarter of [Entity B]'s initial figure ([54\%]). | 句首锚定旧信息/逻辑视角，主干动词紧跟主语，复杂修饰置于宾语之后（符合尾重原则 L011 与旧新信息流 L001）。}
  \item \textbf{S2 [多阶段演进:} \textit{Following five years of marginal fluctuation, electric adoption accelerated sharply after 2006, quadrupling to 38\% by 2018 before leveling off at approximately 41\% through 2025.}\par\noindent{\small 对应步骤：S2 [多阶段演进：缓升转暴涨 Phased Acceleration]：After rising marginally over the firstfive years, [Entity A] surged nearly threefold between [2005] and [2018], peaking at [48\%] before stabilizing on a high plateau. | 句首锚定旧信息/逻辑视角，主干动词紧跟主语，复杂修饰置于宾语之后（符合尾重原则 L011 与旧新信息流 L001）。}
  \item \textbf{S3 [反向曲线对比:} \textit{By contrast, diesel procurement peaked briefly at 67\% in 2004 but subsequently underwent a sustained contraction, plunging to 19\% by the end of the survey.}\par\noindent{\small 对应步骤：S3 [反向曲线对比：冲高回落 Counter-Trajectory]：Conversely, [Entity B] climbed briefly to [61\%] in [2004] but subsequently entered a prolonged decline, dropping by more than half to [24\%] by [2025]. | 句首锚定旧信息/逻辑视角，主干动词紧跟主语，复杂修饰置于宾语之后（符合尾重原则 L011 与旧新信息流 L001）。}
  \item \textbf{S4 [终局交叉总结 Terminal Comparison]:} \textit{This divergent trajectory reversed the initial 58-percentage-point deficit into a more than twofold lead for electric fleets.}\par\noindent{\small 对应步骤：S4 [终局交叉总结 Terminal Comparison]：Consequently, the initial [42]-percentage-point gap between the two categories had completely reversed by the final year. | 句首锚定旧信息/逻辑视角，主干动词紧跟主语，复杂修饰置于宾语之后（符合尾重原则 L011 与旧新信息流 L001）。}
\end{itemize}
\end{TipBox}
\medskip
\subsection*{Paragraph 3: 异常数据点（Outliers）与局部背离现象精准批注段（Outlier \& Anomaly Annotation） (\texttt{BP32})}
\noindent\textbf{Expressive Intent:}\enspace 我这一段想在图表或实验结果段中，专门挑出违背总体规律的 1--2 个异常点（Outliers）或局部反常尖峰进行精准描述与成因剖析。\par
\noindent\textbf{Cohesion Formula:}\enspace {\small\color{AcadSlate}General Rule vs. Notable Departure $\to$ Quantitative Deviation $\to$ Root Cause Attribution $\to$ Calibration Takeaway. 绑定逻辑规则：L007, L013.}\par\vspace{4pt}
\begin{ExampleBox}{Publication-Grade Authentic Text}
Although deep spectral transformers outperform linear baselines across six of the seven benchmarks, a notable departure emerges on the financial Exchange-Rate dataset. Whereas complex attention models reduce MSE by 12\%--18\% on periodic energy and weather streams, they trail a single-layer DLinear baseline by 4.3\% on Exchange-Rate. Autocorrelation analysis reveals that this reversal stems from the near-martingale nature of daily currency returns, where weak seasonal periodicity causes high-capacity attention layers to overfit transient market shocks. This anomaly confirms that model capacity must be matched to the intrinsic signal-to-noise ratio of the underlying stochastic process.
\end{ExampleBox}
\begin{TipBox}{Rhetorical Move Breakdown \& Constituent Analysis}
\begin{itemize}[leftmargin=1.5em]
  \item \textbf{S1 [总体规律背景下的特例锚定 Outlier Identification]:} \textit{Although deep spectral transformers outperform linear baselines across six of the seven benchmarks, a notable departure emerges on the financial Exchange-Rate dataset.}\par\noindent{\small 对应步骤：S1 [总体规律背景下的特例锚定 Outlier Identification]：Although [General Trend holds across most groups/steps], a notable departure from this pattern occurs in [Specific Point/Group]. | 句首锚定旧信息/逻辑视角，主干动词紧跟主语，复杂修饰置于宾语之后（符合尾重原则 L011 与旧新信息流 L001）。}
  \item \textbf{S2 [反常幅度量化 Quantitative Deviation]:} \textit{Whereas complex attention models reduce MSE by 12\%--18\% on periodic energy and weather streams, they trail a single-layer DLinear baseline by 4.3\% on Exchange-Rate.}\par\noindent{\small 对应步骤：S2 [反常幅度量化 Quantitative Deviation]：Whereas peer categories cluster tightly within $[A, B]$, [Outlier Entity] spikes/drops to [Extreme Value], deviating by more than [X\%] from the group median. | 句首锚定旧信息/逻辑视角，主干动词紧跟主语，复杂修饰置于宾语之后（符合尾重原则 L011 与旧新信息流 L001）。}
  \item \textbf{S3 [底层物理/统计成因剖析 Root Cause Analysis]:} \textit{Autocorrelation analysis reveals that this reversal stems from the near-martingale nature of daily currency returns, where weak seasonal periodicity causes high-capacity attention layers to overfit transient market shocks.}\par\noindent{\small 对应步骤：S3 [底层物理/统计成因剖析 Root Cause Analysis]：Closer inspection indicates that this anomaly stems from [Specific Structural/External Factor] rather than random variance. | 句首锚定旧信息/逻辑视角，主干动词紧跟主语，复杂修饰置于宾语之后（符合尾重原则 L011 与旧新信息流 L001）。}
  \item \textbf{S4 [方法论启示 Methodological Implication]:} \textit{This anomaly confirms that model capacity must be matched to the intrinsic signal-to-noise ratio of the underlying stochastic process.}\par\noindent{\small 对应步骤：S4 [方法论启示 Methodological Implication]：Accounting for such localized departures is critical when calibrating [Downstream Model/Policy]. | 句首锚定旧信息/逻辑视角，主干动词紧跟主语，复杂修饰置于宾语之后（符合尾重原则 L011 与旧新信息流 L001）。}
\end{itemize}
\end{TipBox}
\medskip
\clearpage
\section{FT14: 雅思小作文（Task 1）静态柱状图与饼图组合满分范文：云端推理与机载机器人控制系统的能耗结构对比}\label{comp:FT14}
\begin{ConceptBox}{Composition Architecture \& Structural Blueprint Chain}
\textbf{Genre:} \texttt{IELTS\_Task1\_Static} \quad \textbf{Paragraph Count:} 3\par
\textbf{Blueprint Chain:} \texttt{BP25 $\to$ BP27 $\to$ BP28}
\end{ConceptBox}
\subsection*{Paragraph 1: 图表小作文首段：题目改写与总体趋势概览（Paraphrase + Macro Overview） (\texttt{BP25})}
\noindent\textbf{Expressive Intent:}\enspace 我这一段想写雅思 Task 1 或数据报告的第一段：用同义句法改写图表标题，并用两句话高度凝练全图最核心的 2 个宏观趋势（极值、反超拐点或主次格局），绝不堆砌零碎数字。\par
\noindent\textbf{Cohesion Formula:}\enspace {\small\color{AcadSlate}Paraphrase $\to$ Overall Trajectory Contrast $\to$ Dominant Crossover/Extremum (Zero raw numbers in Overview!). 绑定逻辑规则：L013, L015.}\par\vspace{4pt}
\begin{ExampleBox}{Publication-Grade Authentic Text}
The line graph compares the proportion of electricity generated from four fuel sources across three industrial economies between 2000 and 2025. Overall, renewable sources---particularly solar and wind power---followed a pronounced upward trajectory over the 25-year span, whereas coal reliance contracted steadily. Most notably, clean energy generation overtook fossil-fuel consumption around 2018, shifting from a marginal supplement to the dominant power supply by the end of the period.
\end{ExampleBox}
\begin{TipBox}{Rhetorical Move Breakdown \& Constituent Analysis}
\begin{itemize}[leftmargin=1.5em]
  \item \textbf{S1 [图表命题同义改写 Prompt Paraphrase]:} \textit{The line graph compares the proportion of electricity generated from four fuel sources across three industrial economies between 2000 and 2025.}\par\noindent{\small 对应步骤：S1 [图表命题同义改写 Prompt Paraphrase]：The line graph / bar chart compares [Variable/Metric] across [Categories/Regions] over a [N]-year period from [Year A] to [Year B]. | 句首锚定旧信息/逻辑视角，主干动词紧跟主语，复杂修饰置于宾语之后（符合尾重原则 L011 与旧新信息流 L001）。}
  \item \textbf{S2 [宏观趋势 1:} \textit{Overall, renewable sources---particularly solar and wind power---followed a pronounced upward trajectory over the 25-year span, whereas coal reliance contracted steadily.}\par\noindent{\small 对应步骤：S2 [宏观趋势 1：总体走向与反超 Overview Trend 1]：Overall, while [Category A and B] experienced a steady upward trajectory, [Category C] declined markedly, leading to a reversal in dominance around [Mid-Year]. | 句首锚定旧信息/逻辑视角，主干动词紧跟主语，复杂修饰置于宾语之后（符合尾重原则 L011 与旧新信息流 L001）。}
  \item \textbf{S3 [宏观趋势 2:} \textit{Most notably, clean energy generation overtook fossil-fuel consumption around 2018, shifting from a marginal supplement to the dominant power supply by the end of the period.}\par\noindent{\small 对应步骤：S3 [宏观趋势 2：极值或差距特征 Overview Trend 2]：Notably, [Category A] registered the most dramatic growth throughout the period, ultimately overtaking all other sectors to become the primary contributor. | 句首锚定旧信息/逻辑视角，主干动词紧跟主语，复杂修饰置于宾语之后（符合尾重原则 L011 与旧新信息流 L001）。}
\end{itemize}
\end{TipBox}
\medskip
\subsection*{Paragraph 2: 静态多组数据横向对比与倍差/占比凸显段（Cross-Group Proportions） (\texttt{BP27})}
\noindent\textbf{Expressive Intent:}\enspace 我这一段想写静态柱状图或饼图的主体段：横向对比不同国家/不同模型在不同开销维度上的占比差异，灵活运用倍数（double/triple）、分数（nearly a third）与排序结构。\par
\noindent\textbf{Cohesion Formula:}\enspace {\small\color{AcadSlate}Dominant Item with Multiplier $\to$ Sharp Cross-Group Contrast $\to$ Pooled Minor Items $\to$ Structural Synthesis. 绑定逻辑规则：L009, L010.}\par\vspace{4pt}
\begin{ExampleBox}{Publication-Grade Authentic Text}
Within edge-device inference workloads, attention matrix computation consumes the largest fraction of total energy at 48\%, exactly double the power drawn by feed-forward layers (24\%). A strikingly different profile emerges in onboard robotic control loops, where high-frequency sensor preprocessing dominates energy usage at 51\%, while attention overhead contracts to just 16\%. Across both deployment settings, memory offloading and actuation decoding collectively account for under 11\% of total consumption. This contrast highlights that optimization bottlenecks shift from quadratic token mixing on cloud servers to raw sensory I/O on physical robots.
\end{ExampleBox}
\begin{TipBox}{Rhetorical Move Breakdown \& Constituent Analysis}
\begin{itemize}[leftmargin=1.5em]
  \item \textbf{S1 [主导项锁定 Dominant Category]:} \textit{Within edge-device inference workloads, attention matrix computation consumes the largest fraction of total energy at 48\%, exactly double the power drawn by feed-forward layers (24\%).}\par\noindent{\small 对应步骤：S1 [主导项锁定 Dominant Category]：Regarding [Group 1], [Category A] claims the largest share at [46\%], nearly double the proportion allocated to [Category B] ([24\%]). | 句首锚定旧信息/逻辑视角，主干动词紧跟主语，复杂修饰置于宾语之后（符合尾重原则 L011 与旧新信息流 L001）。}
  \item \textbf{S2 [组间横向反差 Cross-Group Contrast]:} \textit{A strikingly different profile emerges in onboard robotic control loops, where high-frequency sensor preprocessing dominates energy usage at 51\%, while attention overhead contracts to just 16\%.}\par\noindent{\small 对应步骤：S2 [组间横向反差 Cross-Group Contrast]：A markedly different allocation pattern emerges in [Group 2], where [Category C] dominates at [52\%], while [Category A] shrinks to just [15\%]. | 句首锚定旧信息/逻辑视角，主干动词紧跟主语，复杂修饰置于宾语之后（符合尾重原则 L011 与旧新信息流 L001）。}
  \item \textbf{S3 [次要项合并归纳 Minor Categories Grouping]:} \textit{Across both deployment settings, memory offloading and actuation decoding collectively account for under 11\% of total consumption.}\par\noindent{\small 对应步骤：S3 [次要项合并归纳 Minor Categories Grouping]：Meanwhile, the remaining categories---[Category D and E]---collectively account for less than [12\%] across both groups. | 句首锚定旧信息/逻辑视角，主干动词紧跟主语，复杂修饰置于宾语之后（符合尾重原则 L011 与旧新信息流 L001）。}
  \item \textbf{S4 [结构性差异提炼 Structural Takeaway]:} \textit{This contrast highlights that optimization bottlenecks shift from quadratic token mixing on cloud servers to raw sensory I/O on physical robots.}\par\noindent{\small 对应步骤：S4 [结构性差异提炼 Structural Takeaway]：These figures indicate that whereas [Group 1] prioritizes [Function A], [Group 2] concentrates resources heavily on [Function C]. | 句首锚定旧信息/逻辑视角，主干动词紧跟主语，复杂修饰置于宾语之后（符合尾重原则 L011 与旧新信息流 L001）。}
\end{itemize}
\end{TipBox}
\medskip
\subsection*{Paragraph 3: 双图关联与互补/因果印证段（Synthesizing Dual Charts） (\texttt{BP28})}
\noindent\textbf{Expressive Intent:}\enspace 我这一段想写双图题（如一张图是算力投入/研发经费变化，另一张图是模型误差/碳排放强度变化）的综合关联段，阐明两张图之间的内在协同或背离关系。\par
\noindent\textbf{Cohesion Formula:}\enspace {\small\color{AcadSlate}Global Correlation $\to$ Co-movement Phase $\to$ Diminishing Returns / Divergence $\to$ Notable Outlier. 绑定逻辑规则：L007, L013.}\par\vspace{4pt}
\begin{ExampleBox}{Publication-Grade Authentic Text}
Cross-referencing both figures reveals a strong inverse relationship between automated inspection investment and manufacturing defect rates up to a \$40 million threshold. Specifically, across plants investing between \$10 million and \$40 million (Figure 1), surface defect incidence dropped steeply from 4.8% to 0.9% (Figure 2). Beyond \$40 million, however, further capital expenditure yielded diminishing returns: doubling investment to \$80 million reduced defects by only an additional 0.1 percentage points. Notably, Plant D defied this expenditure curve, achieving a 0.7% defect rate with just \$22 million through closed-loop sensor calibration.
\end{ExampleBox}
\begin{TipBox}{Rhetorical Move Breakdown \& Constituent Analysis}
\begin{itemize}[leftmargin=1.5em]
  \item \textbf{S1 [双图联动锚定 Cross-Chart Link]:} \textit{Cross-referencing both figures reveals a strong inverse relationship between automated inspection investment and manufacturing defect rates up to a \$40 million threshold.}\par\noindent{\small 对应步骤：S1 [双图联动锚定 Cross-Chart Link]：Synthesizing both charts reveals a clear inverse / positive relationship between [Chart 1 Variable] and [Chart 2 Variable] up to [Threshold Point]. | 句首锚定旧信息/逻辑视角，主干动词紧跟主语，复杂修饰置于宾语之后（符合尾重原则 L011 与旧新信息流 L001）。}
  \item \textbf{S2 [区间 1 协同变化 Phase 1 Co-movement]:} \textit{Specifically, across plants investing between \$10 million and \$40 million (Figure 1), surface defect incidence dropped steeply from 4.8\% to 0.9\% (Figure 2).}\par\noindent{\small 对应步骤：S2 [区间 1 协同变化 Phase 1 Co-movement]：Specifically, as [Variable X in Chart 1] expanded from [A] to [B] during [Period 1], [Variable Y in Chart 2] dropped correspondingly by [C\%]. | 句首锚定旧信息/逻辑视角，主干动词紧跟主语，复杂修饰置于宾语之后（符合尾重原则 L011 与旧新信息流 L001）。}
  \item \textbf{S3 [区间 2 边际递减/背离 Diminishing Returns / Divergence]:} \textit{Beyond \$40 million, however, further capital expenditure yielded diminishing returns: doubling investment to \$80 million reduced defects by only an additional 0.1 percentage points.}\par\noindent{\small 对应步骤：S3 [区间 2 边际递减/背离 Diminishing Returns / Divergence]：However, beyond [Threshold Point], although [Variable X] continued to double, [Variable Y] plateaued around [Value], indicating diminishing marginal returns. | 句首锚定旧信息/逻辑视角，主干动词紧跟主语，复杂修饰置于宾语之后（符合尾重原则 L011 与旧新信息流 L001）。}
  \item \textbf{S4 [例外个体点拨 Exception Highlight]:} \textit{Notably, Plant D defied this expenditure curve, achieving a 0.7\% defect rate with just \$22 million through closed-loop sensor calibration.}\par\noindent{\small 对应步骤：S4 [例外个体点拨 Exception Highlight]：An exception to this pattern is [Entity Z], which achieved the lowest [Variable Y] despite maintaining moderate [Variable X]. | 句首锚定旧信息/逻辑视角，主干动词紧跟主语，复杂修饰置于宾语之后（符合尾重原则 L011 与旧新信息流 L001）。}
\end{itemize}
\end{TipBox}
\medskip
\clearpage
\section{FT15: 雅思小作文（Task 1）闭环工艺与算法流程图满分范文：具身智能遥操作数据采集、双支路增强与主动学习回环}\label{comp:FT15}
\begin{ConceptBox}{Composition Architecture \& Structural Blueprint Chain}
\textbf{Genre:} \texttt{IELTS\_Task1\_Process} \quad \textbf{Paragraph Count:} 3\par
\textbf{Blueprint Chain:} \texttt{BP25 $\to$ BP29 $\to$ BP30}
\end{ConceptBox}
\subsection*{Paragraph 1: 图表小作文首段：题目改写与总体趋势概览（Paraphrase + Macro Overview） (\texttt{BP25})}
\noindent\textbf{Expressive Intent:}\enspace 我这一段想写雅思 Task 1 或数据报告的第一段：用同义句法改写图表标题，并用两句话高度凝练全图最核心的 2 个宏观趋势（极值、反超拐点或主次格局），绝不堆砌零碎数字。\par
\noindent\textbf{Cohesion Formula:}\enspace {\small\color{AcadSlate}Paraphrase $\to$ Overall Trajectory Contrast $\to$ Dominant Crossover/Extremum (Zero raw numbers in Overview!). 绑定逻辑规则：L013, L015.}\par\vspace{4pt}
\begin{ExampleBox}{Publication-Grade Authentic Text}
The line graph compares the proportion of electricity generated from four fuel sources across three industrial economies between 2000 and 2025. Overall, renewable sources---particularly solar and wind power---followed a pronounced upward trajectory over the 25-year span, whereas coal reliance contracted steadily. Most notably, clean energy generation overtook fossil-fuel consumption around 2018, shifting from a marginal supplement to the dominant power supply by the end of the period.
\end{ExampleBox}
\begin{TipBox}{Rhetorical Move Breakdown \& Constituent Analysis}
\begin{itemize}[leftmargin=1.5em]
  \item \textbf{S1 [图表命题同义改写 Prompt Paraphrase]:} \textit{The line graph compares the proportion of electricity generated from four fuel sources across three industrial economies between 2000 and 2025.}\par\noindent{\small 对应步骤：S1 [图表命题同义改写 Prompt Paraphrase]：The line graph / bar chart compares [Variable/Metric] across [Categories/Regions] over a [N]-year period from [Year A] to [Year B]. | 句首锚定旧信息/逻辑视角，主干动词紧跟主语，复杂修饰置于宾语之后（符合尾重原则 L011 与旧新信息流 L001）。}
  \item \textbf{S2 [宏观趋势 1:} \textit{Overall, renewable sources---particularly solar and wind power---followed a pronounced upward trajectory over the 25-year span, whereas coal reliance contracted steadily.}\par\noindent{\small 对应步骤：S2 [宏观趋势 1：总体走向与反超 Overview Trend 1]：Overall, while [Category A and B] experienced a steady upward trajectory, [Category C] declined markedly, leading to a reversal in dominance around [Mid-Year]. | 句首锚定旧信息/逻辑视角，主干动词紧跟主语，复杂修饰置于宾语之后（符合尾重原则 L011 与旧新信息流 L001）。}
  \item \textbf{S3 [宏观趋势 2:} \textit{Most notably, clean energy generation overtook fossil-fuel consumption around 2018, shifting from a marginal supplement to the dominant power supply by the end of the period.}\par\noindent{\small 对应步骤：S3 [宏观趋势 2：极值或差距特征 Overview Trend 2]：Notably, [Category A] registered the most dramatic growth throughout the period, ultimately overtaking all other sectors to become the primary contributor. | 句首锚定旧信息/逻辑视角，主干动词紧跟主语，复杂修饰置于宾语之后（符合尾重原则 L011 与旧新信息流 L001）。}
\end{itemize}
\end{TipBox}
\medskip
\subsection*{Paragraph 2: 流程图/算法流水线与闭环工序衔接段（Process \& Closed-Loop Lifecycle） (\texttt{BP29})}
\noindent\textbf{Expressive Intent:}\enspace 我这一段想写雅思 Task 1 工艺/自然循环流程图（或论文算法数据预处理与闭环训练流水线）：把 5--6 个串联、分叉与回环步骤写得层次分明，绝不机械连用 Then, Then, Then。\par
\noindent\textbf{Cohesion Formula:}\enspace {\small\color{AcadSlate}Pipeline Overview $\to$ Ingestion \& Filtering $\to$ Parallel Bifurcation $\to$ Terminal Merger \& Feedback Loop. 绑定逻辑规则：L003, L005.}\par\vspace{4pt}
\begin{ExampleBox}{Publication-Grade Authentic Text}
The diagram illustrates the closed-loop data curation pipeline for embodied robot learning, spanning four primary stages from teleoperated collection to failure-driven retraining. The cycle begins when raw multimodal demonstrations are recorded via VR teleoperation and filtered to discard kinematic singularities and dropped frames. Once validated, the trajectory stream bifurcates into two parallel pathways: visual frames undergo background augmentation, while proprioceptive states are normalized into local gripper coordinates. These processed streams are subsequently synchronized to train the visuomotor policy, whose real-world rollout failures are logged and fed back into the teleoperation buffer to close the active-learning loop.
\end{ExampleBox}
\begin{TipBox}{Rhetorical Move Breakdown \& Constituent Analysis}
\begin{itemize}[leftmargin=1.5em]
  \item \textbf{S1 [流程总括与起止界定 Lifecycle Overview]:} \textit{The diagram illustrates the closed-loop data curation pipeline for embodied robot learning, spanning four primary stages from teleoperated collection to failure-driven retraining.}\par\noindent{\small 对应步骤：S1 [流程总括与起止界定 Lifecycle Overview]：The diagram illustrates the closed-loop pipeline for [Process Name], comprising [N] main stages from initial [Input Stage] to continuous [Feedback/Output Stage]. | 句首锚定旧信息/逻辑视角，主干动词紧跟主语，复杂修饰置于宾语之后（符合尾重原则 L011 与旧新信息流 L001）。}
  \item \textbf{S2 [前段采集与预处理 Initial Ingestion \& Filtering]:} \textit{The cycle begins when raw multimodal demonstrations are recorded via VR teleoperation and filtered to discard kinematic singularities and dropped frames.}\par\noindent{\small 对应步骤：S2 [前段采集与预处理 Initial Ingestion \& Filtering]：The process begins when raw [Material/Data] is collected via [Source] and passed through [Filter/Chamber] to remove [Impurities/Outliers]. | 句首锚定旧信息/逻辑视角，主干动词紧跟主语，复杂修饰置于宾语之后（符合尾重原则 L011 与旧新信息流 L001）。}
  \item \textbf{S3 [中段并行分流 Parallel Branching]:} \textit{Once validated, the trajectory stream bifurcates into two parallel pathways: visual frames undergo background augmentation, while proprioceptive states are normalized into local gripper coordinates.}\par\noindent{\small 对应步骤：S3 [中段并行分流 Parallel Branching]：Once purified, the stream bifurcates into two parallel pathways: [Fraction A] undergoes [Process A], whereas [Fraction B] is routed to [Process B]. | 句首锚定旧信息/逻辑视角，主干动词紧跟主语，复杂修饰置于宾语之后（符合尾重原则 L011 与旧新信息流 L001）。}
  \item \textbf{S4 [末端合成与闭环回流 Synthesis \& Closed-Loop Recycling]:} \textit{These processed streams are subsequently synchronized to train the visuomotor policy, whose real-world rollout failures are logged and fed back into the teleoperation buffer to close the active-learning loop.}\par\noindent{\small 对应步骤：S4 [末端合成与闭环回流 Synthesis \& Closed-Loop Recycling]：The resulting outputs are subsequently merged in [Unit C], with residual by-products / high-error trajectories recycled back into the initial buffer for iterative refinement. | 句首锚定旧信息/逻辑视角，主干动词紧跟主语，复杂修饰置于宾语之后（符合尾重原则 L011 与旧新信息流 L001）。}
\end{itemize}
\end{TipBox}
\medskip
\subsection*{Paragraph 3: 历史实测与未来预测跨时态衔接段（Historical vs. Projected Trends） (\texttt{BP30})}
\noindent\textbf{Expressive Intent:}\enspace 我这一段想处理同时包含过去历史数据（如 2010--2025）和未来预测数据（如 2026--2040）的图表段落，严格区分过去实测时态与未来审慎预测表达。\par
\noindent\textbf{Cohesion Formula:}\enspace {\small\color{AcadSlate}Historical Past (rose/accounted for) $\to$ Pivot Year $\to$ Future Projection Verbs (is projected/anticipated/expected to) $\to$ Moderating Factor. 绑定逻辑规则：L010, L017.}\par\vspace{4pt}
\begin{ExampleBox}{Publication-Grade Authentic Text}
Between 2010 and 2025, global data-center electricity consumption more than doubled from 200 TWh to 460 TWh. From 2026 onward, total demand is projected to climb at a substantially steeper rate, with baseline forecasts pointing to 1,050 TWh by 2035. Within this expansion, AI model training and inference---which represented under 10\% of server load in 2023---are anticipated to account for nearly half (46\%) of total consumption by 2040. However, widespread deployment of sparsity-aware accelerators is expected to temper the annual growth rate after 2035.
\end{ExampleBox}
\begin{TipBox}{Rhetorical Move Breakdown \& Constituent Analysis}
\begin{itemize}[leftmargin=1.5em]
  \item \textbf{S1 [历史实测轨迹 Historical Trajectory (Past Tense)]:} \textit{Between 2010 and 2025, global data-center electricity consumption more than doubled from 200 TWh to 460 TWh.}\par\noindent{\small 对应步骤：S1 [历史实测轨迹 Historical Trajectory (Past Tense)]：Between [2010] and [2025], global data-center electricity demand rose steadily from [200 TWh] to [460 TWh], driven by cloud migration. | 句首锚定旧信息/逻辑视角，主干动词紧跟主语，复杂修饰置于宾语之后（符合尾重原则 L011 与旧新信息流 L001）。}
  \item \textbf{S2 [当下转折与未来加速 Projected Acceleration (Projection Verbs)]:} \textit{From 2026 onward, total demand is projected to climb at a substantially steeper rate, with baseline forecasts pointing to 1,050 TWh by 2035.}\par\noindent{\small 对应步骤：S2 [当下转折与未来加速 Projected Acceleration (Projection Verbs)]：From [2026] onward, consumption is projected to accelerate sharply, with forecasts indicating a surge to [1,050 TWh] by [2035]. | 句首锚定旧信息/逻辑视角，主干动词紧跟主语，复杂修饰置于宾语之后（符合尾重原则 L011 与旧新信息流 L001）。}
  \item \textbf{S3 [结构占比重塑 Structural Shift in Projections]:} \textit{Within this expansion, AI model training and inference---which represented under 10\% of server load in 2023---are anticipated to account for nearly half (46\%) of total consumption by 2040.}\par\noindent{\small 对应步骤：S3 [结构占比重塑 Structural Shift in Projections]：Within this projected expansion, generative AI training and inference---which accounted for less than [10\%] in [2023]---are anticipated to represent over [45\%] of total load by [2040]. | 句首锚定旧信息/逻辑视角，主干动词紧跟主语，复杂修饰置于宾语之后（符合尾重原则 L011 与旧新信息流 L001）。}
  \item \textbf{S4 [条件式审慎收尾 Conditional Hedging]:} \textit{However, widespread deployment of sparsity-aware accelerators is expected to temper the annual growth rate after 2035.}\par\noindent{\small 对应步骤：S4 [条件式审慎收尾 Conditional Hedging]：Nevertheless, energy-efficient hardware adoption is expected to moderate the growth slope during the final five years. | 句首锚定旧信息/逻辑视角，主干动词紧跟主语，复杂修饰置于宾语之后（符合尾重原则 L011 与旧新信息流 L001）。}
\end{itemize}
\end{TipBox}
\medskip
\clearpage
\chapter{Engineering, Materials \& Physical Sciences / 工程材料与物理科学}\label{ch:anthology_21}
\Needspace{6\baselineskip}
本章精选该学科领域及对应文体的完整学术范文。每一篇范文不仅展示母语级发表标准的正文全貌，更清晰拆解全篇的“段落蓝图链条（Structural Blueprint Chain）”，并对每个自然段逐一剖析其核心交际意图、句际衔接逻辑公式（Cohesion Formula）以及每句话的句法成分与推进机制。

\section{FT16: 托福综合写作（Integrated Writing）满分范文：自主无人机配送能否取代城市货车（阅读三论点 vs. 听力逐点反驳）}\label{comp:FT16}
\begin{ConceptBox}{Composition Architecture \& Structural Blueprint Chain}
\textbf{Genre:} \texttt{TOEFL\_Integrated} \quad \textbf{Paragraph Count:} 4\par
\textbf{Blueprint Chain:} \texttt{BP31 $\to$ BP27 $\to$ BP28 $\to$ BP32}
\end{ConceptBox}
\subsection*{Paragraph 1: 托福综合写作 / 学术争议正反交锋段（Point-by-Point Reading vs. Lecture Synthesis） (\texttt{BP31})}
\noindent\textbf{Expressive Intent:}\enspace 我这一段想写托福综合写作（Integrated Writing）或文献争议交锋段：先用 1 句话精准概括阅读材料（或正方文献）的分论点，再用 2--3 句话详细呈现教授讲座（或反方实证）是如何釜底抽薪推翻该论点的。\par
\noindent\textbf{Cohesion Formula:}\enspace {\small\color{AcadSlate}Reading Claim $\to$ Lecturer's Direct Challenge $\to$ Specific Counter-Mechanism $\to$ Reversed Outcome. 绑定逻辑规则：L008, L019.}\par\vspace{4pt}
\begin{ExampleBox}{Publication-Grade Authentic Text}
The reading passage contends that replacing urban delivery vans with autonomous quadcopter drones will substantially reduce last-mile transportation energy consumption, citing the lighter weight of aerial vehicles. The lecturer directly challenges this premise by examining energy efficiency per kilogram of delivered freight. Specifically, she points out that unlike ground vehicles supported passively by road surfaces, multi-rotor drones must expend continuous battery power simply to counteract gravity and headwinds while carrying only a single parcel per trip. Consequently, servicing a dense apartment block via individual drone flights consumes nearly three times more energy per package than a single consolidated van route.
\end{ExampleBox}
\begin{TipBox}{Rhetorical Move Breakdown \& Constituent Analysis}
\begin{itemize}[leftmargin=1.5em]
  \item \textbf{S1 [正方/阅读论点概括 Reading Claim \& Rationale]:} \textit{The reading passage contends that replacing urban delivery vans with autonomous quadcopter drones will substantially reduce last-mile transportation energy consumption, citing the lighter weight of aerial vehicles.}\par\noindent{\small 对应步骤：S1 [正方/阅读论点概括 Reading Claim \& Rationale]：The reading passage contends that [Claim A], arguing that [Reading's Supporting Mechanism]. | 句首锚定旧信息/逻辑视角，主干动词紧跟主语，复杂修饰置于宾语之后（符合尾重原则 L011 与旧新信息流 L001）。}
  \item \textbf{S2 [反方/听力直接反驳 Lecture Direct Refutation]:} \textit{The lecturer directly challenges this premise by examining energy efficiency per kilogram of delivered freight.}\par\noindent{\small 对应步骤：S2 [反方/听力直接反驳 Lecture Direct Refutation]：The lecturer directly challenges this assertion by pointing out that [Core Flaw / Neglected Factor in Reading's Premise]. | 句首锚定旧信息/逻辑视角，主干动词紧跟主语，复杂修饰置于宾语之后（符合尾重原则 L011 与旧新信息流 L001）。}
  \item \textbf{S3 [反方具体反证展开 Empirical Counter-Evidence]:} \textit{Specifically, she points out that unlike ground vehicles supported passively by road surfaces, multi-rotor drones must expend continuous battery power simply to counteract gravity and headwinds while carrying only a single parcel per trip.}\par\noindent{\small 对应步骤：S3 [反方具体反证展开 Empirical Counter-Evidence]：Specifically, she explains that when [Real-World Condition], [Mechanism B occurs instead], which [Negates the Claimed Benefit]. | 句首锚定旧信息/逻辑视角，主干动词紧跟主语，复杂修饰置于宾语之后（符合尾重原则 L011 与旧新信息流 L001）。}
  \item \textbf{S4 [结论锁定 Refutation Closure]:} \textit{Consequently, servicing a dense apartment block via individual drone flights consumes nearly three times more energy per package than a single consolidated van route.}\par\noindent{\small 对应步骤：S4 [结论锁定 Refutation Closure]：Consequently, rather than achieving [Reading's Goal], the proposed measure actually exacerbates [Negative Outcome]. | 句首锚定旧信息/逻辑视角，主干动词紧跟主语，复杂修饰置于宾语之后（符合尾重原则 L011 与旧新信息流 L001）。}
\end{itemize}
\end{TipBox}
\medskip
\subsection*{Paragraph 2: 静态多组数据横向对比与倍差/占比凸显段（Cross-Group Proportions） (\texttt{BP27})}
\noindent\textbf{Expressive Intent:}\enspace 我这一段想写静态柱状图或饼图的主体段：横向对比不同国家/不同模型在不同开销维度上的占比差异，灵活运用倍数（double/triple）、分数（nearly a third）与排序结构。\par
\noindent\textbf{Cohesion Formula:}\enspace {\small\color{AcadSlate}Dominant Item with Multiplier $\to$ Sharp Cross-Group Contrast $\to$ Pooled Minor Items $\to$ Structural Synthesis. 绑定逻辑规则：L009, L010.}\par\vspace{4pt}
\begin{ExampleBox}{Publication-Grade Authentic Text}
Within edge-device inference workloads, attention matrix computation consumes the largest fraction of total energy at 48\%, exactly double the power drawn by feed-forward layers (24\%). A strikingly different profile emerges in onboard robotic control loops, where high-frequency sensor preprocessing dominates energy usage at 51\%, while attention overhead contracts to just 16\%. Across both deployment settings, memory offloading and actuation decoding collectively account for under 11\% of total consumption. This contrast highlights that optimization bottlenecks shift from quadratic token mixing on cloud servers to raw sensory I/O on physical robots.
\end{ExampleBox}
\begin{TipBox}{Rhetorical Move Breakdown \& Constituent Analysis}
\begin{itemize}[leftmargin=1.5em]
  \item \textbf{S1 [主导项锁定 Dominant Category]:} \textit{Within edge-device inference workloads, attention matrix computation consumes the largest fraction of total energy at 48\%, exactly double the power drawn by feed-forward layers (24\%).}\par\noindent{\small 对应步骤：S1 [主导项锁定 Dominant Category]：Regarding [Group 1], [Category A] claims the largest share at [46\%], nearly double the proportion allocated to [Category B] ([24\%]). | 句首锚定旧信息/逻辑视角，主干动词紧跟主语，复杂修饰置于宾语之后（符合尾重原则 L011 与旧新信息流 L001）。}
  \item \textbf{S2 [组间横向反差 Cross-Group Contrast]:} \textit{A strikingly different profile emerges in onboard robotic control loops, where high-frequency sensor preprocessing dominates energy usage at 51\%, while attention overhead contracts to just 16\%.}\par\noindent{\small 对应步骤：S2 [组间横向反差 Cross-Group Contrast]：A markedly different allocation pattern emerges in [Group 2], where [Category C] dominates at [52\%], while [Category A] shrinks to just [15\%]. | 句首锚定旧信息/逻辑视角，主干动词紧跟主语，复杂修饰置于宾语之后（符合尾重原则 L011 与旧新信息流 L001）。}
  \item \textbf{S3 [次要项合并归纳 Minor Categories Grouping]:} \textit{Across both deployment settings, memory offloading and actuation decoding collectively account for under 11\% of total consumption.}\par\noindent{\small 对应步骤：S3 [次要项合并归纳 Minor Categories Grouping]：Meanwhile, the remaining categories---[Category D and E]---collectively account for less than [12\%] across both groups. | 句首锚定旧信息/逻辑视角，主干动词紧跟主语，复杂修饰置于宾语之后（符合尾重原则 L011 与旧新信息流 L001）。}
  \item \textbf{S4 [结构性差异提炼 Structural Takeaway]:} \textit{This contrast highlights that optimization bottlenecks shift from quadratic token mixing on cloud servers to raw sensory I/O on physical robots.}\par\noindent{\small 对应步骤：S4 [结构性差异提炼 Structural Takeaway]：These figures indicate that whereas [Group 1] prioritizes [Function A], [Group 2] concentrates resources heavily on [Function C]. | 句首锚定旧信息/逻辑视角，主干动词紧跟主语，复杂修饰置于宾语之后（符合尾重原则 L011 与旧新信息流 L001）。}
\end{itemize}
\end{TipBox}
\medskip
\subsection*{Paragraph 3: 双图关联与互补/因果印证段（Synthesizing Dual Charts） (\texttt{BP28})}
\noindent\textbf{Expressive Intent:}\enspace 我这一段想写双图题（如一张图是算力投入/研发经费变化，另一张图是模型误差/碳排放强度变化）的综合关联段，阐明两张图之间的内在协同或背离关系。\par
\noindent\textbf{Cohesion Formula:}\enspace {\small\color{AcadSlate}Global Correlation $\to$ Co-movement Phase $\to$ Diminishing Returns / Divergence $\to$ Notable Outlier. 绑定逻辑规则：L007, L013.}\par\vspace{4pt}
\begin{ExampleBox}{Publication-Grade Authentic Text}
Cross-referencing both figures reveals a strong inverse relationship between automated inspection investment and manufacturing defect rates up to a \$40 million threshold. Specifically, across plants investing between \$10 million and \$40 million (Figure 1), surface defect incidence dropped steeply from 4.8% to 0.9% (Figure 2). Beyond \$40 million, however, further capital expenditure yielded diminishing returns: doubling investment to \$80 million reduced defects by only an additional 0.1 percentage points. Notably, Plant D defied this expenditure curve, achieving a 0.7% defect rate with just \$22 million through closed-loop sensor calibration.
\end{ExampleBox}
\begin{TipBox}{Rhetorical Move Breakdown \& Constituent Analysis}
\begin{itemize}[leftmargin=1.5em]
  \item \textbf{S1 [双图联动锚定 Cross-Chart Link]:} \textit{Cross-referencing both figures reveals a strong inverse relationship between automated inspection investment and manufacturing defect rates up to a \$40 million threshold.}\par\noindent{\small 对应步骤：S1 [双图联动锚定 Cross-Chart Link]：Synthesizing both charts reveals a clear inverse / positive relationship between [Chart 1 Variable] and [Chart 2 Variable] up to [Threshold Point]. | 句首锚定旧信息/逻辑视角，主干动词紧跟主语，复杂修饰置于宾语之后（符合尾重原则 L011 与旧新信息流 L001）。}
  \item \textbf{S2 [区间 1 协同变化 Phase 1 Co-movement]:} \textit{Specifically, across plants investing between \$10 million and \$40 million (Figure 1), surface defect incidence dropped steeply from 4.8\% to 0.9\% (Figure 2).}\par\noindent{\small 对应步骤：S2 [区间 1 协同变化 Phase 1 Co-movement]：Specifically, as [Variable X in Chart 1] expanded from [A] to [B] during [Period 1], [Variable Y in Chart 2] dropped correspondingly by [C\%]. | 句首锚定旧信息/逻辑视角，主干动词紧跟主语，复杂修饰置于宾语之后（符合尾重原则 L011 与旧新信息流 L001）。}
  \item \textbf{S3 [区间 2 边际递减/背离 Diminishing Returns / Divergence]:} \textit{Beyond \$40 million, however, further capital expenditure yielded diminishing returns: doubling investment to \$80 million reduced defects by only an additional 0.1 percentage points.}\par\noindent{\small 对应步骤：S3 [区间 2 边际递减/背离 Diminishing Returns / Divergence]：However, beyond [Threshold Point], although [Variable X] continued to double, [Variable Y] plateaued around [Value], indicating diminishing marginal returns. | 句首锚定旧信息/逻辑视角，主干动词紧跟主语，复杂修饰置于宾语之后（符合尾重原则 L011 与旧新信息流 L001）。}
  \item \textbf{S4 [例外个体点拨 Exception Highlight]:} \textit{Notably, Plant D defied this expenditure curve, achieving a 0.7\% defect rate with just \$22 million through closed-loop sensor calibration.}\par\noindent{\small 对应步骤：S4 [例外个体点拨 Exception Highlight]：An exception to this pattern is [Entity Z], which achieved the lowest [Variable Y] despite maintaining moderate [Variable X]. | 句首锚定旧信息/逻辑视角，主干动词紧跟主语，复杂修饰置于宾语之后（符合尾重原则 L011 与旧新信息流 L001）。}
\end{itemize}
\end{TipBox}
\medskip
\subsection*{Paragraph 4: 异常数据点（Outliers）与局部背离现象精准批注段（Outlier \& Anomaly Annotation） (\texttt{BP32})}
\noindent\textbf{Expressive Intent:}\enspace 我这一段想在图表或实验结果段中，专门挑出违背总体规律的 1--2 个异常点（Outliers）或局部反常尖峰进行精准描述与成因剖析。\par
\noindent\textbf{Cohesion Formula:}\enspace {\small\color{AcadSlate}General Rule vs. Notable Departure $\to$ Quantitative Deviation $\to$ Root Cause Attribution $\to$ Calibration Takeaway. 绑定逻辑规则：L007, L013.}\par\vspace{4pt}
\begin{ExampleBox}{Publication-Grade Authentic Text}
Although deep spectral transformers outperform linear baselines across six of the seven benchmarks, a notable departure emerges on the financial Exchange-Rate dataset. Whereas complex attention models reduce MSE by 12\%--18\% on periodic energy and weather streams, they trail a single-layer DLinear baseline by 4.3\% on Exchange-Rate. Autocorrelation analysis reveals that this reversal stems from the near-martingale nature of daily currency returns, where weak seasonal periodicity causes high-capacity attention layers to overfit transient market shocks. This anomaly confirms that model capacity must be matched to the intrinsic signal-to-noise ratio of the underlying stochastic process.
\end{ExampleBox}
\begin{TipBox}{Rhetorical Move Breakdown \& Constituent Analysis}
\begin{itemize}[leftmargin=1.5em]
  \item \textbf{S1 [总体规律背景下的特例锚定 Outlier Identification]:} \textit{Although deep spectral transformers outperform linear baselines across six of the seven benchmarks, a notable departure emerges on the financial Exchange-Rate dataset.}\par\noindent{\small 对应步骤：S1 [总体规律背景下的特例锚定 Outlier Identification]：Although [General Trend holds across most groups/steps], a notable departure from this pattern occurs in [Specific Point/Group]. | 句首锚定旧信息/逻辑视角，主干动词紧跟主语，复杂修饰置于宾语之后（符合尾重原则 L011 与旧新信息流 L001）。}
  \item \textbf{S2 [反常幅度量化 Quantitative Deviation]:} \textit{Whereas complex attention models reduce MSE by 12\%--18\% on periodic energy and weather streams, they trail a single-layer DLinear baseline by 4.3\% on Exchange-Rate.}\par\noindent{\small 对应步骤：S2 [反常幅度量化 Quantitative Deviation]：Whereas peer categories cluster tightly within $[A, B]$, [Outlier Entity] spikes/drops to [Extreme Value], deviating by more than [X\%] from the group median. | 句首锚定旧信息/逻辑视角，主干动词紧跟主语，复杂修饰置于宾语之后（符合尾重原则 L011 与旧新信息流 L001）。}
  \item \textbf{S3 [底层物理/统计成因剖析 Root Cause Analysis]:} \textit{Autocorrelation analysis reveals that this reversal stems from the near-martingale nature of daily currency returns, where weak seasonal periodicity causes high-capacity attention layers to overfit transient market shocks.}\par\noindent{\small 对应步骤：S3 [底层物理/统计成因剖析 Root Cause Analysis]：Closer inspection indicates that this anomaly stems from [Specific Structural/External Factor] rather than random variance. | 句首锚定旧信息/逻辑视角，主干动词紧跟主语，复杂修饰置于宾语之后（符合尾重原则 L011 与旧新信息流 L001）。}
  \item \textbf{S4 [方法论启示 Methodological Implication]:} \textit{This anomaly confirms that model capacity must be matched to the intrinsic signal-to-noise ratio of the underlying stochastic process.}\par\noindent{\small 对应步骤：S4 [方法论启示 Methodological Implication]：Accounting for such localized departures is critical when calibrating [Downstream Model/Policy]. | 句首锚定旧信息/逻辑视角，主干动词紧跟主语，复杂修饰置于宾语之后（符合尾重原则 L011 与旧新信息流 L001）。}
\end{itemize}
\end{TipBox}
\medskip
\clearpage
\section{FT17: AI 系统评测报告全文：长序列预测与生成模型在历史实测与未来算力扩展下的精度-延迟帕累托报告}\label{comp:FT17}
\begin{ConceptBox}{Composition Architecture \& Structural Blueprint Chain}
\textbf{Genre:} \texttt{Technical\_Benchmark\_Report} \quad \textbf{Paragraph Count:} 5\par
\textbf{Blueprint Chain:} \texttt{BP17 $\to$ BP18 $\to$ BP22 $\to$ BP28 $\to$ BP30}
\end{ConceptBox}
\subsection*{Paragraph 1: 实验设置、数据集特性与强基线选型说明段（Experimental Setup \& Baselines） (\texttt{BP17})}
\noindent\textbf{Expressive Intent:}\enspace 我这一段想在 Experiments 开篇说明：我们为什么选这几类数据集和这几类基线模型，证明我们的实验设计是全面、公平且直击痛点的。\par
\noindent\textbf{Cohesion Formula:}\enspace {\small\color{AcadSlate}Dataset Diversity $\to$ Challenging Regime Highlights $\to$ Categorized Strong Baselines $\to$ Strict Fairness Protocol. 绑定逻辑规则：L009, L010.}\par\vspace{4pt}
\begin{ExampleBox}{Publication-Grade Authentic Text}
To evaluate robustness across diverse non-stationary regimes, we conduct experiments on seven public benchmarks spanning energy grid load (ETT, ECL), meteorological telemetry (Weather), and urban traffic flow (Traffic). Notably, Weather contains abrupt atmospheric pressure transitions that test non-stationary adaptability, whereas Traffic spans 862 synchronized sensors that challenge high-dimensional scalability. We benchmark our approach against 11 state-of-the-art methods across three paradigms: channel-independent patch Transformers (PatchTST), variate-mixing architectures (iTransformer, Crossformer), and frequency-domain forecasters (TimesNet, FEDformer). All baselines are reproduced under identical lookback windows ($L=96$) and prediction horizons ($H \in \{96, 192, 336, 720\}$) using standard MSE and MAE metrics.
\end{ExampleBox}
\begin{TipBox}{Rhetorical Move Breakdown \& Constituent Analysis}
\begin{itemize}[leftmargin=1.5em]
  \item \textbf{S1 [数据集选型逻辑 Dataset Selection Rationale]:} \textit{To evaluate robustness across diverse non-stationary regimes, we conduct experiments on seven public benchmarks spanning energy grid load (ETT, ECL), meteorological telemetry (Weather), and urban traffic flow (Traffic).}\par\noindent{\small 对应步骤：S1 [数据集选型逻辑 Dataset Selection Rationale]：To evaluate performance across diverse non-stationary regimes, we conduct experiments on [N] benchmarks spanning [Domain 1], [Domain 2], and [Domain 3]. | 句首锚定旧信息/逻辑视角，主干动词紧跟主语，复杂修饰置于宾语之后（符合尾重原则 L011 与旧新信息流 L001）。}
  \item \textbf{S2 [极端/困难特性突出 Challenging Dataset Properties]:} \textit{Notably, Weather contains abrupt atmospheric pressure transitions that test non-stationary adaptability, whereas Traffic spans 862 synchronized sensors that challenge high-dimensional scalability.}\par\noindent{\small 对应步骤：S2 [极端/困难特性突出 Challenging Dataset Properties]：Notably, [Dataset A] exhibits [Hard Property 1, e.g., strong covariate shift], whereas [Dataset B] tests scalability across [Hard Property 2, e.g., 862 correlated channels]. | 句首锚定旧信息/逻辑视角，主干动词紧跟主语，复杂修饰置于宾语之后（符合尾重原则 L011 与旧新信息流 L001）。}
  \item \textbf{S3 [基线流派全覆盖 Baseline Family Coverage]:} \textit{We benchmark our approach against 11 state-of-the-art methods across three paradigms: channel-independent patch Transformers (PatchTST), variate-mixing architectures (iTransformer, Crossformer), and frequency-domain forecasters (TimesNet, FEDformer).}\par\noindent{\small 对应步骤：S3 [基线流派全覆盖 Baseline Family Coverage]：We compare our method against [K] competitive baselines grouped into three families: [Family 1], [Family 2], and [Family 3]. | 句首锚定旧信息/逻辑视角，主干动词紧跟主语，复杂修饰置于宾语之后（符合尾重原则 L011 与旧新信息流 L001）。}
  \item \textbf{S4 [公平评测协议 Fair Evaluation Protocol]:} \textit{All baselines are reproduced under identical lookback windows ($L=96$) and prediction horizons ($H \in \{96, 192, 336, 720\}$) using standard MSE and MAE metrics.}\par\noindent{\small 对应步骤：S4 [公平评测协议 Fair Evaluation Protocol]：All models are evaluated under identical lookback windows ($L=96$), prediction horizons ($H \in \{96, 192, 336, 720\}$), and hyperparameter search budgets. | 句首锚定旧信息/逻辑视角，主干动词紧跟主语，复杂修饰置于宾语之后（符合尾重原则 L011 与旧新信息流 L001）。}
\end{itemize}
\end{TipBox}
\medskip
\subsection*{Paragraph 2: 主实验结果整体对比与核心胜出原因剖析段（Main Results Analysis） (\texttt{BP18})}
\noindent\textbf{Expressive Intent:}\enspace 我这一段想解读主实验大表（Table 1）：既报告全局胜率，又挑出最有代表性的对比案例，解释为什么竞争对手在那里跌倒而我们稳赢。\par
\noindent\textbf{Cohesion Formula:}\enspace {\small\color{AcadSlate}Global Win Rate $\to$ Advantage over Family A $\to$ Resilience at Long Horizon vs. Family B $\to$ Mechanistic Takeaway. 绑定逻辑规则：L015, L013, L010.}\par\vspace{4pt}
\begin{ExampleBox}{Publication-Grade Authentic Text}
Table 1 reports the multivariate forecasting results, where DyTS-Net secures the lowest error in 25 out of 28 dataset-horizon configurations. Compared with channel-independent PatchTST, DyTS-Net reduces average MSE by 14.2\% on the 862-channel Traffic benchmark, where modeling spatial-temporal congestion propagation across adjacent road sensors is critical. Furthermore, while variate-mixing iTransformer performs competitively at $H=96$, its MSE escalates by 68\% at $H=720$ on volatile Weather series as unregularized cross-channel attention overfits transient wind gusts. By decoupling stationary phase coupling from neural-ODE scale projection, DyTS-Net suppresses long-horizon drift and maintains an 11.6\% margin over the strongest baseline at $H=720$.
\end{ExampleBox}
\begin{TipBox}{Rhetorical Move Breakdown \& Constituent Analysis}
\begin{itemize}[leftmargin=1.5em]
  \item \textbf{S1 [全局统治力概述 Overall Dominance]:} \textit{Table 1 reports the multivariate forecasting results, where DyTS-Net secures the lowest error in 25 out of 28 dataset-horizon configurations.}\par\noindent{\small 对应步骤：S1 [全局统治力概述 Overall Dominance]：Table 1 summarizes the multivariate forecasting results, where [Our Model] achieves the lowest error in [X] out of [Y] dataset-horizon configurations. | 句首锚定旧信息/逻辑视角，主干动词紧跟主语，复杂修饰置于宾语之后（符合尾重原则 L011 与旧新信息流 L001）。}
  \item \textbf{S2 [对比流派 1 的差距与原因 Contrast vs. Family 1]:} \textit{Compared with channel-independent PatchTST, DyTS-Net reduces average MSE by 14.2\% on the 862-channel Traffic benchmark, where modeling spatial-temporal congestion propagation across adjacent road sensors is critical.}\par\noindent{\small 对应步骤：S2 [对比流派 1 的差距与原因 Contrast vs. Family 1]：Compared with channel-independent baselines such as [Baseline A], our method reduces average MSE by [X\%] on [Dataset], where strong cross-sensor coupling is critical. | 句首锚定旧信息/逻辑视角，主干动词紧跟主语，复杂修饰置于宾语之后（符合尾重原则 L011 与旧新信息流 L001）。}
  \item \textbf{S3 [对比流派 2 在困难场景下的崩溃 Contrast vs. Family 2 under Hard Regime]:} \textit{Furthermore, while variate-mixing iTransformer performs competitively at $H=96$, its MSE escalates by 68\% at $H=720$ on volatile Weather series as unregularized cross-channel attention overfits transient wind gusts.}\par\noindent{\small 对应步骤：S3 [对比流派 2 在困难场景下的崩溃 Contrast vs. Family 2 under Hard Regime]：Furthermore, while [Baseline B] performs competitively at short horizons ($H=96$), its error escalates by [Y\%] at $H=720$ due to [Failure Mechanism]. | 句首锚定旧信息/逻辑视角，主干动词紧跟主语，复杂修饰置于宾语之后（符合尾重原则 L011 与旧新信息流 L001）。}
  \item \textbf{S4 [回扣本文核心机制 Mechanistic Confirmation]:} \textit{By decoupling stationary phase coupling from neural-ODE scale projection, DyTS-Net suppresses long-horizon drift and maintains an 11.6\% margin over the strongest baseline at $H=720$.}\par\noindent{\small 对应步骤：S4 [回扣本文核心机制 Mechanistic Confirmation]：By contrast, our [Core Module] bounds long-range drift, maintaining a [Z\%] margin over the second-best baseline at $H=720$. | 句首锚定旧信息/逻辑视角，主干动词紧跟主语，复杂修饰置于宾语之后（符合尾重原则 L011 与旧新信息流 L001）。}
\end{itemize}
\end{TipBox}
\medskip
\subsection*{Paragraph 3: 计算开销、推理延迟与精度帕累托权衡段（Efficiency-Accuracy Pareto Analysis） (\texttt{BP22})}
\noindent\textbf{Expressive Intent:}\enspace 我这一段想说明我们的模型不仅准，而且算得快、显存省，在“精度-延迟/显存”帕累托前沿（Pareto Frontier）上显著优于现有大模型。\par
\noindent\textbf{Cohesion Formula:}\enspace {\small\color{AcadSlate}Hardware Benchmark $\to$ Complexity Reduction \& Memory Gain $\to$ Real-Time Latency Margin $\to$ Pareto Conclusion. 绑定逻辑规则：L010, L023.}\par\vspace{4pt}
\begin{ExampleBox}{Publication-Grade Authentic Text}
Figure 6 plots the MSE-latency Pareto frontier measured on an onboard NVIDIA Orin AGX module. By replacing $O(L^2)$ point-wise self-attention with $O(L \log L)$ spectral filtering and compact patch tokens, DyTS-Net reduces peak GPU memory by 62\% and training wall-clock time by 3.4x compared with iTransformer. At inference time, generating a 16-step visuomotor trajectory requires only 6.8 ms---well below the 20 ms cycle budget of a 50 Hz industrial robot controller---while outperforming a 200 ms DDPM baseline by 11.4\% in task success rate.
\end{ExampleBox}
\begin{TipBox}{Rhetorical Move Breakdown \& Constituent Analysis}
\begin{itemize}[leftmargin=1.5em]
  \item \textbf{S1 [硬件与测试基准 Hardware Setup]:} \textit{Figure 6 plots the MSE-latency Pareto frontier measured on an onboard NVIDIA Orin AGX module.}\par\noindent{\small 对应步骤：S1 [硬件与测试基准 Hardware Setup]：Figure 6 compares the accuracy-efficiency Pareto frontier across all baselines measured on a single NVIDIA RTX 4090 GPU. | 句首锚定旧信息/逻辑视角，主干动词紧跟主语，复杂修饰置于宾语之后（符合尾重原则 L011 与旧新信息流 L001）。}
  \item \textbf{S2 [复杂度理论与实测对比 Complexity \& Runtime]:} \textit{By replacing $O(L^2)$ point-wise self-attention with $O(L \log L)$ spectral filtering and compact patch tokens, DyTS-Net reduces peak GPU memory by 62\% and training wall-clock time by 3.4x compared with iTransformer.}\par\noindent{\small 对应步骤：S2 [复杂度理论与实测对比 Complexity \& Runtime]：By replacing $O(L^2)$ point-wise attention with $O(L \log L)$ spectral filtering and patch aggregation, our model cuts training time per epoch by [X\%] and peak VRAM by [Y\%]. | 句首锚定旧信息/逻辑视角，主干动词紧跟主语，复杂修饰置于宾语之后（符合尾重原则 L011 与旧新信息流 L001）。}
  \item \textbf{S3 [实时闭环部署可行性 Real-Time Viability]:} \textit{At inference time, generating a 16-step visuomotor trajectory requires only 6.8 ms---well below the 20 ms cycle budget of a 50 Hz industrial robot controller---while outperforming a 200 ms DDPM baseline by 11.4\% in task success rate.}\par\noindent{\small 对应步骤：S3 [实时闭环部署可行性 Real-Time Viability]：At inference time, generating a 720-step forecast (or 16-step robot action chunk) requires only [4.2 ms], well below the [20 ms] real-time control deadline. | 句首锚定旧信息/逻辑视角，主干动词紧跟主语，复杂修饰置于宾语之后（符合尾重原则 L011 与旧新信息流 L001）。}
\end{itemize}
\end{TipBox}
\medskip
\subsection*{Paragraph 4: 双图关联与互补/因果印证段（Synthesizing Dual Charts） (\texttt{BP28})}
\noindent\textbf{Expressive Intent:}\enspace 我这一段想写双图题（如一张图是算力投入/研发经费变化，另一张图是模型误差/碳排放强度变化）的综合关联段，阐明两张图之间的内在协同或背离关系。\par
\noindent\textbf{Cohesion Formula:}\enspace {\small\color{AcadSlate}Global Correlation $\to$ Co-movement Phase $\to$ Diminishing Returns / Divergence $\to$ Notable Outlier. 绑定逻辑规则：L007, L013.}\par\vspace{4pt}
\begin{ExampleBox}{Publication-Grade Authentic Text}
Cross-referencing both figures reveals a strong inverse relationship between automated inspection investment and manufacturing defect rates up to a \$40 million threshold. Specifically, across plants investing between \$10 million and \$40 million (Figure 1), surface defect incidence dropped steeply from 4.8% to 0.9% (Figure 2). Beyond \$40 million, however, further capital expenditure yielded diminishing returns: doubling investment to \$80 million reduced defects by only an additional 0.1 percentage points. Notably, Plant D defied this expenditure curve, achieving a 0.7% defect rate with just \$22 million through closed-loop sensor calibration.
\end{ExampleBox}
\begin{TipBox}{Rhetorical Move Breakdown \& Constituent Analysis}
\begin{itemize}[leftmargin=1.5em]
  \item \textbf{S1 [双图联动锚定 Cross-Chart Link]:} \textit{Cross-referencing both figures reveals a strong inverse relationship between automated inspection investment and manufacturing defect rates up to a \$40 million threshold.}\par\noindent{\small 对应步骤：S1 [双图联动锚定 Cross-Chart Link]：Synthesizing both charts reveals a clear inverse / positive relationship between [Chart 1 Variable] and [Chart 2 Variable] up to [Threshold Point]. | 句首锚定旧信息/逻辑视角，主干动词紧跟主语，复杂修饰置于宾语之后（符合尾重原则 L011 与旧新信息流 L001）。}
  \item \textbf{S2 [区间 1 协同变化 Phase 1 Co-movement]:} \textit{Specifically, across plants investing between \$10 million and \$40 million (Figure 1), surface defect incidence dropped steeply from 4.8\% to 0.9\% (Figure 2).}\par\noindent{\small 对应步骤：S2 [区间 1 协同变化 Phase 1 Co-movement]：Specifically, as [Variable X in Chart 1] expanded from [A] to [B] during [Period 1], [Variable Y in Chart 2] dropped correspondingly by [C\%]. | 句首锚定旧信息/逻辑视角，主干动词紧跟主语，复杂修饰置于宾语之后（符合尾重原则 L011 与旧新信息流 L001）。}
  \item \textbf{S3 [区间 2 边际递减/背离 Diminishing Returns / Divergence]:} \textit{Beyond \$40 million, however, further capital expenditure yielded diminishing returns: doubling investment to \$80 million reduced defects by only an additional 0.1 percentage points.}\par\noindent{\small 对应步骤：S3 [区间 2 边际递减/背离 Diminishing Returns / Divergence]：However, beyond [Threshold Point], although [Variable X] continued to double, [Variable Y] plateaued around [Value], indicating diminishing marginal returns. | 句首锚定旧信息/逻辑视角，主干动词紧跟主语，复杂修饰置于宾语之后（符合尾重原则 L011 与旧新信息流 L001）。}
  \item \textbf{S4 [例外个体点拨 Exception Highlight]:} \textit{Notably, Plant D defied this expenditure curve, achieving a 0.7\% defect rate with just \$22 million through closed-loop sensor calibration.}\par\noindent{\small 对应步骤：S4 [例外个体点拨 Exception Highlight]：An exception to this pattern is [Entity Z], which achieved the lowest [Variable Y] despite maintaining moderate [Variable X]. | 句首锚定旧信息/逻辑视角，主干动词紧跟主语，复杂修饰置于宾语之后（符合尾重原则 L011 与旧新信息流 L001）。}
\end{itemize}
\end{TipBox}
\medskip
\subsection*{Paragraph 5: 历史实测与未来预测跨时态衔接段（Historical vs. Projected Trends） (\texttt{BP30})}
\noindent\textbf{Expressive Intent:}\enspace 我这一段想处理同时包含过去历史数据（如 2010--2025）和未来预测数据（如 2026--2040）的图表段落，严格区分过去实测时态与未来审慎预测表达。\par
\noindent\textbf{Cohesion Formula:}\enspace {\small\color{AcadSlate}Historical Past (rose/accounted for) $\to$ Pivot Year $\to$ Future Projection Verbs (is projected/anticipated/expected to) $\to$ Moderating Factor. 绑定逻辑规则：L010, L017.}\par\vspace{4pt}
\begin{ExampleBox}{Publication-Grade Authentic Text}
Between 2010 and 2025, global data-center electricity consumption more than doubled from 200 TWh to 460 TWh. From 2026 onward, total demand is projected to climb at a substantially steeper rate, with baseline forecasts pointing to 1,050 TWh by 2035. Within this expansion, AI model training and inference---which represented under 10\% of server load in 2023---are anticipated to account for nearly half (46\%) of total consumption by 2040. However, widespread deployment of sparsity-aware accelerators is expected to temper the annual growth rate after 2035.
\end{ExampleBox}
\begin{TipBox}{Rhetorical Move Breakdown \& Constituent Analysis}
\begin{itemize}[leftmargin=1.5em]
  \item \textbf{S1 [历史实测轨迹 Historical Trajectory (Past Tense)]:} \textit{Between 2010 and 2025, global data-center electricity consumption more than doubled from 200 TWh to 460 TWh.}\par\noindent{\small 对应步骤：S1 [历史实测轨迹 Historical Trajectory (Past Tense)]：Between [2010] and [2025], global data-center electricity demand rose steadily from [200 TWh] to [460 TWh], driven by cloud migration. | 句首锚定旧信息/逻辑视角，主干动词紧跟主语，复杂修饰置于宾语之后（符合尾重原则 L011 与旧新信息流 L001）。}
  \item \textbf{S2 [当下转折与未来加速 Projected Acceleration (Projection Verbs)]:} \textit{From 2026 onward, total demand is projected to climb at a substantially steeper rate, with baseline forecasts pointing to 1,050 TWh by 2035.}\par\noindent{\small 对应步骤：S2 [当下转折与未来加速 Projected Acceleration (Projection Verbs)]：From [2026] onward, consumption is projected to accelerate sharply, with forecasts indicating a surge to [1,050 TWh] by [2035]. | 句首锚定旧信息/逻辑视角，主干动词紧跟主语，复杂修饰置于宾语之后（符合尾重原则 L011 与旧新信息流 L001）。}
  \item \textbf{S3 [结构占比重塑 Structural Shift in Projections]:} \textit{Within this expansion, AI model training and inference---which represented under 10\% of server load in 2023---are anticipated to account for nearly half (46\%) of total consumption by 2040.}\par\noindent{\small 对应步骤：S3 [结构占比重塑 Structural Shift in Projections]：Within this projected expansion, generative AI training and inference---which accounted for less than [10\%] in [2023]---are anticipated to represent over [45\%] of total load by [2040]. | 句首锚定旧信息/逻辑视角，主干动词紧跟主语，复杂修饰置于宾语之后（符合尾重原则 L011 与旧新信息流 L001）。}
  \item \textbf{S4 [条件式审慎收尾 Conditional Hedging]:} \textit{However, widespread deployment of sparsity-aware accelerators is expected to temper the annual growth rate after 2035.}\par\noindent{\small 对应步骤：S4 [条件式审慎收尾 Conditional Hedging]：Nevertheless, energy-efficient hardware adoption is expected to moderate the growth slope during the final five years. | 句首锚定旧信息/逻辑视角，主干动词紧跟主语，复杂修饰置于宾语之后（符合尾重原则 L011 与旧新信息流 L001）。}
\end{itemize}
\end{TipBox}
\medskip
\clearpage
\section{FT18: 科研实验反常现象与极端工况诊断报告全文：为何高容量注意力模型在强随机游走序列上退化及改进方案}\label{comp:FT18}
\begin{ConceptBox}{Composition Architecture \& Structural Blueprint Chain}
\textbf{Genre:} \texttt{Failure\_Mode\_Report} \quad \textbf{Paragraph Count:} 5\par
\textbf{Blueprint Chain:} \texttt{BP32 $\to$ BP19 $\to$ BP20 $\to$ BP21 $\to$ BP23}
\end{ConceptBox}
\subsection*{Paragraph 1: 异常数据点（Outliers）与局部背离现象精准批注段（Outlier \& Anomaly Annotation） (\texttt{BP32})}
\noindent\textbf{Expressive Intent:}\enspace 我这一段想在图表或实验结果段中，专门挑出违背总体规律的 1--2 个异常点（Outliers）或局部反常尖峰进行精准描述与成因剖析。\par
\noindent\textbf{Cohesion Formula:}\enspace {\small\color{AcadSlate}General Rule vs. Notable Departure $\to$ Quantitative Deviation $\to$ Root Cause Attribution $\to$ Calibration Takeaway. 绑定逻辑规则：L007, L013.}\par\vspace{4pt}
\begin{ExampleBox}{Publication-Grade Authentic Text}
Although deep spectral transformers outperform linear baselines across six of the seven benchmarks, a notable departure emerges on the financial Exchange-Rate dataset. Whereas complex attention models reduce MSE by 12\%--18\% on periodic energy and weather streams, they trail a single-layer DLinear baseline by 4.3\% on Exchange-Rate. Autocorrelation analysis reveals that this reversal stems from the near-martingale nature of daily currency returns, where weak seasonal periodicity causes high-capacity attention layers to overfit transient market shocks. This anomaly confirms that model capacity must be matched to the intrinsic signal-to-noise ratio of the underlying stochastic process.
\end{ExampleBox}
\begin{TipBox}{Rhetorical Move Breakdown \& Constituent Analysis}
\begin{itemize}[leftmargin=1.5em]
  \item \textbf{S1 [总体规律背景下的特例锚定 Outlier Identification]:} \textit{Although deep spectral transformers outperform linear baselines across six of the seven benchmarks, a notable departure emerges on the financial Exchange-Rate dataset.}\par\noindent{\small 对应步骤：S1 [总体规律背景下的特例锚定 Outlier Identification]：Although [General Trend holds across most groups/steps], a notable departure from this pattern occurs in [Specific Point/Group]. | 句首锚定旧信息/逻辑视角，主干动词紧跟主语，复杂修饰置于宾语之后（符合尾重原则 L011 与旧新信息流 L001）。}
  \item \textbf{S2 [反常幅度量化 Quantitative Deviation]:} \textit{Whereas complex attention models reduce MSE by 12\%--18\% on periodic energy and weather streams, they trail a single-layer DLinear baseline by 4.3\% on Exchange-Rate.}\par\noindent{\small 对应步骤：S2 [反常幅度量化 Quantitative Deviation]：Whereas peer categories cluster tightly within $[A, B]$, [Outlier Entity] spikes/drops to [Extreme Value], deviating by more than [X\%] from the group median. | 句首锚定旧信息/逻辑视角，主干动词紧跟主语，复杂修饰置于宾语之后（符合尾重原则 L011 与旧新信息流 L001）。}
  \item \textbf{S3 [底层物理/统计成因剖析 Root Cause Analysis]:} \textit{Autocorrelation analysis reveals that this reversal stems from the near-martingale nature of daily currency returns, where weak seasonal periodicity causes high-capacity attention layers to overfit transient market shocks.}\par\noindent{\small 对应步骤：S3 [底层物理/统计成因剖析 Root Cause Analysis]：Closer inspection indicates that this anomaly stems from [Specific Structural/External Factor] rather than random variance. | 句首锚定旧信息/逻辑视角，主干动词紧跟主语，复杂修饰置于宾语之后（符合尾重原则 L011 与旧新信息流 L001）。}
  \item \textbf{S4 [方法论启示 Methodological Implication]:} \textit{This anomaly confirms that model capacity must be matched to the intrinsic signal-to-noise ratio of the underlying stochastic process.}\par\noindent{\small 对应步骤：S4 [方法论启示 Methodological Implication]：Accounting for such localized departures is critical when calibrating [Downstream Model/Policy]. | 句首锚定旧信息/逻辑视角，主干动词紧跟主语，复杂修饰置于宾语之后（符合尾重原则 L011 与旧新信息流 L001）。}
\end{itemize}
\end{TipBox}
\medskip
\subsection*{Paragraph 2: 消融实验逐组件归因论证段（Counterfactual Ablation Study） (\texttt{BP19})}
\noindent\textbf{Expressive Intent:}\enspace 我这一段想汇报消融实验（Table 2）：证明我们的模块 A、模块 B 和正则项 C 每一个都不可或缺，且各自负责解决不同的特定子问题。\par
\noindent\textbf{Cohesion Formula:}\enspace {\small\color{AcadSlate}Ablation Overview $\to$ w/o Module A (Specific Failure) $\to$ w/o Module B (Specific Failure) $\to$ Cross-Module Synergy. 绑定逻辑规则：L016, L007.}\par\vspace{4pt}
\begin{ExampleBox}{Publication-Grade Authentic Text}
To isolate the contribution of each design choice, Table 2 compares full DyTS-Net against three ablated variants across ETTm2 and ECL. Removing the neural-ODE scale tracker (w/o Scale-ODE) triggers the sharpest degradation during seasonal transition months on ETTm2, raising 720-step MSE by 19.4\% as the decoder fails to track shifting baseline temperatures. Conversely, replacing phase-aligned graph attention with standard dot-product attention (w/o Phase-Attn) degrades performance most severely on the 321-channel ECL grid (+13.8\% MSE), where unaligned high-frequency load spikes corrupt cross-substation aggregation. Jointly deploying both modules yields complementary gains, as accurate scale tracking prevents amplitude distortion from biasing phase estimation.
\end{ExampleBox}
\begin{TipBox}{Rhetorical Move Breakdown \& Constituent Analysis}
\begin{itemize}[leftmargin=1.5em]
  \item \textbf{S1 [消融实验总括 Ablation Setup]:} \textit{To isolate the contribution of each design choice, Table 2 compares full DyTS-Net against three ablated variants across ETTm2 and ECL.}\par\noindent{\small 对应步骤：S1 [消融实验总括 Ablation Setup]：To isolate the contribution of each architectural component, we evaluate four ablated variants in Table 2 by progressively removing [Mod A], [Mod B], and [Loss C]. | 句首锚定旧信息/逻辑视角，主干动词紧跟主语，复杂修饰置于宾语之后（符合尾重原则 L011 与旧新信息流 L001）。}
  \item \textbf{S2 [移除模块 A 的冲击与归因 Impact of Removing Mod A]:} \textit{Removing the neural-ODE scale tracker (w/o Scale-ODE) triggers the sharpest degradation during seasonal transition months on ETTm2, raising 720-step MSE by 19.4\% as the decoder fails to track shifting baseline temperatures.}\par\noindent{\small 对应步骤：S2 [移除模块 A 的冲击与归因 Impact of Removing Mod A]：Removing [Mod A] (w/o A) causes the largest degradation on [Non-stationary/Volatile Subset], increasing MSE by [X\%] because the model loses [Specific Capability A]. | 句首锚定旧信息/逻辑视角，主干动词紧跟主语，复杂修饰置于宾语之后（符合尾重原则 L011 与旧新信息流 L001）。}
  \item \textbf{S3 [移除模块 B 的冲击与归因 Impact of Removing Mod B]:} \textit{Conversely, replacing phase-aligned graph attention with standard dot-product attention (w/o Phase-Attn) degrades performance most severely on the 321-channel ECL grid (+13.8\% MSE), where unaligned high-frequency load spikes corrupt cross-substation aggregation.}\par\noindent{\small 对应步骤：S3 [移除模块 B 的冲击与归因 Impact of Removing Mod B]：Similarly, replacing [Mod B] with standard [Naive Operator] degrades performance by [Y\%] on high-dimensional benchmarks, confirming its role in [Specific Capability B]. | 句首锚定旧信息/逻辑视角，主干动词紧跟主语，复杂修饰置于宾语之后（符合尾重原则 L011 与旧新信息流 L001）。}
  \item \textbf{S4 [协同效应总结 Synergistic Conclusion]:} \textit{Jointly deploying both modules yields complementary gains, as accurate scale tracking prevents amplitude distortion from biasing phase estimation.}\par\noindent{\small 对应步骤：S4 [协同效应总结 Synergistic Conclusion]：Combining all components yields super-additive gains, as [Mod A] stabilizes input statistics for [Mod B]'s relational reasoning. | 句首锚定旧信息/逻辑视角，主干动词紧跟主语，复杂修饰置于宾语之后（符合尾重原则 L011 与旧新信息流 L001）。}
\end{itemize}
\end{TipBox}
\medskip
\subsection*{Paragraph 3: 超参数敏感性与极端噪声鲁棒性分析段（Sensitivity \& Robustness Analysis） (\texttt{BP20})}
\noindent\textbf{Expressive Intent:}\enspace 我这一段想回应审稿人最常问的：你的模型是不是对超参数很敏感？在缺失值或传感器强噪声下会不会突然崩盘？\par
\noindent\textbf{Cohesion Formula:}\enspace {\small\color{AcadSlate}Goal $\to$ Wide Hyperparameter Plateau $\to$ Stress-Test Comparison $\to$ Mechanistic Explanation. 绑定逻辑规则：L010, L019.}\par\vspace{4pt}
\begin{ExampleBox}{Publication-Grade Authentic Text}
Figure 4 evaluates hyperparameter sensitivity and resilience against severe observation corruption. Varying the phase regularization weight $\lambda$ across $[0.05, 0.50]$ alters validation MSE by less than 1.4\%, indicating a wide, tuning-friendly operating plateau around $\lambda=0.15$. Under stress tests where 30\% of sensor readings are randomly masked or corrupted by impulse noise ($\text{SNR} = 5\text{ dB}$), point-wise Transformer baselines suffer a 34.2\% error surge, whereas DyTS-Net degrades by only 4.8\%. This robustness stems from our spectral energy gate, which automatically attenuates high-entropy corrupted channels prior to cross-variate message passing.
\end{ExampleBox}
\begin{TipBox}{Rhetorical Move Breakdown \& Constituent Analysis}
\begin{itemize}[leftmargin=1.5em]
  \item \textbf{S1 [分析目的 Objective]:} \textit{Figure 4 evaluates hyperparameter sensitivity and resilience against severe observation corruption.}\par\noindent{\small 对应步骤：S1 [分析目的 Objective]：We further examine sensitivity to key hyperparameters and resilience against severe sensor corruption in Figure 4. | 句首锚定旧信息/逻辑视角，主干动词紧跟主语，复杂修饰置于宾语之后（符合尾重原则 L011 与旧新信息流 L001）。}
  \item \textbf{S2 [超参数平坦区间 Plateau Stability]:} \textit{Varying the phase regularization weight $\lambda$ across $[0.05, 0.50]$ alters validation MSE by less than 1.4\%, indicating a wide, tuning-friendly operating plateau around $\lambda=0.15$.}\par\noindent{\small 对应步骤：S2 [超参数平坦区间 Plateau Stability]：Varying [Hyperparameter $\lambda$] across two orders of magnitude ($[0.01, 1.0]$) alters validation MSE by less than [1.8\%], demonstrating a wide stable operating plateau. | 句首锚定旧信息/逻辑视角，主干动词紧跟主语，复杂修饰置于宾语之后（符合尾重原则 L011 与旧新信息流 L001）。}
  \item \textbf{S3 [极端边缘工况压力测试 Stress Test under Corruption]:} \textit{Under stress tests where 30\% of sensor readings are randomly masked or corrupted by impulse noise ($\text{SNR} = 5\text{ dB}$), point-wise Transformer baselines suffer a 34.2\% error surge, whereas DyTS-Net degrades by only 4.8\%.}\par\noindent{\small 对应步骤：S3 [极端边缘工况压力测试 Stress Test under Corruption]：Under synthetic sensor dropouts up to [40\%] and Gaussian noise ($\sigma=0.3$), baseline accuracy plummets by over [35\%], whereas our model degrades by less than [6.5\%]. | 句首锚定旧信息/逻辑视角，主干动词紧跟主语，复杂修饰置于宾语之后（符合尾重原则 L011 与旧新信息流 L001）。}
  \item \textbf{S4 [鲁棒性物理归因 Robustness Mechanism]:} \textit{This robustness stems from our spectral energy gate, which automatically attenuates high-entropy corrupted channels prior to cross-variate message passing.}\par\noindent{\small 对应步骤：S4 [鲁棒性物理归因 Robustness Mechanism]：This resilience stems from [Specific Mechanism], which filters out corrupted channels before global aggregation. | 句首锚定旧信息/逻辑视角，主干动词紧跟主语，复杂修饰置于宾语之后（符合尾重原则 L011 与旧新信息流 L001）。}
\end{itemize}
\end{TipBox}
\medskip
\subsection*{Paragraph 4: 极端/反常情况与失败案例坦诚剖析段（Failure Case \& Outlier Analysis） (\texttt{BP21})}
\noindent\textbf{Expressive Intent:}\enspace 我这一段想通过可视化案例（Case Study / Qualitative Analysis），既展示模型在复杂跳变下的精准跟踪，又坦诚剖析一个模型预测失准的极端案例。\par
\noindent\textbf{Cohesion Formula:}\enspace {\small\color{AcadSlate}Visual Setup $\to$ Success under Shift $\to$ Honest Failure Case under Unobserved Exogenous Shock $\to$ Actionable Insight. 绑定逻辑规则：L017, L018.}\par\vspace{4pt}
\begin{ExampleBox}{Publication-Grade Authentic Text}
Figure 5 visualizes 336-step forecasts during an abrupt voltage drop on the ETTm1 transformer load benchmark. Whereas PatchTST and iTransformer exhibit a 14-step phase lag and severe post-recovery overshoot, DyTS-Net tracks both the sharp downward ramp and subsequent stabilization within 5\% amplitude error. Nevertheless, the failure case in Panel (d) reveals that when an unscheduled manual grid curtailment occurs without preceding spectral precursors, our endogenous model experiences a 6-step delayed response. Anticipating such zero-precursor manual interventions requires integrating exogenous dispatch logs alongside historical telemetry.
\end{ExampleBox}
\begin{TipBox}{Rhetorical Move Breakdown \& Constituent Analysis}
\begin{itemize}[leftmargin=1.5em]
  \item \textbf{S1 [定性对比切入 Qualitative Setup]:} \textit{Figure 5 visualizes 336-step forecasts during an abrupt voltage drop on the ETTm1 transformer load benchmark.}\par\noindent{\small 对应步骤：S1 [定性对比切入 Qualitative Setup]：Figure 5 visualizes predicted trajectories during an abrupt regime shift on [Dataset/Task]. | 句首锚定旧信息/逻辑视角，主干动词紧跟主语，复杂修饰置于宾语之后（符合尾重原则 L011 与旧新信息流 L001）。}
  \item \textbf{S2 [成功案例亮点 Success Highlight]:} \textit{Whereas PatchTST and iTransformer exhibit a 14-step phase lag and severe post-recovery overshoot, DyTS-Net tracks both the sharp downward ramp and subsequent stabilization within 5\% amplitude error.}\par\noindent{\small 对应步骤：S2 [成功案例亮点 Success Highlight]：While [Baseline] lags behind the sudden step change by [K steps] and overshoots, our model rapidly locks onto the new amplitude regime within [M steps]. | 句首锚定旧信息/逻辑视角，主干动词紧跟主语，复杂修饰置于宾语之后（符合尾重原则 L011 与旧新信息流 L001）。}
  \item \textbf{S3 [失败案例坦诚剖析 Honest Failure Mode]:} \textit{Nevertheless, the failure case in Panel (d) reveals that when an unscheduled manual grid curtailment occurs without preceding spectral precursors, our endogenous model experiences a 6-step delayed response.}\par\noindent{\small 对应步骤：S3 [失败案例坦诚剖析 Honest Failure Mode]：Nevertheless, inspecting the bottom-right failure case reveals that when an external manual override occurs without historical precursors, all purely endogenous forecasters---including ours---exhibit a transient lag. | 句首锚定旧信息/逻辑视角，主干动词紧跟主语，复杂修饰置于宾语之后（符合尾重原则 L011 与旧新信息流 L001）。}
  \item \textbf{S4 [改进启示 Constructive Takeaway]:} \textit{Anticipating such zero-precursor manual interventions requires integrating exogenous dispatch logs alongside historical telemetry.}\par\noindent{\small 对应步骤：S4 [改进启示 Constructive Takeaway]：Incorporating exogenous event logs or textual maintenance records is necessary to anticipate such unscheduled interventions. | 句首锚定旧信息/逻辑视角，主干动词紧跟主语，复杂修饰置于宾语之后（符合尾重原则 L011 与旧新信息流 L001）。}
\end{itemize}
\end{TipBox}
\medskip
\subsection*{Paragraph 5: 局限性（Limitations）客观陈述与未来扩展段（Constructive Limitations） (\texttt{BP23})}
\noindent\textbf{Expressive Intent:}\enspace 我这一段想写论文末尾或 NeurIPS/ICLR Checklist 要求的 Limitations 段落：诚实指出当前方法的两个客观边界，但绝不贬损论文的核心贡献。\par
\noindent\textbf{Cohesion Formula:}\enspace {\small\color{AcadSlate}Two Objective Boundaries $\to$ [Boundary 1 + Technical Remedy] $\to$ [Boundary 2 + Stochastic Extension]. 绑定逻辑规则：L017.}\par\vspace{4pt}
\begin{ExampleBox}{Publication-Grade Authentic Text}
Despite its strong empirical performance, our framework exhibits two theoretical boundaries worth noting. First, while our spectral branch scales quasi-linearly with sequence length, computing pairwise Hilbert phase differences scales quadratically with channel count ($O(C^2)$), which increases memory footprint when $C > 2,000$; incorporating anchor-based low-rank graph sparsification provides a direct remedy. Second, our continuous neural-ODE scale tracker assumes piece-wise smooth drift and may smooth over instantaneous step discontinuities caused by sensor recalibration. Extending the continuous drift prior to jump-diffusion stochastic differential equations offers a principled direction for future investigation.
\end{ExampleBox}
\begin{TipBox}{Rhetorical Move Breakdown \& Constituent Analysis}
\begin{itemize}[leftmargin=1.5em]
  \item \textbf{S1 [边界 1:} \textit{Despite its strong empirical performance, our framework exhibits two theoretical boundaries worth noting.}\par\noindent{\small 对应步骤：S1 [边界 1：极端高维或极度稀疏场景 Boundary 1]：Despite its effectiveness across standard benchmarks, our framework currently exhibits two practical boundaries. | 句首锚定旧信息/逻辑视角，主干动词紧跟主语，复杂修饰置于宾语之后（符合尾重原则 L011 与旧新信息流 L001）。}
  \item \textbf{S2 [边界 1 机理分析与解法 Mechanism \& Solution 1]:} \textit{First, while our spectral branch scales quasi-linearly with sequence length, computing pairwise Hilbert phase differences scales quadratically with channel count ($O(C^2)$), which increases memory footprint when $C > 2,000$; incorporating anchor-based low-rank graph sparsification provides a direct remedy.}\par\noindent{\small 对应步骤：S2 [边界 1 机理分析与解法 Mechanism \& Solution 1]：First, when the variate dimension scales beyond $C > 5,000$ unclustered sensors, constructing the cross-channel phase graph incurs noticeable memory overhead; integrating hierarchical sensor clustering offers a natural remedy. | 句首锚定旧信息/逻辑视角，主干动词紧跟主语，复杂修饰置于宾语之后（符合尾重原则 L011 与旧新信息流 L001）。}
  \item \textbf{S3 [边界 2:} \textit{Second, our continuous neural-ODE scale tracker assumes piece-wise smooth drift and may smooth over instantaneous step discontinuities caused by sensor recalibration.}\par\noindent{\small 对应步骤：S3 [边界 2：特定假设限制 Boundary 2]：Second, our continuous scale tracker assumes piece-wise differentiable drift and may under-estimate instantaneous discontinuous jumps caused by hardware resets. | 句首锚定旧信息/逻辑视角，主干动词紧跟主语，复杂修饰置于宾语之后（符合尾重原则 L011 与旧新信息流 L001）。}
  \item \textbf{S4 [未来演进方向 Future Extension]:} \textit{Extending the continuous drift prior to jump-diffusion stochastic differential equations offers a principled direction for future investigation.}\par\noindent{\small 对应步骤：S4 [未来演进方向 Future Extension]：Coupling our neural ODE tracker with jump-diffusion stochastic processes represents a promising avenue for future work. | 句首锚定旧信息/逻辑视角，主干动词紧跟主语，复杂修饰置于宾语之后（符合尾重原则 L011 与旧新信息流 L001）。}
\end{itemize}
\end{TipBox}
\medskip
\clearpage
\section{FT19: 六级写作（CET-6）满分范文：高校教学中生成式 AI 辅助与独立批判性思维培养的平衡之道}\label{comp:FT19}
\begin{ConceptBox}{Composition Architecture \& Structural Blueprint Chain}
\textbf{Genre:} \texttt{CET6\_Essay} \quad \textbf{Paragraph Count:} 5\par
\textbf{Blueprint Chain:} \texttt{BP33 $\to$ BP34 $\to$ BP35 $\to$ BP37 $\to$ BP40}
\end{ConceptBox}
\subsection*{Paragraph 1: 议论文破题首段：具象趋势锚定 + 争议张力 + 明确中心论点（Zero-Template Hook \& Thesis） (\texttt{BP33})}
\noindent\textbf{Expressive Intent:}\enspace 我这一段想写六级、雅思 Task 2 或托福独立写作的第一段（Introduction）：彻底扔掉 'With the rapid development of science and technology, more and more people...' 等烂大街模板，用 3 句话写出考官眼前一亮的母语级破题段。\par
\noindent\textbf{Cohesion Formula:}\enspace {\small\color{AcadSlate}Concrete Phenomenon (No Cliché) $\to$ Balanced Tension $\to$ Conditional \& Specific Thesis Statement. 绑定逻辑规则：L023, L029.}\par\vspace{4pt}
\begin{ExampleBox}{Publication-Grade Authentic Text}
Large language models and automated code assistants have rapidly become routine study aids across university campuses. While proponents contend that these tools accelerate literature synthesis and lower programming barriers, critics warn that unchecked reliance erodes students' capacity for independent problem formulation. In my view, universities should integrate generative tools into advanced coursework only after students master foundational analytical skills, as premature delegation weakens debugging competence and obscures underlying mathematical principles.
\end{ExampleBox}
\begin{TipBox}{Rhetorical Move Breakdown \& Constituent Analysis}
\begin{itemize}[leftmargin=1.5em]
  \item \textbf{S1 [具象现象锚定 Concrete Trend Anchor]:} \textit{Large language models and automated code assistants have rapidly become routine study aids across university campuses.}\par\noindent{\small 对应步骤：S1 [具象现象锚定 Concrete Trend Anchor]：直接从具体技术变革或制度调整切入（如 Generative AI tools now draft code and summarize literature across university campuses），绝不用 With the development of...。 | 句首锚定旧信息/逻辑视角，主干动词紧跟主语，复杂修饰置于宾语之后（符合尾重原则 L011 与旧新信息流 L001）。}
  \item \textbf{S2 [核心争议张力 Core Tension / Debate]:} \textit{While proponents contend that these tools accelerate literature synthesis and lower programming barriers, critics warn that unchecked reliance erodes students' capacity for independent problem formulation.}\par\noindent{\small 对应步骤：S2 [核心争议张力 Core Tension / Debate]：While proponents argue that [View A: efficiency/accessibility], critics warn that [View B: skill erosion/inequality]. | 句首锚定旧信息/逻辑视角，主干动词紧跟主语，复杂修饰置于宾语之后（符合尾重原则 L011 与旧新信息流 L001）。}
  \item \textbf{S3 [带条件的中心论点 Nuanced Thesis Statement]:} \textit{In my view, universities should integrate generative tools into advanced coursework only after students master foundational analytical skills, as premature delegation weakens debugging competence and obscures underlying mathematical principles.}\par\noindent{\small 对应步骤：S3 [带条件的中心论点 Nuanced Thesis Statement]：In my view, rather than banning [X] outright or adopting it uncritically, [Institution/Society] should [Core Stance], because [Primary Reason 1] and [Primary Reason 2]. | 句首锚定旧信息/逻辑视角，主干动词紧跟主语，复杂修饰置于宾语之后（符合尾重原则 L011 与旧新信息流 L001）。}
\end{itemize}
\end{TipBox}
\medskip
\subsection*{Paragraph 2: 核心正方分论点段：主张 $\to$ 机制推演 $\to$ 具象行业微案例 $\to$ 扣题闭环（PEEL 满分段） (\texttt{BP34})}
\noindent\textbf{Expressive Intent:}\enspace 我这一段想写大作文的第一个主体段（Body Paragraph 1）：论证某个核心优势或必要性，做到有理有据、有具象场景支撑，绝不车轱辘话来回说。\par
\noindent\textbf{Cohesion Formula:}\enspace {\small\color{AcadSlate}Topic Claim $\to$ Causal Mechanism $\to$ Concrete Micro-Case (L021) $\to$ Warrant Closure. 绑定逻辑规则：L006, L012, L021.}\par\vspace{4pt}
\begin{ExampleBox}{Publication-Grade Authentic Text}
Delegating repetitive data compilation to automated systems enables professionals to concentrate on high-level diagnostic judgment. In data-intensive fields such as clinical medicine and structural engineering, practitioners traditionally spend hours cross-checking routine telemetry logs before formulating a hypothesis. When hospital radiology departments deploy automated screening algorithms to flag normal chest scans, physicians can redirect their cognitive bandwidth toward ambiguous pathology cases and patient consultation. Consequently, intelligent triage elevates both diagnostic throughput and clinical care quality.
\end{ExampleBox}
\begin{TipBox}{Rhetorical Move Breakdown \& Constituent Analysis}
\begin{itemize}[leftmargin=1.5em]
  \item \textbf{S1 [分论点主题句 Topic Claim]:} \textit{Delegating repetitive data compilation to automated systems enables professionals to concentrate on high-level diagnostic judgment.}\par\noindent{\small 对应步骤：S1 [分论点主题句 Topic Claim]：一句话清晰亮出本段唯一控制中心（[Policy/Factor X] strengthens [Outcome Y] by [Core Mechanism]）。 | 句首锚定旧信息/逻辑视角，主干动词紧跟主语，复杂修饰置于宾语之后（符合尾重原则 L011 与旧新信息流 L001）。}
  \item \textbf{S2 [因果机制推演 Causal Elaboration]:} \textit{In data-intensive fields such as clinical medicine and structural engineering, practitioners traditionally spend hours cross-checking routine telemetry logs before formulating a hypothesis.}\par\noindent{\small 对应步骤：S2 [因果机制推演 Causal Elaboration]：Explain HOW and WHY：阐明如果没有这层干扰或有了这项机制，认知/经济/技术流程会发生什么具体改变。 | 句首锚定旧信息/逻辑视角，主干动词紧跟主语，复杂修饰置于宾语之后（符合尾重原则 L011 与旧新信息流 L001）。}
  \item \textbf{S3 [具象微观案例 Concrete Micro-Case]:} \textit{When hospital radiology departments deploy automated screening algorithms to flag normal chest scans, physicians can redirect their cognitive bandwidth toward ambiguous pathology cases and patient consultation.}\par\noindent{\small 对应步骤：S3 [具象微观案例 Concrete Micro-Case]：For instance, when [Specific Group/Industry] [Specific Action], [Measurable/Observable Consequence]（绝不编造假百分比，用具体场景细节说话）。 | 句首锚定旧信息/逻辑视角，主干动词紧跟主语，复杂修饰置于宾语之后（符合尾重原则 L011 与旧新信息流 L001）。}
  \item \textbf{S4 [反事实对比或扣题升华 Warrant / Link Back]:} \textit{Consequently, intelligent triage elevates both diagnostic throughput and clinical care quality.}\par\noindent{\small 对应步骤：S4 [反事实对比或扣题升华 Warrant / Link Back]：Thus, [Summary Noun] shifts human effort from [Low-value Routine] toward [High-value Judgment]. | 句首锚定旧信息/逻辑视角，主干动词紧跟主语，复杂修饰置于宾语之后（符合尾重原则 L011 与旧新信息流 L001）。}
\end{itemize}
\end{TipBox}
\medskip
\subsection*{Paragraph 3: 深层原因剖析段：表层诱因与结构性根源双层挖掘（Dual-Layer Causal Diagnosis） (\texttt{BP35})}
\noindent\textbf{Expressive Intent:}\enspace 我这一段想在问题分析类作文（Why is this happening?）中，剖析某个社会/科研痛点背后的深层原因（既指出直接表层诱因，又挖出底层的考核激励或结构性根源）。\par
\noindent\textbf{Cohesion Formula:}\enspace {\small\color{AcadSlate}Dual-Cause Opening $\to$ Immediate Operational Driver $\to$ Deeper Institutional Root $\to$ Systemic Compounding. 绑定逻辑规则：L007, L005.}\par\vspace{4pt}
\begin{ExampleBox}{Publication-Grade Authentic Text}
The growing tendency among undergraduates to submit unverified AI-generated essays stems from both lowered friction costs and flawed assessment design. At an immediate level, conversational models produce grammatically polished drafts within seconds, tempting time-pressed students to bypass the arduous process of outlining and primary reading. More fundamentally, many university courses still rely on generic, formulaic essay prompts that reward surface fluency rather than original empirical critique. When effortless text generation meets standardized grading rubrics that fail to penalize hollow reasoning, superficial compliance inevitably crowds out deep inquiry.
\end{ExampleBox}
\begin{TipBox}{Rhetorical Move Breakdown \& Constituent Analysis}
\begin{itemize}[leftmargin=1.5em]
  \item \textbf{S1 [原因总起句 Dual-Cause Thesis]:} \textit{The growing tendency among undergraduates to submit unverified AI-generated essays stems from both lowered friction costs and flawed assessment design.}\par\noindent{\small 对应步骤：S1 [原因总起句 Dual-Cause Thesis]：The widespread emergence of [Problem X] stems from both immediate technological pressures and deeper institutional incentives. | 句首锚定旧信息/逻辑视角，主干动词紧跟主语，复杂修饰置于宾语之后（符合尾重原则 L011 与旧新信息流 L001）。}
  \item \textbf{S2 [第一层:} \textit{At an immediate level, conversational models produce grammatically polished drafts within seconds, tempting time-pressed students to bypass the arduous process of outlining and primary reading.}\par\noindent{\small 对应步骤：S2 [第一层：直接操作/技术诱因 Immediate Driver]：At the operational level, [Factor A] lowers the cost of [Behavior], tempting individuals/organizations to [Short-term Action]. | 句首锚定旧信息/逻辑视角，主干动词紧跟主语，复杂修饰置于宾语之后（符合尾重原则 L011 与旧新信息流 L001）。}
  \item \textbf{S3 [第二层:} \textit{More fundamentally, many university courses still rely on generic, formulaic essay prompts that reward surface fluency rather than original empirical critique.}\par\noindent{\small 对应步骤：S3 [第二层：深层制度/激励根源 Structural Root Cause]：More fundamentally, existing evaluation systems reward [Superficial Metric] rather than [Long-term Quality], creating powerful pressure to prioritize speed over rigor. | 句首锚定旧信息/逻辑视角，主干动词紧跟主语，复杂修饰置于宾语之后（符合尾重原则 L011 与旧新信息流 L001）。}
  \item \textbf{S4 [双因叠加后果 Compounding Effect]:} \textit{When effortless text generation meets standardized grading rubrics that fail to penalize hollow reasoning, superficial compliance inevitably crowds out deep inquiry.}\par\noindent{\small 对应步骤：S4 [双因叠加后果 Compounding Effect]：When effortless generation meets quantity-driven evaluation, [Problem X] becomes systemic rather than isolated. | 句首锚定旧信息/逻辑视角，主干动词紧跟主语，复杂修饰置于宾语之后（符合尾重原则 L011 与旧新信息流 L001）。}
\end{itemize}
\end{TipBox}
\medskip
\subsection*{Paragraph 4: 问题诊断与靶向解决方案段：拒绝“政府应该立法”万金油（Targeted Problem-Solution） (\texttt{BP37})}
\noindent\textbf{Expressive Intent:}\enspace 我这一段想写 Problem-Solution 类作文的解决方案段：针对前文诊断出的两个具体根源，提出两项可操作、有明确执行主体和传导机制的靶向对策，拒绝只喊 'The government should make laws and raise awareness'。\par
\noindent\textbf{Cohesion Formula:}\enspace {\small\color{AcadSlate}Dual-Level Intervention $\to$ Actor 1 + Concrete Mechanism $\to$ Actor 2 + Concrete Standard $\to$ Balanced Closure. 绑定逻辑规则：L002, L010.}\par\vspace{4pt}
\begin{ExampleBox}{Publication-Grade Authentic Text}
Curbing the proliferation of fabricated academic data requires coordinated reforms across both peer-review protocols and institutional promotion criteria. First, journal editorial boards should mandate that empirical submissions include executable code containers and raw sensor logs, enabling automated reproducibility checks before external review. Concurrently, university tenure committees must shift evaluation weight from raw publication counts toward verified open-source contributions and replication studies. Aligning technical verification with career incentives removes the underlying pressure to fabricate incremental results.
\end{ExampleBox}
\begin{TipBox}{Rhetorical Move Breakdown \& Constituent Analysis}
\begin{itemize}[leftmargin=1.5em]
  \item \textbf{S1 [对策总起句 Targeted Solution Thesis]:} \textit{Curbing the proliferation of fabricated academic data requires coordinated reforms across both peer-review protocols and institutional promotion criteria.}\par\noindent{\small 对应步骤：S1 [对策总起句 Targeted Solution Thesis]：Mitigating [Problem X] effectively requires coordinated intervention across both [Level A: Technical/Curriculum Design] and [Level B: Institutional Governance]. | 句首锚定旧信息/逻辑视角，主干动词紧跟主语，复杂修饰置于宾语之后（符合尾重原则 L011 与旧新信息流 L001）。}
  \item \textbf{S2 [对策 1:} \textit{First, journal editorial boards should mandate that empirical submissions include executable code containers and raw sensor logs, enabling automated reproducibility checks before external review.}\par\noindent{\small 对应步骤：S2 [对策 1：操作层可验证改革 Actionable Measure 1]：First, [Specific Actor 1, e.g., university departments] should replace [Vulnerable Practice] with [Verifiable Mechanism, e.g., oral defense and in-class iterative drafting], ensuring that [Specific Effect 1]. | 句首锚定旧信息/逻辑视角，主干动词紧跟主语，复杂修饰置于宾语之后（符合尾重原则 L011 与旧新信息流 L001）。}
  \item \textbf{S3 [对策 2:} \textit{Concurrently, university tenure committees must shift evaluation weight from raw publication counts toward verified open-source contributions and replication studies.}\par\noindent{\small 对应步骤：S3 [对策 2：制度/标准层激励对齐 Actionable Measure 2]：Concurrently, [Specific Actor 2, e.g., regulatory agencies / industry consortia] must establish [Concrete Standard, e.g., mandatory provenance watermarking and audit logs] to penalize [Misconduct]. | 句首锚定旧信息/逻辑视角，主干动词紧跟主语，复杂修饰置于宾语之后（符合尾重原则 L011 与旧新信息流 L001）。}
  \item \textbf{S4 [协同闭环效果 Combined Outcome]:} \textit{Aligning technical verification with career incentives removes the underlying pressure to fabricate incremental results.}\par\noindent{\small 对应步骤：S4 [协同闭环效果 Combined Outcome]：Together, aligning pedagogical verification with technical traceability curbs misuse without stifling legitimate innovation. | 句首锚定旧信息/逻辑视角，主干动词紧跟主语，复杂修饰置于宾语之后（符合尾重原则 L011 与旧新信息流 L001）。}
\end{itemize}
\end{TipBox}
\medskip
\subsection*{Paragraph 5: 议论文收尾段：跳出简单重复，完成核心权衡整合与边界升华（Synthesis Conclusion） (\texttt{BP40})}
\noindent\textbf{Expressive Intent:}\enspace 我这一段想写六级、雅思 Task 2 或托福独立写作的结尾段（Conclusion）：用 2--3 句话收束全文，做到既呼应中心论点、整合正反权衡，又绝不喊假大空口号（Let us work together for a better tomorrow）。\par
\noindent\textbf{Cohesion Formula:}\enspace {\small\color{AcadSlate}Balanced Restatment $\to$ Dual-Pillar Synthesis $\to$ Grounded Governance Principle (Zero Slogans!). 绑定逻辑规则：L023, L029.}\par\vspace{4pt}
\begin{ExampleBox}{Publication-Grade Authentic Text}
In conclusion, while automated decision systems inevitably introduce transitional labor friction and verification challenges, integrating them into industrial and academic workflows yields substantial long-term productivity gains. By pairing algorithmic deployment with transparent audit trails and targeted workforce retraining, institutions can capture efficiency breakthroughs without eroding human analytical oversight. Ultimately, the value of intelligent automation rests not on replacing human judgment, but on elevating the precision of the evidence upon which human judgment acts.
\end{ExampleBox}
\begin{TipBox}{Rhetorical Move Breakdown \& Constituent Analysis}
\begin{itemize}[leftmargin=1.5em]
  \item \textbf{S1 [重申核心权衡与立场 Restate Nuanced Thesis]:} \textit{In conclusion, while automated decision systems inevitably introduce transitional labor friction and verification challenges, integrating them into industrial and academic workflows yields substantial long-term productivity gains.}\par\noindent{\small 对应步骤：S1 [重申核心权衡与立场 Restate Nuanced Thesis]：In conclusion, while [Acknowledged Factor A] introduces real challenges in [Specific Aspect], [Core Thesis Stance] remains the most viable path forward. | 句首锚定旧信息/逻辑视角，主干动词紧跟主语，复杂修饰置于宾语之后（符合尾重原则 L011 与旧新信息流 L001）。}
  \item \textbf{S2 [整合全文两大支柱机制 Synthesize Dual Pillars]:} \textit{By pairing algorithmic deployment with transparent audit trails and targeted workforce retraining, institutions can capture efficiency breakthroughs without eroding human analytical oversight.}\par\noindent{\small 对应步骤：S2 [整合全文两大支柱机制 Synthesize Dual Pillars]：By combining [Pillar 1 from Body 1] with [Pillar 2 / Safeguard from Body 2], [Actors] can harness [Core Benefit] while containing [Primary Risk]. | 句首锚定旧信息/逻辑视角，主干动词紧跟主语，复杂修饰置于宾语之后（符合尾重原则 L011 与旧新信息流 L001）。}
  \item \textbf{S3 [理性边界展望（可选）Grounded Forward Look]:} \textit{Ultimately, the value of intelligent automation rests not on replacing human judgment, but on elevating the precision of the evidence upon which human judgment acts.}\par\noindent{\small 对应步骤：S3 [理性边界展望（可选）Grounded Forward Look]：Ultimately, the success of [Technology/Policy] depends not on its standalone capability, but on the institutional rigor with which it is governed. | 句首锚定旧信息/逻辑视角，主干动词紧跟主语，复杂修饰置于宾语之后（符合尾重原则 L011 与旧新信息流 L001）。}
\end{itemize}
\end{TipBox}
\medskip
\clearpage
\section{FT20: 六级写作（CET-6）满分范文：科研评价体系中抵制数据造假与强化可复现性治理}\label{comp:FT20}
\begin{ConceptBox}{Composition Architecture \& Structural Blueprint Chain}
\textbf{Genre:} \texttt{CET6\_Essay} \quad \textbf{Paragraph Count:} 5\par
\textbf{Blueprint Chain:} \texttt{BP33 $\to$ BP35 $\to$ BP37 $\to$ BP39 $\to$ BP40}
\end{ConceptBox}
\subsection*{Paragraph 1: 议论文破题首段：具象趋势锚定 + 争议张力 + 明确中心论点（Zero-Template Hook \& Thesis） (\texttt{BP33})}
\noindent\textbf{Expressive Intent:}\enspace 我这一段想写六级、雅思 Task 2 或托福独立写作的第一段（Introduction）：彻底扔掉 'With the rapid development of science and technology, more and more people...' 等烂大街模板，用 3 句话写出考官眼前一亮的母语级破题段。\par
\noindent\textbf{Cohesion Formula:}\enspace {\small\color{AcadSlate}Concrete Phenomenon (No Cliché) $\to$ Balanced Tension $\to$ Conditional \& Specific Thesis Statement. 绑定逻辑规则：L023, L029.}\par\vspace{4pt}
\begin{ExampleBox}{Publication-Grade Authentic Text}
Large language models and automated code assistants have rapidly become routine study aids across university campuses. While proponents contend that these tools accelerate literature synthesis and lower programming barriers, critics warn that unchecked reliance erodes students' capacity for independent problem formulation. In my view, universities should integrate generative tools into advanced coursework only after students master foundational analytical skills, as premature delegation weakens debugging competence and obscures underlying mathematical principles.
\end{ExampleBox}
\begin{TipBox}{Rhetorical Move Breakdown \& Constituent Analysis}
\begin{itemize}[leftmargin=1.5em]
  \item \textbf{S1 [具象现象锚定 Concrete Trend Anchor]:} \textit{Large language models and automated code assistants have rapidly become routine study aids across university campuses.}\par\noindent{\small 对应步骤：S1 [具象现象锚定 Concrete Trend Anchor]：直接从具体技术变革或制度调整切入（如 Generative AI tools now draft code and summarize literature across university campuses），绝不用 With the development of...。 | 句首锚定旧信息/逻辑视角，主干动词紧跟主语，复杂修饰置于宾语之后（符合尾重原则 L011 与旧新信息流 L001）。}
  \item \textbf{S2 [核心争议张力 Core Tension / Debate]:} \textit{While proponents contend that these tools accelerate literature synthesis and lower programming barriers, critics warn that unchecked reliance erodes students' capacity for independent problem formulation.}\par\noindent{\small 对应步骤：S2 [核心争议张力 Core Tension / Debate]：While proponents argue that [View A: efficiency/accessibility], critics warn that [View B: skill erosion/inequality]. | 句首锚定旧信息/逻辑视角，主干动词紧跟主语，复杂修饰置于宾语之后（符合尾重原则 L011 与旧新信息流 L001）。}
  \item \textbf{S3 [带条件的中心论点 Nuanced Thesis Statement]:} \textit{In my view, universities should integrate generative tools into advanced coursework only after students master foundational analytical skills, as premature delegation weakens debugging competence and obscures underlying mathematical principles.}\par\noindent{\small 对应步骤：S3 [带条件的中心论点 Nuanced Thesis Statement]：In my view, rather than banning [X] outright or adopting it uncritically, [Institution/Society] should [Core Stance], because [Primary Reason 1] and [Primary Reason 2]. | 句首锚定旧信息/逻辑视角，主干动词紧跟主语，复杂修饰置于宾语之后（符合尾重原则 L011 与旧新信息流 L001）。}
\end{itemize}
\end{TipBox}
\medskip
\subsection*{Paragraph 2: 深层原因剖析段：表层诱因与结构性根源双层挖掘（Dual-Layer Causal Diagnosis） (\texttt{BP35})}
\noindent\textbf{Expressive Intent:}\enspace 我这一段想在问题分析类作文（Why is this happening?）中，剖析某个社会/科研痛点背后的深层原因（既指出直接表层诱因，又挖出底层的考核激励或结构性根源）。\par
\noindent\textbf{Cohesion Formula:}\enspace {\small\color{AcadSlate}Dual-Cause Opening $\to$ Immediate Operational Driver $\to$ Deeper Institutional Root $\to$ Systemic Compounding. 绑定逻辑规则：L007, L005.}\par\vspace{4pt}
\begin{ExampleBox}{Publication-Grade Authentic Text}
The growing tendency among undergraduates to submit unverified AI-generated essays stems from both lowered friction costs and flawed assessment design. At an immediate level, conversational models produce grammatically polished drafts within seconds, tempting time-pressed students to bypass the arduous process of outlining and primary reading. More fundamentally, many university courses still rely on generic, formulaic essay prompts that reward surface fluency rather than original empirical critique. When effortless text generation meets standardized grading rubrics that fail to penalize hollow reasoning, superficial compliance inevitably crowds out deep inquiry.
\end{ExampleBox}
\begin{TipBox}{Rhetorical Move Breakdown \& Constituent Analysis}
\begin{itemize}[leftmargin=1.5em]
  \item \textbf{S1 [原因总起句 Dual-Cause Thesis]:} \textit{The growing tendency among undergraduates to submit unverified AI-generated essays stems from both lowered friction costs and flawed assessment design.}\par\noindent{\small 对应步骤：S1 [原因总起句 Dual-Cause Thesis]：The widespread emergence of [Problem X] stems from both immediate technological pressures and deeper institutional incentives. | 句首锚定旧信息/逻辑视角，主干动词紧跟主语，复杂修饰置于宾语之后（符合尾重原则 L011 与旧新信息流 L001）。}
  \item \textbf{S2 [第一层:} \textit{At an immediate level, conversational models produce grammatically polished drafts within seconds, tempting time-pressed students to bypass the arduous process of outlining and primary reading.}\par\noindent{\small 对应步骤：S2 [第一层：直接操作/技术诱因 Immediate Driver]：At the operational level, [Factor A] lowers the cost of [Behavior], tempting individuals/organizations to [Short-term Action]. | 句首锚定旧信息/逻辑视角，主干动词紧跟主语，复杂修饰置于宾语之后（符合尾重原则 L011 与旧新信息流 L001）。}
  \item \textbf{S3 [第二层:} \textit{More fundamentally, many university courses still rely on generic, formulaic essay prompts that reward surface fluency rather than original empirical critique.}\par\noindent{\small 对应步骤：S3 [第二层：深层制度/激励根源 Structural Root Cause]：More fundamentally, existing evaluation systems reward [Superficial Metric] rather than [Long-term Quality], creating powerful pressure to prioritize speed over rigor. | 句首锚定旧信息/逻辑视角，主干动词紧跟主语，复杂修饰置于宾语之后（符合尾重原则 L011 与旧新信息流 L001）。}
  \item \textbf{S4 [双因叠加后果 Compounding Effect]:} \textit{When effortless text generation meets standardized grading rubrics that fail to penalize hollow reasoning, superficial compliance inevitably crowds out deep inquiry.}\par\noindent{\small 对应步骤：S4 [双因叠加后果 Compounding Effect]：When effortless generation meets quantity-driven evaluation, [Problem X] becomes systemic rather than isolated. | 句首锚定旧信息/逻辑视角，主干动词紧跟主语，复杂修饰置于宾语之后（符合尾重原则 L011 与旧新信息流 L001）。}
\end{itemize}
\end{TipBox}
\medskip
\subsection*{Paragraph 3: 问题诊断与靶向解决方案段：拒绝“政府应该立法”万金油（Targeted Problem-Solution） (\texttt{BP37})}
\noindent\textbf{Expressive Intent:}\enspace 我这一段想写 Problem-Solution 类作文的解决方案段：针对前文诊断出的两个具体根源，提出两项可操作、有明确执行主体和传导机制的靶向对策，拒绝只喊 'The government should make laws and raise awareness'。\par
\noindent\textbf{Cohesion Formula:}\enspace {\small\color{AcadSlate}Dual-Level Intervention $\to$ Actor 1 + Concrete Mechanism $\to$ Actor 2 + Concrete Standard $\to$ Balanced Closure. 绑定逻辑规则：L002, L010.}\par\vspace{4pt}
\begin{ExampleBox}{Publication-Grade Authentic Text}
Curbing the proliferation of fabricated academic data requires coordinated reforms across both peer-review protocols and institutional promotion criteria. First, journal editorial boards should mandate that empirical submissions include executable code containers and raw sensor logs, enabling automated reproducibility checks before external review. Concurrently, university tenure committees must shift evaluation weight from raw publication counts toward verified open-source contributions and replication studies. Aligning technical verification with career incentives removes the underlying pressure to fabricate incremental results.
\end{ExampleBox}
\begin{TipBox}{Rhetorical Move Breakdown \& Constituent Analysis}
\begin{itemize}[leftmargin=1.5em]
  \item \textbf{S1 [对策总起句 Targeted Solution Thesis]:} \textit{Curbing the proliferation of fabricated academic data requires coordinated reforms across both peer-review protocols and institutional promotion criteria.}\par\noindent{\small 对应步骤：S1 [对策总起句 Targeted Solution Thesis]：Mitigating [Problem X] effectively requires coordinated intervention across both [Level A: Technical/Curriculum Design] and [Level B: Institutional Governance]. | 句首锚定旧信息/逻辑视角，主干动词紧跟主语，复杂修饰置于宾语之后（符合尾重原则 L011 与旧新信息流 L001）。}
  \item \textbf{S2 [对策 1:} \textit{First, journal editorial boards should mandate that empirical submissions include executable code containers and raw sensor logs, enabling automated reproducibility checks before external review.}\par\noindent{\small 对应步骤：S2 [对策 1：操作层可验证改革 Actionable Measure 1]：First, [Specific Actor 1, e.g., university departments] should replace [Vulnerable Practice] with [Verifiable Mechanism, e.g., oral defense and in-class iterative drafting], ensuring that [Specific Effect 1]. | 句首锚定旧信息/逻辑视角，主干动词紧跟主语，复杂修饰置于宾语之后（符合尾重原则 L011 与旧新信息流 L001）。}
  \item \textbf{S3 [对策 2:} \textit{Concurrently, university tenure committees must shift evaluation weight from raw publication counts toward verified open-source contributions and replication studies.}\par\noindent{\small 对应步骤：S3 [对策 2：制度/标准层激励对齐 Actionable Measure 2]：Concurrently, [Specific Actor 2, e.g., regulatory agencies / industry consortia] must establish [Concrete Standard, e.g., mandatory provenance watermarking and audit logs] to penalize [Misconduct]. | 句首锚定旧信息/逻辑视角，主干动词紧跟主语，复杂修饰置于宾语之后（符合尾重原则 L011 与旧新信息流 L001）。}
  \item \textbf{S4 [协同闭环效果 Combined Outcome]:} \textit{Aligning technical verification with career incentives removes the underlying pressure to fabricate incremental results.}\par\noindent{\small 对应步骤：S4 [协同闭环效果 Combined Outcome]：Together, aligning pedagogical verification with technical traceability curbs misuse without stifling legitimate innovation. | 句首锚定旧信息/逻辑视角，主干动词紧跟主语，复杂修饰置于宾语之后（符合尾重原则 L011 与旧新信息流 L001）。}
\end{itemize}
\end{TipBox}
\medskip
\subsection*{Paragraph 4: 反事实推理段：如果不采取某项原则，系统将如何一步步滑向失控（Counterfactual Escalation） (\texttt{BP39})}
\noindent\textbf{Expressive Intent:}\enspace 我这一段想在议论文或科研立项书中，通过严密的“反事实推演（If we fail to... / Absent X...）”强化紧迫感与说服力。\par
\noindent\textbf{Cohesion Formula:}\enspace {\small\color{AcadSlate}Positive Safeguard $\to$ Counterfactual Premise (Absent X...) $\to$ Cascading Breakdown $\to$ Precedent Confirmation. 绑定逻辑规则：L020, L003.}\par\vspace{4pt}
\begin{ExampleBox}{Publication-Grade Authentic Text}
Preserving verified human-authored corpora serves as an indispensable anchor for training successive generations of generative models. Absent this human data anchor---if future models were trained recursively on synthetic outputs generated by their predecessors---subtle approximation errors would compound across iterations. In each recursive cycle, low-probability tail events and nuanced stylistic variations are progressively truncated, driving the learned distribution toward high-density mode collapse. This degenerative feedback loop---known empirically as model autophagy---underscores why synthetic data augmentation can complement, but never wholly replace, curated real-world observations.
\end{ExampleBox}
\begin{TipBox}{Rhetorical Move Breakdown \& Constituent Analysis}
\begin{itemize}[leftmargin=1.5em]
  \item \textbf{S1 [正面基准命题 Positive Baseline]:} \textit{Preserving verified human-authored corpora serves as an indispensable anchor for training successive generations of generative models.}\par\noindent{\small 对应步骤：S1 [正面基准命题 Positive Baseline]：Maintaining rigorous [Principle/Standard X] serves as the primary safeguard against [Systemic Risk Y]. | 句首锚定旧信息/逻辑视角，主干动词紧跟主语，复杂修饰置于宾语之后（符合尾重原则 L011 与旧新信息流 L001）。}
  \item \textbf{S2 [反事实假设切入 Counterfactual Premise]:} \textit{Absent this human data anchor---if future models were trained recursively on synthetic outputs generated by their predecessors---subtle approximation errors would compound across iterations.}\par\noindent{\small 对应步骤：S2 [反事实假设切入 Counterfactual Premise]：Absent this safeguard---if [Actors] were to rely exclusively on [Unchecked Alternative]---[Immediate Vulnerability] would quickly emerge. | 句首锚定旧信息/逻辑视角，主干动词紧跟主语，复杂修饰置于宾语之后（符合尾重原则 L011 与旧新信息流 L001）。}
  \item \textbf{S3 [连锁恶化传导 Cascading Deterioration]:} \textit{In each recursive cycle, low-probability tail events and nuanced stylistic variations are progressively truncated, driving the learned distribution toward high-density mode collapse.}\par\noindent{\small 对应步骤：S3 [连锁恶化传导 Cascading Deterioration]：Once [Immediate Vulnerability] takes hold, [Secondary Consequence] compounds across cycles, eventually eroding [Core Foundation]. | 句首锚定旧信息/逻辑视角，主干动词紧跟主语，复杂修饰置于宾语之后（符合尾重原则 L011 与旧新信息流 L001）。}
  \item \textbf{S4 [实证印证或结论锁定 Empirical Echo]:} \textit{This degenerative feedback loop---known empirically as model autophagy---underscores why synthetic data augmentation can complement, but never wholly replace, curated real-world observations.}\par\noindent{\small 对应步骤：S4 [实证印证或结论锁定 Empirical Echo]：This degradation dynamic mirrors [Real-world/Technical Precedent], underscoring why [Principle X] cannot be compromised for short-term convenience. | 句首锚定旧信息/逻辑视角，主干动词紧跟主语，复杂修饰置于宾语之后（符合尾重原则 L011 与旧新信息流 L001）。}
\end{itemize}
\end{TipBox}
\medskip
\subsection*{Paragraph 5: 议论文收尾段：跳出简单重复，完成核心权衡整合与边界升华（Synthesis Conclusion） (\texttt{BP40})}
\noindent\textbf{Expressive Intent:}\enspace 我这一段想写六级、雅思 Task 2 或托福独立写作的结尾段（Conclusion）：用 2--3 句话收束全文，做到既呼应中心论点、整合正反权衡，又绝不喊假大空口号（Let us work together for a better tomorrow）。\par
\noindent\textbf{Cohesion Formula:}\enspace {\small\color{AcadSlate}Balanced Restatment $\to$ Dual-Pillar Synthesis $\to$ Grounded Governance Principle (Zero Slogans!). 绑定逻辑规则：L023, L029.}\par\vspace{4pt}
\begin{ExampleBox}{Publication-Grade Authentic Text}
In conclusion, while automated decision systems inevitably introduce transitional labor friction and verification challenges, integrating them into industrial and academic workflows yields substantial long-term productivity gains. By pairing algorithmic deployment with transparent audit trails and targeted workforce retraining, institutions can capture efficiency breakthroughs without eroding human analytical oversight. Ultimately, the value of intelligent automation rests not on replacing human judgment, but on elevating the precision of the evidence upon which human judgment acts.
\end{ExampleBox}
\begin{TipBox}{Rhetorical Move Breakdown \& Constituent Analysis}
\begin{itemize}[leftmargin=1.5em]
  \item \textbf{S1 [重申核心权衡与立场 Restate Nuanced Thesis]:} \textit{In conclusion, while automated decision systems inevitably introduce transitional labor friction and verification challenges, integrating them into industrial and academic workflows yields substantial long-term productivity gains.}\par\noindent{\small 对应步骤：S1 [重申核心权衡与立场 Restate Nuanced Thesis]：In conclusion, while [Acknowledged Factor A] introduces real challenges in [Specific Aspect], [Core Thesis Stance] remains the most viable path forward. | 句首锚定旧信息/逻辑视角，主干动词紧跟主语，复杂修饰置于宾语之后（符合尾重原则 L011 与旧新信息流 L001）。}
  \item \textbf{S2 [整合全文两大支柱机制 Synthesize Dual Pillars]:} \textit{By pairing algorithmic deployment with transparent audit trails and targeted workforce retraining, institutions can capture efficiency breakthroughs without eroding human analytical oversight.}\par\noindent{\small 对应步骤：S2 [整合全文两大支柱机制 Synthesize Dual Pillars]：By combining [Pillar 1 from Body 1] with [Pillar 2 / Safeguard from Body 2], [Actors] can harness [Core Benefit] while containing [Primary Risk]. | 句首锚定旧信息/逻辑视角，主干动词紧跟主语，复杂修饰置于宾语之后（符合尾重原则 L011 与旧新信息流 L001）。}
  \item \textbf{S3 [理性边界展望（可选）Grounded Forward Look]:} \textit{Ultimately, the value of intelligent automation rests not on replacing human judgment, but on elevating the precision of the evidence upon which human judgment acts.}\par\noindent{\small 对应步骤：S3 [理性边界展望（可选）Grounded Forward Look]：Ultimately, the success of [Technology/Policy] depends not on its standalone capability, but on the institutional rigor with which it is governed. | 句首锚定旧信息/逻辑视角，主干动词紧跟主语，复杂修饰置于宾语之后（符合尾重原则 L011 与旧新信息流 L001）。}
\end{itemize}
\end{TipBox}
\medskip
\clearpage
\chapter{Humanities \& Social Sciences / 人文社科与学术测评}\label{ch:anthology_22}
\Needspace{6\baselineskip}
本章精选该学科领域及对应文体的完整学术范文。每一篇范文不仅展示母语级发表标准的正文全貌，更清晰拆解全篇的“段落蓝图链条（Structural Blueprint Chain）”，并对每个自然段逐一剖析其核心交际意图、句际衔接逻辑公式（Cohesion Formula）以及每句话的句法成分与推进机制。

\section{FT21: 雅思大作文（Task 2 Discuss Both Views）8.5 分范文：工业自动化机器人普及的就业冲击与技能转型}\label{comp:FT21}
\begin{ConceptBox}{Composition Architecture \& Structural Blueprint Chain}
\textbf{Genre:} \texttt{IELTS\_Task2\_DiscussBoth} \quad \textbf{Paragraph Count:} 5\par
\textbf{Blueprint Chain:} \texttt{BP33 $\to$ BP34 $\to$ BP36 $\to$ BP37 $\to$ BP40}
\end{ConceptBox}
\subsection*{Paragraph 1: 议论文破题首段：具象趋势锚定 + 争议张力 + 明确中心论点（Zero-Template Hook \& Thesis） (\texttt{BP33})}
\noindent\textbf{Expressive Intent:}\enspace 我这一段想写六级、雅思 Task 2 或托福独立写作的第一段（Introduction）：彻底扔掉 'With the rapid development of science and technology, more and more people...' 等烂大街模板，用 3 句话写出考官眼前一亮的母语级破题段。\par
\noindent\textbf{Cohesion Formula:}\enspace {\small\color{AcadSlate}Concrete Phenomenon (No Cliché) $\to$ Balanced Tension $\to$ Conditional \& Specific Thesis Statement. 绑定逻辑规则：L023, L029.}\par\vspace{4pt}
\begin{ExampleBox}{Publication-Grade Authentic Text}
Large language models and automated code assistants have rapidly become routine study aids across university campuses. While proponents contend that these tools accelerate literature synthesis and lower programming barriers, critics warn that unchecked reliance erodes students' capacity for independent problem formulation. In my view, universities should integrate generative tools into advanced coursework only after students master foundational analytical skills, as premature delegation weakens debugging competence and obscures underlying mathematical principles.
\end{ExampleBox}
\begin{TipBox}{Rhetorical Move Breakdown \& Constituent Analysis}
\begin{itemize}[leftmargin=1.5em]
  \item \textbf{S1 [具象现象锚定 Concrete Trend Anchor]:} \textit{Large language models and automated code assistants have rapidly become routine study aids across university campuses.}\par\noindent{\small 对应步骤：S1 [具象现象锚定 Concrete Trend Anchor]：直接从具体技术变革或制度调整切入（如 Generative AI tools now draft code and summarize literature across university campuses），绝不用 With the development of...。 | 句首锚定旧信息/逻辑视角，主干动词紧跟主语，复杂修饰置于宾语之后（符合尾重原则 L011 与旧新信息流 L001）。}
  \item \textbf{S2 [核心争议张力 Core Tension / Debate]:} \textit{While proponents contend that these tools accelerate literature synthesis and lower programming barriers, critics warn that unchecked reliance erodes students' capacity for independent problem formulation.}\par\noindent{\small 对应步骤：S2 [核心争议张力 Core Tension / Debate]：While proponents argue that [View A: efficiency/accessibility], critics warn that [View B: skill erosion/inequality]. | 句首锚定旧信息/逻辑视角，主干动词紧跟主语，复杂修饰置于宾语之后（符合尾重原则 L011 与旧新信息流 L001）。}
  \item \textbf{S3 [带条件的中心论点 Nuanced Thesis Statement]:} \textit{In my view, universities should integrate generative tools into advanced coursework only after students master foundational analytical skills, as premature delegation weakens debugging competence and obscures underlying mathematical principles.}\par\noindent{\small 对应步骤：S3 [带条件的中心论点 Nuanced Thesis Statement]：In my view, rather than banning [X] outright or adopting it uncritically, [Institution/Society] should [Core Stance], because [Primary Reason 1] and [Primary Reason 2]. | 句首锚定旧信息/逻辑视角，主干动词紧跟主语，复杂修饰置于宾语之后（符合尾重原则 L011 与旧新信息流 L001）。}
\end{itemize}
\end{TipBox}
\medskip
\subsection*{Paragraph 2: 核心正方分论点段：主张 $\to$ 机制推演 $\to$ 具象行业微案例 $\to$ 扣题闭环（PEEL 满分段） (\texttt{BP34})}
\noindent\textbf{Expressive Intent:}\enspace 我这一段想写大作文的第一个主体段（Body Paragraph 1）：论证某个核心优势或必要性，做到有理有据、有具象场景支撑，绝不车轱辘话来回说。\par
\noindent\textbf{Cohesion Formula:}\enspace {\small\color{AcadSlate}Topic Claim $\to$ Causal Mechanism $\to$ Concrete Micro-Case (L021) $\to$ Warrant Closure. 绑定逻辑规则：L006, L012, L021.}\par\vspace{4pt}
\begin{ExampleBox}{Publication-Grade Authentic Text}
Delegating repetitive data compilation to automated systems enables professionals to concentrate on high-level diagnostic judgment. In data-intensive fields such as clinical medicine and structural engineering, practitioners traditionally spend hours cross-checking routine telemetry logs before formulating a hypothesis. When hospital radiology departments deploy automated screening algorithms to flag normal chest scans, physicians can redirect their cognitive bandwidth toward ambiguous pathology cases and patient consultation. Consequently, intelligent triage elevates both diagnostic throughput and clinical care quality.
\end{ExampleBox}
\begin{TipBox}{Rhetorical Move Breakdown \& Constituent Analysis}
\begin{itemize}[leftmargin=1.5em]
  \item \textbf{S1 [分论点主题句 Topic Claim]:} \textit{Delegating repetitive data compilation to automated systems enables professionals to concentrate on high-level diagnostic judgment.}\par\noindent{\small 对应步骤：S1 [分论点主题句 Topic Claim]：一句话清晰亮出本段唯一控制中心（[Policy/Factor X] strengthens [Outcome Y] by [Core Mechanism]）。 | 句首锚定旧信息/逻辑视角，主干动词紧跟主语，复杂修饰置于宾语之后（符合尾重原则 L011 与旧新信息流 L001）。}
  \item \textbf{S2 [因果机制推演 Causal Elaboration]:} \textit{In data-intensive fields such as clinical medicine and structural engineering, practitioners traditionally spend hours cross-checking routine telemetry logs before formulating a hypothesis.}\par\noindent{\small 对应步骤：S2 [因果机制推演 Causal Elaboration]：Explain HOW and WHY：阐明如果没有这层干扰或有了这项机制，认知/经济/技术流程会发生什么具体改变。 | 句首锚定旧信息/逻辑视角，主干动词紧跟主语，复杂修饰置于宾语之后（符合尾重原则 L011 与旧新信息流 L001）。}
  \item \textbf{S3 [具象微观案例 Concrete Micro-Case]:} \textit{When hospital radiology departments deploy automated screening algorithms to flag normal chest scans, physicians can redirect their cognitive bandwidth toward ambiguous pathology cases and patient consultation.}\par\noindent{\small 对应步骤：S3 [具象微观案例 Concrete Micro-Case]：For instance, when [Specific Group/Industry] [Specific Action], [Measurable/Observable Consequence]（绝不编造假百分比，用具体场景细节说话）。 | 句首锚定旧信息/逻辑视角，主干动词紧跟主语，复杂修饰置于宾语之后（符合尾重原则 L011 与旧新信息流 L001）。}
  \item \textbf{S4 [反事实对比或扣题升华 Warrant / Link Back]:} \textit{Consequently, intelligent triage elevates both diagnostic throughput and clinical care quality.}\par\noindent{\small 对应步骤：S4 [反事实对比或扣题升华 Warrant / Link Back]：Thus, [Summary Noun] shifts human effort from [Low-value Routine] toward [High-value Judgment]. | 句首锚定旧信息/逻辑视角，主干动词紧跟主语，复杂修饰置于宾语之后（符合尾重原则 L011 与旧新信息流 L001）。}
\end{itemize}
\end{TipBox}
\medskip
\subsection*{Paragraph 3: 让步-反驳段：优雅承认反方合理成分并强力扭转权重（Concession \& Refutation） (\texttt{BP36})}
\noindent\textbf{Expressive Intent:}\enspace 我这一段想在议论文中回应反方观点（Some people argue that...）：先大度承认反方担忧在短期或局部确实存在，再用结构性证据或制度设计证明该问题完全可控，从而锁死我方立场。\par
\noindent\textbf{Cohesion Formula:}\enspace {\small\color{AcadSlate}Admittedly (Valid Concern) $\to$ Short-term Reality $\to$ However (Zero-Sum Fallacy Pivot) $\to$ Historical Analogy \& Policy Mitigation. 绑定逻辑规则：L008, L023.}\par\vspace{4pt}
\begin{ExampleBox}{Publication-Grade Authentic Text}
Admittedly, critics raise a legitimate concern that autonomous robotics and generative software can displace routine clerical and assembly-line workers. During initial transition phases, rapid automation does create localized employment friction for personnel whose tasks follow rigid, codifiable rules. However, treating technological upgrading as a fixed-pie zero-sum game overlooks how productivity gains expand downstream service and engineering roles. Much as spreadsheet software eliminated manual ledger keeping while multiplying demand for financial data analysts, coupling industrial automation with subsidized reskilling programs enables displaced workers to transition into robot fleet supervision, calibration, and exception handling.
\end{ExampleBox}
\begin{TipBox}{Rhetorical Move Breakdown \& Constituent Analysis}
\begin{itemize}[leftmargin=1.5em]
  \item \textbf{S1 [呈现反方最有力的担忧 Articulate Opposing Concern]:} \textit{Admittedly, critics raise a legitimate concern that autonomous robotics and generative software can displace routine clerical and assembly-line workers.}\par\noindent{\small 对应步骤：S1 [呈现反方最有力的担忧 Articulate Opposing Concern]：Admittedly, critics raise a valid concern that [Opposing Argument, e.g., automation displaces entry-level clerical and manufacturing workers]. | 句首锚定旧信息/逻辑视角，主干动词紧跟主语，复杂修饰置于宾语之后（符合尾重原则 L011 与旧新信息流 L001）。}
  \item \textbf{S2 [承认局部真实性 Acknowledge Partial Validity]:} \textit{During initial transition phases, rapid automation does create localized employment friction for personnel whose tasks follow rigid, codifiable rules.}\par\noindent{\small 对应步骤：S2 [承认局部真实性 Acknowledge Partial Validity]：In the short term, rapid deployment of [Technology/Policy] can indeed create transitional friction for workers whose routines are easily codified. | 句首锚定旧信息/逻辑视角，主干动词紧跟主语，复杂修饰置于宾语之后（符合尾重原则 L011 与旧新信息流 L001）。}
  \item \textbf{S3 [釜底抽薪式转折 Structural Pivot]:} \textit{However, treating technological upgrading as a fixed-pie zero-sum game overlooks how productivity gains expand downstream service and engineering roles.}\par\noindent{\small 对应步骤：S3 [釜底抽薪式转折 Structural Pivot]：However, viewing [X] purely as a zero-sum threat overlooks how productivity gains generate entirely new categories of skilled employment. | 句首锚定旧信息/逻辑视角，主干动词紧跟主语，复杂修饰置于宾语之后（符合尾重原则 L011 与旧新信息流 L001）。}
  \item \textbf{S4 [历史/现实反证与制度补救 Evidence \& Mitigation]:} \textit{Much as spreadsheet software eliminated manual ledger keeping while multiplying demand for financial data analysts, coupling industrial automation with subsidized reskilling programs enables displaced workers to transition into robot fleet supervision, calibration, and exception handling.}\par\noindent{\small 对应步骤：S4 [历史/现实反证与制度补救 Evidence \& Mitigation]：Just as spreadsheet software reduced manual bookkeeping while expanding demand for financial analysts, pairing automation with targeted vocational retraining allows displaced workers to transition into system supervision and maintenance roles. | 句首锚定旧信息/逻辑视角，主干动词紧跟主语，复杂修饰置于宾语之后（符合尾重原则 L011 与旧新信息流 L001）。}
\end{itemize}
\end{TipBox}
\medskip
\subsection*{Paragraph 4: 问题诊断与靶向解决方案段：拒绝“政府应该立法”万金油（Targeted Problem-Solution） (\texttt{BP37})}
\noindent\textbf{Expressive Intent:}\enspace 我这一段想写 Problem-Solution 类作文的解决方案段：针对前文诊断出的两个具体根源，提出两项可操作、有明确执行主体和传导机制的靶向对策，拒绝只喊 'The government should make laws and raise awareness'。\par
\noindent\textbf{Cohesion Formula:}\enspace {\small\color{AcadSlate}Dual-Level Intervention $\to$ Actor 1 + Concrete Mechanism $\to$ Actor 2 + Concrete Standard $\to$ Balanced Closure. 绑定逻辑规则：L002, L010.}\par\vspace{4pt}
\begin{ExampleBox}{Publication-Grade Authentic Text}
Curbing the proliferation of fabricated academic data requires coordinated reforms across both peer-review protocols and institutional promotion criteria. First, journal editorial boards should mandate that empirical submissions include executable code containers and raw sensor logs, enabling automated reproducibility checks before external review. Concurrently, university tenure committees must shift evaluation weight from raw publication counts toward verified open-source contributions and replication studies. Aligning technical verification with career incentives removes the underlying pressure to fabricate incremental results.
\end{ExampleBox}
\begin{TipBox}{Rhetorical Move Breakdown \& Constituent Analysis}
\begin{itemize}[leftmargin=1.5em]
  \item \textbf{S1 [对策总起句 Targeted Solution Thesis]:} \textit{Curbing the proliferation of fabricated academic data requires coordinated reforms across both peer-review protocols and institutional promotion criteria.}\par\noindent{\small 对应步骤：S1 [对策总起句 Targeted Solution Thesis]：Mitigating [Problem X] effectively requires coordinated intervention across both [Level A: Technical/Curriculum Design] and [Level B: Institutional Governance]. | 句首锚定旧信息/逻辑视角，主干动词紧跟主语，复杂修饰置于宾语之后（符合尾重原则 L011 与旧新信息流 L001）。}
  \item \textbf{S2 [对策 1:} \textit{First, journal editorial boards should mandate that empirical submissions include executable code containers and raw sensor logs, enabling automated reproducibility checks before external review.}\par\noindent{\small 对应步骤：S2 [对策 1：操作层可验证改革 Actionable Measure 1]：First, [Specific Actor 1, e.g., university departments] should replace [Vulnerable Practice] with [Verifiable Mechanism, e.g., oral defense and in-class iterative drafting], ensuring that [Specific Effect 1]. | 句首锚定旧信息/逻辑视角，主干动词紧跟主语，复杂修饰置于宾语之后（符合尾重原则 L011 与旧新信息流 L001）。}
  \item \textbf{S3 [对策 2:} \textit{Concurrently, university tenure committees must shift evaluation weight from raw publication counts toward verified open-source contributions and replication studies.}\par\noindent{\small 对应步骤：S3 [对策 2：制度/标准层激励对齐 Actionable Measure 2]：Concurrently, [Specific Actor 2, e.g., regulatory agencies / industry consortia] must establish [Concrete Standard, e.g., mandatory provenance watermarking and audit logs] to penalize [Misconduct]. | 句首锚定旧信息/逻辑视角，主干动词紧跟主语，复杂修饰置于宾语之后（符合尾重原则 L011 与旧新信息流 L001）。}
  \item \textbf{S4 [协同闭环效果 Combined Outcome]:} \textit{Aligning technical verification with career incentives removes the underlying pressure to fabricate incremental results.}\par\noindent{\small 对应步骤：S4 [协同闭环效果 Combined Outcome]：Together, aligning pedagogical verification with technical traceability curbs misuse without stifling legitimate innovation. | 句首锚定旧信息/逻辑视角，主干动词紧跟主语，复杂修饰置于宾语之后（符合尾重原则 L011 与旧新信息流 L001）。}
\end{itemize}
\end{TipBox}
\medskip
\subsection*{Paragraph 5: 议论文收尾段：跳出简单重复，完成核心权衡整合与边界升华（Synthesis Conclusion） (\texttt{BP40})}
\noindent\textbf{Expressive Intent:}\enspace 我这一段想写六级、雅思 Task 2 或托福独立写作的结尾段（Conclusion）：用 2--3 句话收束全文，做到既呼应中心论点、整合正反权衡，又绝不喊假大空口号（Let us work together for a better tomorrow）。\par
\noindent\textbf{Cohesion Formula:}\enspace {\small\color{AcadSlate}Balanced Restatment $\to$ Dual-Pillar Synthesis $\to$ Grounded Governance Principle (Zero Slogans!). 绑定逻辑规则：L023, L029.}\par\vspace{4pt}
\begin{ExampleBox}{Publication-Grade Authentic Text}
In conclusion, while automated decision systems inevitably introduce transitional labor friction and verification challenges, integrating them into industrial and academic workflows yields substantial long-term productivity gains. By pairing algorithmic deployment with transparent audit trails and targeted workforce retraining, institutions can capture efficiency breakthroughs without eroding human analytical oversight. Ultimately, the value of intelligent automation rests not on replacing human judgment, but on elevating the precision of the evidence upon which human judgment acts.
\end{ExampleBox}
\begin{TipBox}{Rhetorical Move Breakdown \& Constituent Analysis}
\begin{itemize}[leftmargin=1.5em]
  \item \textbf{S1 [重申核心权衡与立场 Restate Nuanced Thesis]:} \textit{In conclusion, while automated decision systems inevitably introduce transitional labor friction and verification challenges, integrating them into industrial and academic workflows yields substantial long-term productivity gains.}\par\noindent{\small 对应步骤：S1 [重申核心权衡与立场 Restate Nuanced Thesis]：In conclusion, while [Acknowledged Factor A] introduces real challenges in [Specific Aspect], [Core Thesis Stance] remains the most viable path forward. | 句首锚定旧信息/逻辑视角，主干动词紧跟主语，复杂修饰置于宾语之后（符合尾重原则 L011 与旧新信息流 L001）。}
  \item \textbf{S2 [整合全文两大支柱机制 Synthesize Dual Pillars]:} \textit{By pairing algorithmic deployment with transparent audit trails and targeted workforce retraining, institutions can capture efficiency breakthroughs without eroding human analytical oversight.}\par\noindent{\small 对应步骤：S2 [整合全文两大支柱机制 Synthesize Dual Pillars]：By combining [Pillar 1 from Body 1] with [Pillar 2 / Safeguard from Body 2], [Actors] can harness [Core Benefit] while containing [Primary Risk]. | 句首锚定旧信息/逻辑视角，主干动词紧跟主语，复杂修饰置于宾语之后（符合尾重原则 L011 与旧新信息流 L001）。}
  \item \textbf{S3 [理性边界展望（可选）Grounded Forward Look]:} \textit{Ultimately, the value of intelligent automation rests not on replacing human judgment, but on elevating the precision of the evidence upon which human judgment acts.}\par\noindent{\small 对应步骤：S3 [理性边界展望（可选）Grounded Forward Look]：Ultimately, the success of [Technology/Policy] depends not on its standalone capability, but on the institutional rigor with which it is governed. | 句首锚定旧信息/逻辑视角，主干动词紧跟主语，复杂修饰置于宾语之后（符合尾重原则 L011 与旧新信息流 L001）。}
\end{itemize}
\end{TipBox}
\medskip
\clearpage
\section{FT22: 雅思大作文（Task 2 To What Extent Agree）8.5 分范文：开源大模型权重的科学透明收益是否远超其滥用风险}\label{comp:FT22}
\begin{ConceptBox}{Composition Architecture \& Structural Blueprint Chain}
\textbf{Genre:} \texttt{IELTS\_Task2\_AgreeDisagree} \quad \textbf{Paragraph Count:} 5\par
\textbf{Blueprint Chain:} \texttt{BP33 $\to$ BP38 $\to$ BP36 $\to$ BP39 $\to$ BP40}
\end{ConceptBox}
\subsection*{Paragraph 1: 议论文破题首段：具象趋势锚定 + 争议张力 + 明确中心论点（Zero-Template Hook \& Thesis） (\texttt{BP33})}
\noindent\textbf{Expressive Intent:}\enspace 我这一段想写六级、雅思 Task 2 或托福独立写作的第一段（Introduction）：彻底扔掉 'With the rapid development of science and technology, more and more people...' 等烂大街模板，用 3 句话写出考官眼前一亮的母语级破题段。\par
\noindent\textbf{Cohesion Formula:}\enspace {\small\color{AcadSlate}Concrete Phenomenon (No Cliché) $\to$ Balanced Tension $\to$ Conditional \& Specific Thesis Statement. 绑定逻辑规则：L023, L029.}\par\vspace{4pt}
\begin{ExampleBox}{Publication-Grade Authentic Text}
Large language models and automated code assistants have rapidly become routine study aids across university campuses. While proponents contend that these tools accelerate literature synthesis and lower programming barriers, critics warn that unchecked reliance erodes students' capacity for independent problem formulation. In my view, universities should integrate generative tools into advanced coursework only after students master foundational analytical skills, as premature delegation weakens debugging competence and obscures underlying mathematical principles.
\end{ExampleBox}
\begin{TipBox}{Rhetorical Move Breakdown \& Constituent Analysis}
\begin{itemize}[leftmargin=1.5em]
  \item \textbf{S1 [具象现象锚定 Concrete Trend Anchor]:} \textit{Large language models and automated code assistants have rapidly become routine study aids across university campuses.}\par\noindent{\small 对应步骤：S1 [具象现象锚定 Concrete Trend Anchor]：直接从具体技术变革或制度调整切入（如 Generative AI tools now draft code and summarize literature across university campuses），绝不用 With the development of...。 | 句首锚定旧信息/逻辑视角，主干动词紧跟主语，复杂修饰置于宾语之后（符合尾重原则 L011 与旧新信息流 L001）。}
  \item \textbf{S2 [核心争议张力 Core Tension / Debate]:} \textit{While proponents contend that these tools accelerate literature synthesis and lower programming barriers, critics warn that unchecked reliance erodes students' capacity for independent problem formulation.}\par\noindent{\small 对应步骤：S2 [核心争议张力 Core Tension / Debate]：While proponents argue that [View A: efficiency/accessibility], critics warn that [View B: skill erosion/inequality]. | 句首锚定旧信息/逻辑视角，主干动词紧跟主语，复杂修饰置于宾语之后（符合尾重原则 L011 与旧新信息流 L001）。}
  \item \textbf{S3 [带条件的中心论点 Nuanced Thesis Statement]:} \textit{In my view, universities should integrate generative tools into advanced coursework only after students master foundational analytical skills, as premature delegation weakens debugging competence and obscures underlying mathematical principles.}\par\noindent{\small 对应步骤：S3 [带条件的中心论点 Nuanced Thesis Statement]：In my view, rather than banning [X] outright or adopting it uncritically, [Institution/Society] should [Core Stance], because [Primary Reason 1] and [Primary Reason 2]. | 句首锚定旧信息/逻辑视角，主干动词紧跟主语，复杂修饰置于宾语之后（符合尾重原则 L011 与旧新信息流 L001）。}
\end{itemize}
\end{TipBox}
\medskip
\subsection*{Paragraph 2: 利弊权衡段：证明结构性长期收益远超短期可管可控成本（Outweigh / Trade-off Analysis） (\texttt{BP38})}
\noindent\textbf{Expressive Intent:}\enspace 我这一段想回应 'Do the advantages outweigh the disadvantages?' 题型：把弊端定性为“暂时的、可通过技术/规则收敛的局部成本”，把优势定性为“不可替代的结构性长期突破”。\par
\noindent\textbf{Cohesion Formula:}\enspace {\small\color{AcadSlate}Asymmetry Thesis $\to$ Bound the Drawback (Temporary \& Mitigable) $\to$ Highlight Compounding Advantage (Irreplaceable) $\to$ Rational Verdict. 绑定逻辑规则：L008, L023.}\par\vspace{4pt}
\begin{ExampleBox}{Publication-Grade Authentic Text}
Weighing both dimensions demonstrates that the structural scientific gains of open-weight foundation models far outweigh their manageable misuse risks. While公开 releasing model weights lowers the barrier for generating deceptive phishing content, such downstream abuse can be substantially contained through cryptographic provenance standards and endpoint spam filters. By contrast, open weights provide academic laboratories and resource-constrained hospitals with transparent, auditable architectures that proprietary black-box APIs cannot substitute---accelerating reproducible breakthroughs in protein folding and rare-disease diagnostics. Because the security risks are mitigable via defense-in-depth whereas the scientific transparency is irreplaceable, open research ecosystems yield a decisive net benefit.
\end{ExampleBox}
\begin{TipBox}{Rhetorical Move Breakdown \& Constituent Analysis}
\begin{itemize}[leftmargin=1.5em]
  \item \textbf{S1 [权衡定性主题句 Asymmetry Thesis]:} \textit{Weighing both dimensions demonstrates that the structural scientific gains of open-weight foundation models far outweigh their manageable misuse risks.}\par\noindent{\small 对应步骤：S1 [权衡定性主题句 Asymmetry Thesis]：Weighing both dimensions reveals that the long-term structural benefits of [X] far outweigh its manageable transitional costs. | 句首锚定旧信息/逻辑视角，主干动词紧跟主语，复杂修饰置于宾语之后（符合尾重原则 L011 与旧新信息流 L001）。}
  \item \textbf{S2 [将弊端限定为可收敛成本 Bounding the Drawback]:} \textit{While公开 releasing model weights lowers the barrier for generating deceptive phishing content, such downstream abuse can be substantially contained through cryptographic provenance standards and endpoint spam filters.}\par\noindent{\small 对应步骤：S2 [将弊端限定为可收敛成本 Bounding the Drawback]：While [Drawback A, e.g., initial capital outlay / training overhead] poses a tangible hurdle, this cost is bounded and diminishes rapidly as [Mitigating Factor, e.g., hardware amortization / open-source standardization] matures. | 句首锚定旧信息/逻辑视角，主干动词紧跟主语，复杂修饰置于宾语之后（符合尾重原则 L011 与旧新信息流 L001）。}
  \item \textbf{S3 [凸显优势的不可替代性 Highlighting Irreversible Gain]:} \textit{By contrast, open weights provide academic laboratories and resource-constrained hospitals with transparent, auditable architectures that proprietary black-box APIs cannot substitute---accelerating reproducible breakthroughs in protein folding and rare-disease diagnostics.}\par\noindent{\small 对应步骤：S3 [凸显优势的不可替代性 Highlighting Irreversible Gain]：By contrast, the core advantage---[Benefit B, e.g., eliminating human exposure to hazardous environments / discovering non-linear scientific laws]---delivers compounding value that manual methods cannot replicate at any cost. | 句首锚定旧信息/逻辑视角，主干动词紧跟主语，复杂修饰置于宾语之后（符合尾重原则 L011 与旧新信息流 L001）。}
  \item \textbf{S4 [权衡结论 Trade-off Verdict]:} \textit{Because the security risks are mitigable via defense-in-depth whereas the scientific transparency is irreplaceable, open research ecosystems yield a decisive net benefit.}\par\noindent{\small 对应步骤：S4 [权衡结论 Trade-off Verdict]：Therefore, addressing the friction through targeted safeguards is far more rational than forgoing the transformative gain. | 句首锚定旧信息/逻辑视角，主干动词紧跟主语，复杂修饰置于宾语之后（符合尾重原则 L011 与旧新信息流 L001）。}
\end{itemize}
\end{TipBox}
\medskip
\subsection*{Paragraph 3: 让步-反驳段：优雅承认反方合理成分并强力扭转权重（Concession \& Refutation） (\texttt{BP36})}
\noindent\textbf{Expressive Intent:}\enspace 我这一段想在议论文中回应反方观点（Some people argue that...）：先大度承认反方担忧在短期或局部确实存在，再用结构性证据或制度设计证明该问题完全可控，从而锁死我方立场。\par
\noindent\textbf{Cohesion Formula:}\enspace {\small\color{AcadSlate}Admittedly (Valid Concern) $\to$ Short-term Reality $\to$ However (Zero-Sum Fallacy Pivot) $\to$ Historical Analogy \& Policy Mitigation. 绑定逻辑规则：L008, L023.}\par\vspace{4pt}
\begin{ExampleBox}{Publication-Grade Authentic Text}
Admittedly, critics raise a legitimate concern that autonomous robotics and generative software can displace routine clerical and assembly-line workers. During initial transition phases, rapid automation does create localized employment friction for personnel whose tasks follow rigid, codifiable rules. However, treating technological upgrading as a fixed-pie zero-sum game overlooks how productivity gains expand downstream service and engineering roles. Much as spreadsheet software eliminated manual ledger keeping while multiplying demand for financial data analysts, coupling industrial automation with subsidized reskilling programs enables displaced workers to transition into robot fleet supervision, calibration, and exception handling.
\end{ExampleBox}
\begin{TipBox}{Rhetorical Move Breakdown \& Constituent Analysis}
\begin{itemize}[leftmargin=1.5em]
  \item \textbf{S1 [呈现反方最有力的担忧 Articulate Opposing Concern]:} \textit{Admittedly, critics raise a legitimate concern that autonomous robotics and generative software can displace routine clerical and assembly-line workers.}\par\noindent{\small 对应步骤：S1 [呈现反方最有力的担忧 Articulate Opposing Concern]：Admittedly, critics raise a valid concern that [Opposing Argument, e.g., automation displaces entry-level clerical and manufacturing workers]. | 句首锚定旧信息/逻辑视角，主干动词紧跟主语，复杂修饰置于宾语之后（符合尾重原则 L011 与旧新信息流 L001）。}
  \item \textbf{S2 [承认局部真实性 Acknowledge Partial Validity]:} \textit{During initial transition phases, rapid automation does create localized employment friction for personnel whose tasks follow rigid, codifiable rules.}\par\noindent{\small 对应步骤：S2 [承认局部真实性 Acknowledge Partial Validity]：In the short term, rapid deployment of [Technology/Policy] can indeed create transitional friction for workers whose routines are easily codified. | 句首锚定旧信息/逻辑视角，主干动词紧跟主语，复杂修饰置于宾语之后（符合尾重原则 L011 与旧新信息流 L001）。}
  \item \textbf{S3 [釜底抽薪式转折 Structural Pivot]:} \textit{However, treating technological upgrading as a fixed-pie zero-sum game overlooks how productivity gains expand downstream service and engineering roles.}\par\noindent{\small 对应步骤：S3 [釜底抽薪式转折 Structural Pivot]：However, viewing [X] purely as a zero-sum threat overlooks how productivity gains generate entirely new categories of skilled employment. | 句首锚定旧信息/逻辑视角，主干动词紧跟主语，复杂修饰置于宾语之后（符合尾重原则 L011 与旧新信息流 L001）。}
  \item \textbf{S4 [历史/现实反证与制度补救 Evidence \& Mitigation]:} \textit{Much as spreadsheet software eliminated manual ledger keeping while multiplying demand for financial data analysts, coupling industrial automation with subsidized reskilling programs enables displaced workers to transition into robot fleet supervision, calibration, and exception handling.}\par\noindent{\small 对应步骤：S4 [历史/现实反证与制度补救 Evidence \& Mitigation]：Just as spreadsheet software reduced manual bookkeeping while expanding demand for financial analysts, pairing automation with targeted vocational retraining allows displaced workers to transition into system supervision and maintenance roles. | 句首锚定旧信息/逻辑视角，主干动词紧跟主语，复杂修饰置于宾语之后（符合尾重原则 L011 与旧新信息流 L001）。}
\end{itemize}
\end{TipBox}
\medskip
\subsection*{Paragraph 4: 反事实推理段：如果不采取某项原则，系统将如何一步步滑向失控（Counterfactual Escalation） (\texttt{BP39})}
\noindent\textbf{Expressive Intent:}\enspace 我这一段想在议论文或科研立项书中，通过严密的“反事实推演（If we fail to... / Absent X...）”强化紧迫感与说服力。\par
\noindent\textbf{Cohesion Formula:}\enspace {\small\color{AcadSlate}Positive Safeguard $\to$ Counterfactual Premise (Absent X...) $\to$ Cascading Breakdown $\to$ Precedent Confirmation. 绑定逻辑规则：L020, L003.}\par\vspace{4pt}
\begin{ExampleBox}{Publication-Grade Authentic Text}
Preserving verified human-authored corpora serves as an indispensable anchor for training successive generations of generative models. Absent this human data anchor---if future models were trained recursively on synthetic outputs generated by their predecessors---subtle approximation errors would compound across iterations. In each recursive cycle, low-probability tail events and nuanced stylistic variations are progressively truncated, driving the learned distribution toward high-density mode collapse. This degenerative feedback loop---known empirically as model autophagy---underscores why synthetic data augmentation can complement, but never wholly replace, curated real-world observations.
\end{ExampleBox}
\begin{TipBox}{Rhetorical Move Breakdown \& Constituent Analysis}
\begin{itemize}[leftmargin=1.5em]
  \item \textbf{S1 [正面基准命题 Positive Baseline]:} \textit{Preserving verified human-authored corpora serves as an indispensable anchor for training successive generations of generative models.}\par\noindent{\small 对应步骤：S1 [正面基准命题 Positive Baseline]：Maintaining rigorous [Principle/Standard X] serves as the primary safeguard against [Systemic Risk Y]. | 句首锚定旧信息/逻辑视角，主干动词紧跟主语，复杂修饰置于宾语之后（符合尾重原则 L011 与旧新信息流 L001）。}
  \item \textbf{S2 [反事实假设切入 Counterfactual Premise]:} \textit{Absent this human data anchor---if future models were trained recursively on synthetic outputs generated by their predecessors---subtle approximation errors would compound across iterations.}\par\noindent{\small 对应步骤：S2 [反事实假设切入 Counterfactual Premise]：Absent this safeguard---if [Actors] were to rely exclusively on [Unchecked Alternative]---[Immediate Vulnerability] would quickly emerge. | 句首锚定旧信息/逻辑视角，主干动词紧跟主语，复杂修饰置于宾语之后（符合尾重原则 L011 与旧新信息流 L001）。}
  \item \textbf{S3 [连锁恶化传导 Cascading Deterioration]:} \textit{In each recursive cycle, low-probability tail events and nuanced stylistic variations are progressively truncated, driving the learned distribution toward high-density mode collapse.}\par\noindent{\small 对应步骤：S3 [连锁恶化传导 Cascading Deterioration]：Once [Immediate Vulnerability] takes hold, [Secondary Consequence] compounds across cycles, eventually eroding [Core Foundation]. | 句首锚定旧信息/逻辑视角，主干动词紧跟主语，复杂修饰置于宾语之后（符合尾重原则 L011 与旧新信息流 L001）。}
  \item \textbf{S4 [实证印证或结论锁定 Empirical Echo]:} \textit{This degenerative feedback loop---known empirically as model autophagy---underscores why synthetic data augmentation can complement, but never wholly replace, curated real-world observations.}\par\noindent{\small 对应步骤：S4 [实证印证或结论锁定 Empirical Echo]：This degradation dynamic mirrors [Real-world/Technical Precedent], underscoring why [Principle X] cannot be compromised for short-term convenience. | 句首锚定旧信息/逻辑视角，主干动词紧跟主语，复杂修饰置于宾语之后（符合尾重原则 L011 与旧新信息流 L001）。}
\end{itemize}
\end{TipBox}
\medskip
\subsection*{Paragraph 5: 议论文收尾段：跳出简单重复，完成核心权衡整合与边界升华（Synthesis Conclusion） (\texttt{BP40})}
\noindent\textbf{Expressive Intent:}\enspace 我这一段想写六级、雅思 Task 2 或托福独立写作的结尾段（Conclusion）：用 2--3 句话收束全文，做到既呼应中心论点、整合正反权衡，又绝不喊假大空口号（Let us work together for a better tomorrow）。\par
\noindent\textbf{Cohesion Formula:}\enspace {\small\color{AcadSlate}Balanced Restatment $\to$ Dual-Pillar Synthesis $\to$ Grounded Governance Principle (Zero Slogans!). 绑定逻辑规则：L023, L029.}\par\vspace{4pt}
\begin{ExampleBox}{Publication-Grade Authentic Text}
In conclusion, while automated decision systems inevitably introduce transitional labor friction and verification challenges, integrating them into industrial and academic workflows yields substantial long-term productivity gains. By pairing algorithmic deployment with transparent audit trails and targeted workforce retraining, institutions can capture efficiency breakthroughs without eroding human analytical oversight. Ultimately, the value of intelligent automation rests not on replacing human judgment, but on elevating the precision of the evidence upon which human judgment acts.
\end{ExampleBox}
\begin{TipBox}{Rhetorical Move Breakdown \& Constituent Analysis}
\begin{itemize}[leftmargin=1.5em]
  \item \textbf{S1 [重申核心权衡与立场 Restate Nuanced Thesis]:} \textit{In conclusion, while automated decision systems inevitably introduce transitional labor friction and verification challenges, integrating them into industrial and academic workflows yields substantial long-term productivity gains.}\par\noindent{\small 对应步骤：S1 [重申核心权衡与立场 Restate Nuanced Thesis]：In conclusion, while [Acknowledged Factor A] introduces real challenges in [Specific Aspect], [Core Thesis Stance] remains the most viable path forward. | 句首锚定旧信息/逻辑视角，主干动词紧跟主语，复杂修饰置于宾语之后（符合尾重原则 L011 与旧新信息流 L001）。}
  \item \textbf{S2 [整合全文两大支柱机制 Synthesize Dual Pillars]:} \textit{By pairing algorithmic deployment with transparent audit trails and targeted workforce retraining, institutions can capture efficiency breakthroughs without eroding human analytical oversight.}\par\noindent{\small 对应步骤：S2 [整合全文两大支柱机制 Synthesize Dual Pillars]：By combining [Pillar 1 from Body 1] with [Pillar 2 / Safeguard from Body 2], [Actors] can harness [Core Benefit] while containing [Primary Risk]. | 句首锚定旧信息/逻辑视角，主干动词紧跟主语，复杂修饰置于宾语之后（符合尾重原则 L011 与旧新信息流 L001）。}
  \item \textbf{S3 [理性边界展望（可选）Grounded Forward Look]:} \textit{Ultimately, the value of intelligent automation rests not on replacing human judgment, but on elevating the precision of the evidence upon which human judgment acts.}\par\noindent{\small 对应步骤：S3 [理性边界展望（可选）Grounded Forward Look]：Ultimately, the success of [Technology/Policy] depends not on its standalone capability, but on the institutional rigor with which it is governed. | 句首锚定旧信息/逻辑视角，主干动词紧跟主语，复杂修饰置于宾语之后（符合尾重原则 L011 与旧新信息流 L001）。}
\end{itemize}
\end{TipBox}
\medskip
\clearpage
\section{FT23: 托福独立写作满分范文：计算机高等教育应孤立偏向工程编码还是将数学基础与项目实践深度融合}\label{comp:FT23}
\begin{ConceptBox}{Composition Architecture \& Structural Blueprint Chain}
\textbf{Genre:} \texttt{TOEFL\_Independent} \quad \textbf{Paragraph Count:} 5\par
\textbf{Blueprint Chain:} \texttt{BP33 $\to$ BP34 $\to$ BP36 $\to$ BP39 $\to$ BP40}
\end{ConceptBox}
\subsection*{Paragraph 1: 议论文破题首段：具象趋势锚定 + 争议张力 + 明确中心论点（Zero-Template Hook \& Thesis） (\texttt{BP33})}
\noindent\textbf{Expressive Intent:}\enspace 我这一段想写六级、雅思 Task 2 或托福独立写作的第一段（Introduction）：彻底扔掉 'With the rapid development of science and technology, more and more people...' 等烂大街模板，用 3 句话写出考官眼前一亮的母语级破题段。\par
\noindent\textbf{Cohesion Formula:}\enspace {\small\color{AcadSlate}Concrete Phenomenon (No Cliché) $\to$ Balanced Tension $\to$ Conditional \& Specific Thesis Statement. 绑定逻辑规则：L023, L029.}\par\vspace{4pt}
\begin{ExampleBox}{Publication-Grade Authentic Text}
Large language models and automated code assistants have rapidly become routine study aids across university campuses. While proponents contend that these tools accelerate literature synthesis and lower programming barriers, critics warn that unchecked reliance erodes students' capacity for independent problem formulation. In my view, universities should integrate generative tools into advanced coursework only after students master foundational analytical skills, as premature delegation weakens debugging competence and obscures underlying mathematical principles.
\end{ExampleBox}
\begin{TipBox}{Rhetorical Move Breakdown \& Constituent Analysis}
\begin{itemize}[leftmargin=1.5em]
  \item \textbf{S1 [具象现象锚定 Concrete Trend Anchor]:} \textit{Large language models and automated code assistants have rapidly become routine study aids across university campuses.}\par\noindent{\small 对应步骤：S1 [具象现象锚定 Concrete Trend Anchor]：直接从具体技术变革或制度调整切入（如 Generative AI tools now draft code and summarize literature across university campuses），绝不用 With the development of...。 | 句首锚定旧信息/逻辑视角，主干动词紧跟主语，复杂修饰置于宾语之后（符合尾重原则 L011 与旧新信息流 L001）。}
  \item \textbf{S2 [核心争议张力 Core Tension / Debate]:} \textit{While proponents contend that these tools accelerate literature synthesis and lower programming barriers, critics warn that unchecked reliance erodes students' capacity for independent problem formulation.}\par\noindent{\small 对应步骤：S2 [核心争议张力 Core Tension / Debate]：While proponents argue that [View A: efficiency/accessibility], critics warn that [View B: skill erosion/inequality]. | 句首锚定旧信息/逻辑视角，主干动词紧跟主语，复杂修饰置于宾语之后（符合尾重原则 L011 与旧新信息流 L001）。}
  \item \textbf{S3 [带条件的中心论点 Nuanced Thesis Statement]:} \textit{In my view, universities should integrate generative tools into advanced coursework only after students master foundational analytical skills, as premature delegation weakens debugging competence and obscures underlying mathematical principles.}\par\noindent{\small 对应步骤：S3 [带条件的中心论点 Nuanced Thesis Statement]：In my view, rather than banning [X] outright or adopting it uncritically, [Institution/Society] should [Core Stance], because [Primary Reason 1] and [Primary Reason 2]. | 句首锚定旧信息/逻辑视角，主干动词紧跟主语，复杂修饰置于宾语之后（符合尾重原则 L011 与旧新信息流 L001）。}
\end{itemize}
\end{TipBox}
\medskip
\subsection*{Paragraph 2: 核心正方分论点段：主张 $\to$ 机制推演 $\to$ 具象行业微案例 $\to$ 扣题闭环（PEEL 满分段） (\texttt{BP34})}
\noindent\textbf{Expressive Intent:}\enspace 我这一段想写大作文的第一个主体段（Body Paragraph 1）：论证某个核心优势或必要性，做到有理有据、有具象场景支撑，绝不车轱辘话来回说。\par
\noindent\textbf{Cohesion Formula:}\enspace {\small\color{AcadSlate}Topic Claim $\to$ Causal Mechanism $\to$ Concrete Micro-Case (L021) $\to$ Warrant Closure. 绑定逻辑规则：L006, L012, L021.}\par\vspace{4pt}
\begin{ExampleBox}{Publication-Grade Authentic Text}
Delegating repetitive data compilation to automated systems enables professionals to concentrate on high-level diagnostic judgment. In data-intensive fields such as clinical medicine and structural engineering, practitioners traditionally spend hours cross-checking routine telemetry logs before formulating a hypothesis. When hospital radiology departments deploy automated screening algorithms to flag normal chest scans, physicians can redirect their cognitive bandwidth toward ambiguous pathology cases and patient consultation. Consequently, intelligent triage elevates both diagnostic throughput and clinical care quality.
\end{ExampleBox}
\begin{TipBox}{Rhetorical Move Breakdown \& Constituent Analysis}
\begin{itemize}[leftmargin=1.5em]
  \item \textbf{S1 [分论点主题句 Topic Claim]:} \textit{Delegating repetitive data compilation to automated systems enables professionals to concentrate on high-level diagnostic judgment.}\par\noindent{\small 对应步骤：S1 [分论点主题句 Topic Claim]：一句话清晰亮出本段唯一控制中心（[Policy/Factor X] strengthens [Outcome Y] by [Core Mechanism]）。 | 句首锚定旧信息/逻辑视角，主干动词紧跟主语，复杂修饰置于宾语之后（符合尾重原则 L011 与旧新信息流 L001）。}
  \item \textbf{S2 [因果机制推演 Causal Elaboration]:} \textit{In data-intensive fields such as clinical medicine and structural engineering, practitioners traditionally spend hours cross-checking routine telemetry logs before formulating a hypothesis.}\par\noindent{\small 对应步骤：S2 [因果机制推演 Causal Elaboration]：Explain HOW and WHY：阐明如果没有这层干扰或有了这项机制，认知/经济/技术流程会发生什么具体改变。 | 句首锚定旧信息/逻辑视角，主干动词紧跟主语，复杂修饰置于宾语之后（符合尾重原则 L011 与旧新信息流 L001）。}
  \item \textbf{S3 [具象微观案例 Concrete Micro-Case]:} \textit{When hospital radiology departments deploy automated screening algorithms to flag normal chest scans, physicians can redirect their cognitive bandwidth toward ambiguous pathology cases and patient consultation.}\par\noindent{\small 对应步骤：S3 [具象微观案例 Concrete Micro-Case]：For instance, when [Specific Group/Industry] [Specific Action], [Measurable/Observable Consequence]（绝不编造假百分比，用具体场景细节说话）。 | 句首锚定旧信息/逻辑视角，主干动词紧跟主语，复杂修饰置于宾语之后（符合尾重原则 L011 与旧新信息流 L001）。}
  \item \textbf{S4 [反事实对比或扣题升华 Warrant / Link Back]:} \textit{Consequently, intelligent triage elevates both diagnostic throughput and clinical care quality.}\par\noindent{\small 对应步骤：S4 [反事实对比或扣题升华 Warrant / Link Back]：Thus, [Summary Noun] shifts human effort from [Low-value Routine] toward [High-value Judgment]. | 句首锚定旧信息/逻辑视角，主干动词紧跟主语，复杂修饰置于宾语之后（符合尾重原则 L011 与旧新信息流 L001）。}
\end{itemize}
\end{TipBox}
\medskip
\subsection*{Paragraph 3: 让步-反驳段：优雅承认反方合理成分并强力扭转权重（Concession \& Refutation） (\texttt{BP36})}
\noindent\textbf{Expressive Intent:}\enspace 我这一段想在议论文中回应反方观点（Some people argue that...）：先大度承认反方担忧在短期或局部确实存在，再用结构性证据或制度设计证明该问题完全可控，从而锁死我方立场。\par
\noindent\textbf{Cohesion Formula:}\enspace {\small\color{AcadSlate}Admittedly (Valid Concern) $\to$ Short-term Reality $\to$ However (Zero-Sum Fallacy Pivot) $\to$ Historical Analogy \& Policy Mitigation. 绑定逻辑规则：L008, L023.}\par\vspace{4pt}
\begin{ExampleBox}{Publication-Grade Authentic Text}
Admittedly, critics raise a legitimate concern that autonomous robotics and generative software can displace routine clerical and assembly-line workers. During initial transition phases, rapid automation does create localized employment friction for personnel whose tasks follow rigid, codifiable rules. However, treating technological upgrading as a fixed-pie zero-sum game overlooks how productivity gains expand downstream service and engineering roles. Much as spreadsheet software eliminated manual ledger keeping while multiplying demand for financial data analysts, coupling industrial automation with subsidized reskilling programs enables displaced workers to transition into robot fleet supervision, calibration, and exception handling.
\end{ExampleBox}
\begin{TipBox}{Rhetorical Move Breakdown \& Constituent Analysis}
\begin{itemize}[leftmargin=1.5em]
  \item \textbf{S1 [呈现反方最有力的担忧 Articulate Opposing Concern]:} \textit{Admittedly, critics raise a legitimate concern that autonomous robotics and generative software can displace routine clerical and assembly-line workers.}\par\noindent{\small 对应步骤：S1 [呈现反方最有力的担忧 Articulate Opposing Concern]：Admittedly, critics raise a valid concern that [Opposing Argument, e.g., automation displaces entry-level clerical and manufacturing workers]. | 句首锚定旧信息/逻辑视角，主干动词紧跟主语，复杂修饰置于宾语之后（符合尾重原则 L011 与旧新信息流 L001）。}
  \item \textbf{S2 [承认局部真实性 Acknowledge Partial Validity]:} \textit{During initial transition phases, rapid automation does create localized employment friction for personnel whose tasks follow rigid, codifiable rules.}\par\noindent{\small 对应步骤：S2 [承认局部真实性 Acknowledge Partial Validity]：In the short term, rapid deployment of [Technology/Policy] can indeed create transitional friction for workers whose routines are easily codified. | 句首锚定旧信息/逻辑视角，主干动词紧跟主语，复杂修饰置于宾语之后（符合尾重原则 L011 与旧新信息流 L001）。}
  \item \textbf{S3 [釜底抽薪式转折 Structural Pivot]:} \textit{However, treating technological upgrading as a fixed-pie zero-sum game overlooks how productivity gains expand downstream service and engineering roles.}\par\noindent{\small 对应步骤：S3 [釜底抽薪式转折 Structural Pivot]：However, viewing [X] purely as a zero-sum threat overlooks how productivity gains generate entirely new categories of skilled employment. | 句首锚定旧信息/逻辑视角，主干动词紧跟主语，复杂修饰置于宾语之后（符合尾重原则 L011 与旧新信息流 L001）。}
  \item \textbf{S4 [历史/现实反证与制度补救 Evidence \& Mitigation]:} \textit{Much as spreadsheet software eliminated manual ledger keeping while multiplying demand for financial data analysts, coupling industrial automation with subsidized reskilling programs enables displaced workers to transition into robot fleet supervision, calibration, and exception handling.}\par\noindent{\small 对应步骤：S4 [历史/现实反证与制度补救 Evidence \& Mitigation]：Just as spreadsheet software reduced manual bookkeeping while expanding demand for financial analysts, pairing automation with targeted vocational retraining allows displaced workers to transition into system supervision and maintenance roles. | 句首锚定旧信息/逻辑视角，主干动词紧跟主语，复杂修饰置于宾语之后（符合尾重原则 L011 与旧新信息流 L001）。}
\end{itemize}
\end{TipBox}
\medskip
\subsection*{Paragraph 4: 反事实推理段：如果不采取某项原则，系统将如何一步步滑向失控（Counterfactual Escalation） (\texttt{BP39})}
\noindent\textbf{Expressive Intent:}\enspace 我这一段想在议论文或科研立项书中，通过严密的“反事实推演（If we fail to... / Absent X...）”强化紧迫感与说服力。\par
\noindent\textbf{Cohesion Formula:}\enspace {\small\color{AcadSlate}Positive Safeguard $\to$ Counterfactual Premise (Absent X...) $\to$ Cascading Breakdown $\to$ Precedent Confirmation. 绑定逻辑规则：L020, L003.}\par\vspace{4pt}
\begin{ExampleBox}{Publication-Grade Authentic Text}
Preserving verified human-authored corpora serves as an indispensable anchor for training successive generations of generative models. Absent this human data anchor---if future models were trained recursively on synthetic outputs generated by their predecessors---subtle approximation errors would compound across iterations. In each recursive cycle, low-probability tail events and nuanced stylistic variations are progressively truncated, driving the learned distribution toward high-density mode collapse. This degenerative feedback loop---known empirically as model autophagy---underscores why synthetic data augmentation can complement, but never wholly replace, curated real-world observations.
\end{ExampleBox}
\begin{TipBox}{Rhetorical Move Breakdown \& Constituent Analysis}
\begin{itemize}[leftmargin=1.5em]
  \item \textbf{S1 [正面基准命题 Positive Baseline]:} \textit{Preserving verified human-authored corpora serves as an indispensable anchor for training successive generations of generative models.}\par\noindent{\small 对应步骤：S1 [正面基准命题 Positive Baseline]：Maintaining rigorous [Principle/Standard X] serves as the primary safeguard against [Systemic Risk Y]. | 句首锚定旧信息/逻辑视角，主干动词紧跟主语，复杂修饰置于宾语之后（符合尾重原则 L011 与旧新信息流 L001）。}
  \item \textbf{S2 [反事实假设切入 Counterfactual Premise]:} \textit{Absent this human data anchor---if future models were trained recursively on synthetic outputs generated by their predecessors---subtle approximation errors would compound across iterations.}\par\noindent{\small 对应步骤：S2 [反事实假设切入 Counterfactual Premise]：Absent this safeguard---if [Actors] were to rely exclusively on [Unchecked Alternative]---[Immediate Vulnerability] would quickly emerge. | 句首锚定旧信息/逻辑视角，主干动词紧跟主语，复杂修饰置于宾语之后（符合尾重原则 L011 与旧新信息流 L001）。}
  \item \textbf{S3 [连锁恶化传导 Cascading Deterioration]:} \textit{In each recursive cycle, low-probability tail events and nuanced stylistic variations are progressively truncated, driving the learned distribution toward high-density mode collapse.}\par\noindent{\small 对应步骤：S3 [连锁恶化传导 Cascading Deterioration]：Once [Immediate Vulnerability] takes hold, [Secondary Consequence] compounds across cycles, eventually eroding [Core Foundation]. | 句首锚定旧信息/逻辑视角，主干动词紧跟主语，复杂修饰置于宾语之后（符合尾重原则 L011 与旧新信息流 L001）。}
  \item \textbf{S4 [实证印证或结论锁定 Empirical Echo]:} \textit{This degenerative feedback loop---known empirically as model autophagy---underscores why synthetic data augmentation can complement, but never wholly replace, curated real-world observations.}\par\noindent{\small 对应步骤：S4 [实证印证或结论锁定 Empirical Echo]：This degradation dynamic mirrors [Real-world/Technical Precedent], underscoring why [Principle X] cannot be compromised for short-term convenience. | 句首锚定旧信息/逻辑视角，主干动词紧跟主语，复杂修饰置于宾语之后（符合尾重原则 L011 与旧新信息流 L001）。}
\end{itemize}
\end{TipBox}
\medskip
\subsection*{Paragraph 5: 议论文收尾段：跳出简单重复，完成核心权衡整合与边界升华（Synthesis Conclusion） (\texttt{BP40})}
\noindent\textbf{Expressive Intent:}\enspace 我这一段想写六级、雅思 Task 2 或托福独立写作的结尾段（Conclusion）：用 2--3 句话收束全文，做到既呼应中心论点、整合正反权衡，又绝不喊假大空口号（Let us work together for a better tomorrow）。\par
\noindent\textbf{Cohesion Formula:}\enspace {\small\color{AcadSlate}Balanced Restatment $\to$ Dual-Pillar Synthesis $\to$ Grounded Governance Principle (Zero Slogans!). 绑定逻辑规则：L023, L029.}\par\vspace{4pt}
\begin{ExampleBox}{Publication-Grade Authentic Text}
In conclusion, while automated decision systems inevitably introduce transitional labor friction and verification challenges, integrating them into industrial and academic workflows yields substantial long-term productivity gains. By pairing algorithmic deployment with transparent audit trails and targeted workforce retraining, institutions can capture efficiency breakthroughs without eroding human analytical oversight. Ultimately, the value of intelligent automation rests not on replacing human judgment, but on elevating the precision of the evidence upon which human judgment acts.
\end{ExampleBox}
\begin{TipBox}{Rhetorical Move Breakdown \& Constituent Analysis}
\begin{itemize}[leftmargin=1.5em]
  \item \textbf{S1 [重申核心权衡与立场 Restate Nuanced Thesis]:} \textit{In conclusion, while automated decision systems inevitably introduce transitional labor friction and verification challenges, integrating them into industrial and academic workflows yields substantial long-term productivity gains.}\par\noindent{\small 对应步骤：S1 [重申核心权衡与立场 Restate Nuanced Thesis]：In conclusion, while [Acknowledged Factor A] introduces real challenges in [Specific Aspect], [Core Thesis Stance] remains the most viable path forward. | 句首锚定旧信息/逻辑视角，主干动词紧跟主语，复杂修饰置于宾语之后（符合尾重原则 L011 与旧新信息流 L001）。}
  \item \textbf{S2 [整合全文两大支柱机制 Synthesize Dual Pillars]:} \textit{By pairing algorithmic deployment with transparent audit trails and targeted workforce retraining, institutions can capture efficiency breakthroughs without eroding human analytical oversight.}\par\noindent{\small 对应步骤：S2 [整合全文两大支柱机制 Synthesize Dual Pillars]：By combining [Pillar 1 from Body 1] with [Pillar 2 / Safeguard from Body 2], [Actors] can harness [Core Benefit] while containing [Primary Risk]. | 句首锚定旧信息/逻辑视角，主干动词紧跟主语，复杂修饰置于宾语之后（符合尾重原则 L011 与旧新信息流 L001）。}
  \item \textbf{S3 [理性边界展望（可选）Grounded Forward Look]:} \textit{Ultimately, the value of intelligent automation rests not on replacing human judgment, but on elevating the precision of the evidence upon which human judgment acts.}\par\noindent{\small 对应步骤：S3 [理性边界展望（可选）Grounded Forward Look]：Ultimately, the success of [Technology/Policy] depends not on its standalone capability, but on the institutional rigor with which it is governed. | 句首锚定旧信息/逻辑视角，主干动词紧跟主语，复杂修饰置于宾语之后（符合尾重原则 L011 与旧新信息流 L001）。}
\end{itemize}
\end{TipBox}
\medskip
\clearpage
\section{FT24: 托福独立写作满分范文：生成模型训练中坚守真实人类语料锚点与防止模型自噬（Model Autophagy）}\label{comp:FT24}
\begin{ConceptBox}{Composition Architecture \& Structural Blueprint Chain}
\textbf{Genre:} \texttt{TOEFL\_Independent} \quad \textbf{Paragraph Count:} 5\par
\textbf{Blueprint Chain:} \texttt{BP33 $\to$ BP35 $\to$ BP39 $\to$ BP38 $\to$ BP40}
\end{ConceptBox}
\subsection*{Paragraph 1: 议论文破题首段：具象趋势锚定 + 争议张力 + 明确中心论点（Zero-Template Hook \& Thesis） (\texttt{BP33})}
\noindent\textbf{Expressive Intent:}\enspace 我这一段想写六级、雅思 Task 2 或托福独立写作的第一段（Introduction）：彻底扔掉 'With the rapid development of science and technology, more and more people...' 等烂大街模板，用 3 句话写出考官眼前一亮的母语级破题段。\par
\noindent\textbf{Cohesion Formula:}\enspace {\small\color{AcadSlate}Concrete Phenomenon (No Cliché) $\to$ Balanced Tension $\to$ Conditional \& Specific Thesis Statement. 绑定逻辑规则：L023, L029.}\par\vspace{4pt}
\begin{ExampleBox}{Publication-Grade Authentic Text}
Large language models and automated code assistants have rapidly become routine study aids across university campuses. While proponents contend that these tools accelerate literature synthesis and lower programming barriers, critics warn that unchecked reliance erodes students' capacity for independent problem formulation. In my view, universities should integrate generative tools into advanced coursework only after students master foundational analytical skills, as premature delegation weakens debugging competence and obscures underlying mathematical principles.
\end{ExampleBox}
\begin{TipBox}{Rhetorical Move Breakdown \& Constituent Analysis}
\begin{itemize}[leftmargin=1.5em]
  \item \textbf{S1 [具象现象锚定 Concrete Trend Anchor]:} \textit{Large language models and automated code assistants have rapidly become routine study aids across university campuses.}\par\noindent{\small 对应步骤：S1 [具象现象锚定 Concrete Trend Anchor]：直接从具体技术变革或制度调整切入（如 Generative AI tools now draft code and summarize literature across university campuses），绝不用 With the development of...。 | 句首锚定旧信息/逻辑视角，主干动词紧跟主语，复杂修饰置于宾语之后（符合尾重原则 L011 与旧新信息流 L001）。}
  \item \textbf{S2 [核心争议张力 Core Tension / Debate]:} \textit{While proponents contend that these tools accelerate literature synthesis and lower programming barriers, critics warn that unchecked reliance erodes students' capacity for independent problem formulation.}\par\noindent{\small 对应步骤：S2 [核心争议张力 Core Tension / Debate]：While proponents argue that [View A: efficiency/accessibility], critics warn that [View B: skill erosion/inequality]. | 句首锚定旧信息/逻辑视角，主干动词紧跟主语，复杂修饰置于宾语之后（符合尾重原则 L011 与旧新信息流 L001）。}
  \item \textbf{S3 [带条件的中心论点 Nuanced Thesis Statement]:} \textit{In my view, universities should integrate generative tools into advanced coursework only after students master foundational analytical skills, as premature delegation weakens debugging competence and obscures underlying mathematical principles.}\par\noindent{\small 对应步骤：S3 [带条件的中心论点 Nuanced Thesis Statement]：In my view, rather than banning [X] outright or adopting it uncritically, [Institution/Society] should [Core Stance], because [Primary Reason 1] and [Primary Reason 2]. | 句首锚定旧信息/逻辑视角，主干动词紧跟主语，复杂修饰置于宾语之后（符合尾重原则 L011 与旧新信息流 L001）。}
\end{itemize}
\end{TipBox}
\medskip
\subsection*{Paragraph 2: 深层原因剖析段：表层诱因与结构性根源双层挖掘（Dual-Layer Causal Diagnosis） (\texttt{BP35})}
\noindent\textbf{Expressive Intent:}\enspace 我这一段想在问题分析类作文（Why is this happening?）中，剖析某个社会/科研痛点背后的深层原因（既指出直接表层诱因，又挖出底层的考核激励或结构性根源）。\par
\noindent\textbf{Cohesion Formula:}\enspace {\small\color{AcadSlate}Dual-Cause Opening $\to$ Immediate Operational Driver $\to$ Deeper Institutional Root $\to$ Systemic Compounding. 绑定逻辑规则：L007, L005.}\par\vspace{4pt}
\begin{ExampleBox}{Publication-Grade Authentic Text}
The growing tendency among undergraduates to submit unverified AI-generated essays stems from both lowered friction costs and flawed assessment design. At an immediate level, conversational models produce grammatically polished drafts within seconds, tempting time-pressed students to bypass the arduous process of outlining and primary reading. More fundamentally, many university courses still rely on generic, formulaic essay prompts that reward surface fluency rather than original empirical critique. When effortless text generation meets standardized grading rubrics that fail to penalize hollow reasoning, superficial compliance inevitably crowds out deep inquiry.
\end{ExampleBox}
\begin{TipBox}{Rhetorical Move Breakdown \& Constituent Analysis}
\begin{itemize}[leftmargin=1.5em]
  \item \textbf{S1 [原因总起句 Dual-Cause Thesis]:} \textit{The growing tendency among undergraduates to submit unverified AI-generated essays stems from both lowered friction costs and flawed assessment design.}\par\noindent{\small 对应步骤：S1 [原因总起句 Dual-Cause Thesis]：The widespread emergence of [Problem X] stems from both immediate technological pressures and deeper institutional incentives. | 句首锚定旧信息/逻辑视角，主干动词紧跟主语，复杂修饰置于宾语之后（符合尾重原则 L011 与旧新信息流 L001）。}
  \item \textbf{S2 [第一层:} \textit{At an immediate level, conversational models produce grammatically polished drafts within seconds, tempting time-pressed students to bypass the arduous process of outlining and primary reading.}\par\noindent{\small 对应步骤：S2 [第一层：直接操作/技术诱因 Immediate Driver]：At the operational level, [Factor A] lowers the cost of [Behavior], tempting individuals/organizations to [Short-term Action]. | 句首锚定旧信息/逻辑视角，主干动词紧跟主语，复杂修饰置于宾语之后（符合尾重原则 L011 与旧新信息流 L001）。}
  \item \textbf{S3 [第二层:} \textit{More fundamentally, many university courses still rely on generic, formulaic essay prompts that reward surface fluency rather than original empirical critique.}\par\noindent{\small 对应步骤：S3 [第二层：深层制度/激励根源 Structural Root Cause]：More fundamentally, existing evaluation systems reward [Superficial Metric] rather than [Long-term Quality], creating powerful pressure to prioritize speed over rigor. | 句首锚定旧信息/逻辑视角，主干动词紧跟主语，复杂修饰置于宾语之后（符合尾重原则 L011 与旧新信息流 L001）。}
  \item \textbf{S4 [双因叠加后果 Compounding Effect]:} \textit{When effortless text generation meets standardized grading rubrics that fail to penalize hollow reasoning, superficial compliance inevitably crowds out deep inquiry.}\par\noindent{\small 对应步骤：S4 [双因叠加后果 Compounding Effect]：When effortless generation meets quantity-driven evaluation, [Problem X] becomes systemic rather than isolated. | 句首锚定旧信息/逻辑视角，主干动词紧跟主语，复杂修饰置于宾语之后（符合尾重原则 L011 与旧新信息流 L001）。}
\end{itemize}
\end{TipBox}
\medskip
\subsection*{Paragraph 3: 反事实推理段：如果不采取某项原则，系统将如何一步步滑向失控（Counterfactual Escalation） (\texttt{BP39})}
\noindent\textbf{Expressive Intent:}\enspace 我这一段想在议论文或科研立项书中，通过严密的“反事实推演（If we fail to... / Absent X...）”强化紧迫感与说服力。\par
\noindent\textbf{Cohesion Formula:}\enspace {\small\color{AcadSlate}Positive Safeguard $\to$ Counterfactual Premise (Absent X...) $\to$ Cascading Breakdown $\to$ Precedent Confirmation. 绑定逻辑规则：L020, L003.}\par\vspace{4pt}
\begin{ExampleBox}{Publication-Grade Authentic Text}
Preserving verified human-authored corpora serves as an indispensable anchor for training successive generations of generative models. Absent this human data anchor---if future models were trained recursively on synthetic outputs generated by their predecessors---subtle approximation errors would compound across iterations. In each recursive cycle, low-probability tail events and nuanced stylistic variations are progressively truncated, driving the learned distribution toward high-density mode collapse. This degenerative feedback loop---known empirically as model autophagy---underscores why synthetic data augmentation can complement, but never wholly replace, curated real-world observations.
\end{ExampleBox}
\begin{TipBox}{Rhetorical Move Breakdown \& Constituent Analysis}
\begin{itemize}[leftmargin=1.5em]
  \item \textbf{S1 [正面基准命题 Positive Baseline]:} \textit{Preserving verified human-authored corpora serves as an indispensable anchor for training successive generations of generative models.}\par\noindent{\small 对应步骤：S1 [正面基准命题 Positive Baseline]：Maintaining rigorous [Principle/Standard X] serves as the primary safeguard against [Systemic Risk Y]. | 句首锚定旧信息/逻辑视角，主干动词紧跟主语，复杂修饰置于宾语之后（符合尾重原则 L011 与旧新信息流 L001）。}
  \item \textbf{S2 [反事实假设切入 Counterfactual Premise]:} \textit{Absent this human data anchor---if future models were trained recursively on synthetic outputs generated by their predecessors---subtle approximation errors would compound across iterations.}\par\noindent{\small 对应步骤：S2 [反事实假设切入 Counterfactual Premise]：Absent this safeguard---if [Actors] were to rely exclusively on [Unchecked Alternative]---[Immediate Vulnerability] would quickly emerge. | 句首锚定旧信息/逻辑视角，主干动词紧跟主语，复杂修饰置于宾语之后（符合尾重原则 L011 与旧新信息流 L001）。}
  \item \textbf{S3 [连锁恶化传导 Cascading Deterioration]:} \textit{In each recursive cycle, low-probability tail events and nuanced stylistic variations are progressively truncated, driving the learned distribution toward high-density mode collapse.}\par\noindent{\small 对应步骤：S3 [连锁恶化传导 Cascading Deterioration]：Once [Immediate Vulnerability] takes hold, [Secondary Consequence] compounds across cycles, eventually eroding [Core Foundation]. | 句首锚定旧信息/逻辑视角，主干动词紧跟主语，复杂修饰置于宾语之后（符合尾重原则 L011 与旧新信息流 L001）。}
  \item \textbf{S4 [实证印证或结论锁定 Empirical Echo]:} \textit{This degenerative feedback loop---known empirically as model autophagy---underscores why synthetic data augmentation can complement, but never wholly replace, curated real-world observations.}\par\noindent{\small 对应步骤：S4 [实证印证或结论锁定 Empirical Echo]：This degradation dynamic mirrors [Real-world/Technical Precedent], underscoring why [Principle X] cannot be compromised for short-term convenience. | 句首锚定旧信息/逻辑视角，主干动词紧跟主语，复杂修饰置于宾语之后（符合尾重原则 L011 与旧新信息流 L001）。}
\end{itemize}
\end{TipBox}
\medskip
\subsection*{Paragraph 4: 利弊权衡段：证明结构性长期收益远超短期可管可控成本（Outweigh / Trade-off Analysis） (\texttt{BP38})}
\noindent\textbf{Expressive Intent:}\enspace 我这一段想回应 'Do the advantages outweigh the disadvantages?' 题型：把弊端定性为“暂时的、可通过技术/规则收敛的局部成本”，把优势定性为“不可替代的结构性长期突破”。\par
\noindent\textbf{Cohesion Formula:}\enspace {\small\color{AcadSlate}Asymmetry Thesis $\to$ Bound the Drawback (Temporary \& Mitigable) $\to$ Highlight Compounding Advantage (Irreplaceable) $\to$ Rational Verdict. 绑定逻辑规则：L008, L023.}\par\vspace{4pt}
\begin{ExampleBox}{Publication-Grade Authentic Text}
Weighing both dimensions demonstrates that the structural scientific gains of open-weight foundation models far outweigh their manageable misuse risks. While公开 releasing model weights lowers the barrier for generating deceptive phishing content, such downstream abuse can be substantially contained through cryptographic provenance standards and endpoint spam filters. By contrast, open weights provide academic laboratories and resource-constrained hospitals with transparent, auditable architectures that proprietary black-box APIs cannot substitute---accelerating reproducible breakthroughs in protein folding and rare-disease diagnostics. Because the security risks are mitigable via defense-in-depth whereas the scientific transparency is irreplaceable, open research ecosystems yield a decisive net benefit.
\end{ExampleBox}
\begin{TipBox}{Rhetorical Move Breakdown \& Constituent Analysis}
\begin{itemize}[leftmargin=1.5em]
  \item \textbf{S1 [权衡定性主题句 Asymmetry Thesis]:} \textit{Weighing both dimensions demonstrates that the structural scientific gains of open-weight foundation models far outweigh their manageable misuse risks.}\par\noindent{\small 对应步骤：S1 [权衡定性主题句 Asymmetry Thesis]：Weighing both dimensions reveals that the long-term structural benefits of [X] far outweigh its manageable transitional costs. | 句首锚定旧信息/逻辑视角，主干动词紧跟主语，复杂修饰置于宾语之后（符合尾重原则 L011 与旧新信息流 L001）。}
  \item \textbf{S2 [将弊端限定为可收敛成本 Bounding the Drawback]:} \textit{While公开 releasing model weights lowers the barrier for generating deceptive phishing content, such downstream abuse can be substantially contained through cryptographic provenance standards and endpoint spam filters.}\par\noindent{\small 对应步骤：S2 [将弊端限定为可收敛成本 Bounding the Drawback]：While [Drawback A, e.g., initial capital outlay / training overhead] poses a tangible hurdle, this cost is bounded and diminishes rapidly as [Mitigating Factor, e.g., hardware amortization / open-source standardization] matures. | 句首锚定旧信息/逻辑视角，主干动词紧跟主语，复杂修饰置于宾语之后（符合尾重原则 L011 与旧新信息流 L001）。}
  \item \textbf{S3 [凸显优势的不可替代性 Highlighting Irreversible Gain]:} \textit{By contrast, open weights provide academic laboratories and resource-constrained hospitals with transparent, auditable architectures that proprietary black-box APIs cannot substitute---accelerating reproducible breakthroughs in protein folding and rare-disease diagnostics.}\par\noindent{\small 对应步骤：S3 [凸显优势的不可替代性 Highlighting Irreversible Gain]：By contrast, the core advantage---[Benefit B, e.g., eliminating human exposure to hazardous environments / discovering non-linear scientific laws]---delivers compounding value that manual methods cannot replicate at any cost. | 句首锚定旧信息/逻辑视角，主干动词紧跟主语，复杂修饰置于宾语之后（符合尾重原则 L011 与旧新信息流 L001）。}
  \item \textbf{S4 [权衡结论 Trade-off Verdict]:} \textit{Because the security risks are mitigable via defense-in-depth whereas the scientific transparency is irreplaceable, open research ecosystems yield a decisive net benefit.}\par\noindent{\small 对应步骤：S4 [权衡结论 Trade-off Verdict]：Therefore, addressing the friction through targeted safeguards is far more rational than forgoing the transformative gain. | 句首锚定旧信息/逻辑视角，主干动词紧跟主语，复杂修饰置于宾语之后（符合尾重原则 L011 与旧新信息流 L001）。}
\end{itemize}
\end{TipBox}
\medskip
\subsection*{Paragraph 5: 议论文收尾段：跳出简单重复，完成核心权衡整合与边界升华（Synthesis Conclusion） (\texttt{BP40})}
\noindent\textbf{Expressive Intent:}\enspace 我这一段想写六级、雅思 Task 2 或托福独立写作的结尾段（Conclusion）：用 2--3 句话收束全文，做到既呼应中心论点、整合正反权衡，又绝不喊假大空口号（Let us work together for a better tomorrow）。\par
\noindent\textbf{Cohesion Formula:}\enspace {\small\color{AcadSlate}Balanced Restatment $\to$ Dual-Pillar Synthesis $\to$ Grounded Governance Principle (Zero Slogans!). 绑定逻辑规则：L023, L029.}\par\vspace{4pt}
\begin{ExampleBox}{Publication-Grade Authentic Text}
In conclusion, while automated decision systems inevitably introduce transitional labor friction and verification challenges, integrating them into industrial and academic workflows yields substantial long-term productivity gains. By pairing algorithmic deployment with transparent audit trails and targeted workforce retraining, institutions can capture efficiency breakthroughs without eroding human analytical oversight. Ultimately, the value of intelligent automation rests not on replacing human judgment, but on elevating the precision of the evidence upon which human judgment acts.
\end{ExampleBox}
\begin{TipBox}{Rhetorical Move Breakdown \& Constituent Analysis}
\begin{itemize}[leftmargin=1.5em]
  \item \textbf{S1 [重申核心权衡与立场 Restate Nuanced Thesis]:} \textit{In conclusion, while automated decision systems inevitably introduce transitional labor friction and verification challenges, integrating them into industrial and academic workflows yields substantial long-term productivity gains.}\par\noindent{\small 对应步骤：S1 [重申核心权衡与立场 Restate Nuanced Thesis]：In conclusion, while [Acknowledged Factor A] introduces real challenges in [Specific Aspect], [Core Thesis Stance] remains the most viable path forward. | 句首锚定旧信息/逻辑视角，主干动词紧跟主语，复杂修饰置于宾语之后（符合尾重原则 L011 与旧新信息流 L001）。}
  \item \textbf{S2 [整合全文两大支柱机制 Synthesize Dual Pillars]:} \textit{By pairing algorithmic deployment with transparent audit trails and targeted workforce retraining, institutions can capture efficiency breakthroughs without eroding human analytical oversight.}\par\noindent{\small 对应步骤：S2 [整合全文两大支柱机制 Synthesize Dual Pillars]：By combining [Pillar 1 from Body 1] with [Pillar 2 / Safeguard from Body 2], [Actors] can harness [Core Benefit] while containing [Primary Risk]. | 句首锚定旧信息/逻辑视角，主干动词紧跟主语，复杂修饰置于宾语之后（符合尾重原则 L011 与旧新信息流 L001）。}
  \item \textbf{S3 [理性边界展望（可选）Grounded Forward Look]:} \textit{Ultimately, the value of intelligent automation rests not on replacing human judgment, but on elevating the precision of the evidence upon which human judgment acts.}\par\noindent{\small 对应步骤：S3 [理性边界展望（可选）Grounded Forward Look]：Ultimately, the success of [Technology/Policy] depends not on its standalone capability, but on the institutional rigor with which it is governed. | 句首锚定旧信息/逻辑视角，主干动词紧跟主语，复杂修饰置于宾语之后（符合尾重原则 L011 与旧新信息流 L001）。}
\end{itemize}
\end{TipBox}
\medskip
\clearpage